\section{The Quot scheme}
\label{mk-chap:Picard_QuotScheme}

\section{Hilbert polynomial of a coherent sheaf}
\label{mk-sec:hilbert_polynomial}

\begin{definition}[Hilbert polynomial]\label{mk-def:hilbert_polynomial}

  \uses{mk-def:fiber_hilbert_function, mk-lem:hilbertPolynomial_unique}
  \dcref{III.5.2,custom}
  \textit{Source: \cite{nitsure-hilbert-quot}, \S 1 (FGA Explained Ch.~5); cf.\ \cite{hartshorne-ag},
  III.5.2.}
  Let \(X \to S\) be a finite-type morphism of noetherian schemes and let \(L\)
  be a line bundle on \(X\). For a coherent sheaf \(\mathcal{F}\) on \(X\) whose
  schematic support is proper over \(S\) and a point \(s \in S\), write \(X_s\)
  for the fibre over \(s\), \(\mathcal{F}_s = \mathcal{F}|_{X_s}\) and
  \(L_s = L|_{X_s}\). The \emph{Hilbert polynomial of \(\mathcal{F}\) at \(s\)
  relative to \(L\)} is the unique polynomial
  \(\Phi_{\mathcal{F},s} \in \mathbb{Q}[\lambda]\) that agrees, for all
  \(m \gg 0\), with the graded Hilbert function
  (\ref{mk-def:fiber_hilbert_function})
  \[
    m \;\longmapsto\; \dim_{\kappa(s)} \Gamma\bigl(X_s,\, \mathcal{F}_s \otimes L_s^{\otimes m}\bigr)
  \]
  of the section module
  \(\bigoplus_{m \ge 0} \Gamma(X_s, \mathcal{F}_s \otimes L_s^{\otimes m})\)
  (\ref{mk-def:sectionGradedModule}), viewed as a finitely generated graded module
  over the homogeneous section ring
  \(\bigoplus_{m \ge 0} \Gamma(X_s, L_s^{\otimes m})\)
  (\ref{mk-def:sectionGradedRing}). Uniqueness of \(\Phi_{\mathcal{F},s}\) is
  unconditional (\ref{mk-lem:hilbertPolynomial_unique}: two polynomials agreeing
  at all sufficiently large integers coincide); existence, for \(X_s\) proper
  over \(\kappa(s)\), \(L_s\) ample and \(\mathcal{F}_s\) coherent, is the
  content of \ref{mk-thm:hilbertPoly_of_sectionModule}.

  This is the classical Hilbert polynomial. When \(\mathcal{F}_s\) is coherent
  and \(L_s\) is ample on the proper \(\kappa(s)\)-scheme \(X_s\), Serre vanishing
  gives \(H^i(X_s, \mathcal{F}_s \otimes L_s^{\otimes m}) = 0\) for all \(i > 0\)
  and \(m \gg 0\), so for \(m \gg 0\)
  \[
    \dim_{\kappa(s)} \Gamma(X_s, \mathcal{F}_s \otimes L_s^{\otimes m})
    \;=\; \chi(X_s,\, \mathcal{F}_s \otimes L_s^{\otimes m})
    \;=\; \sum_{i \ge 0} (-1)^i \dim_{\kappa(s)} H^i(X_s,\, \mathcal{F}_s \otimes L_s^{\otimes m}),
  \]
  and hence \(\Phi_{\mathcal{F},s}\) coincides with the Euler-characteristic
  Hilbert polynomial of Snapper's Lemma -- no numerical invariant is lost. The
  construction above, however, uses only the global sections \(H^0\): the higher
  cohomology enters solely through the classical statement that the two functions
  agree for \(m \gg 0\), never through the definition of \(\Phi_{\mathcal{F},s}\)
  itself. When \(\mathcal{F}\) is \(S\)-flat the function
  \(s \mapsto \Phi_{\mathcal{F},s}\) is locally constant on \(S\), so \(\Phi\) is
  well-defined on connected components.
\end{definition}

\section{Graded Hilbert polynomial}
\label{mk-sec:graded_hilbert_polynomial}

This section gives the graded-Hilbert-function construction that underlies
\ref{mk-def:hilbert_polynomial}. It first records how a fibre \(X_s\) of a
family \(\pi : X \to S\) carries the data the construction consumes: the
restrictions \(\mathcal{F}_s\) and \(L_s\)
(\ref{mk-def:fiber_module_restriction}) and the \(\kappa(s)\)-vector-space
structure on their sections (\ref{mk-def:fiber_sections_module}). Working over
the fibre \(X_s\) with its ample line bundle \(L_s\) and coherent sheaf
\(\mathcal{F}_s\), it builds the section graded ring
(\ref{mk-def:sectionGradedRing}) and module (\ref{mk-def:sectionGradedModule})
together with the graded Hilbert function
(\ref{mk-def:fiber_hilbert_function}), records their finite generation
(\ref{mk-lem:sectionGradedModule_fg}), and extracts the Hilbert polynomial
(\ref{mk-thm:hilbertPoly_of_sectionModule}): the graded Hilbert--Serre
rationality of the Hilbert series (\ref{mk-lem:gradedHilbertSerre_rational})
combines with Mathlib's polynomial-extraction lemma
(\ref{mk-lem:hilbertPoly_exists_mathlib}) in \ref{mk-lem:hilbertPoly_of_isRatHilb},
and the extracted polynomial is unique by
\ref{mk-lem:hilbertPolynomial_unique}. The route uses only \(H^0\)
(global sections), so it is independent of higher coherent cohomology.

\begin{definition}[Restriction of a sheaf of modules to a fibre]\label{mk-def:fiber_module_restriction}

  \dcref{custom}
  \textit{Source: \cite{nitsure-hilbert-quot}, \S 1 (FGA Explained Ch.~5).}
  Let \(\pi : X \to S\) be a morphism of schemes and \(s \in S\) a point with
  residue field \(\kappa(s)\). The \emph{fibre} of \(\pi\) over \(s\) is the
  scheme
  \[
    X_s \;=\; X \times_S \Spec \kappa(s),
  \]
  the base change of \(\pi\) along the canonical morphism
  \(\Spec \kappa(s) \to S\); write \(\iota_s : X_s \to X\) for the projection
  to \(X\) (the \emph{fibre embedding}). For a sheaf of
  \(\mathcal{O}_X\)-modules \(\mathcal{F}\) on \(X\), the \emph{restriction of
  \(\mathcal{F}\) to the fibre} is the module pullback
  \[
    \mathcal{F}_s \;=\; \mathcal{F}|_{X_s} \;:=\; \iota_s^{*} \mathcal{F},
  \]
  a sheaf of \(\mathcal{O}_{X_s}\)-modules. In particular a line bundle \(L\)
  on \(X\) restricts to a line bundle \(L_s = L|_{X_s}\) on \(X_s\).
\end{definition}

\begin{definition}[\(\kappa(s)\)-structure on sections over a fibre]\label{mk-def:fiber_sections_module}

  \uses{mk-def:fiber_module_restriction}
  \dcref{custom}
  \textit{Source: \cite{nitsure-hilbert-quot}, \S 1 (FGA Explained Ch.~5).}
  Let \(\pi : X \to S\) be a morphism of schemes and \(s \in S\). The fibre
  \(X_s = X \times_S \Spec \kappa(s)\) (\ref{mk-def:fiber_module_restriction})
  carries a structural morphism \(p_s : X_s \to \Spec \kappa(s)\), the second
  projection of the fibre product. Taking global sections of \(p_s\) and
  identifying \(\Gamma(\Spec \kappa(s), \mathcal{O}) = \kappa(s)\) yields the
  \emph{structural ring homomorphism of the fibre}
  \[
    \kappa(s) \;\longrightarrow\; \Gamma\bigl(X_s,\, \mathcal{O}_{X_s}\bigr).
  \]
  For any sheaf of \(\mathcal{O}_{X_s}\)-modules \(\mathcal{G}\) on the fibre,
  the global sections \(\Gamma(X_s, \mathcal{G})\) form a module over
  \(\Gamma(X_s, \mathcal{O}_{X_s})\), hence, by restriction of scalars along
  the structural homomorphism, a \(\kappa(s)\)-vector space. All dimensions
  \(\dim_{\kappa(s)}\) of section spaces over a fibre are taken with respect
  to this \(\kappa(s)\)-structure.
\end{definition}

\begin{definition}
[Section graded ring]
  \label{mk-def:sectionGradedRing}
  \dcref{I.7,custom}

  \uses{mk-lem:sectionGradedRing_gcommSemiring}
  \textit{Source: \cite{nitsure-hilbert-quot}, \S 1 (FGA Explained Ch.~5); the section graded
  ring is the homogeneous coordinate ring of \cite{hartshorne-ag}, I.7.}
  Let \(s \in S\) and let \(L_s\) be a line bundle on the fibre \(X_s\), a scheme
  of finite type over the residue field \(\kappa(s)\). The \emph{section graded
  ring} of \(L_s\) is the graded commutative ring
  \[
    R(X_s, L_s) \;=\; \bigoplus_{m \ge 0} \Gamma\bigl(X_s,\, L_s^{\otimes m}\bigr),
  \]
  with degree-\(m\) component \(\Gamma(X_s, L_s^{\otimes m})\) and multiplication
  the tensor product of sections
  \(\Gamma(X_s, L_s^{\otimes m}) \otimes_{\kappa(s)}
  \Gamma(X_s, L_s^{\otimes m'}) \to \Gamma(X_s, L_s^{\otimes (m+m')})\),
  \(\sigma \otimes \tau \mapsto \sigma \cdot \tau\). Its degree-zero component is
  \(\Gamma(X_s, \mathcal{O}_{X_s})\), which contains \(\kappa(s)\). When \(X_s\) is
  proper over \(\kappa(s)\) and \(L_s\) is ample, \(R(X_s, L_s)\) is a finitely
  generated -- hence Noetherian -- graded \(\kappa(s)\)-algebra (a suitable
  Veronese power of \(L_s\) is very ample and embeds \(X_s\) into a projective
  space, exhibiting \(R(X_s, L_s)\) up to a finite extension as a quotient of a
  polynomial ring).
\end{definition}

\begin{definition}
[Section graded module]
  \label{mk-def:sectionGradedModule}
  \dcref{I.7,custom}

  \uses{mk-def:sectionGradedRing, mk-lem:sectionGradedModule_gmodule}
  \textit{Source: \cite{nitsure-hilbert-quot}, \S 1 (FGA Explained Ch.~5); cf.\ \cite{hartshorne-ag},
  I.7.}
  Let \(\mathcal{F}_s\) be a coherent sheaf on the fibre \(X_s\) and \(L_s\) a
  line bundle on \(X_s\). The \emph{section graded module} of \(\mathcal{F}_s\)
  twisted by \(L_s\) is the graded \(R(X_s, L_s)\)-module
  (\ref{mk-def:sectionGradedRing})
  \[
    M(X_s, \mathcal{F}_s, L_s)
    \;=\; \bigoplus_{m \ge 0} \Gamma\bigl(X_s,\, \mathcal{F}_s \otimes L_s^{\otimes m}\bigr),
  \]
  with degree-\(m\) component the \(\kappa(s)\)-vector space
  \(\Gamma(X_s, \mathcal{F}_s \otimes L_s^{\otimes m})\) and scalar action given
  by the section tensor product
  \(\Gamma(X_s, L_s^{\otimes m}) \otimes_{\kappa(s)}
  \Gamma(X_s, \mathcal{F}_s \otimes L_s^{\otimes m'})
  \to \Gamma(X_s, \mathcal{F}_s \otimes L_s^{\otimes (m+m')})\). Its graded Hilbert
  function is \(m \mapsto \dim_{\kappa(s)}
  \Gamma(X_s, \mathcal{F}_s \otimes L_s^{\otimes m})\).
\end{definition}

\begin{definition}[Graded Hilbert function along a fibre]\label{mk-def:fiber_hilbert_function}

  \uses{mk-def:fiber_module_restriction, mk-def:fiber_sections_module,
    mk-def:sheafModuleTwist}
  \dcref{custom}
  \textit{Source: \cite{nitsure-hilbert-quot}, \S 1 (FGA Explained Ch.~5).}
  Let \(S\) be a locally noetherian scheme, \(\pi : X \to S\) a morphism
  locally of finite type, \(L\) a line bundle and \(\mathcal{F}\) a sheaf of
  \(\mathcal{O}_X\)-modules on \(X\), and \(s \in S\). The \emph{graded
  Hilbert function of \(\mathcal{F}\) at \(s\) relative to \(L\)} is the
  function
  \[
    h_{\mathcal{F},s} : \mathbb{N} \longrightarrow \mathbb{N}, \qquad
    h_{\mathcal{F},s}(m) \;=\;
    \dim_{\kappa(s)} \Gamma\bigl(X_s,\, \mathcal{F}_s \otimes L_s^{\otimes m}\bigr),
  \]
  where \(\mathcal{F}_s\) and \(L_s\) are the restrictions to the fibre
  \(X_s\) (\ref{mk-def:fiber_module_restriction}), the twist
  \(\mathcal{F}_s \otimes L_s^{\otimes m}\) is the \(m\)-twist of
  \(\mathcal{F}_s\) by \(L_s\) on \(X_s\) (\ref{mk-def:sheafModuleTwist}), and
  the dimension is taken with respect to the \(\kappa(s)\)-vector-space
  structure on sections over the fibre (\ref{mk-def:fiber_sections_module}).
  This is exactly the graded Hilbert function of the section graded module
  \(M(X_s, \mathcal{F}_s, L_s)\) (\ref{mk-def:sectionGradedModule}); it uses
  only global sections \(H^0\).
\end{definition}

\begin{lemma}
[Finite generation of the section graded module]
  \label{mk-lem:sectionGradedModule_fg}
  \uses{mk-def:sectionGradedModule, mk-def:sectionGradedRing, mk-def:projective_with}
  \dcref{I.7,custom}
  \textit{Source: \cite{nitsure-hilbert-quot}, \S 1 (FGA Explained Ch.~5); Serre's theorem on
  ample line bundles, \cite{hartshorne-ag}, I.7 / II.5.}
  Let \(X_s\) be projective over the field \(\kappa(s)\) carrying the very
  ample line bundle \(L_s\) (\ref{mk-def:projective_with}: a closed immersion
  \(X_s \hookrightarrow \mathbb{P}^N_{\kappa(s)}\) with
  \(L_s \cong \mathcal{O}(1)|_{X_s}\)), and let \(\mathcal{F}_s\) be a
  coherent sheaf on \(X_s\).
  Then the section graded ring \(R(X_s, L_s)\) (\ref{mk-def:sectionGradedRing}) is a
  Noetherian graded ring, and the section graded module
  \(M(X_s, \mathcal{F}_s, L_s)\) (\ref{mk-def:sectionGradedModule}) is a finitely
  generated graded \(R(X_s, L_s)\)-module. In particular each component
  \(\Gamma(X_s, \mathcal{F}_s \otimes L_s^{\otimes m})\) is a finite-dimensional
  \(\kappa(s)\)-vector space.
  (For an ample \(L_s\) the same conclusions follow by passing to a very
  ample Veronese power; the very ample form stated here is the one consumed
  by the Quot-scheme construction.)
\end{lemma}

\begin{proof}
  \uses{mk-def:sectionGradedRing, mk-def:sectionGradedModule}
  The closed immersion \(i : X_s \hookrightarrow \mathbb{P}^N_{\kappa(s)}\)
  with \(L_s \cong i^* \mathcal{O}(1)\) reduces the claims to projective
  space: pushing forward along \(i\) identifies
  \(\Gamma(X_s, \mathcal{F}_s \otimes L_s^{\otimes m})\) with
  \(\Gamma(\mathbb{P}^N, (i_* \mathcal{F}_s)(m))\), and \(i_* \mathcal{F}_s\)
  is coherent. By Serre's theorem the section
  graded ring \(R(X_s, L_s)\) is a finitely generated graded
  \(\kappa(s)\)-algebra (a quotient of the polynomial ring in \(N+1\)
  variables in large degrees), so it is Noetherian. For the coherent sheaf
  \(\mathcal{F}_s\), Serre's theorem yields that
  \(\bigoplus_{m \ge 0} \Gamma(X_s, \mathcal{F}_s \otimes L_s^{\otimes m})\) is a
  finitely generated graded module over \(R(X_s, L_s)\): the twists
  \(\mathcal{F}_s \otimes L_s^{\otimes m}\) are globally generated for
  \(m \gg 0\), and the finitely many generating sections in low degrees together
  with the multiplication maps generate \(M(X_s, \mathcal{F}_s, L_s)\) over
  \(R(X_s, L_s)\). Finite dimensionality of each component over \(\kappa(s)\)
  follows from coherence and properness (finiteness of cohomology of coherent
  sheaves on a proper scheme over a field).
\end{proof}

\begin{lemma}[Existence and uniqueness of the Hilbert polynomial of a rational
  graded series]
  \label{mk-lem:hilbertPoly_exists_mathlib}
  \dcref{custom}

  Let \(F\) be a field of characteristic zero, let \(p \in F[X]\), and let
  \(d \in \mathbb{N}\). Then there is a unique polynomial \(h \in F[X]\) such that
  for some \(N \in \mathbb{N}\) and all \(n > N\) the coefficient of \(X^n\) in the
  power-series expansion of \(p \cdot (1 - X)^{-d} \in F\llbracket X\rrbracket\)
  equals \(h(n)\):
  \[
    \exists!\, h \in F[X],\ \exists N,\ \forall n > N : \;
    \bigl[ X^n \bigr]\bigl( p \cdot (1-X)^{-d} \bigr) \;=\; h(n).
  \]
  Equivalently, the eventual coefficient function of any rational power series
  with denominator a power of \(1 - X\) is, for \(n \gg 0\), given by a unique
  polynomial \(h\) (Mathlib's \(\mathrm{hilbertPoly}\,p\,d\)). The hypothesis
  \([\mathrm{CharZero}\,F]\) is satisfied by \(F = \mathbb{Q}\), which is the
  instance used to obtain \(\Phi_{\mathcal{F},s} \in \mathbb{Q}[\lambda]\).
\end{lemma}

\begin{lemma}[Additivity of \(\dim_\kappa\) on a short exact sequence
  (rank--nullity)]
  \label{mk-lem:finrank_ses_additive_mathlib}
  \dcref{custom}

  Let \(\kappa\) be a field, \(V\) a finite-dimensional \(\kappa\)-vector space and
  \(W \subseteq V\) a subspace. Then
  \[
    \dim_\kappa (V / W) + \dim_\kappa W \;=\; \dim_\kappa V .
  \]
  Equivalently, \(\dim_\kappa = \operatorname{finrank}_\kappa\) is additive on a
  short exact sequence \(0 \to W \to V \to V/W \to 0\) of finite-dimensional
  \(\kappa\)-vector spaces. This is the additivity that converts the degreewise
  short exact sequences of graded components into additive relations among the
  coefficients of the Hilbert series.
\end{lemma}

\begin{lemma}[The rational power series \((1 - X)^{-d}\)]
  \label{mk-lem:invOneSubPow_mathlib}
  \dcref{custom}

  Let \(S\) be a commutative ring and \(d \in \mathbb{N}\). Then \((1 - X)^d\) is a
  unit in the power series ring \(S\llbracket X\rrbracket\), and Mathlib names its
  inverse
  \[
    \operatorname{invOneSubPow}(S, d) \;=\; (1 - X)^{-d}
      \;\in\; S\llbracket X\rrbracket^{\times}.
  \]
  Together with the coercion
  \(\operatorname{toPowerSeries} : F[X] \to F\llbracket X\rrbracket\) of a
  polynomial to its power series, this is the data in which a rational Hilbert
  series \(p \cdot (1 - X)^{-d}\) is expressed; it is the same rational-series shape
  whose eventual coefficients are extracted by
  \ref{mk-lem:hilbertPoly_exists_mathlib}.
\end{lemma}

\begin{lemma}[Graded Hilbert--Serre: rationality of the Hilbert series]\label{mk-lem:gradedHilbertSerre_rational}
  \dcref{00K1,custom}

  \uses{mk-lem:finrank_ses_additive_mathlib, mk-lem:invOneSubPow_mathlib,
    mk-def:ratHilb, mk-lem:ratHilb_ofEventuallyZero, mk-lem:ratHilb_ofDiffEq,
    mk-lem:ratHilb_bump, mk-def:graded_subquotientHilb,
    mk-lem:graded_subquotient_isRatHilb}
  \textit{Source: \cite{stacks-project}, "Commutative Algebra", Tag 00K1
  (proposition on graded Hilbert polynomials); cf.\ \cite{nitsure-hilbert-quot}, \S 1 (Snapper's
  Lemma), \cite{hartshorne-ag}, I.7.}
  Let \(R\) be a Noetherian graded ring whose degree-zero part is a field
  \(\kappa\) and which is generated as a \(\kappa\)-algebra by finitely many
  elements, and let \(M\) be a finitely generated graded \(R\)-module with
  finite-dimensional graded components \(M_n\). Then there exist
  \(p \in \mathbb{Q}[X]\) and \(d \in \mathbb{N}\) and an \(N \in \mathbb{N}\)
  such that for all \(n > N\)
  \[
    \dim_{\kappa} M_n \;=\; \bigl[ X^n \bigr]\bigl( p \cdot (1 - X)^{-d}\bigr);
  \]
  that is, the Hilbert series \(\sum_{n \ge 0} (\dim_\kappa M_n)\, X^n\) is, up to
  a polynomial correction in low degrees, the rational function
  \(p \cdot (1 - X)^{-d}\), expressed in the power-series form
  \(\operatorname{toPowerSeries}(p) \cdot \operatorname{invOneSubPow}(\mathbb{Q}, d)\)
  (\ref{mk-lem:invOneSubPow_mathlib}).

\end{lemma}

\begin{proof}
  \uses{mk-def:ratHilb, mk-lem:ratHilb_ofEventuallyZero, mk-lem:ratHilb_ofDiffEq,
    mk-lem:ratHilb_bump, mk-lem:graded_degreewise_finrank_diff,
    mk-def:graded_subquotientHilb, mk-lem:graded_subquotient_ker_coker,
    mk-lem:graded_subquotient_degreewise_diff, mk-lem:graded_subquotient_finite_transfer,
    mk-lem:graded_subquotient_isRatHilb}
  We argue with the abstract graded \(\kappa\)-algebra and its finitely generated
  graded module. The proof is the ambient subquotient induction described below.

  \emph{Reduction to generation in degree \(1\).} Since \(R\) is Noetherian and
  finitely generated as a \(\kappa\)-algebra, it is generated by finitely many
  homogeneous elements of positive degree. Replacing \(R\) by a Veronese
  subalgebra \(R^{(e)} = \bigoplus_n R_{en}\) for a suitable \(e\) and re-grading
  (and \(M\) by the corresponding Veronese module), we may assume \(R\) is
  generated by its degree-one part, \(\kappa[R_1] = R\). The degree-one part
  \(R_1\) is then a finite-dimensional \(\kappa\)-vector space whose \(\kappa\)-basis
  \(x_1, \dots, x_r\) generates \(R\) as a \(\kappa\)-algebra; \(r = \dim_\kappa R_1\)
  is the number on which we induct.

  \emph{Set-up of the ambient data.} Regard \(M\) as the fixed ambient graded
  \(\kappa\)-module with internal decomposition \(n \mapsto \mathcal{M}_n := M_n\)
  (\ref{mk-def:sectionGradedModule}); the components are finite-dimensional. For each
  basis vector \(x_i \in R_1\) let \(t_i : M \to M\) be multiplication by \(x_i\), a
  \(\kappa\)-linear endomorphism that raises degree by one (it sends \(M_n\) into
  \(M_{n+1}\), \ref{mk-lem:isHomogeneousElem_graded_smul_mathlib}). Since \(R\) is
  commutative the \(t_i\) pairwise commute, and they trivially preserve the pair of
  homogeneous submodules \((N, N') = (\top, \bot)\). Because \(R\) is generated by
  \(x_1, \dots, x_r\) and \(M\) is finitely generated over \(R\), \(M\) is finitely
  generated as a module over the polynomial ring \(\kappa[x_1, \dots, x_r]\) acting
  through the \(t_i\); this is the finiteness condition required below.

  \emph{Invoking the ambient subquotient induction.} Apply
  \ref{mk-lem:graded_subquotient_isRatHilb} to the pair \((N, N') = (\top, \bot)\)
  with the \(r\) commuting degree-\((+1)\) endomorphisms \(t_1, \dots, t_r\). It
  yields \(\mathrm{IsRatHilb}\) of order \(r\) for the ambient subquotient Hilbert
  function of this pair (\ref{mk-def:graded_subquotientHilb}),
  \[
    h(n) \;=\; \dim_\kappa(\top \cap \mathcal{M}_n) - \dim_\kappa(\bot \cap \mathcal{M}_n)
      \;=\; \dim_\kappa M_n - 0 \;=\; \dim_\kappa M_n .
  \]
  Hence \(\mathrm{IsRatHilb}(n \mapsto \dim_\kappa M_n,\, r)\) (\ref{mk-def:ratHilb}):
  there are \(p \in \mathbb{Q}[X]\), \(d := r \in \mathbb{N}\) and \(N \in \mathbb{N}\)
  with
  \[
    \dim_\kappa M_n \;=\; \bigl[X^n\bigr]\bigl(p \cdot (1 - X)^{-d}\bigr)
    \qquad (n > N),
  \]
  which is the claim, with pole order \(d = r\). The internal mechanics of
  \ref{mk-lem:graded_subquotient_isRatHilb} -- the base case via
  \ref{mk-lem:ratHilb_ofEventuallyZero}, and the inductive step assembling the
  kernel/cokernel closure (\ref{mk-lem:graded_subquotient_ker_coker}), the degreewise
  difference identity (\ref{mk-lem:graded_subquotient_degreewise_diff}, built on
  \(D_5\), \ref{mk-lem:graded_degreewise_finrank_diff}) and the finiteness transfer
  (\ref{mk-lem:graded_subquotient_finite_transfer}) into
  \ref{mk-lem:ratHilb_ofDiffEq} and \ref{mk-lem:ratHilb_bump} -- are the Stacks~00K1
  induction realised entirely with ambient homogeneous submodules of \(M\).
\end{proof}

\subsection{The rationality toolkit \texorpdfstring{\(\mathrm{IsRatHilb}\)}{IsRatHilb}}
\label{mk-subsec:isRatHilb}

The proof of \ref{mk-lem:gradedHilbertSerre_rational} runs the Stacks~00K1 induction
on the number of degree-one generators. Each inductive step manipulates the
\emph{shape} of the Hilbert series -- a numerator polynomial over a power of
\((1-X)\) -- rather than the module itself. We isolate that bookkeeping in a
single predicate \(\mathrm{IsRatHilb}\) on functions \(f : \mathbb{N} \to
\mathbb{Q}\), together with the closure lemmas the induction consumes. These are
  The auxiliary lemmas below realise the standard Stacks~00K1 argument at
  the level of power series; the converse extraction is
  (\ref{mk-lem:hilbertPoly_exists_mathlib}).

\begin{lemma}[Coefficients of \((1-X)^{-1}\) times a series are partial sums]\label{mk-lem:coeff_invOneSubPow_one_mul}
  \dcref{custom}

  \uses{mk-lem:invOneSubPow_mathlib}
  Let \(F \in \mathbb{Q}\llbracket X\rrbracket\) be a power series. Then for every
  \(n \in \mathbb{N}\) the \(n\)-th coefficient of
  \(\operatorname{invOneSubPow}(\mathbb{Q}, 1)\cdot F = (1-X)^{-1}\cdot F\)
  (\ref{mk-lem:invOneSubPow_mathlib}) is the partial sum of the coefficients of
  \(F\):
  \[
    \bigl[X^n\bigr]\bigl((1-X)^{-1}\cdot F\bigr)
    \;=\; \sum_{k=0}^{n} \bigl[X^k\bigr] F .
  \]
  Equivalently, multiplication by \((1-X)^{-1} = \sum_{m\ge 0} X^m\) is the
  summation operator on coefficient sequences, by the Cauchy product.
\end{lemma}

\begin{lemma}[Antidifference of a rational Hilbert series]\label{mk-lem:rationalHilbert_antidiff}
  \dcref{custom}

  \uses{mk-lem:coeff_invOneSubPow_one_mul, mk-lem:invOneSubPow_mathlib}
  Let \(H, \delta : \mathbb{N} \to \mathbb{Q}\), let \(q \in \mathbb{Q}[X]\) and
  \(e, N \in \mathbb{N}\). Suppose that for all \(n > N\) the sequence \(\delta\)
  is the coefficient sequence of the rational series \(q\cdot(1-X)^{-e}\),
  \[
    \delta(n) \;=\; \bigl[X^n\bigr]\bigl(q\cdot(1-X)^{-e}\bigr),
  \]
  and that \(\delta\) is the first difference of \(H\) beyond \(N\), i.e.
  \(H(n+1) - H(n) = \delta(n+1)\) for all \(n > N\). Then there is a numerator
  polynomial \(p \in \mathbb{Q}[X]\) such that for all \(n > N\)
  \[
    H(n) \;=\; \bigl[X^n\bigr]\bigl(p\cdot(1-X)^{-(e+1)}\bigr).
  \]
  That is, summing a coefficient sequence (anti-differencing) raises the pole
  order by one: the partial-sum series of \(q\cdot(1-X)^{-e}\) equals
  \(p\cdot(1-X)^{-(e+1)}\) up to a polynomial correction in low degrees. This is
  the power-series heart of the Stacks~00K1 inductive step.
\end{lemma}

\begin{proof}
  \uses{mk-lem:coeff_invOneSubPow_one_mul, mk-lem:invOneSubPow_mathlib}
  Write \(F = q\cdot(1-X)^{-e}\). Multiplying by
  \((1-X)^{-1}\) and using the factorisation
  \((1-X)^{-(e+1)} = (1-X)^{-1}\cdot(1-X)^{-e}\)
  (\ref{mk-lem:invOneSubPow_mathlib}), \ref{mk-lem:coeff_invOneSubPow_one_mul} gives
  \(\bigl[X^m\bigr]\bigl(q\cdot(1-X)^{-(e+1)}\bigr) = \sum_{k=0}^{m}[X^k]F\), the
  partial sums of \(F\). On the other hand, telescoping the first-difference
  hypothesis from \(N+1\) expresses \(H(N+1+j)\) as \(H(N+1)\) plus a window sum
  of the coefficients of \(F\). Splitting the partial sum
  \(\sum_{k=0}^{n}[X^k]F\) at \(N+2\) and absorbing the resulting constant
  discrepancy into a numerator \(C\,(1-X)^{e}\) (a constant function is the
  order-\((e+1)\) coefficient sequence of \(C\,(1-X)^{e}\cdot(1-X)^{-(e+1)} =
  C\,(1-X)^{-1}\)), the explicit numerator
  \(p = C\,(1-X)^{e} + q\) realises \(H(n)\) as
  \([X^n]\bigl(p\cdot(1-X)^{-(e+1)}\bigr)\) for all \(n > N\).
\end{proof}

\begin{definition}[Rational Hilbert function]\label{mk-def:ratHilb}
  \dcref{custom}

  \uses{mk-lem:invOneSubPow_mathlib}
  \textit{Source: Stacks 00K1.}
  For a function \(f : \mathbb{N} \to \mathbb{Q}\) and \(d \in \mathbb{N}\), say
  that \(f\) is a \emph{rational Hilbert function of order \(d\)}, written
  \(\mathrm{IsRatHilb}(f, d)\), when there exist a numerator polynomial
  \(p \in \mathbb{Q}[X]\) and an \(N \in \mathbb{N}\) such that for all \(n > N\)
  \[
    f(n) \;=\; \bigl[X^n\bigr]\bigl(p\cdot(1-X)^{-d}\bigr)
    \;=\; \bigl[X^n\bigr]\bigl(\operatorname{toPowerSeries}(p)\cdot
      \operatorname{invOneSubPow}(\mathbb{Q}, d)\bigr)
  \]
  (\ref{mk-lem:invOneSubPow_mathlib}). This is the ``eventually equal to the
  coefficient sequence of a rational series \(p\cdot(1-X)^{-d}\)'' predicate that
  the graded Hilbert--Serre induction tracks; its conclusion shape is exactly the
  one consumed by \ref{mk-lem:hilbertPoly_exists_mathlib}.
\end{definition}

\begin{lemma}[Base case: eventually-zero functions are rational of order \(0\)]\label{mk-lem:ratHilb_ofEventuallyZero}
  \dcref{custom}

  \uses{mk-def:ratHilb}
  If \(f : \mathbb{N} \to \mathbb{Q}\) is eventually zero -- there is
  \(N\) with \(f(n) = 0\) for all \(n > N\) -- then \(\mathrm{IsRatHilb}(f, 0)\)
  (\ref{mk-def:ratHilb}). Indeed \(f\) is then the coefficient sequence of the
  numerator polynomial \(0\) over \((1-X)^{0} = 1\). This is the base case of the
  Stacks~00K1 induction (no degree-one generators).
\end{lemma}

\begin{lemma}[Raising the pole order]\label{mk-lem:ratHilb_bump}
  \dcref{custom}

  \uses{mk-def:ratHilb, mk-lem:invOneSubPow_mathlib}
  If \(\mathrm{IsRatHilb}(f, d)\) (\ref{mk-def:ratHilb}) then
  \(\mathrm{IsRatHilb}(f, d+1)\): multiplying the numerator by \((1-X)\) cancels
  one factor of \((1-X)^{-1}\), so the same coefficient sequence is realised at
  pole order \(d+1\). This lets two rational Hilbert functions of different
  orders be brought to a common denominator, using
  \((1-X)\cdot(1-X)^{-(d+1)} = (1-X)^{-d}\) (\ref{mk-lem:invOneSubPow_mathlib}).
\end{lemma}

\begin{lemma}[Closure under difference]\label{mk-lem:ratHilb_sub}
  \dcref{custom}

  \uses{mk-def:ratHilb}
  If \(\mathrm{IsRatHilb}(f, d)\) and \(\mathrm{IsRatHilb}(g, d)\)
  (\ref{mk-def:ratHilb}) at the same order \(d\), then
  \(\mathrm{IsRatHilb}(n \mapsto f(n) - g(n),\, d)\): the numerator of the
  difference is the difference of the numerators (over the common denominator
  \((1-X)^{-d}\)).
\end{lemma}

\begin{lemma}[Closure under right shift]\label{mk-lem:ratHilb_shiftRight}
  \dcref{custom}

  \uses{mk-def:ratHilb}
  If \(\mathrm{IsRatHilb}(f, d)\) (\ref{mk-def:ratHilb}) then
  \(\mathrm{IsRatHilb}(n \mapsto f(n-1),\, d)\): the right shift on coefficient
  sequences is multiplication of the series by \(X\), so the numerator
  \(p\) is replaced by \(X\cdot p\) at the same pole order.
\end{lemma}

\begin{lemma}[Antidifference step for the predicate]\label{mk-lem:ratHilb_antidiff}
  \dcref{custom}

  \uses{mk-def:ratHilb, mk-lem:rationalHilbert_antidiff}
  Let \(H, g : \mathbb{N} \to \mathbb{Q}\) and \(e, N \in \mathbb{N}\). If
  \(\mathrm{IsRatHilb}(g, e)\) (\ref{mk-def:ratHilb}) and \(g\) is the first
  difference of \(H\) beyond \(N\) (\(H(n+1) - H(n) = g(n+1)\) for all
  \(n > N\)), then \(\mathrm{IsRatHilb}(H, e+1)\). This is the predicate-level
  packaging of \ref{mk-lem:rationalHilbert_antidiff}: a finite difference raises
  the pole order by one, the inductive heart of graded Hilbert--Serre.
\end{lemma}

\begin{proof}
  \uses{mk-def:ratHilb, mk-lem:rationalHilbert_antidiff}
  Unfolding \(\mathrm{IsRatHilb}(g, e)\) supplies a numerator \(q\) and a
  threshold \(N_g\) with \(g(n) = [X^n](q\cdot(1-X)^{-e})\) for \(n > N_g\).
  Applying \ref{mk-lem:rationalHilbert_antidiff} with this \(q\) and threshold
  \(\max(N, N_g)\) -- whose two hypotheses are exactly the rational form of \(g\)
  and the first-difference identity for \(H\) -- yields a numerator \(p\) with
  \(H(n) = [X^n](p\cdot(1-X)^{-(e+1)})\) for \(n > \max(N, N_g)\), i.e.
  \(\mathrm{IsRatHilb}(H, e+1)\).
\end{proof}

\begin{lemma}[Inductive-step engine from a degreewise short exact sequence]\label{mk-lem:ratHilb_ofDiffEq}
  \dcref{custom}

  \uses{mk-def:ratHilb, mk-lem:ratHilb_antidiff, mk-lem:ratHilb_sub, mk-lem:ratHilb_shiftRight}
  Let \(h_M, h_C, h_K : \mathbb{N} \to \mathbb{Q}\) and \(d, N \in \mathbb{N}\).
  Suppose \(\mathrm{IsRatHilb}(h_C, d)\) and \(\mathrm{IsRatHilb}(h_K, d)\)
  (\ref{mk-def:ratHilb}), and that for all \(n > N\)
  \[
    h_M(n+1) - h_M(n) \;=\; h_C(n+1) - h_K(n).
  \]
  Then \(\mathrm{IsRatHilb}(h_M, d+1)\). This packages the entire power-series
  side of the Stacks~00K1 inductive step: \(h_M\) is the Hilbert function of
  \(M\), and \(h_C, h_K\) are the Hilbert functions of the cokernel
  \(C = M/xM\) and kernel \(K = \ker(x : M \to M(1))\) of multiplication by a
  degree-one element \(x\); the displayed identity is the alternating-sum
  consequence of the degreewise short exact sequence
  \(0 \to K_n \to M_n \xrightarrow{x} M_{n+1} \to C_{n+1} \to 0\). Given the
  inductive hypothesis that \(h_C, h_K\) are rational of order \(d\), the
  conclusion is rationality of \(h_M\) at order \(d+1\).
\end{lemma}

\begin{proof}
  \uses{mk-lem:ratHilb_antidiff, mk-lem:ratHilb_sub, mk-lem:ratHilb_shiftRight}
  Set \(g(n) = h_C(n) - h_K(n-1)\). By closure under right shift
  (\ref{mk-lem:ratHilb_shiftRight}) \(n \mapsto h_K(n-1)\) is rational of order
  \(d\), and by closure under difference (\ref{mk-lem:ratHilb_sub}) so is \(g\). The
  hypothesis rewrites as \(h_M(n+1) - h_M(n) = g(n+1)\) for \(n > N\) (since
  \(g(n+1) = h_C(n+1) - h_K(n)\)). Applying the antidifference step
  (\ref{mk-lem:ratHilb_antidiff}) to \(H = h_M\) and this \(g\) gives
  \(\mathrm{IsRatHilb}(h_M, d+1)\).
\end{proof}

\begin{lemma}[Uniqueness of the Hilbert polynomial]\label{mk-lem:hilbertPolynomial_unique}
  \dcref{custom}

  \uses{mk-def:fiber_hilbert_function}
  Let \(h_{\mathcal{F},s}\) be the graded Hilbert function of
  \ref{mk-def:fiber_hilbert_function}. If \(\Phi, \Psi \in \mathbb{Q}[\lambda]\)
  both agree with \(h_{\mathcal{F},s}\) at all sufficiently large integers --
  there are \(N_1, N_2 \in \mathbb{N}\) with
  \(\Phi(m) = h_{\mathcal{F},s}(m)\) for all \(m > N_1\) and
  \(\Psi(m) = h_{\mathcal{F},s}(m)\) for all \(m > N_2\) -- then
  \(\Phi = \Psi\). In particular a polynomial agreeing with
  \(h_{\mathcal{F},s}\) for all \(m \gg 0\) is unique when it exists.
\end{lemma}

\begin{proof}
  The difference \(\Phi - \Psi \in \mathbb{Q}[\lambda]\) vanishes at every
  integer \(m > \max(N_1, N_2)\), hence at infinitely many points of
  \(\mathbb{Q}\). A nonzero polynomial over a field has at most
  \(\deg(\Phi - \Psi)\) roots, so \(\Phi - \Psi = 0\), that is
  \(\Phi = \Psi\).
\end{proof}

\begin{lemma}[Hilbert polynomial from a rational Hilbert function]\label{mk-lem:hilbertPoly_of_isRatHilb}
  \dcref{custom}

  \uses{mk-def:fiber_hilbert_function, mk-def:ratHilb,
    mk-lem:hilbertPoly_exists_mathlib, mk-lem:hilbertPolynomial_unique}
  Let \(h_{\mathcal{F},s}\) be the graded Hilbert function of
  \ref{mk-def:fiber_hilbert_function}, and suppose \(h_{\mathcal{F},s}\) is a
  rational Hilbert function of some order \(d\) (\ref{mk-def:ratHilb}): there
  are \(p \in \mathbb{Q}[X]\) and \(N \in \mathbb{N}\) with
  \(h_{\mathcal{F},s}(m) = [X^m]\bigl(p \cdot (1-X)^{-d}\bigr)\) for all
  \(m > N\). Then there exists a unique polynomial
  \(\Phi \in \mathbb{Q}[\lambda]\) such that
  \(\Phi(m) = h_{\mathcal{F},s}(m)\) for all \(m \gg 0\).
\end{lemma}

\begin{proof}
  \uses{mk-def:ratHilb, mk-lem:hilbertPoly_exists_mathlib,
    mk-lem:hilbertPolynomial_unique}
  Unfolding \ref{mk-def:ratHilb} supplies \(p \in \mathbb{Q}[X]\),
  \(d \in \mathbb{N}\) and \(N\) with
  \(h_{\mathcal{F},s}(m) = [X^m](p\cdot(1-X)^{-d})\) for all \(m > N\). By
  \ref{mk-lem:hilbertPoly_exists_mathlib} over the characteristic-zero field
  \(\mathbb{Q}\), applied to the pair \((p, d)\), there are a polynomial
  \(h \in \mathbb{Q}[X]\) and an \(N'\) with
  \(h(n) = [X^n](p\cdot(1-X)^{-d})\) for all \(n > N'\). Hence
  \(\Phi := h\) satisfies \(\Phi(m) = h_{\mathcal{F},s}(m)\) for all
  \(m > \max(N, N')\), proving existence. Uniqueness is
  \ref{mk-lem:hilbertPolynomial_unique}: any two polynomials agreeing with
  \(h_{\mathcal{F},s}\) at all sufficiently large integers are equal.
\end{proof}

\begin{corollary}[Defining property of the Hilbert polynomial]\label{mk-cor:hilbertPolynomial_eval_eventually}
  \dcref{custom}

  \uses{mk-def:hilbert_polynomial, mk-def:fiber_hilbert_function,
    mk-lem:hilbertPolynomial_unique}
  Let \(h_{\mathcal{F},s}\) be the graded Hilbert function of
  \ref{mk-def:fiber_hilbert_function}. If \emph{some} polynomial in
  \(\mathbb{Q}[\lambda]\) agrees with \(h_{\mathcal{F},s}\) at all
  sufficiently large integers, then the Hilbert polynomial
  \(\Phi_{\mathcal{F},s}\) of \ref{mk-def:hilbert_polynomial} does:
  \(\Phi_{\mathcal{F},s}(m) = h_{\mathcal{F},s}(m)\) for all \(m \gg 0\).
\end{corollary}

\begin{proof}
  \uses{mk-def:hilbert_polynomial, mk-lem:hilbertPolynomial_unique}
  Let \(\Psi \in \mathbb{Q}[\lambda]\) agree with \(h_{\mathcal{F},s}\) at all
  \(m > N\). By \ref{mk-def:hilbert_polynomial} the Hilbert polynomial
  \(\Phi_{\mathcal{F},s}\) is a polynomial agreeing with
  \(h_{\mathcal{F},s}\) for all \(m \gg 0\) whenever one exists, and by
  \ref{mk-lem:hilbertPolynomial_unique} any such polynomial is unique, so
  \(\Phi_{\mathcal{F},s} = \Psi\); hence
  \(\Phi_{\mathcal{F},s}(m) = h_{\mathcal{F},s}(m)\) for all \(m > N\).
\end{proof}

\begin{corollary}[The Hilbert polynomial under a rational Hilbert function]\label{mk-cor:hilbertPolynomial_eval_of_isRatHilb}
  \dcref{custom}

  \uses{mk-def:hilbert_polynomial, mk-def:ratHilb, mk-lem:hilbertPoly_of_isRatHilb,
    mk-cor:hilbertPolynomial_eval_eventually}
  If the graded Hilbert function \(h_{\mathcal{F},s}\) is a rational Hilbert
  function of some order \(d\) (\ref{mk-def:ratHilb}), then the Hilbert
  polynomial \(\Phi_{\mathcal{F},s}\) of \ref{mk-def:hilbert_polynomial}
  satisfies \(\Phi_{\mathcal{F},s}(m) = h_{\mathcal{F},s}(m)\) for all
  \(m \gg 0\).
\end{corollary}

\begin{proof}
  \uses{mk-lem:hilbertPoly_of_isRatHilb, mk-cor:hilbertPolynomial_eval_eventually}
  By \ref{mk-lem:hilbertPoly_of_isRatHilb} some polynomial agrees with
  \(h_{\mathcal{F},s}\) at all sufficiently large integers; conclude by
  \ref{mk-cor:hilbertPolynomial_eval_eventually}.
\end{proof}

\begin{corollary}[The Hilbert polynomial computes as the extracted numerator polynomial]\label{mk-cor:hilbertPolynomial_eq_hilbertPoly}
  \dcref{custom}

  \uses{mk-def:hilbert_polynomial, mk-def:ratHilb, mk-lem:hilbertPoly_exists_mathlib,
    mk-lem:hilbertPolynomial_unique}
  Suppose there are \(p \in \mathbb{Q}[X]\) and \(d, N \in \mathbb{N}\) with
  \(h_{\mathcal{F},s}(m) = [X^m]\bigl(p \cdot (1-X)^{-d}\bigr)\) for all
  \(m > N\) (\ref{mk-def:ratHilb}). Then the Hilbert polynomial of
  \ref{mk-def:hilbert_polynomial} is the polynomial extracted from the pair
  \((p, d)\) by \ref{mk-lem:hilbertPoly_exists_mathlib}:
  \(\Phi_{\mathcal{F},s} = \mathrm{hilbertPoly}\,p\,d\).
\end{corollary}

\begin{proof}
  \uses{mk-lem:hilbertPoly_exists_mathlib, mk-lem:hilbertPolynomial_unique}
  By \ref{mk-lem:hilbertPoly_exists_mathlib} there is an \(N'\) with
  \((\mathrm{hilbertPoly}\,p\,d)(n) = [X^n](p\cdot(1-X)^{-d})\) for all
  \(n > N'\), so \(\mathrm{hilbertPoly}\,p\,d\) agrees with
  \(h_{\mathcal{F},s}\) at all \(m > \max(N, N')\). Both
  \(\Phi_{\mathcal{F},s}\) and \(\mathrm{hilbertPoly}\,p\,d\) therefore agree
  with \(h_{\mathcal{F},s}\) at all sufficiently large integers, and
  \ref{mk-lem:hilbertPolynomial_unique} gives
  \(\Phi_{\mathcal{F},s} = \mathrm{hilbertPoly}\,p\,d\).
\end{proof}

\begin{theorem}[Hilbert polynomial from the section module]
  \label{mk-thm:hilbertPoly_of_sectionModule}

  \uses{mk-def:hilbert_polynomial, mk-def:fiber_hilbert_function,
    mk-lem:sectionGradedModule_fg, mk-lem:gradedHilbertSerre_rational,
    mk-lem:hilbertPoly_of_isRatHilb}
  \dcref{custom}
  \textit{Source: \cite{nitsure-hilbert-quot}, \S 1 (FGA Explained Ch.~5).}
  Let \(X_s\) be proper over the field \(\kappa(s)\), let \(L_s\) be an ample line
  bundle on \(X_s\), and let \(\mathcal{F}_s\) be a coherent sheaf on \(X_s\).
  Then there is a unique polynomial \(\Phi_{\mathcal{F},s} \in \mathbb{Q}[\lambda]\)
  such that for all \(m \gg 0\)
  \[
    \Phi_{\mathcal{F},s}(m)
    \;=\; \dim_{\kappa(s)} \Gamma\bigl(X_s,\, \mathcal{F}_s \otimes L_s^{\otimes m}\bigr).
  \]
  This \(\Phi_{\mathcal{F},s}\) is the Hilbert polynomial of
  \ref{mk-def:hilbert_polynomial}.

\end{theorem}

\begin{proof}
  \uses{mk-lem:sectionGradedModule_fg, mk-lem:gradedHilbertSerre_rational,
    mk-lem:hilbertPoly_of_isRatHilb}
  By \ref{mk-lem:sectionGradedModule_fg} the section graded module
  \(M(X_s, \mathcal{F}_s, L_s)\) is a finitely generated graded module over the
  Noetherian graded ring \(R(X_s, L_s)\), with finite-dimensional components and
  degree-zero part the field \(\kappa(s)\). Applying the abstract graded
  Hilbert--Serre rationality lemma (\ref{mk-lem:gradedHilbertSerre_rational}) to the
  graded \(\kappa(s)\)-algebra \(R = R(X_s, L_s)\) and the finitely generated
  graded module \(M = M(X_s, \mathcal{F}_s, L_s)\) exhibits the graded Hilbert
  function
  \(m \mapsto \dim_{\kappa(s)} \Gamma(X_s, \mathcal{F}_s \otimes L_s^{\otimes m})\)
  (\ref{mk-def:fiber_hilbert_function}) as a rational Hilbert function of some
  order \(d\) (\ref{mk-def:ratHilb}): there are \(p \in \mathbb{Q}[X]\),
  \(d \in \mathbb{N}\) and an \(N\) with
  \(\dim_{\kappa(s)} \Gamma(X_s, \mathcal{F}_s \otimes L_s^{\otimes m})
  = [X^m]\,(p \cdot (1 - X)^{-d})\) for \(m > N\). The extraction lemma
  \ref{mk-lem:hilbertPoly_of_isRatHilb} then yields a unique polynomial
  \(\Phi_{\mathcal{F},s} \in \mathbb{Q}[\lambda]\) with
  \(\Phi_{\mathcal{F},s}(m) = \dim_{\kappa(s)}
  \Gamma(X_s, \mathcal{F}_s \otimes L_s^{\otimes m})\) for all \(m \gg 0\),
  as claimed.
\end{proof}

\subsection{The ambient subquotient induction for the Stacks~00K1 step}
\label{mk-subsec:gradedModuleApi}

The power-series side of \ref{mk-lem:gradedHilbertSerre_rational} is fully isolated
in the \(\mathrm{IsRatHilb}\) toolkit (\ref{mk-subsec:isRatHilb}). What remains is
the \emph{graded-module} side of the Stacks~00K1 induction. Rather than build the
kernel and cokernel of multiplication by a degree-one element as graded modules
over a smaller quotient ring \(R/(x)\) -- which would force a graded structure on a
derived carrier -- we range the induction over \emph{pairs of homogeneous
submodules of a single fixed ambient graded module} \(M\). A pair \((N, N')\) with
\(N' \le N\) homogeneous submodules of \(M\) stands for the subquotient \(N/N'\); its
ambient Hilbert function is the difference \(n \mapsto \dim_\kappa(N \cap M_n) -
\dim_\kappa(N' \cap M_n)\). The class of such pairs, equipped with finitely many
pairwise-commuting degree-one endomorphisms preserving \(N\) and \(N'\), is
\emph{closed} under the kernel and cokernel of one of those endomorphisms, both
again presented as ambient subquotient pairs. Consequently no graded quotient ring,
graded quotient module, or regrading of a subtype/quotient carrier is ever formed:
every object that appears is an ambient family \(n \mapsto N_{\mathrm{aux}} \cap M_n\)
of submodules of the fixed \(M\).

This subsection records that machinery as a thin layer over existing Mathlib
scaffold: Mathlib already provides the notion of a homogeneous submodule
(\ref{mk-lem:submodule_isHomogeneous_mathlib}), the homogeneity of a span of
homogeneous elements (\ref{mk-lem:ideal_homogeneous_span_mathlib}), the degree-shift
of multiplication by a homogeneous element
(\ref{mk-lem:isHomogeneousElem_graded_smul_mathlib}), the rank--nullity identity
(\ref{mk-lem:finrank_range_add_finrank_ker_mathlib}), and the finite-generation
transfer along a surjection (\ref{mk-lem:fg_restrictScalars_of_surjective_mathlib}).
What is supplied here is the internal direct-sum decomposition of a homogeneous
submodule as an ambient family (the split \(G_1\)), the subquotient Hilbert data,
the kernel/cokernel closure, the degreewise difference identity, the ambient
finiteness transfer, and the induction itself.

\emph{Setting for this subsection.} Let \(R = \bigoplus_{n \ge 0} R_n\) be a graded
ring whose degree-zero part \(R_0 = \kappa\) is a field and which is generated as a
\(\kappa\)-algebra in degree one, and let \(M = \bigoplus_{n \ge 0} M_n\) be a graded
\(\kappa\)-vector space (the fixed ambient module) with finite-dimensional
components \(M_n\). All submodules \(N, N', \dots\) below are homogeneous submodules
of \(M\), and the relevant endomorphisms are \(\kappa\)-linear maps \(M \to M\)
raising degree by one -- typically multiplication by a homogeneous element
\(x \in R_1\).

\subsubsection*{Mathlib dependency anchors}

\begin{lemma}[Homogeneous submodule]
  \label{mk-lem:submodule_isHomogeneous_mathlib}
  \dcref{custom}

  Let \(R = \bigoplus_n R_n\) be a graded ring and \(M = \bigoplus_n M_n\) an
  internally graded \(R\)-module. A submodule \(p \subseteq M\) is
  \emph{homogeneous} when it is closed under taking homogeneous components:
  \(m \in p\) implies each graded piece \(\mathrm{proj}_n(m) \in p\). Equivalently
  \(p = \bigoplus_n (p \cap M_n)\) as a set. Mathlib bundles this into a homogeneous
  submodule type providing order and extensionality structure; the ambient internal
  direct-sum decomposition \(p = \bigoplus_n (M_n \cap p)\) is supplied by the split
  \(G_1\) (\ref{mk-lem:graded_homogeneousSubmodule_iSupIndep},
  \ref{mk-lem:graded_homogeneousSubmodule_iSup_eq}).
\end{lemma}

\begin{lemma}[A span of homogeneous elements is homogeneous]
  \label{mk-lem:ideal_homogeneous_span_mathlib}
  \dcref{custom}

  Let \(R = \bigoplus_n R_n\) be a graded ring and \(s \subseteq R\) a set each of
  whose elements is homogeneous. Then the ideal \(\mathrm{span}\,s\) is a
  homogeneous ideal. In particular, for \(x \in R_1\) the principal ideal \((x)\)
  is homogeneous, hence \(xM\) and the auxiliary submodule \(N' + x \cdot N\) are
  homogeneous; this homogeneity feeds the kernel/cokernel closure
  (\ref{mk-lem:graded_subquotient_ker_coker}).
\end{lemma}

\begin{lemma}[Rank--nullity]
  \label{mk-lem:finrank_range_add_finrank_ker_mathlib}
  \dcref{custom}

  Let \(\kappa\) be a field, \(V\) a finite-dimensional \(\kappa\)-vector space and
  \(\varphi : V \to W\) a \(\kappa\)-linear map. Then
  \[
    \dim_\kappa(\operatorname{range}\varphi) + \dim_\kappa(\ker\varphi)
      \;=\; \dim_\kappa V .
  \]
  This is the rank--nullity identity underlying the degreewise difference identity
  \(D_5\) (\ref{mk-lem:graded_degreewise_finrank_diff}).
\end{lemma}

\begin{lemma}[Finite generation transfers along a surjection]
  \label{mk-lem:fg_restrictScalars_of_surjective_mathlib}
  \dcref{custom}

  Let \(\varphi : R \to A\) be a surjective ring homomorphism and \(M\) an
  \(A\)-module that is finitely generated when viewed as an \(R\)-module via
  \(\varphi\). Then \(M\) is finitely generated as an \(A\)-module. Applied to the
  surjection \(\kappa[t_0, \dots, t_{r-1}] \twoheadrightarrow \kappa[t_0, \dots,
  t_{r-2}]\) killing the endomorphism \(x = t_{r-1}\) that annihilates the two
  subquotients, this transfers their finite generation down one generator
  (\ref{mk-lem:graded_subquotient_finite_transfer}).
\end{lemma}

\begin{lemma}[Multiplication by a homogeneous element shifts degree]
  \label{mk-lem:isHomogeneousElem_graded_smul_mathlib}
  \dcref{custom}

  Let \(\mathcal{A}\) be a graded family acting on a graded family \(\mathcal{M}\)
  via a graded scalar multiplication. If \(a \in R\) is homogeneous of degree
  \(i\) -- in particular \(x \in R_1\) -- then for \(m \in M_n\) one has
  \(a \cdot m \in M_{i+n}\). This is the degree-shifting property that makes \(xM\)
  a homogeneous submodule (with degree-\((n+1)\) part \(x \cdot M_n\)) and that
  underlies the ambient identity \(x \cdot N \cap M_{n+1} = x \cdot (N \cap M_n)\)
  used in the subquotient closure (\ref{mk-lem:graded_subquotient_ker_coker}) and
  difference identity (\ref{mk-lem:graded_subquotient_degreewise_diff}).
\end{lemma}

\begin{lemma}[Commutative subalgebra generated by commuting elements]
  \label{mk-lem:adjoinCommRingOfComm_mathlib}
  \dcref{custom}

  Let \(R\) be a commutative semiring and \(A\) an \(R\)-algebra, possibly
  noncommutative. If \(s \subseteq A\) is a set of pairwise-commuting elements
  (\(a b = b a\) for all \(a, b \in s\)), then the \(R\)-subalgebra
  \(\mathrm{adjoin}_R(s)\) generated by \(s\) has commutative multiplication; the
  commutative-ring structure on it follows. Applied to the \(r\) pairwise-commuting degree-raising endomorphisms
  \(x_0, \dots, x_{r-1} \in \operatorname{End}_\kappa(M)\), it yields a
  \emph{commutative} \(\kappa\)-subalgebra \(A = \mathrm{adjoin}_\kappa\{x_0, \dots,
  x_{r-1}\} \subseteq \operatorname{End}_\kappa(M)\); the polynomial action of
  \(\kappa[t_0, \dots, t_{r-1}]\) on \(M\) factors through this commutative \(A\)
  (\ref{mk-lem:graded_subquotient_finite_transfer}).
\end{lemma}

\begin{lemma}[Finiteness transfers along a surjective linear map]
  \label{mk-lem:module_finite_of_surjective_mathlib}
  \dcref{custom}

  Let \(A\) be a ring, \(M, N\) two \(A\)-modules and \(\varphi : M \to N\) a
  surjective \(A\)-linear map. If \(M\) is a finite \(A\)-module, then so is \(N\).
  This is the finiteness-transfer step at the heart of the ambient finiteness
  transfer (\ref{mk-lem:graded_subquotient_finite_transfer}) and of the
  numerator/denominator finiteness lemmas
  (\ref{mk-lem:graded_polyQuot_finite_of_le_numerator},
  \ref{mk-lem:graded_polyQuot_finite_of_le_denominator}).
\end{lemma}

\begin{lemma}[Finiteness descends along an injective linear map into a Noetherian module]
  \label{mk-lem:module_finite_of_injective_mathlib}
  \dcref{custom}

  Let \(A\) be a ring, \(M, N\) two \(A\)-modules with \(N\) Noetherian (in
  particular, finite over a Noetherian ring \(A\)), and \(\varphi : M \to N\) an
  injective \(A\)-linear map. Then \(M\) is a finite \(A\)-module. This is the
  submodule-finiteness step in \ref{mk-lem:graded_polyQuot_finite_of_le_numerator}.
\end{lemma}

\begin{lemma}[The lift of a linear map through a quotient]
  \label{mk-lem:submodule_liftQ_mathlib}
  \dcref{custom}

  Let \(A\) be a ring, \(M, N\) two \(A\)-modules, \(P \subseteq M\) a submodule and
  \(\varphi : M \to N\) an \(A\)-linear map with \(P \subseteq \ker\varphi\). Then
  \(\varphi\) factors uniquely through a linear map \(M/P \to N\) on the quotient.
  This realises the descent of a (\(\sigma\)-semilinear) map to the subquotient
  quotient in the ambient finiteness transfer
  (\ref{mk-lem:graded_subquotient_finite_transfer}) and the denominator lemma
  (\ref{mk-lem:graded_polyQuot_finite_of_le_denominator}).
\end{lemma}

\begin{lemma}[Only finitely many members of an independent family are nonzero]
  \label{mk-lem:finite_ne_bot_of_iSupIndep_mathlib}
  \dcref{custom}

  Let \(\kappa\) be a field (more generally a division ring) and \(V\) a
  finite-dimensional \(\kappa\)-vector space. If \((p_i)_{i \in \iota}\) is an
  \(\mathrm{iSupIndep}\) family of subspaces of \(V\), then the set
  \(\{\,i : p_i \ne \bot\,\}\) is finite. This is the finiteness input that turns
  degreewise independence into eventual vanishing in the base case
  (\ref{mk-lem:graded_subquotient_base_eventuallyZero}).
\end{lemma}

\begin{lemma}[The polynomial ring in finitely many variables is Noetherian]
  \label{mk-lem:mvPolynomial_isNoetherianRing_fin_mathlib}
  \dcref{custom}

  Let \(R\) be a Noetherian commutative ring and \(r \in \mathbb{N}\). Then the
  multivariate polynomial ring \(\operatorname{MvPolynomial}(\{0, \dots, r-1\}, R)\)
  is a Noetherian ring (Hilbert's basis theorem in finitely many variables). With
  \(R = \kappa\) a field this makes \(\kappa[t_0, \dots, t_{r-1}]\) Noetherian, the
  base ring over which the subquotient finiteness is measured
  (\ref{mk-lem:graded_polyQuot_finite_of_le_numerator}).
\end{lemma}

\begin{lemma}[A finite module over a Noetherian ring is Noetherian]
  \label{mk-lem:isNoetherian_of_isNoetherianRing_of_finite_mathlib}
  \dcref{custom}

  Let \(A\) be a Noetherian ring and \(M\) a finite \(A\)-module. Then \(M\) is a
  Noetherian \(A\)-module. Combined with
  \ref{mk-lem:mvPolynomial_isNoetherianRing_fin_mathlib} this makes a finite
  \(\kappa[t_0, \dots, t_{r-1}]\)-module Noetherian, so its submodules are again
  finite (\ref{mk-lem:graded_polyQuot_finite_of_le_numerator}).
\end{lemma}

\subsubsection*{The internal grading of a homogeneous submodule (split $G_1$)}

\begin{lemma}[$G_1$ (independence) -- the family $n \mapsto M_n \cap p$ is independent]\label{mk-lem:graded_homogeneousSubmodule_iSupIndep}
  \dcref{custom}

  \uses{mk-lem:submodule_isHomogeneous_mathlib}
  Let \(M = \bigoplus_n M_n\) be an internally graded \(\kappa\)-module and let
  \(p \subseteq M\) be a homogeneous submodule
  (\ref{mk-lem:submodule_isHomogeneous_mathlib}). Then the family
  \(n \mapsto M_n \cap p\) of \(\kappa\)-subspaces of \(M\) is independent: for each
  \(n\),
  \[
    (M_n \cap p) \;\cap\; \Bigl(\bigsqcup_{m \ne n} (M_m \cap p)\Bigr) \;=\; 0 .
  \]
  The whole statement lives in the ambient module \(M\): the subspaces
  \(M_n \cap p\) are submodules of \(M\), never of a derived carrier \(\smash{\hat p}\).
\end{lemma}

\begin{proof}
  \uses{mk-lem:submodule_isHomogeneous_mathlib}
  The components \(M_n\) of the internal decomposition of \(M\) are independent. The
  subfamily \(n \mapsto M_n \cap p\) consists of submodules with \(M_n \cap p \le M_n\),
  so independence is inherited: a relation among the \(M_n \cap p\) is in particular a
  relation among the \(M_n\), forcing each term to vanish.
\end{proof}

\begin{lemma}[$G_1$ (supremum) -- $\bigsqcup_n (M_n \cap p) = p$ for $p$ homogeneous]\label{mk-lem:graded_homogeneousSubmodule_iSup_eq}
  \dcref{custom}

  \uses{mk-lem:submodule_isHomogeneous_mathlib}
  Let \(p \subseteq M\) be a homogeneous submodule
  (\ref{mk-lem:submodule_isHomogeneous_mathlib}). Then
  \[
    \bigsqcup_{n} (M_n \cap p) \;=\; p .
  \]
  Together with the independence of
  \ref{mk-lem:graded_homogeneousSubmodule_iSupIndep}, this is the ambient internal
  direct sum \(p = \bigoplus_n (M_n \cap p)\), stated entirely inside \(M\) without
  requiring a direct-sum decomposition of the subtype \(\smash{\hat p}\).
\end{lemma}

\begin{proof}
  \uses{mk-lem:submodule_isHomogeneous_mathlib}
  The inclusion \(\bigsqcup_n (M_n \cap p) \le p\) is clear, since each
  \(M_n \cap p \le p\). For the reverse inclusion, take \(m \in p\) and write
  \(m = \sum_n \mathrm{proj}_n(m)\) with \(\mathrm{proj}_n(m) \in M_n\) using the
  internal decomposition of \(M\). Homogeneity of \(p\) gives
  \(\mathrm{proj}_n(m) \in p\), hence \(\mathrm{proj}_n(m) \in M_n \cap p\); thus
  \(m\) is a finite sum of elements of the \(M_n \cap p\), i.e.
  \(m \in \bigsqcup_n (M_n \cap p)\).
\end{proof}

\subsubsection*{The ambient subquotient class and its Hilbert function}

\begin{definition}[Subquotient Hilbert data]\label{mk-def:graded_subquotientHilb}
  \dcref{custom}



  \uses{mk-lem:graded_homogeneousSubmodule_iSupIndep,
    mk-lem:graded_homogeneousSubmodule_iSup_eq}
  Fix a graded \(\kappa\)-module \(\mathcal{M}\), i.e.\ \(M\) together with an
  internal decomposition \(M = \bigoplus_n \mathcal{M}_n\) whose components
  \(\mathcal{M}_n\) are finite-dimensional over \(\kappa\). A \emph{subquotient
  datum of length \(r\)} consists of:
  \begin{itemize}
    \item a pair of homogeneous submodules \(N' \le N \subseteq M\); and
    \item a family \(t : \{0, \dots, r-1\} \to (M \to_\kappa M)\) of pairwise
      commuting \(\kappa\)-linear endomorphisms, each raising degree by one
      (\(t_i(\mathcal{M}_n) \subseteq \mathcal{M}_{n+1}\)), each preserving \(N\)
      and \(N'\) (\(t_i N \subseteq N\), \(t_i N' \subseteq N'\)); together with
    \item a finiteness condition: the subquotient \(N/N'\) is a finite module over the
      polynomial ring \(\kappa[t_0, \dots, t_{r-1}]\) acting through the \(t_i\) (via the
      ambient \(t\)-stable carriers \(\mathrm{polySubmodule}\,N\),
      \(\mathrm{polySubmodule}\,N'\) of \ref{mk-def:graded_polySubmodule}).
  \end{itemize}
  The pair \((N, N')\) represents the subquotient \(N/N'\). Its \emph{ambient
  Hilbert function} is
  \[
    \mathrm{hilb}(n) \;=\; \Bigl(\dim_\kappa (N \cap \mathcal{M}_n)
      \;-\; \dim_\kappa (N' \cap \mathcal{M}_n)\Bigr) \;\in\; \mathbb{Q},
  \]
  the difference of the (finite) \(\kappa\)-dimensions of the ambient intersections,
  computed in \(\mathbb{Z}\) and cast to \(\mathbb{Q}\). Since \(N' \le N\) one has
  \(N' \cap \mathcal{M}_n \le N \cap \mathcal{M}_n\), so
  \(\mathrm{hilb}(n) = \dim_\kappa\bigl((N \cap \mathcal{M}_n)/(N' \cap \mathcal{M}_n)\bigr)\)
  is the dimension of the degree-\(n\) piece of \(N/N'\). The split \(G_1\)
  (\ref{mk-lem:graded_homogeneousSubmodule_iSupIndep},
  \ref{mk-lem:graded_homogeneousSubmodule_iSup_eq}) guarantees that these ambient
  intersections assemble the homogeneous submodules \(N, N'\), so the data depends
  only on the ambient families \(n \mapsto N \cap \mathcal{M}_n\),
  \(n \mapsto N' \cap \mathcal{M}_n\).
\end{definition}

\subsubsection*{Ambient homogeneity calculus}

The kernel/cokernel closure and the degreewise difference identity rest on a small
calculus of how a degree-raising endomorphism interacts with the internal grading and
with the lattice of homogeneous submodules. Throughout, \(M = \bigoplus_n \mathcal{M}_n\)
is the fixed internally graded \(\kappa\)-module and \(x : M \to M\) is a
\(\kappa\)-linear endomorphism raising degree by one. None of these statements forms a
derived carrier: every object is a submodule of the fixed \(M\).

\begin{definition}[Degree-raising endomorphism]\label{mk-def:graded_raisesDegree}
  \dcref{custom}

  Let \(M = \bigoplus_n \mathcal{M}_n\) be an internally graded \(\kappa\)-module. A
  \(\kappa\)-linear endomorphism \(x : M \to M\) \emph{raises degree by one} when it
  carries each graded piece into the next: \(x(\mathcal{M}_n) \subseteq
  \mathcal{M}_{n+1}\) for every \(n\). This is the grading-ring-free form of
  ``multiplication by a degree-one element'': it abstracts each of the \(r\)
  commuting degree-one generators \(x \in R_1\) acting on \(M\) into a single
  endomorphism, with no reference to the grading ring \(R\).
\end{definition}

\begin{lemma}[Membership form of degree-raising]\label{mk-lem:graded_raisesDegree_mem}
  \dcref{custom}

  \uses{mk-def:graded_raisesDegree}
  If \(x\) raises degree by one (\ref{mk-def:graded_raisesDegree}) and \(m \in
  \mathcal{M}_n\), then \(x m \in \mathcal{M}_{n+1}\).
\end{lemma}

\begin{proof}
  \uses{mk-def:graded_raisesDegree}
  Immediate from the definition: \(x m \in x(\mathcal{M}_n) \subseteq
  \mathcal{M}_{n+1}\).
\end{proof}

\begin{lemma}[Degree-shift commutation]\label{mk-lem:graded_decompose_raisesDegree}
  \dcref{custom}

  \uses{mk-def:graded_raisesDegree, mk-lem:graded_raisesDegree_mem}
  Let \(x\) raise degree by one. Then for every \(m \in M\) and every \(i\) the
  degree-\((i+1)\) homogeneous component of \(x m\) equals \(x\) applied to the
  degree-\(i\) component of \(m\):
  \[
    \mathrm{proj}_{i+1}(x m) \;=\; x\bigl(\mathrm{proj}_i(m)\bigr) .
  \]
  This load-bearing commutation says a degree-raising endomorphism shifts the
  internal decomposition by exactly one slot; it is what makes preimages and images
  of homogeneous submodules under \(x\) homogeneous.
\end{lemma}

\begin{proof}
  \uses{mk-def:graded_raisesDegree, mk-lem:graded_raisesDegree_mem}
  Decompose \(m = \sum_j \mathrm{proj}_j(m)\) into its finitely many homogeneous
  components and apply \(x\): \(x m = \sum_j x(\mathrm{proj}_j(m))\) with
  \(x(\mathrm{proj}_j(m)) \in \mathcal{M}_{j+1}\) by
  \ref{mk-lem:graded_raisesDegree_mem}. Reading off the degree-\((i+1)\) component, all
  terms with \(j+1 \ne i+1\) drop out, leaving \(x(\mathrm{proj}_i(m))\).
\end{proof}

\begin{lemma}[Degree-raising kills the bottom]\label{mk-lem:graded_decompose_raisesDegree_zero}
  \dcref{custom}

  \uses{mk-def:graded_raisesDegree}
  Let \(x\) raise degree by one. Then the degree-\(0\) component of \(x m\) vanishes
  for every \(m\): \(\mathrm{proj}_0(x m) = 0\).
\end{lemma}

\begin{proof}
  \uses{mk-def:graded_raisesDegree}
  Writing \(x m = \sum_j x(\mathrm{proj}_j(m))\) with each summand in
  \(\mathcal{M}_{j+1}\), no summand has degree \(0\) (every \(j+1 \ge 1\)), so the
  degree-\(0\) component is \(0\).
\end{proof}

\begin{lemma}[Preimage of a homogeneous submodule is homogeneous]\label{mk-lem:graded_comap_isHomogeneous}
  \dcref{custom}

  \uses{mk-lem:submodule_isHomogeneous_mathlib, mk-lem:graded_decompose_raisesDegree}
  Let \(x\) raise degree by one and let \(N' \subseteq M\) be a homogeneous submodule
  (\ref{mk-lem:submodule_isHomogeneous_mathlib}). Then the preimage \(x^{-1}(N') =
  \{m : x m \in N'\}\) is a homogeneous submodule of \(M\).
\end{lemma}

\begin{proof}
  \uses{mk-lem:submodule_isHomogeneous_mathlib, mk-lem:graded_decompose_raisesDegree}
  Let \(m \in x^{-1}(N')\) and fix a degree \(i\). By the degree-shift commutation
  (\ref{mk-lem:graded_decompose_raisesDegree}), \(x(\mathrm{proj}_i(m)) =
  \mathrm{proj}_{i+1}(x m)\), which lies in \(N'\) because \(x m \in N'\) and \(N'\)
  is homogeneous. Hence \(\mathrm{proj}_i(m) \in x^{-1}(N')\), so \(x^{-1}(N')\) is
  homogeneous.
\end{proof}

\begin{lemma}[Image of a homogeneous submodule is homogeneous]\label{mk-lem:graded_map_isHomogeneous}
  \dcref{custom}

  \uses{mk-lem:submodule_isHomogeneous_mathlib, mk-lem:graded_decompose_raisesDegree,
    mk-lem:graded_decompose_raisesDegree_zero}
  Let \(x\) raise degree by one and let \(N \subseteq M\) be a homogeneous submodule.
  Then the image \(x \cdot N = \{x m : m \in N\}\) is a homogeneous submodule of
  \(M\), with degree-\((n+1)\) piece \(x(N \cap \mathcal{M}_n)\) and zero
  degree-\(0\) piece.
\end{lemma}

\begin{proof}
  \uses{mk-lem:submodule_isHomogeneous_mathlib, mk-lem:graded_decompose_raisesDegree,
    mk-lem:graded_decompose_raisesDegree_zero}
    Take \(x m \in x \cdot N\) with \(m \in N\) and fix a degree. In degree \(0\) the
  component \(\mathrm{proj}_0(x m) = 0 \in x \cdot N\)
  (\ref{mk-lem:graded_decompose_raisesDegree_zero}). In degree \(i+1\),
  \(\mathrm{proj}_{i+1}(x m) = x(\mathrm{proj}_i(m))\)
  (\ref{mk-lem:graded_decompose_raisesDegree}); since \(N\) is homogeneous,
  \(\mathrm{proj}_i(m) \in N\), so this component lies in \(x \cdot N\). Thus every
  homogeneous component of \(x m\) lies in \(x \cdot N\), which is therefore
  homogeneous.
\end{proof}

\begin{lemma}[Intersection of homogeneous submodules is homogeneous]\label{mk-lem:graded_inf_isHomogeneous}
  \dcref{custom}

  \uses{mk-lem:submodule_isHomogeneous_mathlib}
  If \(p, q \subseteq M\) are homogeneous submodules
  (\ref{mk-lem:submodule_isHomogeneous_mathlib}) then so is their intersection \(p
  \cap q\). Mathlib provides no lattice-closure lemmas for homogeneous submodules, so
  this is recorded here.
\end{lemma}

\begin{proof}
  \uses{mk-lem:submodule_isHomogeneous_mathlib}
  For \(m \in p \cap q\) and any degree \(i\), homogeneity of \(p\) and of \(q\)
  gives \(\mathrm{proj}_i(m) \in p\) and \(\mathrm{proj}_i(m) \in q\), hence
  \(\mathrm{proj}_i(m) \in p \cap q\).
\end{proof}

\begin{lemma}[Sum of homogeneous submodules is homogeneous]\label{mk-lem:graded_sup_isHomogeneous}
  \dcref{custom}

  \uses{mk-lem:submodule_isHomogeneous_mathlib}
  If \(p, q \subseteq M\) are homogeneous submodules then so is their sum \(p + q\).
\end{lemma}

\begin{proof}
  \uses{mk-lem:submodule_isHomogeneous_mathlib}
  Write an element of \(p + q\) as \(a + b\) with \(a \in p\), \(b \in q\). For any
  degree \(i\), \(\mathrm{proj}_i(a+b) = \mathrm{proj}_i(a) + \mathrm{proj}_i(b)\)
  with \(\mathrm{proj}_i(a) \in p\) and \(\mathrm{proj}_i(b) \in q\) by homogeneity,
  so \(\mathrm{proj}_i(a+b) \in p + q\).
\end{proof}

\begin{lemma}[Ambient image identity]\label{mk-lem:graded_map_inf_degree_eq}
  \dcref{custom}

  \uses{mk-lem:graded_decompose_raisesDegree, mk-lem:graded_map_isHomogeneous}
  Let \(x\) raise degree by one and let \(N \subseteq M\) be homogeneous. Then the
  degree-\((n+1)\) part of the image \(x \cdot N\) is exactly \(x\) of the
  degree-\(n\) part of \(N\):
  \[
    (x \cdot N) \cap \mathcal{M}_{n+1} \;=\; x \cdot (N \cap \mathcal{M}_n) .
  \]
\end{lemma}

\begin{proof}
  \uses{mk-lem:graded_decompose_raisesDegree, mk-lem:graded_map_isHomogeneous}
  \(\supseteq\): for \(m \in N \cap \mathcal{M}_n\), \(x m \in x \cdot N\) and \(x m
  \in \mathcal{M}_{n+1}\) since \(x\) raises degree. \(\subseteq\): an element \(x m
  \in x \cdot N\) lying in \(\mathcal{M}_{n+1}\) equals its own degree-\((n+1)\)
  component \(\mathrm{proj}_{n+1}(x m) = x(\mathrm{proj}_n(m))\)
  (\ref{mk-lem:graded_decompose_raisesDegree}); homogeneity of \(N\) puts
  \(\mathrm{proj}_n(m) \in N \cap \mathcal{M}_n\), so \(x m \in x \cdot (N \cap
  \mathcal{M}_n)\). The image \(x \cdot N\) is homogeneous by
  \ref{mk-lem:graded_map_isHomogeneous}.
\end{proof}

\begin{lemma}[Ambient distributive law]\label{mk-lem:graded_sup_inf_degree_eq}
  \dcref{custom}

  \uses{mk-lem:graded_inf_isHomogeneous, mk-lem:graded_sup_isHomogeneous,
    mk-lem:submodule_isHomogeneous_mathlib}
  For homogeneous submodules \(P, Q \subseteq M\) and any degree \(k\), intersecting
  the sum with a graded piece distributes:
  \[
    (P + Q) \cap \mathcal{M}_k \;=\; (P \cap \mathcal{M}_k) + (Q \cap \mathcal{M}_k) .
  \]
\end{lemma}

\begin{proof}
  \uses{mk-lem:graded_inf_isHomogeneous, mk-lem:graded_sup_isHomogeneous,
    mk-lem:submodule_isHomogeneous_mathlib}
    \(\supseteq\) is the general lattice inequality. For \(\subseteq\), take \(z = p +
  q \in (P + Q) \cap \mathcal{M}_k\) with \(p \in P\), \(q \in Q\). Since \(z \in
  \mathcal{M}_k\), \(z = \mathrm{proj}_k(z) = \mathrm{proj}_k(p) +
  \mathrm{proj}_k(q)\); homogeneity of \(P\) and \(Q\)
  (\ref{mk-lem:submodule_isHomogeneous_mathlib}) puts \(\mathrm{proj}_k(p) \in P \cap
  \mathcal{M}_k\) and \(\mathrm{proj}_k(q) \in Q \cap \mathcal{M}_k\), so \(z \in (P
  \cap \mathcal{M}_k) + (Q \cap \mathcal{M}_k)\). The intersection and sum stay in
  the homogeneous lattice by \ref{mk-lem:graded_inf_isHomogeneous} and
  \ref{mk-lem:graded_sup_isHomogeneous}.
\end{proof}

\subsubsection*{Closure under kernel and cokernel}

The kernel/cokernel closure of the subquotient class is realised by eight ambient
component facts: for a degree-raising endomorphism \(x\) and a homogeneous pair
\(N' \le N\), the kernel pair is \((N \cap x^{-1}(N'),\, N')\) and the cokernel pair is
\((N,\, N' + x \cdot N)\). The lemmas below record that both new pairs are homogeneous,
nest correctly, are annihilated by \(x\), and are preserved by any endomorphism \(t\)
commuting with \(x\) that preserves the original pair.

\begin{lemma}[Homogeneity of the kernel pair's lower module]\label{mk-lem:graded_ker_isHomogeneous}
  \dcref{custom}

  \uses{mk-lem:graded_inf_isHomogeneous, mk-lem:graded_comap_isHomogeneous}
  Let \(x\) raise degree by one (\ref{mk-def:graded_raisesDegree}) and let \(N, N'\) be
  homogeneous submodules of \(M\). Then the kernel pair's lower module
  \(N \cap x^{-1}(N')\) is homogeneous. Indeed it is the intersection
  (\ref{mk-lem:graded_inf_isHomogeneous}) of the homogeneous \(N\) with the homogeneous
  preimage \(x^{-1}(N')\) (\ref{mk-lem:graded_comap_isHomogeneous}).
\end{lemma}

\begin{lemma}[Homogeneity of the cokernel pair's upper module]\label{mk-lem:graded_coker_isHomogeneous}
  \dcref{custom}

  \uses{mk-lem:graded_sup_isHomogeneous, mk-lem:graded_map_isHomogeneous}
  Under the same hypotheses, the cokernel pair's upper module \(N' + x \cdot N\) is
  homogeneous: it is the sum (\ref{mk-lem:graded_sup_isHomogeneous}) of the homogeneous
  \(N'\) with the homogeneous image \(x \cdot N\)
  (\ref{mk-lem:graded_map_isHomogeneous}).
\end{lemma}

\begin{lemma}[The kernel pair nests]\label{mk-lem:graded_ker_le}
  \dcref{custom}

  If \(N' \le N\) and \(x\) preserves \(N'\) (\(x \cdot N' \subseteq N'\)), then
  \(N' \le N \cap x^{-1}(N')\). Indeed \(N' \le N\) by hypothesis, and \(x \cdot N'
  \subseteq N'\) is exactly \(N' \subseteq x^{-1}(N')\); so \(N'\) is below the
  intersection. Hence the kernel pair \((N \cap x^{-1}(N'),\, N')\) is a valid
  subquotient pair (\ref{mk-def:graded_subquotientHilb}).
\end{lemma}

\begin{lemma}[The cokernel pair nests]\label{mk-lem:graded_coker_le}
  \dcref{custom}

  If \(N' \le N\) and \(x\) preserves \(N\) (\(x \cdot N \subseteq N\)), then
  \(N' + x \cdot N \le N\): both summands lie in \(N\). Hence the cokernel pair
  \((N,\, N' + x \cdot N)\) is a valid subquotient pair
  (\ref{mk-def:graded_subquotientHilb}).
\end{lemma}

\begin{lemma}[$x$ annihilates the kernel subquotient]\label{mk-lem:graded_ker_annihilate}
  \dcref{custom}

  The endomorphism \(x\) carries the kernel pair's lower module into \(N'\):
  \(x \cdot (N \cap x^{-1}(N')) \subseteq N'\). Indeed \(N \cap x^{-1}(N') \subseteq
  x^{-1}(N')\), and applying \(x\) to a preimage of \(N'\) lands in \(N'\). Thus \(x\)
  acts as zero on the kernel subquotient.
\end{lemma}

\begin{lemma}[$x$ annihilates the cokernel subquotient]\label{mk-lem:graded_coker_annihilate}
  \dcref{custom}

  The endomorphism \(x\) carries \(N\) into the cokernel pair's upper module:
  \(x \cdot N \subseteq N' + x \cdot N\), trivially (the right summand). Thus \(x\)
  acts as zero on the cokernel subquotient \(N/(N' + x \cdot N)\).
\end{lemma}

\begin{lemma}[Commuting endomorphisms preserve the preimage]\label{mk-lem:graded_comap_map_le_of_commute}
  \dcref{custom}

  Let \(t\) commute with \(x\) and preserve \(N'\) (\(t \cdot N' \subseteq N'\)). Then
  \(t\) preserves the preimage: \(t \cdot x^{-1}(N') \subseteq x^{-1}(N')\). Indeed for
  \(m\) with \(x m \in N'\) one has \(x(t m) = t(x m) \in t \cdot N' \subseteq N'\), so
  \(t m \in x^{-1}(N')\). This is what shows the remaining \(t_i\) preserve the kernel
  pair's lower module \(N \cap x^{-1}(N')\).
\end{lemma}

\begin{lemma}[Commuting endomorphisms preserve the image]\label{mk-lem:graded_map_map_le_of_commute}
  \dcref{custom}

  Let \(t\) commute with \(x\) and preserve \(N\) (\(t \cdot N \subseteq N\)). Then
  \(t\) preserves the image: \(t \cdot (x \cdot N) \subseteq x \cdot N\). Indeed for
  \(m \in N\) one has \(t(x m) = x(t m)\) with \(t m \in N\), so \(t(x m) \in x \cdot
  N\). This is what shows the remaining \(t_i\) preserve the cokernel pair's upper
  module \(N' + x \cdot N\).
\end{lemma}

\begin{lemma}[Kernel/cokernel closure of the subquotient class]\label{mk-lem:graded_subquotient_ker_coker}
  \dcref{custom}

  \uses{mk-lem:graded_ker_isHomogeneous, mk-lem:graded_coker_isHomogeneous,
    mk-lem:graded_ker_le, mk-lem:graded_coker_le, mk-lem:graded_ker_annihilate,
    mk-lem:graded_coker_annihilate, mk-lem:graded_comap_map_le_of_commute,
    mk-lem:graded_map_map_le_of_commute,
    mk-def:graded_subquotientHilb, mk-lem:graded_homogeneousSubmodule_iSupIndep,
    mk-lem:graded_homogeneousSubmodule_iSup_eq, mk-lem:ideal_homogeneous_span_mathlib,
    mk-lem:isHomogeneousElem_graded_smul_mathlib, mk-lem:graded_comap_isHomogeneous,
    mk-lem:graded_map_isHomogeneous, mk-lem:graded_inf_isHomogeneous,
    mk-lem:graded_sup_isHomogeneous}
  Let \((N, N')\) be a subquotient datum (\ref{mk-def:graded_subquotientHilb}) with
  endomorphisms \(t_0, \dots, t_{r-1}\), and set \(x = t_{r-1}\). Form the two new
  pairs of homogeneous submodules of \(M\)
  \[
    K \;=\; \bigl(\,N \cap x^{-1}(N'),\ N'\,\bigr),
    \qquad
    C \;=\; \bigl(\,N,\ N' + x \cdot N\,\bigr),
  \]
  where \(x^{-1}(N') = \{\,m \in M : x \cdot m \in N'\,\}\) is the preimage and
  \(x \cdot N\) is the image of \(N\) under \(x\). Then:
  \begin{enumerate}
    \item \(N \cap x^{-1}(N')\) and \(N' + x \cdot N\) are homogeneous submodules of
      \(M\), with \(N' \le N \cap x^{-1}(N')\) and \(N' + x \cdot N \le N\), so both
      \(K\) and \(C\) are subquotient pairs;
    \item the remaining endomorphisms \(t_0, \dots, t_{r-2}\) pairwise commute,
      raise degree by one, and preserve both members of \(K\) and of \(C\); and
    \item \(x\) annihilates both subquotients: \(x \cdot (N \cap x^{-1}(N')) \subseteq N'\)
      and \(x \cdot N \subseteq N' + x \cdot N\).
  \end{enumerate}
  Consequently \(K\) and \(C\) are again subquotient data, now of length \(r-1\)
  (the endomorphisms \(t_0, \dots, t_{r-2}\)), and \(K = \ker x\) and \(C =
  \operatorname{coker} x\) on the subquotient \(N/N'\): \(K\) represents
  \(\ker\bigl(x : N/N' \to N/N'\bigr)\) and \(C\) represents
  \(\operatorname{coker}\bigl(x : N/N' \to N/N'\bigr)\). No quotient or subtype
  carrier is constructed; \(K\) and \(C\) are ambient pairs of homogeneous
  submodules of the fixed \(M\).
\end{lemma}

\begin{proof}
  \uses{mk-lem:graded_ker_isHomogeneous, mk-lem:graded_coker_isHomogeneous,
    mk-lem:graded_ker_le, mk-lem:graded_coker_le, mk-lem:graded_ker_annihilate,
    mk-lem:graded_coker_annihilate, mk-lem:graded_comap_map_le_of_commute,
    mk-lem:graded_map_map_le_of_commute,
    mk-lem:graded_homogeneousSubmodule_iSupIndep,
    mk-lem:graded_homogeneousSubmodule_iSup_eq, mk-lem:ideal_homogeneous_span_mathlib,
    mk-lem:isHomogeneousElem_graded_smul_mathlib, mk-lem:graded_comap_isHomogeneous,
    mk-lem:graded_map_isHomogeneous, mk-lem:graded_inf_isHomogeneous,
    mk-lem:graded_sup_isHomogeneous}
    \emph{Homogeneity.} The preimage \(x^{-1}(N')\) is homogeneous: if
  \(m = \sum_n \mathrm{proj}_n(m)\) with \(x \cdot m \in N'\), then since \(x\) raises
  degree by one (\ref{mk-lem:isHomogeneousElem_graded_smul_mathlib}) the components
  \(x \cdot \mathrm{proj}_n(m) \in \mathcal{M}_{n+1}\) lie in distinct degrees, and
  homogeneity of \(N'\) (split \(G_1\),
  \ref{mk-lem:graded_homogeneousSubmodule_iSup_eq}) forces each
  \(x \cdot \mathrm{proj}_n(m) \in N'\), i.e.\ \(\mathrm{proj}_n(m) \in x^{-1}(N')\).
  Thus \(N \cap x^{-1}(N')\) is an intersection of homogeneous submodules, hence
  homogeneous. The image \(x \cdot N\) is homogeneous with degree-\((n+1)\) piece
  \(x \cdot (N \cap \mathcal{M}_n)\) by
  \ref{mk-lem:isHomogeneousElem_graded_smul_mathlib}, and a sum of homogeneous
  submodules is homogeneous (\ref{mk-lem:ideal_homogeneous_span_mathlib} applied at the
  module level), so \(N' + x \cdot N\) is homogeneous. The inclusions
  \(N' \le N \cap x^{-1}(N')\) (as \(x \cdot N' \subseteq N'\)) and
  \(N' + x \cdot N \le N\) (as \(N' \le N\) and \(x \cdot N \subseteq N\)) are
  immediate. In the language of the ambient homogeneity calculus, this paragraph is
  exactly: \(x^{-1}(N')\) is homogeneous (\ref{mk-lem:graded_comap_isHomogeneous}),
  \(x \cdot N\) is homogeneous (\ref{mk-lem:graded_map_isHomogeneous}), and the
  homogeneous lattice is closed under intersection
  (\ref{mk-lem:graded_inf_isHomogeneous}) and sum
  (\ref{mk-lem:graded_sup_isHomogeneous}).

  \emph{Preservation.} Each \(t_i\) with \(i \le r-2\) commutes with \(x\) and
  preserves \(N, N'\); hence it preserves \(x^{-1}(N')\) (if \(x t_i m = t_i x m \in
  N'\) whenever \(x m \in N'\)), so it preserves \(N \cap x^{-1}(N')\), and it
  preserves \(x \cdot N\) (as \(t_i x N = x t_i N \subseteq x N\)), so it preserves
  \(N' + x \cdot N\). Degree-raising and commutativity are inherited.

  \emph{Annihilation by \(x\).} For \(m \in N \cap x^{-1}(N')\) one has
  \(x \cdot m \in N'\) by definition, so \(x\) carries the first member of \(K\) into
  its second; and \(x \cdot N \subseteq N' + x \cdot N\) trivially. Thus \(x\) acts as
  zero on \(N/N'\) restricted to \(K\) and on the cokernel quotient \(C\).

  \emph{Identification with \(\ker x\) and \(\operatorname{coker} x\).} On the
  subquotient \(N/N'\), the class of \(m \in N\) is killed by \(x\) iff
  \(x \cdot m \in N'\), i.e.\ iff \(m \in N \cap x^{-1}(N')\); so \(\ker(x : N/N' \to
  N/N')\) is exactly \((N \cap x^{-1}(N'))/N'\), represented by \(K\). The image of
  \(x\) on \(N/N'\) is \((N' + x \cdot N)/N'\), so the cokernel
  \((N/N')/\operatorname{im} x = N/(N' + x \cdot N)\) is represented by \(C\).
\end{proof}

\subsubsection*{The subquotient degreewise difference identity ($D_6$)}

\begin{lemma}[preimage dimension equals ambient-intersection dimension]\label{mk-lem:graded_finrank_comap_subtype}
  \dcref{custom}

  For submodules \(p, S\) of a \(\kappa\)-vector space \(M\), the dimension of the
  preimage of \(S\) under the inclusion \(p \hookrightarrow M\) equals the dimension of
  the ambient intersection: \(\dim_\kappa(\iota_p^{-1} S) = \dim_\kappa(p \cap S)\),
  where \(\iota_p : p \to M\) is the canonical injection. Project-local helper consumed
  by the two kernel-dimension computations of \ref{mk-lem:graded_subquotient_degreewise_diff}.
\end{lemma}
\begin{proof}

The inclusion \(\iota_p\) is injective, so it restricts to a \(\kappa\)-linear
  isomorphism from \(\iota_p^{-1} S\) onto its image \(\iota_p(\iota_p^{-1} S) = p \cap S\)
  (image of a comap under an injective map). An isomorphism preserves dimension, giving
  the claimed equality.
\end{proof}

\begin{lemma}[$D_6$ -- subquotient degreewise difference]\label{mk-lem:graded_subquotient_degreewise_diff}
  \dcref{00K1,custom}

  \uses{mk-def:graded_subquotientHilb, mk-lem:graded_subquotient_ker_coker,
    mk-lem:graded_degreewise_finrank_diff, mk-lem:graded_homogeneousSubmodule_iSupIndep,
    mk-lem:graded_homogeneousSubmodule_iSup_eq, mk-lem:graded_finrank_comap_subtype,
    mk-lem:isHomogeneousElem_graded_smul_mathlib, mk-lem:graded_map_inf_degree_eq,
    mk-lem:graded_sup_inf_degree_eq}
  \textit{Source: \cite{stacks-project}, Tag 00K1 (degreewise short exact sequence).}
  Let \((N, N')\) be a subquotient datum (\ref{mk-def:graded_subquotientHilb}) with
  Hilbert function \(h_M\), and let \(x = t_{r-1}\). Let \(h_K\) and \(h_C\) be the
  Hilbert functions of the kernel and cokernel subquotients
  \(K = (N \cap x^{-1}(N'),\, N')\) and \(C = (N,\, N' + x \cdot N)\) of
  \ref{mk-lem:graded_subquotient_ker_coker}. Then for every \(n\)
  \[
    h_M(n+1) - h_M(n) \;=\; h_C(n+1) - h_K(n).
  \]
  Writing \(Q_n = (N \cap \mathcal{M}_n)/(N' \cap \mathcal{M}_n)\) for the
  degree-\(n\) piece of \(N/N'\) (so \(\dim_\kappa Q_n = h_M(n)\)), multiplication by
  \(x\) is a \(\kappa\)-linear map \(x : Q_n \to Q_{n+1}\) with
  \(\ker(x : Q_n \to Q_{n+1})\) the degree-\(n\) piece of \(K\) and
  \(Q_{n+1}/\operatorname{im} x\) the degree-\((n+1)\) piece of \(C\); the identity is
  then the degreewise rank--nullity difference \(D_5\)
  (\ref{mk-lem:graded_degreewise_finrank_diff}). Specialising to \((N, N') = (\top,
  \bot)\) recovers the classical \(0 \to \mathcal{M}_n \xrightarrow{x}
  \mathcal{M}_{n+1} \to \mathcal{M}_{n+1}/x\mathcal{M}_n \to 0\).
\end{lemma}
\begin{proof}
  \uses{mk-lem:graded_subquotient_ker_coker, mk-lem:graded_degreewise_finrank_diff,
    mk-lem:graded_homogeneousSubmodule_iSup_eq,
    mk-lem:isHomogeneousElem_graded_smul_mathlib, mk-lem:graded_map_inf_degree_eq,
    mk-lem:graded_sup_inf_degree_eq}
    Since \(N' \le N\) are homogeneous, the split \(G_1\)
  (\ref{mk-lem:graded_homogeneousSubmodule_iSup_eq}) identifies the degree-\(n\) piece
  of \(N/N'\) with \(Q_n = (N \cap \mathcal{M}_n)/(N' \cap \mathcal{M}_n)\), a
  finite-dimensional \(\kappa\)-vector space of dimension \(h_M(n)\). Because \(x\)
  raises degree by one and preserves \(N\) and \(N'\)
  (\ref{mk-lem:isHomogeneousElem_graded_smul_mathlib}), it induces a \(\kappa\)-linear
  map \(x : Q_n \to Q_{n+1}\).

  \emph{Kernel.} A class \([m] \in Q_n\) (with \(m \in N \cap \mathcal{M}_n\)) lies in
  \(\ker(x : Q_n \to Q_{n+1})\) iff \(x \cdot m \in N' \cap \mathcal{M}_{n+1}\), i.e.\
  iff \(m \in (N \cap x^{-1}(N')) \cap \mathcal{M}_n\). Hence \(\ker(x : Q_n \to
  Q_{n+1}) = \bigl((N \cap x^{-1}(N')) \cap \mathcal{M}_n\bigr)/(N' \cap \mathcal{M}_n)\)
  is exactly the degree-\(n\) piece of \(K\), of dimension \(h_K(n)\)
  (\ref{mk-lem:graded_subquotient_ker_coker}).

  \emph{Cokernel.} The image of \(x : Q_n \to Q_{n+1}\) is
  \(\bigl(x \cdot (N \cap \mathcal{M}_n) + (N' \cap \mathcal{M}_{n+1})\bigr)/(N' \cap
  \mathcal{M}_{n+1})\). Using the ambient identity \((N' + x \cdot N) \cap
  \mathcal{M}_{n+1} = (N' \cap \mathcal{M}_{n+1}) + x \cdot (N \cap \mathcal{M}_n)\)
  -- which is the ambient distributive law (\ref{mk-lem:graded_sup_inf_degree_eq})
  applied to \(P = N'\), \(Q = x \cdot N\), combined with the ambient image identity
  \((x \cdot N) \cap \mathcal{M}_{n+1} = x \cdot (N \cap \mathcal{M}_n)\)
  (\ref{mk-lem:graded_map_inf_degree_eq}) -- the cokernel \(Q_{n+1}/\operatorname{im} x = (N \cap
  \mathcal{M}_{n+1})/\bigl((N' + x \cdot N) \cap \mathcal{M}_{n+1}\bigr)\) is the
  degree-\((n+1)\) piece of \(C\), of dimension \(h_C(n+1)\).

  \emph{Difference.} Apply the degreewise rank--nullity difference \(D_5\)
  (\ref{mk-lem:graded_degreewise_finrank_diff}) to \(\varphi = (x : Q_n \to Q_{n+1})\),
  with \(V = Q_n\), \(W = Q_{n+1}\):
  \[
    h_M(n+1) - h_M(n)
      = \dim_\kappa Q_{n+1} - \dim_\kappa Q_n
      = \dim_\kappa(Q_{n+1}/\operatorname{im}\varphi) - \dim_\kappa(\ker\varphi)
      = h_C(n+1) - h_K(n).
  \]
  This is the \(\mathrm{hdiff}\) hypothesis consumed by \ref{mk-lem:ratHilb_ofDiffEq}.
\end{proof}

\subsubsection*{The free polynomial-ring module structure}

The finiteness condition of a subquotient datum (\ref{mk-def:graded_subquotientHilb}) and
its transfer down one generator (\ref{mk-lem:graded_subquotient_finite_transfer}) are
phrased over the \emph{free} polynomial ring \(\kappa[t_0, \dots, t_{r-1}] =
\operatorname{MvPolynomial}(\{0, \dots, r-1\}, \kappa)\). The following blocks build the
action of that ring on the ambient module \(M\) from the \(r\) pairwise-commuting
degree-raising endomorphisms \(t_i\), keeping every carrier an ambient submodule of
\(M\). The free polynomial ring -- rather than the subalgebra
\(\mathrm{adjoin}_\kappa\{t_0, \dots, t_{r-1}\} \subseteq \operatorname{End}_\kappa(M)\)
-- is used precisely so the inductive transfer has the ring surjection
\(\kappa[t_0, \dots, t_{r-1}] \twoheadrightarrow \kappa[t_0, \dots, t_{r-2}]\) available
(\ref{mk-lem:graded_subquotient_finite_transfer}).

\begin{definition}[Polynomial evaluation ring homomorphism]\label{mk-def:graded_polyEndHom}
  \dcref{custom}

  \uses{mk-lem:adjoinCommRingOfComm_mathlib}
  Let \(t : \{0, \dots, r-1\} \to \operatorname{End}_\kappa(M)\) be a family of
  pairwise-commuting \(\kappa\)-linear endomorphisms. The \emph{polynomial evaluation
  ring homomorphism}
  \[
    \mathrm{polyEndHom}(t) : \kappa[t_0, \dots, t_{r-1}] \longrightarrow
      \operatorname{End}_\kappa(M)
  \]
  sends each variable \(X_i\) to \(t_i\). It is constructed by evaluating into the
  \emph{commutative} \(\kappa\)-subalgebra \(A = \mathrm{adjoin}_\kappa(\{t_i\})
  \subseteq \operatorname{End}_\kappa(M)\) -- commutative by
  \ref{mk-lem:adjoinCommRingOfComm_mathlib} because the generators commute -- and then
  including \(A \hookrightarrow \operatorname{End}_\kappa(M)\); evaluation into the
  noncommutative \(\operatorname{End}_\kappa(M)\) directly is not an algebra
  homomorphism, which is why the commutative \(A\) is interposed.
\end{definition}

\begin{lemma}[Evaluation sends $X_i$ to $t_i$]\label{mk-lem:graded_polyEndHom_X}
  \dcref{custom}

  \uses{mk-def:graded_polyEndHom}
  \(\mathrm{polyEndHom}(t)(X_i) = t_i\) for each \(i\). Immediate from the construction
  (\ref{mk-def:graded_polyEndHom}), as evaluation sends \(X_i\) to the chosen image
  \(t_i\).
\end{lemma}

\begin{lemma}[Evaluation sends a constant to a scalar]\label{mk-lem:graded_polyEndHom_C}
  \dcref{custom}

  \uses{mk-def:graded_polyEndHom}
  \(\mathrm{polyEndHom}(t)(C\,c) = c \cdot \mathrm{id}_M\) for each \(c \in \kappa\):
  the evaluation homomorphism (\ref{mk-def:graded_polyEndHom}) is \(\kappa\)-linear and
  carries the constant polynomial \(C\,c\) to the scalar endomorphism \(c \cdot 1\).
\end{lemma}

\begin{definition}[Polynomial-ring module structure on $M$]\label{mk-def:graded_polyModule}
  \dcref{custom}

  \uses{mk-def:graded_polyEndHom}
  The \emph{polynomial-module structure} is the \(\kappa[t_0, \dots, t_{r-1}]\)-module
  structure on \(M\) obtained by restricting scalars along
  \(\mathrm{polyEndHom}(t)\) (\ref{mk-def:graded_polyEndHom}): a polynomial \(p\) acts on
  \(m \in M\) by \(p \cdot m = \mathrm{polyEndHom}(t)(p)(m)\). This is the module over the
  free polynomial ring on which the ambient subquotient induction's finiteness condition
  is measured.
\end{definition}

\begin{lemma}[$X_i$ acts as $t_i$]\label{mk-lem:graded_polyModule_X_smul}
  \dcref{custom}

  \uses{mk-def:graded_polyModule, mk-lem:graded_polyEndHom_X}
  In the polynomial-module structure (\ref{mk-def:graded_polyModule}), \(X_i \cdot m =
  t_i(m)\) for all \(m \in M\). Immediate from
  \(\mathrm{polyEndHom}(t)(X_i) = t_i\) (\ref{mk-lem:graded_polyEndHom_X}).
\end{lemma}

\begin{lemma}[A constant acts as a scalar]\label{mk-lem:graded_polyModule_C_smul}
  \dcref{custom}

  \uses{mk-def:graded_polyModule, mk-lem:graded_polyEndHom_C}
  In the polynomial-module structure (\ref{mk-def:graded_polyModule}), \((C\,c) \cdot m =
  c \cdot m\) for all \(c \in \kappa\), \(m \in M\). Immediate from
  \(\mathrm{polyEndHom}(t)(C\,c) = c \cdot \mathrm{id}_M\)
  (\ref{mk-lem:graded_polyEndHom_C}).
\end{lemma}


\begin{definition}[Polynomial-ring submodule from a $t$-stable submodule]\label{mk-def:graded_polySubmodule}
  \dcref{custom}

  \uses{mk-def:graded_polyModule, mk-lem:graded_polyModule_X_smul,
    mk-lem:graded_polyModule_C_smul}
  Let \(N \subseteq M\) be a \(\kappa\)-submodule that is stable under each \(t_i\)
  (\(t_i \cdot N \subseteq N\)). Then \(N\), with the \emph{same} ambient carrier, is a
  \(\kappa[t_0, \dots, t_{r-1}]\)-submodule of \(M\) for the polynomial-module structure
  (\ref{mk-def:graded_polyModule}). Stability under the action of an arbitrary polynomial
  follows by induction on monomials: constants act as scalars
  (\ref{mk-lem:graded_polyModule_C_smul}) and multiplication by \(X_i\) acts as \(t_i\)
  (\ref{mk-lem:graded_polyModule_X_smul}), both of which preserve \(N\).
\end{definition}


\subsubsection*{Finiteness-transfer infrastructure}

The ambient finiteness transfer (\ref{mk-lem:graded_subquotient_finite_transfer})
factors the polynomial action of the larger ring through a surjection onto the
smaller one. The following blocks build that surjection and the semilinearity
identity it rests on, and record the two elementary finiteness manipulations
(shrinking a numerator, enlarging a denominator) used to feed the kernel and
cokernel constructors.

\begin{definition}[The last-variable surjection of polynomial rings]\label{mk-lem:graded_lastVarAlgHom}
  \dcref{custom}







  \uses{mk-lem:fg_restrictScalars_of_surjective_mathlib}
  For a field \(\kappa\) and \(r \in \mathbb{N}\), the \emph{last-variable
  homomorphism} is the \(\kappa\)-algebra map
  \[
    \mathrm{lastVarAlgHom} : \kappa[t_0, \dots, t_r]
      = \operatorname{MvPolynomial}(\{0,\dots,r\}, \kappa)
      \;\longrightarrow\;
      \operatorname{MvPolynomial}(\{0,\dots,r-1\}, \kappa) = \kappa[t_0, \dots, t_{r-1}]
  \]
  obtained by evaluation: it sends the last variable \(X_r\) to \(0\) and each
  earlier variable \(X_i\) (for \(0 \le i < r\)) to \(X_i\), and fixes constants.
  It is a left inverse of the algebra map
  \(\kappa[t_0, \dots, t_{r-1}] \to \kappa[t_0, \dots, t_r]\)
  induced by the canonical inclusion \(\{0, \dots, r-1\} \hookrightarrow \{0, \dots, r\}\)
  on variable indices, and is therefore surjective. This is the concrete
  ring surjection \(\kappa[t_0, \dots, t_r] \twoheadrightarrow
  \kappa[t_0, \dots, t_{r-1}]\) along which finiteness is transferred
  (\ref{mk-lem:fg_restrictScalars_of_surjective_mathlib}).
\end{definition}

\begin{lemma}[A $t$-stable submodule is closed under the polynomial action]\label{mk-lem:graded_polyEndHom_mem_of_stable}
  \dcref{custom}

  \uses{mk-def:graded_polyEndHom, mk-lem:graded_polyEndHom_X, mk-lem:graded_polyEndHom_C}
  Let \(t_0, \dots, t_{r-1}\) be pairwise-commuting \(\kappa\)-linear endomorphisms
  of \(M\) and let \(P' \subseteq M\) be a \(\kappa\)-submodule stable under each
  \(t_i\) (\(t_i \cdot P' \subseteq P'\)). Then \(P'\) is closed under the action of
  every polynomial: for all \(p \in \kappa[t_0, \dots, t_{r-1}]\) and \(m \in P'\),
  \(\mathrm{polyEndHom}(t)(p)(m) \in P'\) (\ref{mk-def:graded_polyEndHom}). The proof
  is induction on \(p\): a constant acts as a scalar (\ref{mk-lem:graded_polyEndHom_C}),
  multiplication by a variable acts as the stabilising \(t_i\)
  (\ref{mk-lem:graded_polyEndHom_X}), and sums are handled additively.
\end{lemma}

\begin{lemma}[Mod-$P'$ semilinearity of the polynomial action]\label{mk-lem:graded_polyEndHom_lastVar_sub_mem}
  \dcref{custom}

  \uses{mk-def:graded_polyEndHom, mk-lem:graded_polyEndHom_X, mk-lem:graded_polyEndHom_C,
    mk-lem:graded_polyEndHom_mem_of_stable, mk-lem:graded_lastVarAlgHom}
  Let \(t_0, \dots, t_r\) be pairwise-commuting \(\kappa\)-linear endomorphisms of
  \(M\), and let \(P' \le P\) be \(\kappa\)-submodules each stable under every
  \(t_i\), with the last endomorphism \(x = t_r\) carrying \(P\) into \(P'\)
  (\(x \cdot P \subseteq P'\)). Then for every polynomial
  \(s \in \kappa[t_0, \dots, t_r]\) and every \(m \in P\),
  \[
    \mathrm{polyEndHom}(t)(s)(m)
      \;-\; \mathrm{polyEndHom}(t_0, \dots, t_{r-1})
        \bigl(\mathrm{lastVarAlgHom}(s)\bigr)(m)
      \;\in\; P' .
  \]
  In words: acting by \(s\) through the full family \(t\) agrees, modulo \(P'\), with
  first projecting away the last variable via \(\mathrm{lastVarAlgHom}\)
  (\ref{mk-lem:graded_lastVarAlgHom}) and acting through the reduced family
  \((t_0, \dots, t_{r-1})\). This is the semilinearity identity that makes the
  \(\sigma\)-semilinear quotient map of \ref{mk-lem:graded_subquotient_finite_transfer}
  well-defined. The proof is induction on \(s\): the case of an earlier variable
  \(s = X_i\) (\(i < r\)) is the inductive hypothesis after \(\mathrm{lastVarAlgHom}\)
  fixes it; the case of the last variable \(s = X_r\) annihilates the reduced term,
  and the full term lands in \(P'\) because \(x = t_r\) carries \(P\) into the
  \(t\)-stable \(P'\) (\ref{mk-lem:graded_polyEndHom_mem_of_stable}).
\end{lemma}

\begin{lemma}[Shrinking the numerator preserves finiteness]\label{mk-lem:graded_polyQuot_finite_of_le_numerator}
  \dcref{custom}

  \uses{mk-def:graded_polyModule, mk-def:graded_polySubmodule,
    mk-lem:module_finite_of_injective_mathlib,
    mk-lem:mvPolynomial_isNoetherianRing_fin_mathlib,
    mk-lem:isNoetherian_of_isNoetherianRing_of_finite_mathlib}
  Let \(N_1 \le N_2\) and \(P'\) be \(\kappa\)-submodules of \(M\), each stable under
  the commuting family \(t_0, \dots, t_{r-1}\), and suppose the subquotient
  \(N_2/P'\) is finite over \(\kappa[t_0, \dots, t_{r-1}]\)
  (\ref{mk-def:graded_polyModule}, \ref{mk-def:graded_polySubmodule}). Then \(N_1/P'\) is
  also finite over \(\kappa[t_0, \dots, t_{r-1}]\). Indeed the inclusion
  \(N_1 \subseteq N_2\) induces an injective \(\kappa[t_0, \dots, t_{r-1}]\)-linear
  map \(N_1/P' \hookrightarrow N_2/P'\); since
  \(\kappa[t_0, \dots, t_{r-1}]\) is Noetherian
  (\ref{mk-lem:mvPolynomial_isNoetherianRing_fin_mathlib}), the finite module
  \(N_2/P'\) is Noetherian
  (\ref{mk-lem:isNoetherian_of_isNoetherianRing_of_finite_mathlib}), so its submodule
  \(N_1/P'\) is finite (\ref{mk-lem:module_finite_of_injective_mathlib}). This supplies
  the finiteness witness of the kernel constructor
  (\ref{mk-def:graded_subquotientDatum_ker}).
\end{lemma}

\begin{lemma}[Enlarging the denominator preserves finiteness]\label{mk-lem:graded_polyQuot_finite_of_le_denominator}
  \dcref{custom}

  \uses{mk-def:graded_polyModule, mk-def:graded_polySubmodule,
    mk-lem:module_finite_of_surjective_mathlib, mk-lem:submodule_liftQ_mathlib}
  Let \(N\) and \(P_1 \le P_2\) be \(\kappa\)-submodules of \(M\), each stable under
  the commuting family \(t_0, \dots, t_{r-1}\), and suppose the subquotient
  \(N/P_1\) is finite over \(\kappa[t_0, \dots, t_{r-1}]\)
  (\ref{mk-def:graded_polyModule}, \ref{mk-def:graded_polySubmodule}). Then \(N/P_2\) is
  also finite over \(\kappa[t_0, \dots, t_{r-1}]\). Indeed the inclusion
  \(P_1 \subseteq P_2\) gives a surjection \(N/P_1 \twoheadrightarrow N/P_2\)
  (\ref{mk-lem:submodule_liftQ_mathlib}), and finiteness transfers along it
  (\ref{mk-lem:module_finite_of_surjective_mathlib}). This supplies the
  finiteness witness of the cokernel constructor
  (\ref{mk-def:graded_subquotientDatum_coker}).
\end{lemma}

\subsubsection*{Ambient finiteness transfer}

\begin{lemma}[Abstract single-pair finiteness transfer down one variable]\label{mk-lem:graded_subquotient_finite_transfer}
  \dcref{custom}

  \uses{mk-def:graded_polyModule, mk-def:graded_polySubmodule,
    mk-lem:graded_lastVarAlgHom, mk-lem:graded_polyEndHom_lastVar_sub_mem,
    mk-lem:module_finite_of_surjective_mathlib, mk-lem:submodule_liftQ_mathlib,
    mk-lem:fg_restrictScalars_of_surjective_mathlib}
  Let \(t_0, \dots, t_r\) be pairwise-commuting \(\kappa\)-linear endomorphisms of
  \(M\), and let \(P' \le P\) be \(\kappa\)-submodules of \(M\), each stable under
  every \(t_i\). Suppose the last endomorphism \(x = t_r\) carries \(P\) into \(P'\)
  (\(x \cdot P \subseteq P'\)), and that the subquotient \(P/P'\) is finite over the
  free polynomial ring \(\kappa[t_0, \dots, t_r]\) for the polynomial-module structure
  of \ref{mk-def:graded_polyModule} (with carriers the ambient \(t\)-stable submodules
  \(\mathrm{polySubmodule}\,P\), \(\mathrm{polySubmodule}\,P'\),
  \ref{mk-def:graded_polySubmodule}). Then \(P/P'\) is finite over
  \(\kappa[t_0, \dots, t_{r-1}]\) for the reduced family \((t_0, \dots, t_{r-1})\).

  This is the abstract single-pair transfer: it does not mention the kernel/cokernel
  pairs, but is the engine applied to each of them by the kernel and cokernel
  constructors (\ref{mk-def:graded_subquotientDatum_ker},
  \ref{mk-def:graded_subquotientDatum_coker}). Because the last variable
  annihilates the subquotient, the \(\kappa[t_0, \dots, t_r]\)-action factors through
  the last-variable surjection \(\kappa[t_0, \dots, t_r] \twoheadrightarrow
  \kappa[t_0, \dots, t_{r-1}]\), which is what makes the transfer drop one variable.
\end{lemma}

\begin{proof}
  \uses{mk-def:graded_polyModule, mk-def:graded_polySubmodule,
    mk-lem:graded_lastVarAlgHom, mk-lem:graded_polyEndHom_lastVar_sub_mem,
    mk-lem:module_finite_of_surjective_mathlib, mk-lem:submodule_liftQ_mathlib}
    Write \(\sigma = \mathrm{lastVarAlgHom} : \kappa[t_0, \dots, t_r]
  \twoheadrightarrow \kappa[t_0, \dots, t_{r-1}]\) for the surjective ring
  homomorphism sending the last variable to \(0\) and each earlier variable \(X_i\)
  (\(i < r\)) to \(X_i\) (\ref{mk-lem:graded_lastVarAlgHom}). Consider the map on numerators
  \[
    g : \mathrm{polySubmodule}\,P \;\longrightarrow\;
      \bigl(\mathrm{polySubmodule}\,P\bigr)\big/
        \bigl(\mathrm{polySubmodule}\,P'\bigr), \qquad y \longmapsto [\,y\,],
  \]
  where the source carries the \(\kappa[t_0, \dots, t_r]\)-module structure and the
  target the reduced \(\kappa[t_0, \dots, t_{r-1}]\)-module structure (both with the
  same ambient carrier \(P\), \ref{mk-def:graded_polySubmodule}). The mod-\(P'\)
  semilinearity identity (\ref{mk-lem:graded_polyEndHom_lastVar_sub_mem}) -- acting by
  \(s\) through the full family and acting by \(\sigma(s)\) through the reduced family
  agree modulo \(P'\), precisely because \(x = t_r\) carries \(P\) into \(P'\) --
  says exactly that \(g\) is \(\sigma\)-semilinear. The map \(g\) is surjective (it is
  the quotient projection on the common carrier), and it kills the denominator
  \(\mathrm{polySubmodule}\,P'\), so it descends through that submodule
  (\ref{mk-lem:submodule_liftQ_mathlib}) to a surjective
  \(\kappa[t_0, \dots, t_{r-1}]\)-linear map
  \[
    \bigl(\mathrm{polySubmodule}\,P\bigr)\big/\bigl(\mathrm{polySubmodule}\,P'\bigr)
      \;\twoheadrightarrow\;
    \bigl(\mathrm{polySubmodule}\,P\bigr)\big/\bigl(\mathrm{polySubmodule}\,P'\bigr)
  \]
  from the \(\kappa[t_0, \dots, t_r]\)-finite subquotient onto the reduced one.
  Finiteness transfers along this surjection
  (\ref{mk-lem:module_finite_of_surjective_mathlib}), giving finiteness of \(P/P'\)
  over \(\kappa[t_0, \dots, t_{r-1}]\).

  \emph{Why the free polynomial ring.} The smaller ring must be the \emph{free}
  polynomial ring \(\kappa[t_0, \dots, t_{r-1}]\), not the subalgebra
  \(\mathrm{adjoin}_\kappa\{t_0, \dots, t_{r-1}\} \subseteq
  \operatorname{End}_\kappa(M)\): relations among the endomorphisms inside
  \(\operatorname{End}_\kappa(M)\) would obstruct the polynomial-ring surjection
  \(\sigma\) that the scalar transfer requires
  (\ref{mk-lem:fg_restrictScalars_of_surjective_mathlib}). Keeping free polynomial
  rings at every step keeps \(\sigma\) -- and hence the transfer -- available
  throughout. Everything is finite generation of ambient submodules of the fixed
  \(M\) and their quotients.
\end{proof}

\subsubsection*{Subquotient constructors}

The kernel/cokernel closure (\ref{mk-lem:graded_subquotient_ker_coker}) and the
abstract finiteness transfer (\ref{mk-lem:graded_subquotient_finite_transfer}) are
bundled into two constructors that take a length-\((r+1)\) subquotient datum to the
length-\(r\) kernel and cokernel data consumed by the induction. The two stability
lemmas below first record that the new pairs are stable under the \emph{whole}
endomorphism family (not merely the reduced one), which is what lets the over-\(\kappa[t_0,\dots,t_r]\)
finiteness inputs be stated before the variable is dropped.

\begin{lemma}[The kernel pair's lower module is stable under the full family]\label{mk-lem:graded_ker_stable_full}
  \dcref{custom}

  \uses{mk-lem:graded_comap_map_le_of_commute}
  Let \((N, N')\) be a length-\((r+1)\) subquotient datum
  (\ref{mk-def:graded_subquotientHilb}) with endomorphisms \(t_0, \dots, t_r\) and
  \(x = t_r\). Then the kernel pair's lower module \(N \cap x^{-1}(N')\) is stable
  under \emph{every} \(t_i\) (\(0 \le i \le r\)), not only the reduced family: each
  \(t_i\) preserves \(N\) and preserves the preimage \(x^{-1}(N')\) because it
  commutes with \(x\) and preserves \(N'\)
  (\ref{mk-lem:graded_comap_map_le_of_commute}). This full stability is what allows the
  kernel subquotient's finiteness to be measured over \(\kappa[t_0, \dots, t_r]\)
  before the transfer drops the last variable.
\end{lemma}

\begin{lemma}[The cokernel pair's upper module is stable under the full family]\label{mk-lem:graded_coker_stable_full}
  \dcref{custom}

  \uses{mk-lem:graded_map_map_le_of_commute}
  Under the same hypotheses, the cokernel pair's upper module \(N' + x \cdot N\) is
  stable under every \(t_i\) (\(0 \le i \le r\)): each \(t_i\) preserves \(N'\) and
  preserves the image \(x \cdot N\) because it commutes with \(x\) and preserves
  \(N\) (\ref{mk-lem:graded_map_map_le_of_commute}), and a map preserves a sum of
  preserved submodules. This full stability supports the over-\(\kappa[t_0, \dots,
  t_r]\) finiteness input of the cokernel constructor.
\end{lemma}

\begin{definition}[Kernel constructor of a subquotient datum]\label{mk-def:graded_subquotientDatum_ker}
  \dcref{custom}

  \uses{mk-def:graded_subquotientHilb, mk-lem:graded_subquotient_ker_coker,
    mk-lem:graded_ker_isHomogeneous, mk-lem:graded_ker_le, mk-lem:graded_ker_annihilate,
    mk-lem:graded_comap_map_le_of_commute, mk-lem:graded_ker_stable_full,
    mk-lem:graded_subquotient_finite_transfer,
    mk-lem:graded_polyQuot_finite_of_le_numerator}
  From a length-\((r+1)\) subquotient datum \((N, N')\) with endomorphisms
  \(t_0, \dots, t_r\) and \(x = t_r\) (\ref{mk-def:graded_subquotientHilb}), the
  \emph{kernel constructor} produces the length-\(r\) datum
  \[
    \ker D \;=\; \bigl(\,N \cap x^{-1}(N'),\ N'\,\bigr)
  \]
  on the reduced family \((t_0, \dots, t_{r-1})\). Its
  non-finiteness fields are supplied by the ambient kernel calculus: homogeneity of
  \(N \cap x^{-1}(N')\) (\ref{mk-lem:graded_ker_isHomogeneous}), the nesting
  \(N' \le N \cap x^{-1}(N')\) (\ref{mk-lem:graded_ker_le}), and stability under the
  reduced family (\ref{mk-lem:graded_ker_stable_full}). The finiteness witness
  is obtained by first shrinking the numerator from \(N\) to
  \(N \cap x^{-1}(N')\) (\ref{mk-lem:graded_polyQuot_finite_of_le_numerator}) to get
  \(\kappa[t_0, \dots, t_r]\)-finiteness of the kernel subquotient, then transferring
  down one variable (\ref{mk-lem:graded_subquotient_finite_transfer}), valid because
  \(x\) annihilates the kernel subquotient (\ref{mk-lem:graded_ker_annihilate}). This
  realises the kernel half of \ref{mk-lem:graded_subquotient_ker_coker} as a genuine
  length-\(r\) datum.
\end{definition}

\begin{definition}[Cokernel constructor of a subquotient datum]\label{mk-def:graded_subquotientDatum_coker}
  \dcref{custom}

  \uses{mk-def:graded_subquotientHilb, mk-lem:graded_subquotient_ker_coker,
    mk-lem:graded_coker_isHomogeneous, mk-lem:graded_coker_le, mk-lem:graded_coker_annihilate,
    mk-lem:graded_map_map_le_of_commute, mk-lem:graded_coker_stable_full,
    mk-lem:graded_subquotient_finite_transfer,
    mk-lem:graded_polyQuot_finite_of_le_denominator}
  From the same length-\((r+1)\) datum, the \emph{cokernel constructor} produces the
  length-\(r\) datum
  \[
    \operatorname{coker} D \;=\; \bigl(\,N,\ N' + x \cdot N\,\bigr)
  \]
  on the reduced family. Its non-finiteness fields come from the ambient cokernel
  calculus: homogeneity of \(N' + x \cdot N\) (\ref{mk-lem:graded_coker_isHomogeneous}),
  the nesting \(N' + x \cdot N \le N\) (\ref{mk-lem:graded_coker_le}), and stability
  under the reduced family (\ref{mk-lem:graded_coker_stable_full}). The finiteness field
  is obtained by enlarging the denominator from \(N'\) to \(N' + x \cdot N\)
  (\ref{mk-lem:graded_polyQuot_finite_of_le_denominator}) to get
  \(\kappa[t_0, \dots, t_r]\)-finiteness of the cokernel subquotient, then
  transferring down one variable (\ref{mk-lem:graded_subquotient_finite_transfer}),
  valid because \(x\) annihilates the cokernel subquotient
  (\ref{mk-lem:graded_coker_annihilate}). This realises the cokernel half of
  \ref{mk-lem:graded_subquotient_ker_coker} as a genuine length-\(r\) datum.
\end{definition}

\subsubsection*{The degreewise difference identity}

\begin{lemma}[$D_5$ -- degreewise rank--nullity difference]\label{mk-lem:graded_degreewise_finrank_diff}
  \dcref{custom}

  \uses{mk-lem:finrank_range_add_finrank_ker_mathlib,
    mk-lem:finrank_ses_additive_mathlib}
  Let \(\varphi : V \to W\) be a \(\kappa\)-linear map between finite-dimensional
  \(\kappa\)-vector spaces. Then
  \[
    \dim_\kappa W - \dim_\kappa V
      \;=\; \dim_\kappa(W/\operatorname{range}\varphi) - \dim_\kappa(\ker\varphi).
  \]
  Applied degreewise with \(\varphi = (x : \mathcal{M}_n \to \mathcal{M}_{n+1})\),
  \(V = \mathcal{M}_n\), \(W = \mathcal{M}_{n+1}\) -- so that \(\ker\varphi\) is the
  degree-\(n\) piece \(K_n\) of the kernel subquotient and
  \(W/\operatorname{range}\varphi = \mathcal{M}_{n+1}/x \mathcal{M}_n\) the
  degree-\((n+1)\) piece \(C_{n+1}\) of the cokernel subquotient
  (\ref{mk-lem:graded_subquotient_degreewise_diff}, via the split \(G_1\)
  \ref{mk-lem:graded_homogeneousSubmodule_iSup_eq} and the degree shift
  \ref{mk-lem:isHomogeneousElem_graded_smul_mathlib}) -- this gives, for every \(n\),
  \[
    \dim_\kappa \mathcal{M}_{n+1} - \dim_\kappa \mathcal{M}_n
      \;=\; \dim_\kappa C_{n+1} - \dim_\kappa K_n,
  \]
  which is the difference identity \(\mathrm{hdiff}\) consumed by
  \ref{mk-lem:ratHilb_ofDiffEq}. No graded structure is used in proving the
  underlying vector-space identity -- it is pure rank--nullity.
\end{lemma}

\begin{proof}
  \uses{mk-lem:finrank_range_add_finrank_ker_mathlib,
    mk-lem:finrank_ses_additive_mathlib}
    By rank--nullity (\ref{mk-lem:finrank_range_add_finrank_ker_mathlib}),
  \(\dim_\kappa V = \dim_\kappa(\ker\varphi) + \dim_\kappa(\operatorname{range}
  \varphi)\), so \(\dim_\kappa(\operatorname{range}\varphi) = \dim_\kappa V -
  \dim_\kappa(\ker\varphi)\). By additivity on the short exact sequence
  \(0 \to \operatorname{range}\varphi \to W \to W/\operatorname{range}\varphi \to 0\)
  (\ref{mk-lem:finrank_ses_additive_mathlib}),
  \(\dim_\kappa(W/\operatorname{range}\varphi) = \dim_\kappa W -
  \dim_\kappa(\operatorname{range}\varphi) = \dim_\kappa W - \dim_\kappa V +
  \dim_\kappa(\ker\varphi)\). Rearranging gives
  \(\dim_\kappa W - \dim_\kappa V = \dim_\kappa(W/\operatorname{range}\varphi) -
  \dim_\kappa(\ker\varphi)\). The degreewise application identifies
  \(\ker(x : \mathcal{M}_n \to \mathcal{M}_{n+1})\) and
  \(\mathcal{M}_{n+1}/\operatorname{range}(x : \mathcal{M}_n \to \mathcal{M}_{n+1})\)
  with the degree-\(n\) piece of the kernel subquotient and the degree-\((n+1)\)
  piece of the cokernel subquotient (\ref{mk-lem:graded_subquotient_degreewise_diff}),
  the latter because the degree-\((n+1)\) graded piece of the homogeneous submodule
  \(x\mathcal{M}\) is exactly \(x \mathcal{M}_n\)
  (\ref{mk-lem:isHomogeneousElem_graded_smul_mathlib}).
\end{proof}

\subsubsection*{The base case ($r = 0$)}

\begin{lemma}[Base-case finiteness over the empty polynomial ring]\label{mk-lem:graded_finiteDimensional_of_mvPolynomial_isEmpty_finite}
  \dcref{custom}

  Let \(Q\) be a \(\kappa\)-vector space that is also a module over
  \(\operatorname{MvPolynomial}(\sigma, \kappa)\) for an \emph{empty} index set
  \(\sigma\) -- in particular \(\sigma = \{0, \dots, r-1\}\) with \(r = 0\) -- in a way
  compatible with the \(\kappa\)-action (scalar tower). If \(Q\) is module-finite over
  \(\operatorname{MvPolynomial}(\sigma, \kappa)\), then \(Q\) is finite-dimensional over
  \(\kappa\). Indeed for empty \(\sigma\) the polynomial ring is
  \(\operatorname{MvPolynomial}(\sigma, \kappa) \cong \kappa\), a finite \(\kappa\)-module,
  so finiteness of \(Q\) over it transfers to finiteness over \(\kappa\) along the scalar
  tower. This is the length-zero input to the ambient subquotient induction
  (\ref{mk-lem:graded_subquotient_isRatHilb}): when \(r = 0\) the subquotient \(N/N'\) is a
  finite-dimensional \(\kappa\)-vector space, hence \(\mathrm{hilb}\) is eventually zero.
\end{lemma}

\begin{lemma}[independence of images modulo the kernel]\label{mk-lem:graded_iSupIndep_map_of_mem_ker_sup}
  \dcref{custom}

  Let \(\pi : A \to Q\) be a \(\kappa\)-linear map and \((g_i)_{i \in \iota}\) a family of
  subspaces of \(A\) that is \emph{independent modulo \(\ker\pi\)}: for each \(i\), every
  \(a \in g_i\) that also lies in \(\ker\pi + \sum_{j \ne i} g_j\) already lies in
  \(\ker\pi\). Then the image family \((\pi(g_i))_{i \in \iota}\) is independent
  (\(\mathrm{iSupIndep}\)) in \(Q\). Ring-agnostic lattice core of the base-case
  independence step below; applied to the quotient map \(N \twoheadrightarrow N/N'\).
\end{lemma}
\begin{proof}

By the disjointness criterion for independence it suffices to show, for each \(i\),
  that \(\pi(g_i) \cap \sum_{j \ne i} \pi(g_j) = 0\). Take \(q\) in that intersection;
  write \(q = \pi(a)\) with \(a \in g_i\), and (pushing the sum through \(\pi\), since
  \(\pi(\sum_{j\ne i} g_j) = \sum_{j\ne i}\pi(g_j)\)) \(q = \pi(b)\) with
  \(b \in \sum_{j \ne i} g_j\). Then \(a - b \in \ker\pi\), so
  \(a = (a-b) + b \in \ker\pi + \sum_{j \ne i} g_j\). By the hypothesis applied at \(i\),
  \(a \in \ker\pi\), whence \(q = \pi(a) = 0\).
\end{proof}

\begin{lemma}[Base case: the length-zero subquotient Hilbert function vanishes eventually]\label{mk-lem:graded_subquotient_base_eventuallyZero}
  \dcref{00K1,custom}

  \uses{mk-def:graded_subquotientHilb,
    mk-lem:graded_finiteDimensional_of_mvPolynomial_isEmpty_finite,
    mk-lem:graded_homogeneousSubmodule_iSupIndep, mk-lem:graded_iSupIndep_map_of_mem_ker_sup,
    mk-lem:finite_ne_bot_of_iSupIndep_mathlib, mk-lem:ratHilb_ofEventuallyZero}
  \textit{Source: \cite{stacks-project}, Tag 00K1 (base case of the induction).}
  Let \((N, N')\) be a subquotient datum of length \(0\)
  (\ref{mk-def:graded_subquotientHilb}): a homogeneous pair \(N' \le N\) of the fixed
  ambient graded module \(M = \bigoplus_n \mathcal{M}_n\) with \emph{no} endomorphisms,
  whose finiteness witness says the subquotient \(Q_0 := N/N'\) is module-finite over
  \(\operatorname{MvPolynomial}(\varnothing, \kappa)\). Then the ambient Hilbert
  function vanishes eventually: there is \(K \in \mathbb{N}\) with
  \(\mathrm{hilb}(n) = 0\) for all \(n > K\). Consequently
  \(\mathrm{IsRatHilb}(\mathrm{hilb}, 0)\) (\ref{mk-lem:ratHilb_ofEventuallyZero}),
  which is the base case of the ambient subquotient induction
  (\ref{mk-lem:graded_subquotient_isRatHilb}).
\end{lemma}

\begin{proof}
  \uses{mk-lem:graded_finiteDimensional_of_mvPolynomial_isEmpty_finite,
    mk-lem:graded_homogeneousSubmodule_iSupIndep,
    mk-lem:finite_ne_bot_of_iSupIndep_mathlib, mk-lem:ratHilb_ofEventuallyZero}
    Since the index set of variables is empty,
  \(\operatorname{MvPolynomial}(\varnothing, \kappa) \cong \kappa\), so the finiteness
  witness makes the subquotient \(Q_0 = N/N'\) a \emph{finite-dimensional}
  \(\kappa\)-vector space
  (\ref{mk-lem:graded_finiteDimensional_of_mvPolynomial_isEmpty_finite}). For each
  degree \(n\), composing the inclusion of the homogeneous piece
  \(N \cap \mathcal{M}_n \hookrightarrow N\) with the quotient projection
  \(N \twoheadrightarrow Q_0\) gives a \(\kappa\)-linear map
  \(\psi_n : N \cap \mathcal{M}_n \to Q_0\), \(m \mapsto [m]\), whose image is the
  class of the degree-\(n\) homogeneous piece inside the subquotient.

  \emph{Step 1: independence of the images.} The family
  \(\bigl(\operatorname{range}(\psi_n)\bigr)_{n}\) is \(\mathrm{iSupIndep}\) in
  \(Q_0\). The reason is the ambient direct-sum grading of \(M\): the homogeneous
  pieces \(\mathcal{M}_n\) form an independent family
  (\ref{mk-lem:graded_homogeneousSubmodule_iSupIndep}). An element of
  \(\bigsqcup_{j \ne n} \operatorname{range}(\psi_j)\) is a finite sum of classes of
  homogeneous elements of degrees \(j \ne n\); since \(N'\) is homogeneous, the
  degree-\(n\) homogeneous component of any representative of such a class lies in
  \(N'\), so the class has zero degree-\(n\) component in \(Q_0\). A nonzero class in
  \(\operatorname{range}(\psi_n)\), by contrast, is represented by a homogeneous
  element of degree \(n\) not in \(N'\), so has nonzero degree-\(n\) component.
  Comparing degree-\(n\) components forces the intersection
  \(\operatorname{range}(\psi_n) \cap \bigsqcup_{j \ne n} \operatorname{range}(\psi_j)
  = 0\); independence follows.

  \emph{Step 2: finiteness of the support.} Because \(Q_0\) is finite-dimensional and
  the ranges \(\operatorname{range}(\psi_n)\) are independent, only finitely many of
  them are nonzero: the set \(\{\,n : \operatorname{range}(\psi_n) \ne 0\,\}\) is
  finite (\ref{mk-lem:finite_ne_bot_of_iSupIndep_mathlib}). In particular it is bounded
  above, by some \(K \in \mathbb{N}\).

  \emph{Step 3: eventual vanishing.} For \(n > K\) we have
  \(\operatorname{range}(\psi_n) = 0\), i.e.\ every homogeneous element of
  \(N \cap \mathcal{M}_n\) has class \(0\) in \(Q_0 = N/N'\); equivalently
  \(N \cap \mathcal{M}_n \le N'\), hence \(N \cap \mathcal{M}_n = N' \cap \mathcal{M}_n\)
  and
  \[
    \mathrm{hilb}(n) \;=\; \dim_\kappa(N \cap \mathcal{M}_n)
      - \dim_\kappa(N' \cap \mathcal{M}_n) \;=\; 0 .
  \]
  Thus \(\mathrm{hilb}\) is eventually zero, and
  \ref{mk-lem:ratHilb_ofEventuallyZero} gives \(\mathrm{IsRatHilb}(\mathrm{hilb}, 0)\).
\end{proof}

\subsubsection*{The ambient subquotient induction}

\begin{lemma}[Rationality of the ambient subquotient Hilbert function]\label{mk-lem:graded_subquotient_isRatHilb}
  \dcref{00K1,custom}

  \uses{mk-def:graded_subquotientHilb, mk-lem:graded_subquotient_ker_coker,
    mk-def:graded_subquotientDatum_ker, mk-def:graded_subquotientDatum_coker,
    mk-lem:graded_subquotient_degreewise_diff,
    mk-lem:graded_subquotient_finite_transfer, mk-lem:ratHilb_ofEventuallyZero,
    mk-lem:ratHilb_ofDiffEq, mk-lem:ratHilb_bump,
    mk-lem:graded_subquotient_base_eventuallyZero,
    mk-lem:graded_finiteDimensional_of_mvPolynomial_isEmpty_finite}
  \textit{Source: \cite{stacks-project}, Tag 00K1 (induction on the number of
  degree-one generators).}
  Let \((N, N')\) be a subquotient datum of length \(r\)
  (\ref{mk-def:graded_subquotientHilb}) with ambient Hilbert function
  \(\mathrm{hilb}\). Then \(\mathrm{IsRatHilb}(\mathrm{hilb}, r)\)
  (\ref{mk-def:ratHilb}): the difference \(n \mapsto \dim_\kappa(N \cap \mathcal{M}_n)
  - \dim_\kappa(N' \cap \mathcal{M}_n)\) is, for \(n \gg 0\), the coefficient
  sequence of a rational series \(p \cdot (1 - X)^{-r}\) with \(p \in
  \mathbb{Q}[X]\). The induction is on the length \(r\), and every object that
  appears is an ambient pair of homogeneous submodules of the fixed graded module
  \(M\).
\end{lemma}
\begin{proof}
  \uses{mk-lem:graded_subquotient_ker_coker,
    mk-def:graded_subquotientDatum_ker, mk-def:graded_subquotientDatum_coker,
    mk-lem:graded_subquotient_degreewise_diff,
    mk-lem:graded_subquotient_finite_transfer, mk-lem:ratHilb_ofEventuallyZero,
    mk-lem:ratHilb_ofDiffEq, mk-lem:ratHilb_bump,
    mk-lem:graded_subquotient_base_eventuallyZero,
    mk-lem:graded_finiteDimensional_of_mvPolynomial_isEmpty_finite}
    We induct on the length \(r\).

  \emph{Base case \(r = 0\).} A length-zero datum has a finiteness witness making the
  subquotient \(N/N'\) a finite-dimensional \(\kappa\)-vector space
  (\ref{mk-lem:graded_finiteDimensional_of_mvPolynomial_isEmpty_finite}); its degreewise
  images form an independent family inside that finite-dimensional space, so only
  finitely many are nonzero and \(\mathrm{hilb}(n) = 0\) for \(n \gg 0\)
  (\ref{mk-lem:graded_subquotient_base_eventuallyZero}). By
  \ref{mk-lem:ratHilb_ofEventuallyZero} an eventually-zero function is rational of order
  \(0\), giving \(\mathrm{IsRatHilb}(\mathrm{hilb}, 0)\).

  \emph{Inductive step \(r \ge 1\).} Set \(x = t_{r-1}\). By the kernel/cokernel
  closure (\ref{mk-lem:graded_subquotient_ker_coker}) the kernel subquotient
  \(K = (N \cap x^{-1}(N'),\, N')\) and the cokernel subquotient
  \(C = (N,\, N' + x \cdot N)\) are subquotient data of length \(r-1\), assembled by
  the kernel and cokernel constructors (\ref{mk-def:graded_subquotientDatum_ker},
  \ref{mk-def:graded_subquotientDatum_coker}) whose finiteness witnesses come from the
  abstract finiteness transfer (\ref{mk-lem:graded_subquotient_finite_transfer}). The
  induction hypothesis applied
  to \(K\) and \(C\) gives \(\mathrm{IsRatHilb}(h_K, r-1)\) and
  \(\mathrm{IsRatHilb}(h_C, r-1)\) (\ref{mk-def:ratHilb}). By the degreewise
  difference identity (\ref{mk-lem:graded_subquotient_degreewise_diff}),
  \[
    \mathrm{hilb}(n+1) - \mathrm{hilb}(n) \;=\; h_C(n+1) - h_K(n)
    \qquad \text{for all } n .
  \]
  Feeding \(\mathrm{IsRatHilb}(h_C, r-1)\), \(\mathrm{IsRatHilb}(h_K, r-1)\) and this
  difference identity into the inductive-step engine
  (\ref{mk-lem:ratHilb_ofDiffEq}) yields \(\mathrm{IsRatHilb}(\mathrm{hilb}, r)\). (The
  common order \(r-1\) for \(h_C\) and \(h_K\) is arranged, if necessary, by raising
  the lower one with \ref{mk-lem:ratHilb_bump}.) This completes the induction.
\end{proof}

\section{Support and freeness predicates}
\label{mk-sec:quot_predicates}

The Quot functor (\ref{mk-def:quot_functor}) and the Grassmannian
(\ref{mk-def:grassmannian_scheme}) classify quotient sheaves subject to two
geometric conditions: a coherent quotient must have \emph{schematic support
proper over the base}, and the Grassmannian's quotients must be \emph{locally
free of a fixed rank}. This section introduces both predicates.

\begin{definition}[Annihilator ideal sheaf of a sheaf of modules]\label{mk-def:modules_annihilator}
  \dcref{custom}

  \textit{Source: \cite{nitsure-hilbert-quot}, \S 1 (FGA Explained Ch.~5).}
  Let \(X\) be a scheme and let \(\mathcal{F}\) be a sheaf of
  \(\mathcal{O}_X\)-modules that is \emph{quasi-coherent and of finite type}
  (over a locally noetherian base, so that on each affine open \(U\) the section
  module \(\mathcal{F}(U)\) is a finitely generated \(\mathcal{O}_X(U)\)-module).
  The \emph{annihilator ideal sheaf} of \(\mathcal{F}\) is the ideal sheaf
  \(\mathrm{Ann}(\mathcal{F})\) defined on affine opens: for each affine open
  \(U \subseteq X\),
  \[
    \mathrm{Ann}(\mathcal{F})(U)
    \;=\; \mathrm{Ann}_{\mathcal{O}_X(U)}\bigl(\mathcal{F}(U)\bigr)
    \;\subseteq\; \mathcal{O}_X(U),
  \]
  the annihilator of the \(\mathcal{O}_X(U)\)-module \(\mathcal{F}(U)\). These
  local data patch to a coherent ideal sheaf on \(X\) whose vanishing locus is
  the set-theoretic support of \(\mathcal{F}\).

  \emph{Well-definedness on basic opens.} The coherence condition for an ideal
  sheaf requires that, for an affine open \(U\) and \(f \in \mathcal{O}_X(U)\),
  the section \(\mathrm{Ann}(\mathcal{F})(U)\) restricts correctly to the basic
  open \(D(f) \subseteq U\), i.e. that
  \(\mathrm{Ann}_{\mathcal{O}_X(D(f))}(\mathcal{F}(D(f)))\) is the extension of
  \(\mathrm{Ann}_{\mathcal{O}_X(U)}(\mathcal{F}(U))\) along the localization map
  \(\mathcal{O}_X(U) \to \mathcal{O}_X(D(f))\). This is exactly where the two
  hypotheses enter. The bridge
  \ref{mk-lem:qcoh_section_localization_basicOpen} supplies, for a quasi-coherent
  finite-type \(\mathcal{F}\), that \(\mathcal{O}_X(D(f)) = (\mathcal{O}_X(U))_f\)
  and that the restriction \(\mathcal{F}(U) \to \mathcal{F}(D(f))\) exhibits
  \(\mathcal{F}(D(f))\) as the localized module \((\mathrm{powers}\,f)^{-1}
  \mathcal{F}(U)\); the algebra engine
  \ref{mk-lem:annihilator_localization_eq_map} then yields, using that
  \(\mathcal{F}(U)\) is finitely generated, the required identity
  \(\mathrm{Ann}_{(\mathcal{O}_X(U))_f}\bigl((\mathrm{powers}\,f)^{-1}
  \mathcal{F}(U)\bigr) = \mathrm{Ann}_{\mathcal{O}_X(U)}(\mathcal{F}(U)) \cdot
  (\mathcal{O}_X(U))_f\). Thus the local annihilators are compatible with basic-open
  restriction, which is the patching datum that makes \(\mathrm{Ann}(\mathcal{F})\)
  a well-defined ideal sheaf. The finite-type hypothesis is genuinely required:
  without finite generation of \(\mathcal{F}(U)\) the annihilator need not commute
  with localization, and the local data fail to patch.

  \emph{Note.} Mathlib does not carry an annihilator for sheaves of modules on a scheme.
\end{definition}

\begin{lemma}[Annihilator ideal bounded by the module annihilator on affine opens]\label{mk-lem:modules_annihilator_ideal_le}
  \dcref{custom}

  \uses{mk-def:modules_annihilator}
  Let \(X\) be a scheme, \(\mathcal{F}\) a sheaf of \(\mathcal{O}_X\)-modules and
  \(U \subseteq X\) an affine open. Then the section over \(U\) of the annihilator
  ideal sheaf \(\mathrm{Ann}(\mathcal{F})\) (\ref{mk-def:modules_annihilator}) is
  contained in the module-annihilator of the sections of \(\mathcal{F}\) over
  \(U\):
  \[
    \mathrm{Ann}(\mathcal{F})(U)
    \;\subseteq\; \mathrm{Ann}_{\mathcal{O}_X(U)}\bigl(\mathcal{F}(U)\bigr).
  \]
  This is the inclusion direction from the definition of
  \(\mathrm{Ann}(\mathcal{F})\); the reverse inclusion, under finiteness of the
  section modules, is \ref{mk-lem:modules_annihilator_ideal}. It supports the
  well-definedness of the schematic support
  (\ref{mk-def:schematic_support}).
\end{lemma}

\begin{lemma}[Annihilator ideal sheaf equals the module annihilator on affine opens]\label{mk-lem:modules_annihilator_ideal}
  \dcref{custom}

  \uses{mk-def:modules_annihilator}
  \textit{Source: Stacks, Algebra, Lemma \texttt{annihilator-flat-base-change}
  (localization instance);
  Stacks~\href{https://stacks.math.columbia.edu/tag/01IB}{01IB}.}
  Let \(X\) be a scheme and \(\mathcal{F}\) a quasi-coherent sheaf of
  \(\mathcal{O}_X\)-modules such that \(\Gamma(\mathcal{F}, V)\) is a finitely
  generated \(\Gamma(X, V)\)-module for every affine open \(V \subseteq X\).
  Then for every affine open \(U \subseteq X\),
  \[
    \mathrm{Ann}(\mathcal{F})(U)
    \;=\; \mathrm{Ann}_{\mathcal{O}_X(U)}\bigl(\mathcal{F}(U)\bigr).
  \]
\end{lemma}
\begin{proof}\uses{mk-lem:modules_annihilator_ideal_le, mk-lem:annihilator_localization_eq_map,
    mk-lem:qcoh_section_localization_basicOpen, mk-lem:isLocalization_basicOpen_mathlib}
  The family \(I(V) := \mathrm{Ann}_{\mathcal{O}_X(V)}(\mathcal{F}(V))\),
  indexed by the affine opens of \(X\), satisfies the basic-open coherence
  required of an ideal sheaf: for an affine open \(V\) and
  \(f \in \mathcal{O}_X(V)\), the restriction
  \(\mathcal{O}_X(V) \to \mathcal{O}_X(D(f))\) is the localization away from
  \(f\) (\ref{mk-lem:isLocalization_basicOpen_mathlib}), the section restriction
  \(\mathcal{F}(V) \to \mathcal{F}(D(f))\) is the corresponding module
  localization (\ref{mk-lem:qcoh_section_localization_basicOpen}), and — because
  \(\mathcal{F}(V)\) is finitely generated — the annihilator commutes with this
  localization (\ref{mk-lem:annihilator_localization_eq_map}):
  \(I(D(f)) = I(V) \cdot \mathcal{O}_X(D(f))\).  Hence \(I\) is itself an
  ideal sheaf.  Since \(\mathrm{Ann}(\mathcal{F})\) is by definition the
  largest ideal sheaf whose affine-open values are contained in the family
  \(I\) (\ref{mk-def:modules_annihilator}), and \(I\) is such an ideal sheaf,
  maximality gives \(I(U) \subseteq \mathrm{Ann}(\mathcal{F})(U)\) for every
  affine \(U\); the reverse inclusion is
  \ref{mk-lem:modules_annihilator_ideal_le}.
\end{proof}


\begin{lemma}[Sections on a basic open localize the affine ring]
  \label{mk-lem:isLocalization_basicOpen_mathlib}
  \dcref{custom}

  Let \(X\) be a scheme, \(U \subseteq X\) an affine open, and
  \(f \in \Gamma(X, U)\). Then the restriction map
  \(\Gamma(X, U) \to \Gamma\bigl(X,\, D(f)\bigr)\) on the basic open
  \(D(f) = X.\mathrm{basicOpen}\,f\) exhibits \(\Gamma(X, D(f))\) as the
  localization of \(\Gamma(X, U)\) away from \(f\); that is, it is
  \(\mathrm{IsLocalization.Away}\,f\), so
  \(\Gamma(X, D(f)) = (\Gamma(X, U))_f\) as a \(\Gamma(X,U)\)-algebra.
\end{lemma}

\begin{lemma}[Annihilator of a finitely generated module commutes with localization]
  \label{mk-lem:annihilator_localization_eq_map}
  \dcref{custom}

  \textit{Source: Stacks, Algebra, Lemma \texttt{annihilator-flat-base-change}
  (localization instance for finite modules).}
  Let \(R\) be a commutative ring, \(S \subseteq R\) a multiplicative submonoid, and
  \(R_S\) the localization of \(R\) at \(S\). Let \(M\) be a \emph{finitely
  generated} \(R\)-module and let \(g : M \to M_S\) exhibit \(M_S\) as the
  localization of \(M\) at \(S\) (an \(R_S\)-module whose \(R\)-action is the
  restriction of the \(R_S\)-action along \(R \to R_S\)). Then
  \[
    \mathrm{Ann}_{R_S}(M_S) \;=\; \mathrm{Ann}_R(M) \cdot R_S,
  \]
  the extension of the annihilator ideal along the canonical map
  \(R \to R_S\). Finite generation is essential: for a general module only the
  inclusion \(\supseteq\) holds.
\end{lemma}
\begin{proof}
  Fix a finite generating set \(m_1, \dots, m_n\) of \(M\). An element of \(R\)
  annihilating every \(m_i\) annihilates all of \(M\): the elements it kills form
  a submodule (they are closed under sums and scalar multiples), and that
  submodule contains the generators.

  (\(\supseteq\)) It suffices that the image of \(\mathrm{Ann}_R(M)\)
  annihilates \(M_S\). Every element of \(M_S\) has the form \(g(m)/s\) with
  \(m \in M\), \(s \in S\); for \(a \in \mathrm{Ann}_R(M)\),
  \((a/1)\cdot(g(m)/s) = g(a \cdot m)/s = 0\).

  (\(\subseteq\)) Let \(y \in \mathrm{Ann}_{R_S}(M_S)\) and write \(y = a/s\)
  with \(a \in R\), \(s \in S\). For each \(i\) the hypothesis gives
  \(y \cdot (g(m_i)/1) = g(a \cdot m_i)/s = 0\) in \(M_S\), so there is
  \(u_i \in S\) with \(u_i \cdot a \cdot m_i = 0\) in \(M\). Put
  \(u := u_1 \cdots u_n \in S\); since each \(u_i\) divides \(u\) in \(S\),
  \((u a) \cdot m_i = 0\) for every \(i\), hence
  \(u a \in \mathrm{Ann}_R(M)\) by the generating-set criterion. Finally
  \[
    \frac{a}{s} \;=\; \frac{ua}{1} \cdot \frac{1}{us}
  \]
  exhibits \(y\) as an \(R_S\)-multiple of the image of
  \(ua \in \mathrm{Ann}_R(M)\), so \(y \in \mathrm{Ann}_R(M) \cdot R_S\).
\end{proof}


\begin{lemma}[Sheaf Mayer--Vietoris: the pullback cone over a two-open cover is a limit]
  \label{mk-lem:isLimitPullbackCone_mathlib}
  \dcref{custom}

  For a sheaf \(F\) on a topological space and two opens \(U, V\), the square relating
  \(F(U \cup V)\), \(F(U)\), \(F(V)\), \(F(U \cap V)\) is a pullback (limit). This is the generic
  module Mayer--Vietoris used in the inductive step of the descent.
\end{lemma}


\begin{lemma}[Separatedness: a section is determined by its restrictions to a cover]
  \label{mk-lem:eq_of_locally_eq_mathlib}
  \dcref{custom}

  Let \(F\) be a sheaf on a topological space, \(V\) an open, and \(\{U_i\}\) a family of opens with
  \(\bigsqcup_i U_i = V\). If two sections \(s, t \in F(V)\) have equal restrictions
  \(s|_{U_i} = t|_{U_i}\) for every \(i\), then \(s = t\). This separatedness over a cover is the
  uniqueness half of the sheaf condition; it discharges the \(\mathrm{exists\_of\_eq}\) obligation
  (a global section vanishing on \(D(f)\) is killed by a power of \(f\)) and the final
  ``conclude on \(D(f)\)'' step of the cover-form descent
  (\ref{mk-lem:section_localization_descent_of_cover}).
\end{lemma}

\begin{lemma}[Pullback of sheaves of modules along a morphism of schemes]
  \label{mk-lem:modules_pullback_mathlib}
  \dcref{custom}

  For a morphism of schemes \(f : X \to Y\), Mathlib supplies the pullback functor
  \(f^{*} : Y.\mathrm{Modules} \to X.\mathrm{Modules}\) on sheaves of modules, left adjoint to
  pushforward, hence colimit-preserving, and carrying the unit sheaf to the unit sheaf. It is the
  geometric restriction
  through which a presentation is transported to an open subscheme.
\end{lemma}

\begin{lemma}[Restriction of sheaves of modules along an open immersion]
  \label{mk-lem:modules_restrictFunctor_mathlib}
  \dcref{custom}

  For an open immersion \(f : X \to Y\), Mathlib supplies the restriction functor
  \(f^{\mathrm{res}} : Y.\mathrm{Modules} \to X.\mathrm{Modules}\), with
  \(\Gamma(f^{\mathrm{res}} M, U) = \Gamma(M, f(U))\) definitionally; it is a left adjoint (colimit
  preserving) and carries the unit sheaf to the unit sheaf. This is the geometric
  ``restrict to open \(U\)'' functor on \((\Spec R).\mathrm{Modules}\).
\end{lemma}

\begin{lemma}[Restriction is isomorphic to pullback]
  \label{mk-lem:modules_restrictFunctorIsoPullback_mathlib}
  \dcref{custom}

  For an open immersion \(f\), the restriction functor (\ref{mk-lem:modules_restrictFunctor_mathlib})
  is naturally isomorphic to the pullback functor (\ref{mk-lem:modules_pullback_mathlib}):
  \(f^{\mathrm{res}} \cong f^{*}\). This identifies the restriction functor with the
  universal-property pullback, so a presentation transported along one transports along the other.
\end{lemma}

\begin{lemma}[Slice of opens equivalent to opens of the open subspace]
  \label{mk-lem:opens_overEquivalence_mathlib}
  \dcref{custom}

  For a topological space \(X\) and an open \(U \subseteq X\), the slice category of opens of
  \(X\) over \(U\) (the opens of \(X\) contained in \(U\), with inclusions) is equivalent to the
  opens of the subspace \(U\):
  \((\mathrm{Opens}\,X)_{/U} \simeq \mathrm{Opens}(U)\). The corresponding
  equivalence of sheaf categories is obtained by transport along this site equivalence.
\end{lemma}

\begin{lemma}[Sheaf categories of equivalent sites are equivalent]
  \label{mk-lem:equivalence_sheafCongr_mathlib}
  \dcref{custom}

  Let \(e : C \simeq D\) be an equivalence of categories carrying Grothendieck topologies \(J\)
  on \(C\) and \(K\) on \(D\), and suppose \(e.\mathrm{inverse}\) is a dense subsite for
  \((K, J)\). Then for any target category \(A\) the sheaf categories are equivalent:
  \(\mathrm{Sheaf}(J, A) \simeq \mathrm{Sheaf}(K, A)\).
\end{lemma}

\begin{lemma}[The over-equivalence functor is cocontinuous]\label{mk-lem:overEquivalence_functor_isCocontinuous}
  \dcref{custom}

  \uses{mk-lem:opens_overEquivalence_mathlib}
  The functor of \(\mathrm{Opens.overEquivalence}\,U\) (\ref{mk-lem:opens_overEquivalence_mathlib})
  is cocontinuous (cover-lifting) from the \(U\)-slice \(J_X|_U\) of the opens Grothendieck
  topology of \(X\) to the opens Grothendieck topology \(J_U\) of the subspace \(U\).
\end{lemma}
\begin{proof}
  \uses{mk-lem:opens_overEquivalence_mathlib}
  A sieve covers in the slice topology \(J_X|_U\) iff its image sieve covers in \(J_X\), and
  covering for the opens topology is pointwise: a sieve covers iff its members' opens cover the
  base set-theoretically. Given a point \(x\) of a slice object and a cover member in \(J_U\)
  containing the image of \(x\), transport it across the open embedding \(U \hookrightarrow X\),
  which is injective: the set-theoretic image of an open \(V \subseteq U\) is an open of \(X\)
  contained in the base, a neighbourhood of \(x\) lying in the image sieve. This witnesses the
  cover-lifting condition.
\end{proof}

\begin{lemma}[The over-equivalence inverse is cocontinuous]\label{mk-lem:overEquivalence_inverse_isCocontinuous}
  \dcref{custom}

  \uses{mk-lem:opens_overEquivalence_mathlib}
  Symmetrically, the inverse of \(\mathrm{Opens.overEquivalence}\,U\) is cocontinuous from the
  opens topology \(J_U\) of the subspace \(U\) to the \(U\)-slice \(J_X|_U\).
\end{lemma}
\begin{proof}
  \uses{mk-lem:opens_overEquivalence_mathlib}
  As in \ref{mk-lem:overEquivalence_functor_isCocontinuous}, reduce to the pointwise covering
  criterion and the over-sieve image criterion; given a point and a slice-cover member, pull the
  member back along the embedding \(U \hookrightarrow X\) (preimage of an open of \(X\) is an open
  of the subspace) to obtain a cover member of \(J_U\) witnessing the lift.
\end{proof}

\begin{lemma}[The over-equivalence inverse is a dense subsite]\label{mk-lem:overEquivalence_inverse_isDenseSubsite}
  \dcref{custom}

  \uses{mk-lem:overEquivalence_functor_isCocontinuous, mk-lem:overEquivalence_inverse_isCocontinuous}
  The inverse of \(\mathrm{Opens.overEquivalence}\,U\) is a dense subsite for \((J_U, J_X|_U)\).
  This is formal: an equivalence both of whose legs are cocontinuous
  (\ref{mk-lem:overEquivalence_functor_isCocontinuous},
  \ref{mk-lem:overEquivalence_inverse_isCocontinuous}) yields a dense subsite on its inverse.
\end{lemma}

\begin{lemma}[The over-equivalence functor is continuous]\label{mk-lem:overEquivalence_functor_isContinuous}
  \dcref{custom}

  \uses{mk-lem:overEquivalence_inverse_isCocontinuous}
  The functor of \(\mathrm{Opens.overEquivalence}\,U\) is continuous (cover-preserving) for
  \((J_X|_U, J_U)\). It is obtained from the cocontinuity of the inverse
  (\ref{mk-lem:overEquivalence_inverse_isCocontinuous}) through the equivalence adjunction: the
  left adjoint of a cocontinuous functor is continuous.
\end{lemma}

\begin{lemma}[The over-equivalence inverse is continuous]\label{mk-lem:overEquivalence_inverse_isContinuous}
  \dcref{custom}

  \uses{mk-lem:overEquivalence_functor_isCocontinuous}
  Symmetrically, the inverse of \(\mathrm{Opens.overEquivalence}\,U\) is continuous for
  \((J_U, J_X|_U)\), obtained from the cocontinuity of the functor
  (\ref{mk-lem:overEquivalence_functor_isCocontinuous}) through the equivalence adjunction.
\end{lemma}

\begin{definition}[The over-site/open-subspace sheaf equivalence]\label{mk-def:overEquivalence_sheafCongr}
  \dcref{custom}

  \uses{mk-lem:opens_overEquivalence_mathlib, mk-lem:equivalence_sheafCongr_mathlib,
    mk-lem:overEquivalence_inverse_isDenseSubsite}
  \textit{Keystone of the topological layer of bridge C.}
  For any category \(A\), the site equivalence \(\mathrm{Opens.overEquivalence}\,U\)
  (\ref{mk-lem:opens_overEquivalence_mathlib}) lifts to an equivalence of sheaf categories
  \[
    \mathrm{Sheaf}(J_X|_U,\, A) \;\simeq\; \mathrm{Sheaf}(J_U,\, A),
  \]
  obtained by feeding the dense-subsite witness
  (\ref{mk-lem:overEquivalence_inverse_isDenseSubsite}) to the general sheaf-congruence
  construction (\ref{mk-lem:equivalence_sheafCongr_mathlib}). This is exactly the equivalence left
  and it is the object across which \ref{mk-lem:over_restrict_iso} transports the
  structure sheaf and the module \(M\).
\end{definition}

\begin{definition}[The slice-to-geometric equivalence of module categories (gap1, C, step 3)]
  \label{mk-def:over_restrict_equiv}
  \dcref{custom}

  \uses{mk-lem:opens_overEquivalence_mathlib, mk-lem:pushforwardPushforwardEquivalence_mathlib,
    mk-lem:overEquivalence_functor_isContinuous, mk-lem:overEquivalence_inverse_isContinuous}
  Let \(X\) be a scheme and \(U \subseteq X\) an open. The category of sheaves of modules
  over the sliced structure sheaf \(\mathcal{O}_X.\mathrm{over}\,U\) on the slice site
  \(J_X|_U\) is equivalent to the category of sheaves of modules on the open subscheme
  \(U.\mathrm{toScheme}\):
  \[
    \mathrm{SheafOfModules}\bigl(\mathcal{O}_X.\mathrm{over}\,U\bigr)
    \;\simeq\; U.\mathrm{toScheme}.\mathrm{Modules}.
  \]
  The equivalence is obtained by lifting the site equivalence
  \(\mathrm{Opens.overEquivalence}\,U\) (\ref{mk-lem:opens_overEquivalence_mathlib}) — both of
  whose legs are continuous (\ref{mk-lem:overEquivalence_functor_isContinuous},
  \ref{mk-lem:overEquivalence_inverse_isContinuous}) — to module categories through the
  ringed-site transport of \ref{mk-lem:pushforwardPushforwardEquivalence_mathlib}. The two
  ring-sheaf comparison morphisms that transport consumes are the unit of the site
  equivalence whiskered into the sliced structure presheaf, and the identity: the structure
  sheaf of the open subscheme is precisely the transport of the sliced structure sheaf
  across the site equivalence.
\end{definition}

\begin{lemma}[The Grothendieck-slice restriction equals the geometric restriction (gap1, C)]\label{mk-lem:over_restrict_iso}
  \dcref{custom}

  \uses{mk-lem:modules_restrictFunctor_mathlib, mk-lem:modules_restrictFunctorIsoPullback_mathlib,
    mk-lem:isLocalization_basicOpen_mathlib, mk-def:overEquivalence_sheafCongr,
    mk-def:over_restrict_equiv,
    mk-lem:overEquivalence_functor_isContinuous, mk-lem:overEquivalence_inverse_isContinuous,
    mk-lem:opens_overEquivalence_mathlib, mk-lem:pushforwardPushforwardEquivalence_mathlib}
  \textit{The one lemma that must touch the Grothendieck slice site; everything downstream of it is
  geometric.}
  Let \(X\) be a scheme and \(U \subseteq X\) an open. The \emph{abstract Grothendieck-slice}
  restriction \(M.\mathrm{over}\,U\) — an object of the category of sheaves of modules over the
  sliced structure sheaf \(\mathcal{O}_X.\mathrm{over}\,U\) on the slice site \(J.\mathrm{over}\,U\) —
  is canonically isomorphic to the \emph{geometric} restriction
  \((\mathrm{restrictFunctor}\,U.\iota).\mathrm{obj}\,M\)
  (\ref{mk-lem:modules_restrictFunctor_mathlib}), equivalently to the pullback \(U.\iota^{*} M\)
  (\ref{mk-lem:modules_restrictFunctorIsoPullback_mathlib}), as a sheaf of modules on the small Zariski
  site of the open subscheme \(U.\mathrm{toScheme}\). The isomorphism is induced by the equivalence
  of sites \(J.\mathrm{over}\,U \simeq \mathrm{Opens}(U.\mathrm{toScheme})\) coming from the opens
  functor of the open immersion \(U.\iota\), under which the sliced sheaf of rings
  \(\mathcal{O}_X.\mathrm{over}\,U\) corresponds to the structure sheaf of \(U.\mathrm{toScheme}\).
  For \(X = \Spec R\) and \(U = D(r)\) a basic open, the target category is
  \((\Spec R_r).\mathrm{Modules}\) under the affine identification \(D(r) \cong \Spec R_r\)
  (\ref{mk-lem:isLocalization_basicOpen_mathlib}).
\end{lemma}
\begin{proof}
  \uses{mk-def:overEquivalence_sheafCongr, mk-def:over_restrict_equiv,
    mk-lem:overEquivalence_functor_isContinuous,
    mk-lem:overEquivalence_inverse_isContinuous, mk-lem:opens_overEquivalence_mathlib,
    mk-lem:pushforwardPushforwardEquivalence_mathlib, mk-lem:modules_restrictFunctor_mathlib,
    mk-lem:modules_restrictFunctorIsoPullback_mathlib, mk-lem:isLocalization_basicOpen_mathlib}
  The isomorphism is built in four steps, ascending from the purely topological layer to the
  module-level statement.

  \emph{Step 1 (topological layer --- done).} The opens functor of the open immersion
  \(U.\iota\) realizes the slice of the opens site of \(X\) over \(U\) as the opens site of the
  subspace \(U\); on sheaf categories this is the equivalence
  \(\mathrm{Sheaf}(J_X|_U,\,A) \simeq \mathrm{Sheaf}(J_U,\,A)\) of
  \ref{mk-def:overEquivalence_sheafCongr}, for every target category \(A\). This is the
  topos-theoretic layer, assembled from the cocontinuity/continuity of both legs of
  \(\mathrm{Opens.overEquivalence}\,U\) (\ref{mk-lem:opens_overEquivalence_mathlib}).

  \emph{Step 2 (geometric ring-sheaf identification --- the current obstacle).} Identify the
  sliced structure sheaf \(\mathcal{O}_X.\mathrm{over}\,U\) with the structure sheaf
  \(\mathcal{O}_{U}\) of the open subscheme \(U.\mathrm{toScheme}\), transported across the
  \(\mathrm{RingCat}\)-valued instance of \ref{mk-def:overEquivalence_sheafCongr}; this is an
  isomorphism of sheaves of rings. The bridge is the factorization of the opens functor of the
  immersion,
  \[
    U.\iota.\mathrm{opensFunctor}
      \;=\; (\mathrm{Opens.overEquivalence}\,U).\mathrm{inverse}
        \;\circ\; \mathrm{Over.forget}\,U ,
  \]
  under which the structure-sheaf comparison reduces to the standard open-immersion
  structure-sheaf identifications; recall that \(\mathrm{restrictFunctor}\,U.\iota\)
  (\ref{mk-lem:modules_restrictFunctor_mathlib}) is precisely pushforward along
  \(U.\iota.\mathrm{opensFunctor}\). This ring-sheaf identification is the remaining geometric
  obstacle.

  \emph{Step 3 (lift to modules).} Given the step-2 identification of the underlying ring sheaves
  and the continuity of both legs (\ref{mk-lem:overEquivalence_functor_isContinuous},
  \ref{mk-lem:overEquivalence_inverse_isContinuous}), lift the sheaf-of-rings comparison to an
  equivalence of categories of sheaves of modules via
  \(\mathrm{pushforwardPushforwardEquivalence}\)
  (\ref{mk-lem:pushforwardPushforwardEquivalence_mathlib}). Under this equivalence the sliced
  module \(M.\mathrm{over}\,U\) corresponds to the geometric restriction
  \((\mathrm{restrictFunctor}\,U.\iota).\mathrm{obj}\,M\).

  \emph{Step 4 (compose).} Compose with
  \(\mathrm{restrictFunctorIsoPullback}\) (\ref{mk-lem:modules_restrictFunctorIsoPullback_mathlib})
  to land the isomorphism \(M.\mathrm{over}\,U \cong M.\mathrm{restrict}\,U.\iota \cong
  U.\iota^{*}M\). The two sides live a priori in different module categories, so the literal
  statement of \(\mathrm{overRestrictIso}\) is phrased through the step-3 equivalence functor; for
  \(X = \Spec R\) and \(U = D(r)\) the target is \((\Spec R_r).\mathrm{Modules}\) under
  \(D(r) \cong \Spec R_r\) (\ref{mk-lem:isLocalization_basicOpen_mathlib}).
\end{proof}

\begin{lemma}[The slice-to-geometric isomorphism in pullback form (gap1, C, step 4')]\label{mk-lem:over_restrict_pullback_iso}
  \dcref{custom}

  \uses{mk-lem:over_restrict_iso, mk-def:over_restrict_equiv,
    mk-lem:modules_restrictFunctorIsoPullback_mathlib,
    mk-lem:modules_pullback_mathlib}
  \textit{The pullback packaging of the slice-touching bridge; the form consumed by the P1 transport.}
  With \(X\) and \(U \subseteq X\) as above and \(M\) a sheaf of modules on \(X\), the transport of the
  abstract Grothendieck-slice restriction \(M.\mathrm{over}\,U\) under the step-3 equivalence functor
  (\ref{mk-def:over_restrict_equiv}) is canonically isomorphic to the inverse-image (pullback) of \(M\)
  along the open immersion \(U.\iota\):
  \[
    (\mathrm{overRestrictEquiv}\,U).\mathrm{functor}.\mathrm{obj}\,(M.\mathrm{over}\,U)
      \;\cong\; (U.\iota^{*})\,M
  \]
  (\ref{mk-lem:modules_pullback_mathlib}). This is the pullback packaging of the slice-touching bridge
  \ref{mk-lem:over_restrict_iso}; it is exactly the form in which a presentation of
  \(M.\mathrm{over}\,U\) is transported into a presentation of the geometric pullback \(U.\iota^{*}M\),
  and it is the ingredient consumed by the per-element local-tilde transport
  \ref{mk-lem:isIso_fromTildeΓ_basicOpen_of_quasicoherent}.
\end{lemma}
\begin{proof}
  \uses{mk-lem:over_restrict_iso, mk-lem:modules_restrictFunctorIsoPullback_mathlib}
  Compose the slice-to-geometric isomorphism \ref{mk-lem:over_restrict_iso}, which identifies the
  transported slice restriction with the geometric restriction
  \((\mathrm{restrictFunctor}\,U.\iota).\mathrm{obj}\,M\), with the component at \(M\) of the natural
  isomorphism \(\mathrm{restrictFunctorIsoPullback}\,U.\iota\)
  (\ref{mk-lem:modules_restrictFunctorIsoPullback_mathlib}) identifying the geometric restriction with
  the pullback \(U.\iota^{*}M\).
\end{proof}

\begin{lemma}[Quasi-coherent modules restrict to a tilde on basic opens of the cover
(gap1, P1)]
  \label{mk-lem:isIso_fromTildeΓ_basicOpen_of_quasicoherent}
  \dcref{custom}

  \uses{mk-def:over_restrict_equiv, mk-lem:over_restrict_pullback_iso,
    mk-lem:modules_pullback_mathlib, mk-lem:tilde_gamma_adjunction_mathlib,
    mk-lem:isLocalization_basicOpen_mathlib}
  Let \(M\) be a sheaf of modules on \(\Spec R\) equipped with quasi-coherence data: an
  open cover \(\{X_i\}\) of \(\Spec R\) together with, for each \(i\), a presentation of
  the slice restriction \(M.\mathrm{over}\,X_i\) — an exhibition of it as the cokernel of
  a map between coproducts of copies of the unit module. Let \(r \in R\) with
  \(D(r) \subseteq X_i\) for some \(i\), and let \(W := \iota_{X_i}^{-1}(D(r))\) be the
  corresponding affine open of the subscheme \(Z := X_i\). Then the restriction of \(M\)
  to \(D(r)\) — the pullback of \(M\) to \(Z\), then to \(W\), transported across the
  affine identification \(W \cong \Spec\Gamma(Z, W)\) — has invertible
  tilde--\(\Gamma\) counit (\ref{mk-lem:tilde_gamma_adjunction_mathlib}); that is,
  \(M|_{D(r)}\) is the tilde of its section module over
  \(\Gamma(Z, W) \cong R_r\) (\ref{mk-lem:isLocalization_basicOpen_mathlib}).
\end{lemma}
\begin{proof}
  \uses{mk-def:over_restrict_equiv, mk-lem:over_restrict_pullback_iso,
    mk-lem:modules_pullback_mathlib, mk-lem:tilde_gamma_adjunction_mathlib}
  \emph{Step 1 (geometrize the datum).} The quasi-coherence datum supplies a presentation
  of the slice restriction \(M.\mathrm{over}\,X_i\). The slice-to-geometric equivalence
  (\ref{mk-def:over_restrict_equiv}) preserves colimits and carries the unit module to the
  unit module, so the image of this presentation, transported along the isomorphism of
  \ref{mk-lem:over_restrict_pullback_iso}, is a presentation of the geometric pullback
  \(N := \iota_{X_i}^{*} M\) on the scheme \(Z\).

  \emph{Step 2 (restrict to \(W\)).} Restriction to the slice site over \(W\) again
  preserves colimits and the unit, so the global presentation of \(N\) slices to a
  presentation of \(N.\mathrm{over}\,W\); geometrizing across the slice-to-geometric
  equivalence at \(W\) (\ref{mk-def:over_restrict_equiv},
  \ref{mk-lem:over_restrict_pullback_iso}) yields a presentation of the geometric
  restriction \(\iota_W^{*} N\) on the open subscheme \(W\).

  \emph{Step 3 (affine identification).} \(W\) is affine: \(D(r)\) is a basic open of the
  affine scheme \(\Spec R\), and its preimage under the open immersion \(\iota_{X_i}\) is
  again affine because \(D(r)\) is contained in the immersion's range. Pull the
  presentation back across the isomorphism \(W \cong \Spec\Gamma(Z, W)\): the pullback of
  sheaves of modules along an isomorphism of schemes
  (\ref{mk-lem:modules_pullback_mathlib}) preserves colimits, being a left adjoint, and
  carries the unit module to the unit module, since the induced functor on opens sites is
  an equivalence. This lands a global presentation of the transported module on the
  affine scheme \(\Spec\Gamma(Z, W)\).

  \emph{Step 4 (a presented module on \(\Spec S\) is a tilde).} Let \(P\) be a sheaf of
  modules on \(\Spec S\) with a presentation
  \(\mathcal{O}^{\oplus J} \to \mathcal{O}^{\oplus K} \to P \to 0\). The map of coproducts
  of units corresponds under the tilde--\(\Gamma\) adjunction
  (\ref{mk-lem:tilde_gamma_adjunction_mathlib}) to a map of free modules
  \(S^{\oplus J} \to S^{\oplus K}\), because \(\mathcal{O}^{\oplus K}\) is the tilde of
  \(S^{\oplus K}\) and the unit of the adjunction is an isomorphism. Since the tilde
  functor preserves cokernels (it is a left adjoint), \(P\) is isomorphic to the tilde of
  the cokernel of that module map. Finally, a module isomorphic to a tilde has invertible
  counit: on a tilde \(\widetilde{N}\) the triangle identity of the adjunction shows the
  counit is inverse to the tilde of the unit isomorphism \(N \cong
  \Gamma(\widetilde{N}, \top)\), and invertibility of the counit transports across any
  isomorphism of modules by naturality. Applying this to the transported module of
  Step~3 concludes.
\end{proof}

\begin{lemma}[Section-localization descent from per-cover localization data (gap1, D,
cover form)]
  \label{mk-lem:section_localization_descent_of_cover}
  \dcref{custom}

  \uses{mk-lem:eq_of_locally_eq_mathlib}
  Let \(M\) be a sheaf of modules on \(\Spec R\), \(f \in R\), and
  \(t \subseteq R\) a finite subset generating the unit ideal, so that the basic opens
  \(\{D(r)\}_{r \in t}\) cover \(\Spec R\). Assume the localization datum: for every open
  \(U\) contained in some \(D(r)\) with \(r \in t\), the section restriction
  \[
    \Gamma(M, U) \longrightarrow \Gamma\bigl(M, D(f) \cap U\bigr)
  \]
  exhibits the target as the localization of \(\Gamma(M, U)\) at the powers of \(f\),
  over \(R\). Then the global section restriction
  \(\Gamma(M, \top) \to \Gamma(M, D(f))\) exhibits \(\Gamma(M, D(f))\) as the
  localization \(\Gamma(M, \top)[1/f]\) over \(R\).
\end{lemma}
\begin{proof}
  \uses{mk-lem:eq_of_locally_eq_mathlib}
  We verify the three conditions characterizing the localization map.

  \emph{Units.} Every power of \(f\) acts invertibly on \(\Gamma(M, D(f))\): the action
  of \(f\) on sections over \(D(f)\) is through the restriction of \(f\) to
  \(\Gamma(\Spec R, D(f))\), which is a unit there (\(f\) is invertible on its own basic
  open), and multiplication by a unit of the acting ring is bijective. This holds for an
  arbitrary \(M\), without the covering hypothesis.

  \emph{Torsion.} Let \(x \in \Gamma(M, \top)\) restrict to \(0\) on \(D(f)\); we find
  \(n\) with \(f^n \cdot x = 0\) (applied to a difference \(x_1 - x_2\), this is the
  uniqueness condition). For each \(r \in t\), the restriction \(x|_{D(r)}\) restricts to
  \(0\) on \(D(f) \cap D(r)\) (functoriality of restriction through \(D(f)\)); the
  localization datum at \(U = D(r)\) then provides \(k_r\) with
  \(f^{k_r} \cdot x|_{D(r)} = 0\). Take \(n := \max_{r \in t} k_r\); then
  \((f^n \cdot x)|_{D(r)} = f^{\,n-k_r} \cdot \bigl(f^{k_r} \cdot x|_{D(r)}\bigr) = 0\)
  for every \(r\), and since the \(D(r)\) cover \(\Spec R\), sheaf separatedness
  (\ref{mk-lem:eq_of_locally_eq_mathlib}) gives \(f^n \cdot x = 0\).

  \emph{Surjectivity.} Let \(y \in \Gamma(M, D(f))\); we find a global \(x\) and \(n\)
  with \(f^n \cdot y = x|_{D(f)}\). For each \(r \in t\), the localization datum at
  \(U = D(r)\) writes \(f^{m_r} \cdot y|_{D(f) \cap D(r)}\) as the restriction of a
  section \(b_r \in \Gamma(M, D(r))\); normalizing with the common exponent
  \(N := \max_r m_r\) (replace \(b_r\) by \(f^{\,N-m_r} \cdot b_r\)) we may assume
  \(b_r|_{D(f) \cap D(r)} = f^{N} \cdot y|_{D(f) \cap D(r)}\) for every \(r\). On a
  pairwise overlap \(D(r) \cap D(r')\) — an open contained in \(D(r)\) — the difference
  \(b_r - b_{r'}\) restricts to \(0\) on \(D(f) \cap D(r) \cap D(r')\), since both terms
  restrict to \(f^N \cdot y\) there; the localization datum at
  \(U = D(r) \cap D(r')\) provides a power \(f^{P_{r,r'}}\) killing the difference, and
  taking \(P := \max P_{r,r'}\) over the finitely many pairs, the family
  \(\{f^P \cdot b_r\}_{r \in t}\) agrees on all overlaps. By the sheaf axiom (unique
  gluing over the finite cover \(\{D(r)\}_{r \in t}\) of \(\Spec R\)) the family glues to
  a global section \(x \in \Gamma(M, \top)\). Finally
  \(x|_{D(f)} = f^{N+P} \cdot y\): the two sides agree on each member of the cover
  \(\{D(f) \cap D(r)\}_{r \in t}\) of \(D(f)\), so they are equal by separatedness
  (\ref{mk-lem:eq_of_locally_eq_mathlib}).
\end{proof}

\begin{lemma}[Section-localization descent, basic-open hypothesis form (gap1, D)]
  \label{mk-lem:section_localization_descent_of_basicOpen_cover}
  \dcref{custom}

  \uses{mk-lem:section_localization_descent_of_cover, mk-lem:eq_of_locally_eq_mathlib}
  Let \(M\), \(f\), and \(t\) be as in \ref{mk-lem:section_localization_descent_of_cover},
  but assume the localization datum only for \emph{basic} opens: for every \(s \in R\)
  with \(D(s) \subseteq D(r)\) for some \(r \in t\), the restriction
  \(\Gamma(M, D(s)) \to \Gamma(M, D(f) \cap D(s))\) is the localization of
  \(\Gamma(M, D(s))\) at the powers of \(f\) over \(R\). Then the conclusion of
  \ref{mk-lem:section_localization_descent_of_cover} still holds: the global restriction
  \(\Gamma(M, \top) \to \Gamma(M, D(f))\) is the localization at the powers of \(f\)
  over \(R\).
\end{lemma}
\begin{proof}
  \uses{mk-lem:section_localization_descent_of_cover, mk-lem:eq_of_locally_eq_mathlib}
  The proof of \ref{mk-lem:section_localization_descent_of_cover} consults the
  localization datum only at the cover members \(U = D(r)\) (\(r \in t\)) and at their
  pairwise intersections \(U = D(r) \cap D(r') = D(r r')\) — all of them \emph{basic}
  opens contained in a cover member. The datum restricted to basic opens therefore
  suffices, and the same three-part verification applies verbatim. (The restriction of
  the hypothesis matters: the geometric producer below furnishes the localization datum
  only at basic opens, a general open of \(\Spec R\) being neither affine nor
  quasi-compact in general.)
\end{proof}


\subsection{Section-transport producer for the basic-open Hfr}

This subsection builds the geometric input of the gap1 keystone descent
\ref{mk-lem:section_localization_descent}: the geometric \emph{producer}
\ref{mk-lem:section_localization_hfr_basicOpen} that manufactures, for a quasi-coherent \(M\) on
\(\Spec R\) and a basic open \(D(s)\) inside a quasi-coherence cover member, the per-cover-element
localization datum \(\mathrm{Hfr}\) the basic-open cover form
\ref{mk-lem:section_localization_descent_of_basicOpen_cover} consumes. The construction transports the
\(R_s\)-level localization living on the affine slice \(D(s) \cong \Spec R_s\) back to an
\(R\)-level statement through the combiner \ref{mk-lem:isLocalizedModule_powers_transport}; the
geometric inputs --- the iterated-pullback tilde--\(\Gamma\) isomorphism, the identification of the
immersion's image opens, and the semilinear section comparisons --- enter in the proof of the
producer.

\begin{lemma}[Localization transport across a semilinear section comparison]
  \label{mk-lem:isLocalizedModule_powers_transport}
  \dcref{custom}

  Let \(R\) be a commutative ring, \(A\) an \(R\)-algebra, \(S\) a commutative ring, and
  \(\sigma : S \xrightarrow{\ \sim\ } A\) a ring isomorphism; let \(f \in R\) and
  \(f' \in S\) with \(\sigma(f')\) equal to the image of \(f\) in \(A\). Let
  \(g : M_1 \to M_2\) be an \(S\)-linear map exhibiting \(M_2\) as the localization of
  \(M_1\) at the powers of \(f'\); let \(N_1, N_2\) be \(A\)-modules, regarded also as
  \(R\)-modules through \(R \to A\); let \(e_1 : M_1 \to N_1\) and \(e_2 : M_2 \to N_2\)
  be additive isomorphisms that are \(\sigma\)-semilinear, i.e.
  \(e_i(a \cdot x) = \sigma(a) \cdot e_i(x)\) for \(a \in S\); and let
  \(h : N_1 \to N_2\) be an \(A\)-linear map intertwining \(g\), i.e.
  \(h \circ e_1 = e_2 \circ g\). Then \(h\), regarded as an \(R\)-linear map, exhibits
  \(N_2\) as the localization of \(N_1\) at the powers of \(f\).
\end{lemma}
\begin{proof}
  We check the three localization conditions in two stages.

  \emph{Stage I: crossing \(\sigma\).} The map \(h\) is a localization at the powers of
  \(\sigma(f')\) over \(A\).
  \emph{Units:} for \(k \ge 0\), the semilinearity of \(e_2\) gives that the action of
  \(\sigma(f')^k = \sigma(f'^k)\) on \(N_2\) equals
  \(e_2 \circ (\text{action of } f'^k \text{ on } M_2) \circ e_2^{-1}\), a composite of
  bijections because \(f'^k\) acts invertibly on the localization \(M_2\).
  \emph{Surjectivity:} given \(y \in N_2\), choose \(x \in M_1\) and \(k\) with
  \(f'^k \cdot e_2^{-1}(y) = g(x)\); applying \(e_2\) and semilinearity,
  \(\sigma(f')^k \cdot y = e_2(g(x)) = h(e_1(x))\).
  \emph{Uniqueness:} if \(h(y_1) = h(y_2)\) then, by the intertwining relation,
  \(e_2\bigl(g(e_1^{-1} y_1)\bigr) = e_2\bigl(g(e_1^{-1} y_2)\bigr)\), so
  \(g(e_1^{-1} y_1) = g(e_1^{-1} y_2)\); hence
  \(f'^k \cdot e_1^{-1} y_1 = f'^k \cdot e_1^{-1} y_2\) for some \(k\), and applying
  \(e_1\), \(\sigma(f')^k \cdot y_1 = \sigma(f')^k \cdot y_2\).

  \emph{Stage II: descending to \(R\).} By hypothesis \(\sigma(f')\) is the image of
  \(f\) under \(R \to A\), and by compatibility of the two module structures the action
  of \(f^k\) through \(R\) coincides with the action of the image of \(f^k\) through
  \(A\), on \(N_1\) and on \(N_2\). Each of the three conditions at the powers of \(f\)
  over \(R\) is therefore literally the corresponding Stage-I condition at the powers of
  \(\sigma(f')\) over \(A\): the units condition transfers because the two actions
  coincide; surjectivity reads \(f^k \cdot y = \sigma(f')^k \cdot y = h(x)\); and
  uniqueness likewise.
\end{proof}

\begin{lemma}[Section-transport producer: basic-open localization inside a cover member
(gap1, Hfr)]
  \label{mk-lem:section_localization_hfr_basicOpen}
  \dcref{custom}

  \uses{mk-lem:isIso_fromTildeΓ_basicOpen_of_quasicoherent,
    mk-lem:isLocalizedModule_restrict_of_isIso_fromTildeΓ,
    mk-lem:isLocalizedModule_powers_transport,
    mk-lem:modules_restrictFunctorIsoPullback_mathlib, mk-lem:modules_pullback_mathlib}
  Let \(M\) be a sheaf of modules on \(\Spec R\) with quasi-coherence data (an open cover
  \(\{X_i\}\) with per-member presentations, as in
  \ref{mk-lem:isIso_fromTildeΓ_basicOpen_of_quasicoherent}), let \(f, s \in R\), and
  suppose \(D(s) \subseteq X_i\) for some \(i\). Then the section restriction
  \[
    \Gamma\bigl(M, D(s)\bigr) \longrightarrow \Gamma\bigl(M, D(f) \cap D(s)\bigr)
  \]
  exhibits the target as the localization of \(\Gamma(M, D(s))\) at the powers of \(f\),
  over \(R\).
\end{lemma}
\begin{proof}
  \uses{mk-lem:isIso_fromTildeΓ_basicOpen_of_quasicoherent,
    mk-lem:isLocalizedModule_restrict_of_isIso_fromTildeΓ,
    mk-lem:isLocalizedModule_powers_transport,
    mk-lem:modules_restrictFunctorIsoPullback_mathlib, mk-lem:modules_pullback_mathlib}
  Let \(W := \iota_{X_i}^{-1}(D(s))\) be the affine open of the subscheme \(Z := X_i\)
  corresponding to \(D(s)\), put \(S := \Gamma(Z, W)\), and let
  \(j : \Spec S \to \Spec R\) be the composite open immersion
  \(\Spec S \cong W \hookrightarrow Z \hookrightarrow \Spec R\) (the affine
  identification of \(W\) followed by the two canonical immersions). Its image is
  \(D(s)\): the first leg is an isomorphism onto \(W\), and
  \(\iota_{X_i}(W) = X_i \cap D(s) = D(s)\).

  Write \(M' := j^{*} M\) (\ref{mk-lem:modules_pullback_mathlib}). The per-element
  transport \ref{mk-lem:isIso_fromTildeΓ_basicOpen_of_quasicoherent} gives invertibility of
  the tilde--\(\Gamma\) counit for the iterated pullback of \(M\) along the three legs of
  \(j\); pullback along a composite agrees with the composite of the pullbacks up to
  canonical natural isomorphism, and invertibility of the counit is invariant under
  isomorphism of modules (the essential image of the tilde functor is closed under
  isomorphism). Hence the counit of \(M'\) is invertible, and by
  \ref{mk-lem:isLocalizedModule_restrict_of_isIso_fromTildeΓ}, for every \(f' \in S\) the
  restriction \(\Gamma(M', \top) \to \Gamma(M', D(f'))\) is the localization at the
  powers of \(f'\) over \(S\).

  Let \(\sigma : S \xrightarrow{\ \sim\ } \Gamma(\Spec R, D(s))\) be the ring comparison
  of the open immersion \(j\) on top sections
  (\(S \cong \Gamma(\Spec S, \top) \cong \Gamma(\Spec R, j(\top))\), and
  \(j(\top) = D(s)\)), and choose \(f' := \sigma^{-1}\bigl(f|_{D(s)}\bigr)\). The image
  \(j\bigl(D(f')\bigr)\) is the locus of \(D(s)\) where \(\sigma(f') = f|_{D(s)}\) is
  invertible, that is, \(j\bigl(D(f')\bigr) = D(f) \cap D(s)\).

  Since \(j\) is an open immersion, applying the section functor to the identification of
  the pullback with the restriction (\ref{mk-lem:modules_restrictFunctorIsoPullback_mathlib})
  gives additive isomorphisms
  \(e_1 : \Gamma(M', \top) \cong \Gamma\bigl(M, D(s)\bigr)\) and
  \(e_2 : \Gamma(M', D(f')) \cong \Gamma\bigl(M, D(f) \cap D(s)\bigr)\) which intertwine
  the two restriction maps (naturality of the identification in the open) and are
  \(\sigma\)-semilinear over the structure-sheaf comparisons: the action on the sections
  of the pullback is through the immersion's section-ring isomorphism.

  Finally apply the combiner \ref{mk-lem:isLocalizedModule_powers_transport} with
  \(A := \Gamma(\Spec R, D(s))\) — an \(R\)-algebra through restriction of global
  sections — the isomorphism \(\sigma\), the elements \(f\) and \(f'\) (so that
  \(\sigma(f')\) is the image of \(f\) in \(A\)), the localization
  \(g := \bigl(\Gamma(M', \top) \to \Gamma(M', D(f'))\bigr)\), the comparisons
  \(e_1, e_2\), and \(h := \bigl(\Gamma(M, D(s)) \to \Gamma(M, D(f) \cap D(s))\bigr)\):
  the conclusion is that \(h\) is the localization at the powers of \(f\) over \(R\).
\end{proof}

\begin{lemma}[Section-localization descent for quasi-coherent modules (gap1 keystone)]
  \label{mk-lem:section_localization_descent}
  \dcref{custom}

  \uses{mk-lem:section_localization_descent_of_basicOpen_cover,
    mk-lem:section_localization_hfr_basicOpen}
  For a quasi-coherent sheaf of modules \(M\) on \(\Spec R\) and \(f \in R\), the section
  restriction \(\Gamma(M, \top) \to \Gamma(M, D(f))\) exhibits the target as the
  localization \(\Gamma(M, \top)[1/f]\) over \(R\). This is the statement of
  \ref{mk-lem:qcoh_section_localization_basicOpen} at the affine open \(\top\) of
  \(\Spec R\), proved here by finite-cover descent --- not through the affine
  equivalence of quasi-coherent modules with modules, which is instead \emph{derived}
  from this lemma in \ref{mk-lem:qcoh_affine_isIso_fromTildeΓ}.
\end{lemma}
\begin{proof}
  \uses{mk-lem:section_localization_descent_of_basicOpen_cover,
    mk-lem:section_localization_hfr_basicOpen}
  Quasi-coherence provides a cover \(\{X_i\}\) of \(\Spec R\) with per-member
  presentations. Refine it to a finite basic-open cover: the basic opens form a basis of
  \(\Spec R\), so every point has a basic-open neighbourhood contained in some \(X_i\);
  quasi-compactness of \(\Spec R\) extracts a finite family \(t \subseteq R\) with each
  \(D(r)\), \(r \in t\), contained in some \(X_i\), and covering is equivalent to \(t\)
  generating the unit ideal. For every basic open \(D(u)\) contained in some \(D(r)\)
  with \(r \in t\) — hence contained in a cover member \(X_i\) — the producer
  \ref{mk-lem:section_localization_hfr_basicOpen} supplies the localization datum at
  \(D(u)\). The basic-open cover form
  \ref{mk-lem:section_localization_descent_of_basicOpen_cover} then yields the conclusion.
\end{proof}


\subsection{G1-assemble: from per-basic-open section localization to the
\texorpdfstring{$\widetilde{(-)}$}{tilde}-counit}

This subsection records the bridge \emph{section localization on every basic open}
\(\Rightarrow\) \emph{invertible tilde--\(\Gamma\) counit}: given that the section
restriction localizes on every basic open, the tilde--\(\Gamma\) counit is an
isomorphism. Together with its converse
\ref{mk-lem:isLocalizedModule_restrict_of_isIso_fromTildeΓ} it assembles the equivalence
\ref{mk-lem:isIso_fromTildeΓ_iff_isLocalizedModule_restrict}, which is exactly what makes G1-core
(\ref{mk-lem:qcoh_affine_section_localization}) and gap1
(\ref{mk-lem:qcoh_affine_isIso_fromTildeΓ}) interderivable -- the route by which G1-core is obtained as
a corollary of gap1. The supporting Mathlib facts enter through the tilde--\(\Gamma\)
adjunction anchor (\ref{mk-lem:tilde_gamma_adjunction_mathlib}).

\begin{lemma}[Basic-open section localization for a module with invertible
tilde--\(\Gamma\) counit]
  \label{mk-lem:isLocalizedModule_restrict_of_isIso_fromTildeΓ}
  \dcref{custom}

  \uses{mk-lem:tilde_gamma_adjunction_mathlib}
  \textit{Source: Stacks, Schemes, Lemma \texttt{widetilde-pullback} (sections of a
  tilde over a basic open are the localized module).}
  Let \(M\) be a sheaf of modules on \(\Spec R\) whose tilde--\(\Gamma\) counit
  \(\widetilde{\Gamma(M, \top)} \to M\) (\ref{mk-lem:tilde_gamma_adjunction_mathlib}) is an
  isomorphism. Then for every \(f \in R\) the section restriction
  \(\Gamma(M, \top) \to \Gamma(M, D(f))\) exhibits the target as the localization of
  \(\Gamma(M, \top)\) at the powers of \(f\), over \(R\).
\end{lemma}
\begin{proof}
  \uses{mk-lem:tilde_gamma_adjunction_mathlib}
  First let \(M = \widetilde{N}\) be a tilde. The unit
  \(N \to \Gamma(\widetilde{N}, \top)\) of the adjunction
  (\ref{mk-lem:tilde_gamma_adjunction_mathlib}) is an isomorphism, and the canonical map
  \(N \to \Gamma(\widetilde{N}, D(f))\) exhibits the target as the localized module
  \(N_f\) — the sections of a tilde over a basic open are the localization. The
  canonical map factors as the unit followed by the section restriction, so the
  restriction \(\Gamma(\widetilde{N}, \top) \to \Gamma(\widetilde{N}, D(f))\) is the
  composite of the inverse unit isomorphism with a localization map, hence itself a
  localization at the powers of \(f\).

  For general \(M\), put \(N := \Gamma(M, \top)\). The counit isomorphism
  \(\widetilde{N} \cong M\) induces, naturally in the open, \(R\)-linear isomorphisms
  \(\Gamma(\widetilde{N}, \top) \cong \Gamma(M, \top)\) and
  \(\Gamma(\widetilde{N}, D(f)) \cong \Gamma(M, D(f))\) intertwining the two restriction
  maps. Pre- and post-composing a localization map with \(R\)-linear isomorphisms
  compatible in this way preserves the localization property.
\end{proof}

\begin{lemma}[Characterization of the invertible tilde--\(\Gamma\) counit by
section localization]
  \label{mk-lem:isIso_fromTildeΓ_iff_isLocalizedModule_restrict}
  \dcref{custom}

  \uses{mk-lem:isLocalizedModule_restrict_of_isIso_fromTildeΓ,
    mk-lem:tilde_gamma_adjunction_mathlib}
  For a sheaf of modules \(M\) on \(\Spec R\), the tilde--\(\Gamma\) counit
  \(\widetilde{\Gamma(M, \top)} \to M\) is an isomorphism if and only if for every
  \(f \in R\) the section restriction \(\Gamma(M, \top) \to \Gamma(M, D(f))\) exhibits
  the target as the localization of \(\Gamma(M, \top)\) at the powers of \(f\) over
  \(R\).
\end{lemma}
\begin{proof}
  \uses{mk-lem:isLocalizedModule_restrict_of_isIso_fromTildeΓ,
    mk-lem:tilde_gamma_adjunction_mathlib}
  The forward direction is \ref{mk-lem:isLocalizedModule_restrict_of_isIso_fromTildeΓ}.

  For the converse, write \(N := \Gamma(M, \top)\) and let
  \(\varepsilon : \widetilde{N} \to M\) be the counit. Fix \(f \in R\). The
  \(D(f)\)-component of \(\varepsilon\) is a map between two localizations of \(N\) at
  the powers of \(f\): the source \(\Gamma(\widetilde{N}, D(f))\) receives the canonical
  map from \(N\) (sections of a tilde over a basic open, as in the forward direction),
  and the target \(\Gamma(M, D(f))\) receives the restriction
  \(N = \Gamma(M, \top) \to \Gamma(M, D(f))\), a localization by hypothesis. The
  component intertwines these two maps out of \(N\): the counit composed with the
  canonical section map is the section restriction, by the compatibility of the counit
  with the adjunction unit (\ref{mk-lem:tilde_gamma_adjunction_mathlib}) and naturality.
  A map between two localizations of the same module at the same multiplicative set
  commuting with the two localization maps is the canonical isomorphism between them,
  hence bijective. Thus \(\varepsilon\) is an isomorphism on every basic open.

  Since the basic opens form a basis of \(\Spec R\), \(\varepsilon\) is an isomorphism
  on every stalk: injective, because a germ in the kernel is represented by a section
  over some basic neighbourhood killed by \(\varepsilon\) there, hence zero; surjective,
  because every germ of \(M\) is the germ of a section over some basic open, which is
  hit by the corresponding section of \(\widetilde{N}\). A morphism of sheaves that is
  an isomorphism on all stalks is an isomorphism, and a morphism of sheaves of
  modules whose underlying morphism of sheaves is an isomorphism is itself an
  isomorphism.
\end{proof}

\begin{lemma}[A quasi-coherent module on an affine scheme is the tilde of its sections
(gap1, Stacks 01I8)]
  \label{mk-lem:qcoh_affine_isIso_fromTildeΓ}
  \dcref{custom}

  \uses{mk-lem:section_localization_descent,
    mk-lem:isIso_fromTildeΓ_iff_isLocalizedModule_restrict}
  \textit{Source: Stacks~\href{https://stacks.math.columbia.edu/tag/01I8}{01I8}
  (\texttt{lemma-quasi-coherent-affine}).}
  For a quasi-coherent sheaf of modules \(M\) on \(\Spec R\), the tilde--\(\Gamma\)
  counit \(\widetilde{\Gamma(M, \top)} \to M\) is an isomorphism; equivalently, \(M\)
  lies in the essential image of the tilde functor, i.e. \(M\) is isomorphic to the
  tilde of an \(R\)-module.
\end{lemma}
\begin{proof}
  \uses{mk-lem:section_localization_descent,
    mk-lem:isIso_fromTildeΓ_iff_isLocalizedModule_restrict}
  By the keystone descent \ref{mk-lem:section_localization_descent}, every basic-open
  section restriction \(\Gamma(M, \top) \to \Gamma(M, D(f))\) is the localization at the
  powers of \(f\). The reverse implication of
  \ref{mk-lem:isIso_fromTildeΓ_iff_isLocalizedModule_restrict} concludes.
\end{proof}

\begin{lemma}[Affine section localization for quasi-coherent modules (G1-core)]
  \label{mk-lem:qcoh_affine_section_localization}
  \dcref{custom}

  \uses{mk-lem:qcoh_affine_isIso_fromTildeΓ,
    mk-lem:isLocalizedModule_restrict_of_isIso_fromTildeΓ}
  For a quasi-coherent sheaf of modules \(M\) on \(\Spec R\) and \(f \in R\), the
  section restriction \(\Gamma(M, \top) \to \Gamma(M, D(f))\) exhibits the target as the
  localization \(\Gamma(M, \top)[1/f]\) over \(R\). This is the statement of
  \ref{mk-lem:section_localization_descent} in its clean named form — obtained here as a
  corollary of gap1 rather than by re-running the descent — and the affine instance
  (\(X = \Spec R\), \(U = \top\)) of the general-scheme keystone
  \ref{mk-lem:qcoh_section_localization_basicOpen}.
\end{lemma}
\begin{proof}
  \uses{mk-lem:qcoh_affine_isIso_fromTildeΓ,
    mk-lem:isLocalizedModule_restrict_of_isIso_fromTildeΓ}
  Gap1 (\ref{mk-lem:qcoh_affine_isIso_fromTildeΓ}) makes the tilde--\(\Gamma\) counit of
  \(M\) an isomorphism, and
  \ref{mk-lem:isLocalizedModule_restrict_of_isIso_fromTildeΓ} applies.
\end{proof}


\begin{definition}[Schematic support of a sheaf of modules]\label{mk-def:schematic_support}
  \dcref{custom}

  \uses{mk-def:modules_annihilator, mk-lem:modules_annihilator_ideal_le}
  \textit{Source: \cite{nitsure-hilbert-quot}, \S 1 (FGA Explained Ch.~5).}
  Let \(\mathcal{F}\) be a sheaf of modules on a scheme \(X\). The
  \emph{schematic support} of \(\mathcal{F}\) is the closed subscheme of \(X\)
  cut out by its annihilator ideal sheaf (\ref{mk-def:modules_annihilator}):
  \[
    \mathrm{schematicSupport}(\mathcal{F})
    \;=\; V\bigl(\mathrm{Ann}(\mathcal{F})\bigr).
  \]
  This realises Nitsure's ``schematic support of \(\mathcal{F}\)'' as a closed
  subscheme of \(X\).
\end{definition}

\begin{definition}[Schematic-support immersion]\label{mk-def:schematic_support_immersion}
  \dcref{custom}

  \uses{mk-def:schematic_support, mk-def:modules_annihilator}
  Let \(\mathcal{F}\) be a sheaf of modules on a scheme \(X\). The
  \emph{schematic-support immersion} is the canonical closed immersion
  \[
    \mathrm{schematicSupport}(\mathcal{F}) \;\hookrightarrow\; X
  \]
  of the schematic support (\ref{mk-def:schematic_support}) into the ambient
  scheme, namely the subscheme immersion of the annihilator ideal sheaf
  \(\mathrm{Ann}(\mathcal{F})\) (\ref{mk-def:modules_annihilator}). It is a
  quasi-compact preimmersion onto the support, and it is the morphism whose
  composite with a base morphism \(f : X \to S\) defines proper support
  (\ref{mk-def:has_proper_support}).
\end{definition}

\begin{lemma}[Properness as a morphism property]
  \label{mk-lem:isProper_mathlib}
  \dcref{custom}

  A morphism \(f : X \to S\) of schemes is \emph{proper} iff it is separated,
  universally closed, and locally of finite type. Properness is stable under base
  change: if \(f : X \to S\) is proper and \(g : T \to S\) is any morphism then
  the base change \(X \times_S T \to T\) is proper. Consequently the base-change
  clause that Nitsure invokes (``properness \ldots\ preserved by base-change'')
  applies directly to the proper-support condition below.  Moreover properness
  is stable under composition; closed immersions are proper (they are finite
  morphisms); and cancellation holds along separated morphisms: if a composite
  \(g \circ f\) is proper and \(g\) is separated, then \(f\) is proper.
\end{lemma}

\begin{definition}[Proper support over the base]\label{mk-def:has_proper_support}
  \dcref{custom}

  \uses{mk-def:schematic_support, mk-def:schematic_support_immersion, mk-lem:isProper_mathlib}
  \textit{Source: \cite{nitsure-hilbert-quot}, \S 1 (FGA Explained Ch.~5).}
  Let \(\mathcal{F}\) be a sheaf of modules on a scheme \(X\) and let
  \(f : X \to S\) be a morphism. We say \(\mathcal{F}\) \emph{has proper support
  over \(S\) (along \(f\))} when the composite morphism
  \[
    \mathrm{schematicSupport}(\mathcal{F}) \hookrightarrow X \xrightarrow{f} S
  \]
  is proper (\ref{mk-lem:isProper_mathlib}). Since properness is stable under base
  change (\ref{mk-lem:isProper_mathlib}), this condition is preserved by pull-back
  along any \(S'\to S\), which is exactly what the well-definedness of the Quot
  functor's pullback action requires.
\end{definition}

\begin{definition}[Locally free of rank \(d\)]\label{mk-def:is_locally_free_of_rank}
  \dcref{custom}

  \textit{Source: \cite{nitsure-hilbert-quot}, \S 1, Exercise (2) (FGA Explained Ch.~5).}
  Let \(\mathcal{M}\) be a sheaf of modules on a scheme \(X\) and let
  \(d \in \mathbb{N}\). We say \(\mathcal{M}\) is \emph{locally free of rank
  \(d\)} when there exists an open cover \(\{U_i\}\) of \(X\) together with
  isomorphisms \(\mathcal{M}|_{U_i} \cong \mathcal{O}_{U_i}^{\oplus d}\) for
  each \(i\).

  \emph{Note.} Mathlib does not provide this predicate.
\end{definition}

\section{The Quot functor}
\label{mk-sec:quot_functor}

\begin{definition}[Quot functor]\label{mk-def:quot_functor}
  \dcref{custom}

  \uses{mk-def:hilbert_polynomial, mk-def:has_proper_support, mk-def:coherent_sheaf_flat,
    mk-lem:pullback_finite_presentation, mk-lem:coherent_flat_base_change,
    mk-lem:proper_support_base_change, mk-lem:hilbert_fibre_base_change}
  \textit{Source: \cite{nitsure-hilbert-quot}, \S 1 (FGA Explained Ch.~5).}
  Let \(S\) be a locally noetherian scheme, \(\pi : X \to S\) a morphism
  locally of finite type,
  \(L\) a line bundle on \(X\), \(E\) a coherent sheaf on \(X\), and
  \(\Phi \in \mathbb{Q}[\lambda]\). The \emph{Quot functor with Hilbert
  polynomial \(\Phi\)}
  \[
    \Quot^{\Phi,L}_{E/X/S} : (\mathrm{Sch}/S)^{op} \longrightarrow \mathbf{Set}
  \]
  sends an \(S\)-scheme \(T \to S\) to the set of equivalence classes
  \(\langle \mathcal{F}, q\rangle\) of pairs \((\mathcal{F}, q)\) where
  \begin{enumerate}
    \item \(\mathcal{F}\) is a coherent sheaf on \(X_T = X \times_S T\) whose
      schematic support is proper over \(T\) (\ref{mk-def:has_proper_support}) and
      which is flat over \(T\);
    \item \(q : E_T \twoheadrightarrow \mathcal{F}\) is a surjective
      \(\mathcal{O}_{X_T}\)-linear homomorphism, where \(E_T\) is the pull-back
      of \(E\) along \(X_T \to X\);
    \item for every \(t \in T\) the Hilbert polynomial of \(\mathcal{F}|_{X_t}\)
      with respect to \(L|_{X_t}\) (\ref{mk-def:hilbert_polynomial}) equals
      \(\Phi\).
  \end{enumerate}
  Two pairs \((\mathcal{F}, q)\) and \((\mathcal{F}', q')\) are equivalent iff
  \(\ker(q) = \ker(q')\). On morphisms \(f : T' \to T\) over \(S\), the functor
  acts by pulling back the quotient along \(X_{T'} \to X_T\) (well-defined
  since tensor product is right-exact and preserves both flatness and
  properness). The Hilbert scheme is the special case
  \(\Hilb^{\Phi,L}_{X/S} := \Quot^{\Phi,L}_{\mathcal{O}_X/X/S}\).
\end{definition}

\subsection{Base-change well-definedness of the Quot functor}
\label{mk-sec:quot_base_change}

The pullback action of the Quot functor (\ref{mk-def:quot_functor}) along an
\(S\)-morphism \(T' \to T\) is well defined because each of the four
conditions on a family of quotients is preserved by base change; this is the
content of the following four lemmas.  Throughout, write
\(g' : X_{T'} \to X_T\) for the induced morphism on relative products, which
is cartesian over \(T' \to T\).

\begin{lemma}[Pullback preserves finite presentation]\label{mk-lem:pullback_finite_presentation}

  \dcref{custom}
  \textit{Source: \cite{nitsure-hilbert-quot}, \S 1 (FGA Explained Ch.~5).}
  Let \(g : Y \to X\) be a morphism of schemes and \(\mathcal{F}\) a finitely
  presented sheaf of \(\mathcal{O}_X\)-modules (locally the cokernel of a map
  of finite free modules).  Then \(g^* \mathcal{F}\) is a finitely presented
  sheaf of \(\mathcal{O}_Y\)-modules.
\end{lemma}
\begin{proof}Choose an open cover \(\{U_i\}\) of \(X\) on whose members \(\mathcal{F}\)
  admits presentations
  \(\mathcal{O}_{U_i}^{\oplus b_i} \to \mathcal{O}_{U_i}^{\oplus a_i} \to
  \mathcal{F}|_{U_i} \to 0\) with \(a_i, b_i\) finite.  The preimages
  \(g^{-1}(U_i)\) cover \(Y\).  The pullback functor is a left adjoint, hence
  right exact, and it sends \(\mathcal{O}_X\) to \(\mathcal{O}_Y\) and finite
  free modules to finite free modules of the same rank; moreover restriction
  to \(g^{-1}(U_i)\) of the pullback is the pullback of the restriction along
  \(g^{-1}(U_i) \to U_i\).  Applying the restricted pullback to the chosen
  presentation therefore yields a presentation
  \(\mathcal{O}^{\oplus b_i} \to \mathcal{O}^{\oplus a_i} \to
  (g^*\mathcal{F})|_{g^{-1}(U_i)} \to 0\) with the same finite index sets.
\end{proof}

\begin{lemma}[Flatness over the base is stable under base change]\label{mk-lem:coherent_flat_base_change}

  \uses{mk-def:coherent_sheaf_flat}
  \dcref{custom}
  \textit{Source: \cite{nitsure-hilbert-quot}, \S 1 (FGA Explained Ch.~5); Stacks 01U9.}
  Let
  \[
    \begin{array}{ccc}
      X' & \xrightarrow{\ g'\ } & X \\
      \downarrow f' & & \downarrow f \\
      S' & \xrightarrow{\ g\ } & S
    \end{array}
  \]
  be a cartesian square of schemes and \(\mathcal{F}\) a quasi-coherent sheaf
  of \(\mathcal{O}_X\)-modules that is flat over \(S\) in the sense of
  \ref{mk-def:coherent_sheaf_flat}.  Then \(g'^* \mathcal{F}\) is flat over
  \(S'\).
\end{lemma}
\begin{proof}\uses{mk-lem:pullback_isQuasicoherent_hom}
  The pullback \(g'^*\mathcal{F}\) is quasi-coherent
  (\ref{mk-lem:pullback_isQuasicoherent_hom}).  For a quasi-coherent module,
  flatness over the base in the sense of \ref{mk-def:coherent_sheaf_flat} is an
  affine-local condition: if every point of \(X'\) has an affine chart
  \(W \subseteq f'^{-1}(U_t)\) (with \(U_t \subseteq S'\) affine) on which
  \(\Gamma(g'^*\mathcal{F}, W)\) is flat over \(\Gamma(S', U_t)\), then the
  same holds at \emph{every} affine pair
  \(V' \subseteq f'^{-1}(U')\).  Indeed, given such a pair, cover \(V'\) by
  opens that are simultaneously basic in \(V'\) and in one of the charts
  \(W\); on such an open the sections of the quasi-coherent module
  \(g'^*\mathcal{F}\) are a localization of \(\Gamma(g'^*\mathcal{F}, W)\)
  along the algebra \(\Gamma(X', W) \to \Gamma(X', W_b)\), and flatness over
  the base ring is preserved by such localizations of the module and by the
  base exchanges \(\Gamma(S', U_t) \rightsquigarrow \Gamma(S', U')\) along
  basic-open localizations of the base; the local-to-global flatness
  criterion along a spanning family of \(\Gamma(X', V')\) then yields
  flatness of \(\Gamma(g'^*\mathcal{F}, V')\) over \(\Gamma(S', U')\)
  (Stacks 00HT-style affine locality).

  It remains to produce the charts.  Given \(x' \in X'\), choose an affine
  \(U \subseteq S\) around the common image \(f(g'(x')) = g(f'(x'))\), an
  affine \(V \subseteq f^{-1}(U)\) around \(g'(x')\), and an affine
  \(U_t \subseteq g^{-1}(U)\) around \(f'(x')\).  The piece
  \(W := g'^{-1}(V) \cap f'^{-1}(U_t)\) is an affine open of \(X'\): it is
  the fibre product \(V \times_U U_t\) of the restricted cartesian square.
  Taking sections, \(\Gamma(X', W) =
  \Gamma(X, V) \otimes_{\Gamma(S, U)} \Gamma(S', U_t)\) is a pushout of
  rings, and by quasi-coherence of \(\mathcal{F}\) the sections of the
  pullback on the piece are the base change
  \[
    \Gamma(g'^*\mathcal{F}, W)
      \;=\; \Gamma(X', W) \otimes_{\Gamma(X, V)} \Gamma(\mathcal{F}, V)
      \;=\; \Gamma(\mathcal{F}, V) \otimes_{\Gamma(S, U)} \Gamma(S', U_t).
  \]
  The hypothesis at the affine pair \((U, V)\) gives
  \(\Gamma(\mathcal{F}, V)\) flat over \(\Gamma(S, U)\), and flatness is
  stable under base change of the ring, so \(\Gamma(g'^*\mathcal{F}, W)\) is
  flat over \(\Gamma(S', U_t)\), as required.
\end{proof}

\begin{lemma}[Module finiteness is local along a spanning family]
  \label{mk-lem:module_finite_of_localizationSpan_mathlib}
  \dcref{custom}

  Let \(R\) be a commutative ring, \(M\) an \(R\)-module, and
  \(t \subseteq R\) a subset generating the unit ideal.  If for every
  \(g \in t\) the localized module \(M_g\) is a finite \(R_g\)-module, then
  \(M\) is a finite \(R\)-module.
\end{lemma}

\begin{lemma}[Finitely presented modules have finite sections on every affine open]\label{mk-lem:finite_sections_affine_all}

  \dcref{custom}
  \textit{Source: Stacks~\href{https://stacks.math.columbia.edu/tag/01PC}{01PC}
  (finite-type half).}
  Let \(\mathcal{F}\) be a finitely presented sheaf of
  \(\mathcal{O}_X\)-modules on a scheme \(X\).  Then for \emph{every} affine
  open \(V \subseteq X\) the section module \(\Gamma(\mathcal{F}, V)\) is a
  finite \(\Gamma(X, V)\)-module.
\end{lemma}
\begin{proof}\uses{mk-lem:affine_finite_sections_nhds, mk-lem:isLocalization_basicOpen_mathlib,
    mk-lem:qcoh_section_localization_basicOpen,
    mk-lem:module_finite_of_localizationSpan_mathlib}
  Let \(t \subseteq \Gamma(X, V)\) be the set of all \(f\) such that
  \(D(f) \subseteq V\) and \(\Gamma(\mathcal{F}, D(f))\) is a finite
  \(\Gamma(X, D(f))\)-module.

  \emph{The basic opens of \(t\) cover \(V\).}  Fix \(x \in V\).  By
  \ref{mk-lem:affine_finite_sections_nhds} there is an affine open chart \(W\)
  with \(x \in W \subseteq V\) and \(\Gamma(\mathcal{F}, W)\) finite over
  \(\Gamma(X, W)\).  Since the basic opens of the affine \(V\) form a basis of
  its topology, there is \(f \in \Gamma(X, V)\) with
  \(x \in D(f) \subseteq W\).  Writing \(f|_W\) for the restriction of \(f\)
  to \(W\), one has \(D(f) = D(f|_W)\) as opens contained in \(W\); the
  restriction \(\Gamma(X, W) \to \Gamma(X, D(f))\) is the localization away
  from \(f|_W\) (\ref{mk-lem:isLocalization_basicOpen_mathlib}), and
  \(\Gamma(\mathcal{F}, W) \to \Gamma(\mathcal{F}, D(f))\) is the
  corresponding module localization
  (\ref{mk-lem:qcoh_section_localization_basicOpen}, applicable because a
  finitely presented module is quasi-coherent).  The images of a finite
  generating set of \(\Gamma(\mathcal{F}, W)\) generate the localized module
  over the localized ring, so \(\Gamma(\mathcal{F}, D(f))\) is finite over
  \(\Gamma(X, D(f))\) and \(f \in t\), with \(x \in D(f)\).

  \emph{Conclusion.}  The basic opens \(D(f)\), \(f \in t\), cover the affine
  \(V\), so \(t\) generates the unit ideal of \(\Gamma(X, V)\).  For each
  \(f \in t\) the restriction
  \(\Gamma(\mathcal{F}, V) \to \Gamma(\mathcal{F}, D(f))\) is the localization
  of \(\Gamma(\mathcal{F}, V)\) away from \(f\)
  (\ref{mk-lem:qcoh_section_localization_basicOpen}), with target a finite
  \(\Gamma(X, D(f)) = \Gamma(X, V)_f\)-module by membership in \(t\);
  \ref{mk-lem:module_finite_of_localizationSpan_mathlib} glues these into
  finiteness of \(\Gamma(\mathcal{F}, V)\) over \(\Gamma(X, V)\).
\end{proof}

\begin{lemma}[Annihilators are stable under pullback: section level]\label{mk-lem:annihilator_pullback_sections}

  \dcref{custom}
  \textit{Source:
  Stacks~\href{https://stacks.math.columbia.edu/tag/056H}{056H};
  \cite{nitsure-hilbert-quot}, \S 1.}
  Let \(g' : X' \to X\) be a morphism of schemes, \(\mathcal{F}\) a
  quasi-coherent sheaf of \(\mathcal{O}_X\)-modules, \(U \subseteq X\) an
  affine open, and \(V' \subseteq g'^{-1}(U)\) any open of \(X'\).  If
  \(r \in \mathcal{O}_X(U)\) annihilates \(\Gamma(\mathcal{F}, U)\), then its
  image in \(\mathcal{O}_{X'}(V')\) under \(g'^\sharp\) and restriction
  annihilates \(\Gamma(g'^* \mathcal{F}, V')\).
\end{lemma}
\begin{proof}\uses{mk-lem:modules_pullback_mathlib, mk-lem:pullback_isQuasicoherent_hom,
    mk-lem:qcoh_section_localization_basicOpen}
  For a sheaf of modules \(\mathcal{M}\) on a scheme \(Z\) and a global
  section \(w \in \mathcal{O}_Z(Z)\), scalar multiplication by \(w\) is an
  endomorphism \(w \cdot : \mathcal{M} \to \mathcal{M}\): the componentwise
  multiplications by the restrictions of \(w\) are linear because the section
  rings are commutative, and they commute with the restriction maps because
  restriction is semilinear.  The proof transports the vanishing of this
  endomorphism.

  \emph{Pullback intertwines the multiplication endomorphisms.}  For a
  morphism \(h : Z' \to Z\) one has \(h^*(w \cdot) = (h^\sharp w) \cdot\) as
  endomorphisms of \(h^* \mathcal{M}\).  Indeed, under the adjunction between
  \(h^*\) and \(h_*\) (\ref{mk-lem:modules_pullback_mathlib}) both sides
  transpose to morphisms \(\mathcal{M} \to h_* h^* \mathcal{M}\), and these
  transposes agree: the adjunction unit
  \(\eta : \mathcal{M} \to h_* h^* \mathcal{M}\) is a morphism of
  \(\mathcal{O}_Z\)-modules, and the \(\mathcal{O}_Z\)-module structure of
  \(h_* h^* \mathcal{M}\) is restriction of scalars along \(h^\sharp\), so
  both transposes send a section \(m\) over \(O \subseteq Z\) to
  \(\bigl(h^\sharp(w|_O)\bigr) \cdot \eta(m) = \bigl(w|_O\bigr) \cdot
  \eta(m)\).

  \emph{On an affine scheme, killing global sections kills all sections.}
  Let \(Z\) be affine, \(\mathcal{M}\) quasi-coherent, and suppose
  \(w \cdot \Gamma(\mathcal{M}, Z) = 0\).  On a basic open \(D(b)\) the
  restriction \(\Gamma(\mathcal{M}, Z) \to \Gamma(\mathcal{M}, D(b))\) is the
  localization away from \(b\)
  (\ref{mk-lem:qcoh_section_localization_basicOpen}); for
  \(y \in \Gamma(\mathcal{M}, D(b))\) there are \(n\) and
  \(m \in \Gamma(\mathcal{M}, Z)\) with \(b^n y = m|_{D(b)}\), whence
  \(b^n (w y) = w|_{D(b)} \cdot m|_{D(b)} = (w \cdot m)|_{D(b)} = 0\); since
  \(b\) acts invertibly on the localized sections, \(w y = 0\).  A general
  open of \(Z\) is covered by basic opens and the abelian sheaf underlying
  \(\mathcal{M}\) is separated, so multiplication by (the restriction of)
  \(w\) kills all sections of \(\mathcal{M}\).

  \emph{Conclusion.}  Restrict \(\mathcal{F}\) to the affine open subscheme
  \(U\), i.e.\ form \(\iota_U^* \mathcal{F}\) for the open immersion
  \(\iota_U : U \to X\); it is quasi-coherent
  (\ref{mk-lem:pullback_isQuasicoherent_hom}), and the global section
  \(r_0 \in \mathcal{O}_U(U)\) corresponding to \(r\) kills
  \(\Gamma(\iota_U^* \mathcal{F}, U) = \Gamma(\mathcal{F}, U)\).  By the
  affine step, \(r_0 \cdot\) is the zero endomorphism of
  \(\iota_U^* \mathcal{F}\).  Let \(h : V' \to U\) be the restriction of
  \(g'\), defined since \(V' \subseteq g'^{-1}(U)\); by the intertwining step
  \((h^\sharp r_0) \cdot\) is the zero endomorphism of
  \(h^* \iota_U^* \mathcal{F}\), which is identified with
  \(\iota_{V'}^* (g'^* \mathcal{F})\) by pseudofunctoriality of pullback
  along \(\iota_U \circ h = g' \circ \iota_{V'}\)
  (\ref{mk-lem:modules_pullback_mathlib}).  Reading off global sections over
  \(V'\) — the identification
  \(\Gamma(\iota_{V'}^* g'^* \mathcal{F}, V') = \Gamma(g'^* \mathcal{F}, V')\)
  is equivariant for the ring identification
  \(\mathcal{O}_{V'}(V') = \mathcal{O}_{X'}(V')\), which carries
  \(h^\sharp r_0\) to the image of \(r\) — the image of \(r\) annihilates
  \(\Gamma(g'^* \mathcal{F}, V')\).
\end{proof}

\begin{lemma}[The annihilator ideal sheaf pulls back]\label{mk-lem:annihilator_le_map_pullback}

  \uses{mk-def:modules_annihilator, mk-def:schematic_support_immersion}
  \dcref{custom}
  \textit{Source:
  Stacks~\href{https://stacks.math.columbia.edu/tag/056H}{056H}.}
  Let \(g' : X' \to X\) be a morphism of schemes and \(\mathcal{F}\) a
  finitely presented sheaf of \(\mathcal{O}_X\)-modules.  Write
  \(\iota' : Z' \to X'\) for the schematic-support immersion of
  \(g'^* \mathcal{F}\) (\ref{mk-def:schematic_support_immersion}).  Then
  \(\mathrm{Ann}(\mathcal{F})\) is contained in the kernel ideal sheaf of
  \(\mathcal{O}_X \to (\iota' \circ g')_* \mathcal{O}_{Z'}\); equivalently,
  \(g' \circ \iota'\) factors through the schematic support
  \(Z \subseteq X\) of \(\mathcal{F}\) by a morphism \(l : Z' \to Z\) over
  \(g'\).
\end{lemma}
\begin{proof}\uses{mk-lem:annihilator_pullback_sections, mk-lem:modules_annihilator_ideal,
    mk-lem:finite_sections_affine_all, mk-lem:pullback_finite_presentation,
    mk-lem:modules_annihilator_ideal_le}
  The kernel ideal sheaf of \((\iota' \circ g')^\sharp\) is the largest ideal
  sheaf whose value on each affine open \(U \subseteq X\) is contained in
  \(\ker\bigl((\iota' \circ g')^\sharp_U\bigr)\), so it suffices to prove,
  for each affine open \(U\) and each \(r \in \mathrm{Ann}(\mathcal{F})(U)\),
  that \((\iota' \circ g')^\sharp_U(r) = 0\) in
  \(\Gamma\bigl(Z', (\iota' \circ g')^{-1}(U)\bigr)\).  By
  \ref{mk-lem:modules_annihilator_ideal_le}, \(r\) annihilates
  \(\Gamma(\mathcal{F}, U)\).

  The open \((\iota')^{-1}(g'^{-1}(U))\) of \(Z'\) is covered by the
  preimages \((\iota')^{-1}(V')\) of the affine opens
  \(V' \subseteq g'^{-1}(U)\) of \(X'\), and the structure sheaf of \(Z'\) is
  separated, so it suffices that the restriction of
  \((\iota' \circ g')^\sharp_U(r)\) to each \((\iota')^{-1}(V')\) vanishes.
  By naturality of \(\iota'^\sharp\), that restriction is
  \(\iota'^\sharp_{V'}(\rho)\), where \(\rho \in \mathcal{O}_{X'}(V')\) is the
  image of \(r\).  By \ref{mk-lem:annihilator_pullback_sections}, \(\rho\)
  annihilates \(\Gamma(g'^* \mathcal{F}, V')\).  The pullback
  \(g'^* \mathcal{F}\) is finitely presented
  (\ref{mk-lem:pullback_finite_presentation}), hence quasi-coherent with finite
  section modules on all affine opens
  (\ref{mk-lem:finite_sections_affine_all}); therefore
  \(\mathrm{Ann}_{\mathcal{O}_{X'}(V')}\Gamma(g'^* \mathcal{F}, V')
  = \mathrm{Ann}(g'^* \mathcal{F})(V')\)
  (\ref{mk-lem:modules_annihilator_ideal}), and on affine opens the subscheme
  immersion of an ideal sheaf kills exactly that ideal, so
  \(\iota'^\sharp_{V'}(\rho) = 0\), as required.  Finally, a morphism to
  \(X\) under whose structure map the defining ideal sheaf of the closed
  subscheme \(Z\) dies factors through \(Z\); this produces \(l : Z' \to Z\)
  with \(\iota \circ l = g' \circ \iota'\).
\end{proof}

\begin{lemma}[Proper support is stable under base change]\label{mk-lem:proper_support_base_change}

  \uses{mk-def:has_proper_support, mk-def:modules_annihilator}
  \dcref{custom}
  \textit{Source: \cite{nitsure-hilbert-quot}, \S 1 (FGA Explained Ch.~5); Stacks 01W4, 056H.}
  In the cartesian square of \ref{mk-lem:coherent_flat_base_change}, let
  \(\mathcal{F}\) be a finitely presented sheaf of
  \(\mathcal{O}_X\)-modules whose schematic support is proper over \(S\)
  (\ref{mk-def:has_proper_support}).  Then the schematic support of
  \(g'^* \mathcal{F}\) is proper over \(S'\).
\end{lemma}
\begin{proof}\uses{mk-lem:annihilator_le_map_pullback, mk-lem:isProper_mathlib,
    mk-def:schematic_support_immersion}
  Write \(\iota : Z \to X\) and \(\iota' : Z' \to X'\) for the
  schematic-support immersions (\ref{mk-def:schematic_support_immersion}) of
  \(\mathcal{F}\) and of \(g'^* \mathcal{F}\).  By
  \ref{mk-lem:annihilator_le_map_pullback} there is a morphism
  \(l : Z' \to Z\) with \(\iota \circ l = g' \circ \iota'\).

  Form the fibre product \(W := Z \times_X X'\) with projections
  \(p : W \to Z\) and \(q : W \to X'\).  Pasting its defining cartesian
  square vertically with the given cartesian square exhibits \(W\), together
  with \(p\) and \(f' \circ q\), as the fibre product of
  \(Z \xrightarrow{\,f \circ \iota\,} S \xleftarrow{\,g\,} S'\).  The
  hypothesis says \(f \circ \iota\) is proper, and properness is stable under
  base change (\ref{mk-lem:isProper_mathlib}), so \(f' \circ q : W \to S'\) is
  proper.

  The pair \((l, \iota')\) induces \(c : Z' \to W\) with
  \(q \circ c = \iota'\).  The immersion \(\iota'\) is a closed immersion,
  hence proper; and \(q\) is a closed immersion (the base change of the
  closed immersion \(\iota\) along \(g'\)), hence separated.  By cancellation
  of properness along separated morphisms applied to \(q \circ c = \iota'\)
  (\ref{mk-lem:isProper_mathlib}), \(c\) is proper.  Finally
  \[
    f' \circ \iota' \;=\; (f' \circ q) \circ c
  \]
  is a composite of proper morphisms, hence proper
  (\ref{mk-lem:isProper_mathlib}); that is, the schematic support of
  \(g'^* \mathcal{F}\) is proper over \(S'\).
\end{proof}

\begin{lemma}[Pullback commutes with tensor product and twists]\label{mk-lem:pullback_moduleTensorPow}

  \uses{mk-def:sheafTensorObj, mk-def:sheafTensorPow, mk-def:sheafModuleTwist,
    mk-lem:modules_pullback_mathlib, mk-lem:pullback_tensor_map_isiso}
  \dcref{01CD,custom}
  \textit{Source: \cite{stacks-project}, "Sheaves of Modules", Tag 01CD (pullback
  and tensor products); the stalkwise description is Tags 01CB and 0098.}
  Let \(f : Y \to X\) be a morphism of schemes and let \(\mathcal{F}\), \(L\)
  be quasi-coherent sheaves of \(\mathcal{O}_X\)-modules.  For every
  \(m \in \mathbb{N}\) there is an isomorphism of \(\mathcal{O}_Y\)-modules
  \[
    f^*\bigl(\mathcal{F} \otimes_{\mathcal{O}_X} L^{\otimes m}\bigr)
    \;\cong\;
    f^*\mathcal{F} \otimes_{\mathcal{O}_Y} \bigl(f^*L\bigr)^{\otimes m}.
  \]
  (The fact holds for arbitrary sheaves of modules on ringed spaces, with a
  canonical comparison isomorphism; the quasi-coherent case stated here is
  the one consumed by \ref{mk-lem:hilbert_fibre_base_change}.)
\end{lemma}
\begin{proof}
  The twist \(\mathcal{F} \otimes L^{\otimes m}\) is built from the binary
  tensor product (\ref{mk-def:sheafTensorObj}) and the unit
  \(\mathcal{O}_X\) by the recursion \(L^{\otimes 0} = \mathcal{O}_X\),
  \(L^{\otimes (m+1)} = L^{\otimes m} \otimes L\)
  (\ref{mk-def:sheafTensorPow}), so by induction on \(m\) it suffices that the
  pullback functor carries \(\mathcal{O}_X\) to \(\mathcal{O}_Y\) — the
  unitality of \ref{mk-lem:modules_pullback_mathlib} — and commutes with the
  binary tensor product up to natural isomorphism, which is
  \ref{mk-lem:pullback_tensor_map_isiso}.  The induction is carried out in the formal proof
  (the base and step assembled from those two isomorphisms and the tensor
  congruence \(\texttt{sheafTensorObjCongr}\)); only the binary comparison of
  \ref{mk-lem:pullback_tensor_map_isiso} remains as an open leaf.
\end{proof}

\begin{lemma}[The canonical pullback--tensor comparison is an isomorphism]\label{mk-lem:pullback_tensor_map_isiso}

  \uses{mk-def:sheafTensorObj, mk-lem:modules_pullback_mathlib}
  \dcref{01CD,custom}
  \textit{Source: \cite{stacks-project}, "Sheaves of Modules", Tag 01CD (pullback
  and tensor products); the stalkwise description is Tags 01CB and 0098.}
  Let \(f : Y \to X\) be a morphism of schemes and let \(\mathcal{F}\),
  \(\mathcal{G}\) be sheaves of \(\mathcal{O}_X\)-modules.  The canonical
  comparison morphism
  \[
    f^*\bigl(\mathcal{F} \otimes_{\mathcal{O}_X} \mathcal{G}\bigr)
    \;\longrightarrow\;
    f^*\mathcal{F} \otimes_{\mathcal{O}_Y} f^*\mathcal{G}
  \]
  (the sheafified mate of the lax monoidal structure of \(f_*\)) is an
  isomorphism.  The statement holds for arbitrary sheaves of modules on
  ringed spaces; no quasi-coherence hypothesis is needed.
\end{lemma}
\begin{proof}
  For every \(\mathcal{O}_Y\)-module \(\mathcal{H}\) there are isomorphisms,
  natural in all three variables:
  \begin{align*}
    \operatorname{Hom}_{\mathcal{O}_Y}\!\bigl(f^*(\mathcal{F} \otimes
      \mathcal{G}), \mathcal{H}\bigr)
    &\cong \operatorname{Hom}_{\mathcal{O}_X}\!\bigl(\mathcal{F} \otimes
      \mathcal{G}, f_*\mathcal{H}\bigr)
    \cong \operatorname{Hom}_{\mathcal{O}_X}\!\bigl(\mathcal{F},
      \mathcal{H}om_{\mathcal{O}_X}(\mathcal{G}, f_*\mathcal{H})\bigr) \\
    &\cong \operatorname{Hom}_{\mathcal{O}_X}\!\bigl(\mathcal{F},
      f_*\,\mathcal{H}om_{\mathcal{O}_Y}(f^*\mathcal{G}, \mathcal{H})\bigr)
    \cong \operatorname{Hom}_{\mathcal{O}_Y}\!\bigl(f^*\mathcal{F},
      \mathcal{H}om_{\mathcal{O}_Y}(f^*\mathcal{G}, \mathcal{H})\bigr) \\
    &\cong \operatorname{Hom}_{\mathcal{O}_Y}\!\bigl(f^*\mathcal{F} \otimes
      f^*\mathcal{G}, \mathcal{H}\bigr),
  \end{align*}
  by, in order: the pullback--pushforward adjunction; the tensor--hom
  adjunction for \(\mathcal{O}_X\)-modules; the isomorphism
  \(\mathcal{H}om_{\mathcal{O}_X}(\mathcal{G}, f_*\mathcal{H}) \cong
  f_*\,\mathcal{H}om_{\mathcal{O}_Y}(f^*\mathcal{G}, \mathcal{H})\) — over
  an open \(V \subseteq X\) both sides compute
  \(\operatorname{Hom}_{\mathcal{O}_{f^{-1}V}}\bigl((f^*\mathcal{G})|_{f^{-1}V},
  \mathcal{H}|_{f^{-1}V}\bigr)\) by the pullback--pushforward adjunction for
  the restricted morphism \(f^{-1}V \to V\), since pullback commutes with
  restriction to opens, and these identifications are compatible with
  further restriction; the pullback--pushforward adjunction again; and the
  tensor--hom adjunction for \(\mathcal{O}_Y\)-modules.  By the Yoneda lemma
  the composite isomorphism of functors in \(\mathcal{H}\) is induced by a
  unique isomorphism \(f^*(\mathcal{F} \otimes \mathcal{G}) \cong
  f^*\mathcal{F} \otimes_{\mathcal{O}_Y} f^*\mathcal{G}\), and it is inverse
  to the canonical comparison map.
\end{proof}


\begin{lemma}[Annihilator monotonicity under tensoring]\label{mk-lem:annihilator_le_annihilator_tensorProduct}
  \dcref{custom}

  Let \(R\) be a commutative ring and \(M, N\) two \(R\)-modules. Then
  \[
    \mathrm{Ann}_R(M) \;\subseteq\; \mathrm{Ann}_R(M \otimes_R N),
    \qquad
    \mathrm{Ann}_R(N) \;\subseteq\; \mathrm{Ann}_R(M \otimes_R N).
  \]
\end{lemma}
\begin{proof}A scalar \(a \in \mathrm{Ann}_R(M)\) kills every elementary tensor via
  \(a \cdot (m \otimes n) = (a \cdot m) \otimes n = 0\), hence the whole tensor
  product by additivity; symmetrically for \(N\).
\end{proof}

\begin{remark}
  \label{mk-rem:annihilator_tensorProduct_support}
  \dcref{custom}
  \uses{mk-lem:annihilator_le_annihilator_tensorProduct, mk-def:modules_annihilator}
  On an affine open, the section module of a twist \(\mathcal F \otimes
  L^{\otimes m}\) of quasi-coherent sheaves is the tensor product of the
  section modules, so \ref{mk-lem:annihilator_le_annihilator_tensorProduct}
  is the sections-level algebraic content behind the monotonicity
  \(\mathrm{Ann}(\mathcal F) \subseteq \mathrm{Ann}(\mathcal F \otimes
  L^{\otimes m})\) of the annihilator ideal sheaves
  (\ref{mk-def:modules_annihilator}) --- and hence, via
  \(\mathrm{IdealSheafData.ofIdeals\_mono}\), the containment of schematic
  supports \(\mathrm{Supp}(\mathcal F \otimes L^{\otimes m}) \subseteq
  \mathrm{Supp}(\mathcal F)\) that would give properness of the twisted
  fibre's support from properness of the support of \(\mathcal F\). Lifting
  this pure-algebra monotonicity to the sheaf-level inclusion needs an
  affine \emph{tensor-section comparison}
  \(\Gamma(\mathcal F \otimes \mathcal G, U) \cong \Gamma(\mathcal F, U)
  \otimes_{\Gamma(X, U)} \Gamma(\mathcal G, U)\) for quasi-coherent
  \(\mathcal F, \mathcal G\) on an affine open \(U\) (Stacks 01CC); this
  comparison is the subject of the tensor-section bricks below, and its
  affine-open case (the genuinely open step) is recorded in
  \ref{mk-rem:tensor_section_formula_gap}.
\end{remark}

\begin{lemma}[Dimension transport across a base-change linear equivalence]\label{mk-lem:finrank_eq_of_baseChange_linearEquiv}
  \dcref{custom}

  Let \(\kappa \to \kappa'\) be a field extension, \(V\) a \(\kappa\)-vector
  space, and \(W\) a \(\kappa'\)-vector space. Every \(\kappa'\)-linear
  equivalence \(\kappa' \otimes_\kappa V \simeq_{\kappa'} W\) forces
  \(\dim_{\kappa'} W = \dim_\kappa V\).
\end{lemma}
\begin{proof}Over fields the base-change identity
  \(\dim_{\kappa'}(\kappa' \otimes_\kappa V) = \dim_\kappa V\) holds
  unconditionally (with the junk value \(0\) shared correctly in the
  infinite-dimensional case), and a linear equivalence transports dimension.
\end{proof}

\begin{remark}
  \label{mk-rem:finrank_baseChange_linearEquiv_role}
  \dcref{custom}
  \uses{mk-lem:finrank_eq_of_baseChange_linearEquiv}
  This is the terminal dimension-count packaging of
  \ref{mk-lem:gamma_fiber_baseChange_field}'s alternative route: flat base
  change over the residue field extension \(\kappa(t) \to \kappa(t')\)
  supplies a \(\kappa(t')\)-linear equivalence between
  \(\kappa(t') \otimes_{\kappa(t)} \Gamma(Z, \mathcal N)\) and
  \(\Gamma(Z', v^*\mathcal N)\) for the schematic support \(Z\) of the
  twisted fibre module and its base change \(Z'\), and
  \ref{mk-lem:finrank_eq_of_baseChange_linearEquiv} converts this directly into
  the dimension equality, without any finiteness case split.
\end{remark}

\subsection{The affine tensor-section comparison}
\label{mk-subsec:tensor_section_formula}

The map-half of the affine tensor-section comparison flagged in
\ref{mk-rem:annihilator_tensorProduct_support} is recorded here: for sheaves
of \(\mathcal O_X\)-modules \(A, B\), the sheafification unit of the
presheaf-of-modules tensor product supplies, on every open \(V\), a
canonical \(\Gamma(X,V)\)-linear comparison map into the sections of the
substrate tensor \(\mathrm{tensorObj}(A,B)\) (\ref{mk-def:scheme_modules_tensorobj}).

\begin{definition}[The presheaf-of-modules tensor and its section comparison]\label{mk-def:tensor_section_hom}
  \dcref{custom}

  \uses{mk-def:scheme_modules_tensorobj}
  Let \(A, B\) be sheaves of \(\mathcal O_X\)-modules. Write \(A \otimes^p B\)
  for the presheaf-of-modules tensor product of the underlying presheaves of
  \(A\) and \(B\); objectwise, \((A \otimes^p B)(V) = \Gamma(A,V)
  \otimes_{\Gamma(X,V)} \Gamma(B,V)\) for every open \(V\). The
  \emph{section comparison} at \(V\) is the \(\Gamma(X,V)\)-linear map
  \[
    \tau_{A,B,V} \;:\; \Gamma(A,V) \otimes_{\Gamma(X,V)} \Gamma(B,V)
    \;\longrightarrow\; \Gamma(\mathrm{tensorObj}(A,B), V)
  \]
  given by the \(V\)-component of the sheafification-adjunction unit at
  \(A \otimes^p B\) (\(\mathrm{tensorObj}(A,B)\) being, by definition, the
  sheafification of \(A \otimes^p B\)).
\end{definition}
\begin{proof}\uses{mk-def:scheme_modules_tensorobj}
  The presheaf-of-modules tensor product and its objectwise formula are
  Mathlib's monoidal structure on presheaves of modules; the section
  comparison is the sheafification-adjunction unit, restricted to the open
  \(V\).
\end{proof}

\begin{lemma}[Naturality and normalisation of the section comparison]\label{mk-lem:tensor_section_hom_naturality}
  \dcref{custom}

  \uses{mk-def:tensor_section_hom, mk-def:sectionMul}
  The section comparison \(\tau_{A,B,-}\) of \ref{mk-def:tensor_section_hom}
  is natural in restriction: for opens \(V \subseteq W\) it commutes with
  the restriction maps of \(A \otimes^p B\) and of \(\mathrm{tensorObj}(A,B)\).
  Over the total space, \(\tau_{A,B,\top}\) is the section multiplication
  \(\mathrm{sectionsMul}\) of \ref{mk-def:sectionMul}.
\end{lemma}
\begin{proof}\uses{mk-def:sectionMul}
  Naturality is naturality of the sheafification-adjunction unit as a
  morphism of presheaves. At \(V = \top\) the section comparison is by
  definition the global-sections component of that same unit evaluated at
  \(A \otimes^p B\), which is exactly how \(\mathrm{sectionsMul}\) is
  constructed in \ref{mk-def:sectionMul}.
\end{proof}

\begin{lemma}[The substrate tensor and the section-graded tensor coincide]\label{mk-lem:tensorObj_iso_sheafTensorObj}
  \dcref{custom}

  \uses{mk-def:scheme_modules_tensorobj, mk-def:sheafTensorObj}
  The substrate tensor \(\mathrm{tensorObj}(A,B)\)
  (\ref{mk-def:scheme_modules_tensorobj}) and the section-graded tensor
  \(\mathcal F \otimes_{\mathcal O_X} \mathcal G\) of \ref{mk-def:sheafTensorObj}
  are the same object: both are the sheafification of the presheaf-of-modules
  tensor of the underlying presheaves.
\end{lemma}
\begin{proof}Definitional: both constructions sheafify \(A \otimes^p B\).
\end{proof}

\begin{lemma}[Sheafifying the tensor comparison unit is an isomorphism]\label{mk-lem:isIso_sheafification_tensorSectionUnit}
  \dcref{custom}

  \uses{mk-def:tensor_section_hom}
  The sheafification of the presheaf comparison unit
  \(A \otimes^p B \to \mathrm{tensorObj}(A,B)\) is an isomorphism of sheaves
  of modules.
\end{lemma}
\begin{proof}The sheafification functor is a reflective localization, so it inverts the
  localization unit of any presheaf --- here, of \(A \otimes^p B\).
\end{proof}

\begin{remark}
[The remaining affine step, and a narrowing recommendation]
  \label{mk-rem:tensor_section_formula_gap}
  \dcref{custom}
  \uses{mk-def:tensor_section_hom, mk-lem:isIso_sheafification_tensorSectionUnit,
    mk-lem:pullback_tensor_map_isiso}
  What is \emph{not} yet established is that \(\tau_{A,B,V}\)
  (\ref{mk-def:tensor_section_hom}) is an isomorphism for \(A, B\)
  quasi-coherent and \(V\) affine --- the Stacks 01CC affine tensor-section
  formula flagged in \ref{mk-rem:annihilator_tensorProduct_support}. The route
  is: identify \(\Gamma(A,V)_f \otimes_{\Gamma(X,V)_f} \Gamma(B,V)_f\) with
  \((\Gamma(A,V) \otimes_{\Gamma(X,V)} \Gamma(B,V))_f\) on a basic open
  \(D(f) \subseteq V\) (a localization-tensor identity), transport this
  through the tilde-comparison engine for quasi-coherent sheaves, and
  globalize over an affine cover; \ref{mk-lem:isIso_sheafification_tensorSectionUnit}
  is the categorical fact that makes the sheafification step in this route
  exact rather than merely a comparison map. This affine case is also all
  that is needed by the sole consumer of the general pullback--tensor
  comparison \ref{mk-lem:pullback_tensor_map_isiso}
  (\ref{mk-lem:pullback_moduleTensorPow}, stated for quasi-coherent sheaves),
  so narrowing \ref{mk-lem:pullback_tensor_map_isiso} to the quasi-coherent
  case would let this affine formula suffice in place of the general
  ringed-space stalk argument sketched there.
\end{remark}

\begin{lemma}[Sections of the twisted fibre module under residue field extension]\label{mk-lem:gamma_fiber_baseChange_field}

  \uses{mk-def:fiber_hilbert_function, mk-def:has_proper_support,
    mk-def:modules_annihilator}
  \dcref{02KH,custom}
  \textit{Source: \cite{nitsure-hilbert-quot}, \S 1 (FGA Explained Ch.~5); \cite{stacks-project},
  Tag 02KH (flat base change of cohomology), degree \(0\).}
  Let \(\pi : X \to S\) be a morphism of schemes, \(L\) a quasi-coherent
  sheaf of \(\mathcal{O}_X\)-modules, \(\psi : T' \to T\) a morphism of
  \(S\)-schemes, and \(\mathcal{F}\) a finitely presented sheaf on \(X_T\)
  whose schematic support is proper over \(T\)
  (\ref{mk-def:has_proper_support}).  Let \(t' \in T'\) with image
  \(t \in T\), let
  \[
    u : (X_{T'})_{t'} \longrightarrow (X_T)_t
  \]
  be the morphism exhibiting the fibre of \(X_{T'} \to T'\) at \(t'\) as the
  base change of the fibre of \(X_T \to T\) at \(t\) along the residue field
  extension \(\kappa(t) \to \kappa(t')\), and let \(m \in \mathbb{N}\).
  Then
  \[
    \dim_{\kappa(t')} \Gamma\bigl((X_{T'})_{t'},\,
      u^*(\mathcal{F}_t \otimes L_t^{\otimes m})\bigr)
    \;=\;
    \dim_{\kappa(t)} \Gamma\bigl((X_T)_t,\,
      \mathcal{F}_t \otimes L_t^{\otimes m}\bigr),
  \]
  where \(\mathcal{F}_t\) and \(L_t\) denote the restrictions of
  \(\mathcal{F}\) and of the pullback of \(L\) to the fibre \((X_T)_t\), and
  the right-hand side is the graded Hilbert function of
  \ref{mk-def:fiber_hilbert_function} at \((t, m)\).
\end{lemma}
\begin{proof}
  \uses{mk-lem:proper_support_base_change}
  Write \(\mathcal{M} := \mathcal{F}_t \otimes L_t^{\otimes m}\); it is
  quasi-coherent (\(\mathcal{F}_t\) is finitely presented and \(L_t\)
  quasi-coherent, both being restrictions of such along the fibre
  embedding), and its pullback \(u^*\mathcal{M}\) is quasi-coherent as well.
  Write \(i : Z \hookrightarrow (X_T)_t\) for the schematic support of
  \(\mathcal{M}\), the closed subscheme cut out by the annihilator ideal
  \(\mathrm{Ann}(\mathcal{M})\) (\ref{mk-def:modules_annihilator}).

  \emph{The support \(Z\) is proper over \(\kappa(t)\).}  The fibre
  \((X_T)_t = X_T \times_T \operatorname{Spec} \kappa(t)\) is a base change
  of \(X_T\), so by \ref{mk-lem:proper_support_base_change} the schematic
  support of \(\mathcal{F}_t\) is proper over
  \(\operatorname{Spec} \kappa(t)\).  A local section annihilating
  \(\mathcal{F}_t\) annihilates \(\mathcal{M} = \mathcal{F}_t \otimes
  L_t^{\otimes m}\) (it annihilates every elementary tensor of the tensor
  product presheaf, hence the sheafification), so
  \(\mathrm{Ann}(\mathcal{F}_t) \subseteq \mathrm{Ann}(\mathcal{M})\) and
  \(Z\) is a closed subscheme of the schematic support of
  \(\mathcal{F}_t\); a closed immersion followed by a proper morphism is
  proper, so \(Z \to \operatorname{Spec} \kappa(t)\) is proper — in
  particular \(Z\) is quasi-compact and separated.

  \emph{Reduction to the support and flat base change.}  Since
  \(\mathrm{Ann}(\mathcal{M})\) annihilates \(\mathcal{M}\), the module
  \(\mathcal{M}\) is \(i_*\) of a quasi-coherent module \(\mathcal{N}\) on
  \(Z\).  Let \(i' : Z' \hookrightarrow (X_{T'})_{t'}\) be the base change of
  \(i\) along \(\operatorname{Spec} \kappa(t') \to
  \operatorname{Spec} \kappa(t)\), where
  \(Z' = Z \times_{\operatorname{Spec} \kappa(t)}
  \operatorname{Spec} \kappa(t')\) with projection \(v : Z' \to Z\);
  pushforward along a closed immersion commutes with flat base change of
  quasi-coherent modules, so \(u^*\mathcal{M} = u^* i_* \mathcal{N} =
  i'_* (v^* \mathcal{N})\).  Since \(\Gamma \circ i_* = \Gamma\) and
  \(\Gamma \circ i'_* = \Gamma\), both displayed section spaces are global
  sections over \(Z\), respectively over its base change \(Z'\).  The
  extension \(\kappa(t) \to \kappa(t')\) of fields is flat and \(Z\) is
  quasi-compact and separated over \(\kappa(t)\), so flat base change for
  the cohomology of quasi-coherent modules (\cite{stacks-project}, Tag 02KH, in
  degree \(0\)) gives a \(\kappa(t')\)-linear isomorphism
  \[
    \Gamma(Z', v^*\mathcal{N})
    \;\cong\;
    \Gamma(Z, \mathcal{N}) \otimes_{\kappa(t)} \kappa(t').
  \]

  \emph{Dimension count.}  For any \(\kappa(t)\)-vector space \(V\), a
  \(\kappa(t)\)-basis of \(V\) induces a \(\kappa(t')\)-basis of
  \(V \otimes_{\kappa(t)} \kappa(t')\), so
  \(\dim_{\kappa(t')}(V \otimes_{\kappa(t)} \kappa(t')) =
  \dim_{\kappa(t)} V\) — in the infinite-dimensional case both sides are
  infinite.  Applying this to \(V = \Gamma(Z, \mathcal{N}) =
  \Gamma\bigl((X_T)_t, \mathcal{M}\bigr)\) yields the claim.
\end{proof}

\begin{lemma}[The fibrewise Hilbert function is invariant under base change]\label{mk-lem:hilbert_fibre_base_change}

  \uses{mk-def:fiber_hilbert_function, mk-def:has_proper_support}
  \dcref{custom}
  \textit{Source: \cite{nitsure-hilbert-quot}, \S 1 (FGA Explained Ch.~5).}
  Let \(\pi : X \to S\) be a morphism of schemes, \(L\) a quasi-coherent
  sheaf of \(\mathcal{O}_X\)-modules, \(T' \to T\) a morphism of
  \(S\)-schemes, and
  \(\mathcal{F}\) a finitely presented sheaf on \(X_T\) whose schematic
  support is proper over \(T\).  For \(t' \in T'\) with image \(t \in T\)
  and every \(m \in \mathbb{N}\),
  \[
    \dim_{\kappa(t')} \Gamma\bigl((X_{T'})_{t'},\,
      (g'^*\mathcal{F})_{t'} \otimes L_{t'}^{\otimes m}\bigr)
    \;=\;
    \dim_{\kappa(t)} \Gamma\bigl((X_T)_t,\,
      \mathcal{F}_t \otimes L_t^{\otimes m}\bigr),
  \]
  where \(g' : X_{T'} \to X_T\) is the induced morphism.  In particular the
  graded Hilbert function of \ref{mk-def:fiber_hilbert_function} is invariant
  under base change of the parameter scheme, and so is the Hilbert polynomial
  of \ref{mk-def:hilbert_polynomial}.
\end{lemma}
\begin{proof}
  \uses{mk-lem:pullback_moduleTensorPow, mk-lem:gamma_fiber_baseChange_field}
  \emph{Fibre pasting.}  Since \(X_{T'} = X_T \times_T T'\), the fibre of
  \(X_{T'} \to T'\) at \(t'\) is the base change of the fibre of
  \(X_T \to T\) at \(t\) along the residue field extension
  \(\kappa(t) \to \kappa(t')\):
  \[
    (X_{T'})_{t'}
    = X_{T'} \times_{T'} \operatorname{Spec} \kappa(t')
    = X_T \times_T \operatorname{Spec} \kappa(t')
    = (X_T)_t \times_{\operatorname{Spec} \kappa(t)}
      \operatorname{Spec} \kappa(t'),
  \]
  by pasting the cartesian square \(X_{T'} = X_T \times_T T'\) with the
  squares defining the two fibres, using the factorisation
  \(\operatorname{Spec}\kappa(t') \to \operatorname{Spec}\kappa(t) \to T\)
  of \(\operatorname{Spec}\kappa(t') \to T' \to T\).  Write
  \(u : (X_{T'})_{t'} \to (X_T)_t\) for the top morphism of the pasted
  cartesian square.

  \emph{Matching the twisted modules.}  Restriction to the fibre and
  pullback along \(g'\) (resp.\ along the projections to \(X\)) are
  pullbacks along two factorisations of one and the same composite
  morphism, so \((g'^*\mathcal{F})_{t'} = u^* (\mathcal{F}_t)\) and
  \(L_{t'} = u^* (L_t)\).  All modules in sight are quasi-coherent
  (\(\mathcal{F}\) is finitely presented, \(L\) is quasi-coherent, and
  quasi-coherence is preserved by module pullback), so by
  \ref{mk-lem:pullback_moduleTensorPow} applied to \(u\),
  \[
    (g'^*\mathcal{F})_{t'} \otimes L_{t'}^{\otimes m}
    \;\cong\;
    u^*(\mathcal{F}_t) \otimes \bigl(u^* (L_t)\bigr)^{\otimes m}
    \;\cong\;
    u^*\bigl(\mathcal{F}_t \otimes L_t^{\otimes m}\bigr).
  \]
  An isomorphism of sheaves of modules induces a \(\kappa(t')\)-linear
  isomorphism on global sections (the \(\kappa(t')\)-action on either side
  is restriction of the \(\Gamma((X_{T'})_{t'}, \mathcal{O})\)-action along
  the structural map \(\kappa(t') \to \Gamma((X_{T'})_{t'}, \mathcal{O})\)),
  so the left-hand side of the claim equals
  \(\dim_{\kappa(t')} \Gamma\bigl((X_{T'})_{t'},
  u^*(\mathcal{F}_t \otimes L_t^{\otimes m})\bigr)\).

  \emph{Base change of sections.}  By
  \ref{mk-lem:gamma_fiber_baseChange_field},
  \[
    \dim_{\kappa(t')} \Gamma\bigl((X_{T'})_{t'},\,
      u^*(\mathcal{F}_t \otimes L_t^{\otimes m})\bigr)
    =
    \dim_{\kappa(t)} \Gamma\bigl((X_T)_t,\,
      \mathcal{F}_t \otimes L_t^{\otimes m}\bigr),
  \]
  which is the claim.
\end{proof}

\section{The relative-divisor functor}
\label{mk-sec:div_functor}

In Grothendieck's quotient route to the Picard scheme --- the route set beside
the Milne--Koll\'ar one in \ref{mk-sec:fga_pic_setup}, and not the one the
the construction follows --- the Picard scheme arises from a family of effective
divisors modulo linear equivalence.  This section defines the parameter functor
of that family: the relative-effective-divisor functor \(\Div_{X/S}\), encoded
inside the Hilbert functor \(\Hilb_{X/S} = \Quot_{\mathcal O_X/X/S}\) (before the
Hilbert-polynomial decomposition) as the quotients of
\(\mathcal O_{X_T}\) whose kernel ideal is invertible.  The functor is a piece of
mathematics in its own right, and the divisor substrate it rests on is shared
with the committed route.

\begin{definition}[Family of relative effective divisors]\label{mk-def:div_family}

  \uses{mk-def:has_proper_support, mk-def:coherent_sheaf_flat, mk-def:IsLocallyTrivial}
  \dcref{custom}
  \textit{Source: \cite{kleiman-picard}, \S 3, Defs.~df:red and df:div (FGA Explained
  Ch.~9).}
  Let \(\pi : X \to S\) be a morphism of schemes and \(T \to S\) an
  \(S\)-scheme; write \(X_T := X \times_S T\).  A \emph{family of relative
  effective divisors on \(X_T/T\)} is a pair \((\mathcal F, q)\) where
  \begin{enumerate}
    \item \(\mathcal F\) is a finitely presented sheaf of
      \(\mathcal O_{X_T}\)-modules, flat over \(T\)
      (\ref{mk-def:coherent_sheaf_flat}) and with schematic support proper
      over \(T\) (\ref{mk-def:has_proper_support});
    \item \(q : \mathcal O_{X_T} \twoheadrightarrow \mathcal F\) is a
      surjective \(\mathcal O_{X_T}\)-linear homomorphism;
    \item the kernel ideal \(\mathcal I := \ker q\) is invertible, i.e.\
      locally trivial of rank one (\ref{mk-def:IsLocallyTrivial}).
  \end{enumerate}
  Two families \((\mathcal F, q)\) and \((\mathcal F', q')\) are
  \emph{equivalent} iff \(\ker q = \ker q'\), i.e.\ iff they cut out the same
  closed subscheme.

  Such a pair is precisely the structure-sheaf quotient
  \(q : \mathcal O_{X_T} \to \mathcal O_D\) of a relative effective divisor
  \(D \subseteq X_T\) in the sense of \cite{kleiman-picard} df:red --- a closed
  subscheme, flat over \(T\), whose ideal is invertible --- subject to two
  conditions that are automatic in the regime of interest: finite
  presentation of \(\mathcal O_D\) always holds (\(\mathcal O_D\) is the
  cokernel of \(\mathcal I \hookrightarrow \mathcal O_{X_T}\) with
  \(\mathcal I\) locally free of rank one), and properness of the support
  \(D\) over \(T\) holds whenever \(\pi\) is proper (a closed subscheme of
  the \(T\)-proper \(X_T\) is proper over \(T\)).
\end{definition}

\begin{lemma}[Invertibility is invariant under isomorphism]\label{mk-lem:locally_trivial_of_iso}
  \dcref{custom}

  \uses{mk-def:IsLocallyTrivial}
  Let \(M \cong N\) be isomorphic sheaves of \(\mathcal O_X\)-modules on a
  scheme \(X\).  If \(M\) is locally trivial of rank one
  (\ref{mk-def:IsLocallyTrivial}), then so is \(N\).
\end{lemma}
\begin{proof}Fix \(x \in X\) and choose an affine open \(U \ni x\) with an isomorphism
  \(M|_U \cong \mathcal O_U\).  Restriction along the open immersion
  \(U \subseteq X\) is functorial, so the given isomorphism \(M \cong N\)
  restricts to \(N|_U \cong M|_U \cong \mathcal O_U\).
\end{proof}

\begin{lemma}[A trivialising chart restricts to any smaller open]\label{mk-lem:IsLocallyTrivial_trivialization_of_le}
  \dcref{custom}

  \uses{mk-def:IsLocallyTrivial}
  \textit{Source: Stacks~\href{https://stacks.math.columbia.edu/tag/01HH}{01HH}
  (locally trivial modules), specialised to a chart shrink.}
  Let \(\mathcal M\) be an \(\mathcal O_X\)-module, \(V \subseteq U\) opens of
  \(X\), and \(t : \mathcal M|_U \cong \mathcal O_U\) a trivialisation. Then
  \(\mathcal M|_V \cong \mathcal O_V\).
\end{lemma}
\begin{proof}Restriction to \(V\) factors as restriction to \(U\) followed by
  restriction along the open immersion \(V \subseteq U\); applying \(t\) and
  using that the restriction of the structure sheaf along an open immersion
  is again the structure sheaf gives \(\mathcal M|_V \cong \mathcal O_U|_V
  \cong \mathcal O_V\).
\end{proof}

\begin{lemma}[Affine trivialising charts are cofinal]\label{mk-cor:IsLocallyTrivial_exists_affine_trivializing_le}
  \dcref{custom}

  \uses{mk-def:IsLocallyTrivial, mk-lem:IsLocallyTrivial_trivialization_of_le}
  Let \(\mathcal M\) be locally trivial of rank one
  (\ref{mk-def:IsLocallyTrivial}), \(x \in X\), and \(W\) an open neighbourhood
  of \(x\). Then there is an affine open \(V\) with \(x \in V \subseteq W\)
  and a trivialisation \(\mathcal M|_V \cong \mathcal O_V\).
\end{lemma}
\begin{proof}By local triviality there is an affine open \(U \ni x\) with a
  trivialisation of \(\mathcal M|_U\). Choose an affine open \(V\) with
  \(x \in V \subseteq U \cap W\) (affine opens form a basis of the topology
  of a scheme); restricting the trivialisation to \(V\) via
  \ref{mk-lem:IsLocallyTrivial_trivialization_of_le} gives the claim.
\end{proof}

\begin{lemma}[Flat cokernel keeps a tensored kernel inclusion injective]\label{mk-lem:flat_cokernel_rTensor_injective}

  \dcref{custom}
  \textit{Source: Stacks~\href{https://stacks.math.columbia.edu/tag/00HL}{00HL}.}
  Let \(R\) be a commutative ring and \(0 \to A \xrightarrow{f} B
  \xrightarrow{g} C \to 0\) an exact sequence of \(R\)-modules with \(C\)
  flat. Then for every \(R\)-module \(N\) the map
  \(f \otimes \mathrm{id}_N : A \otimes_R N \to B \otimes_R N\) is injective.
\end{lemma}
\begin{proof}Choose a free presentation \(0 \to K \xrightarrow{\iota} F_0
  \xrightarrow{\pi} N \to 0\) of \(N\); \(F_0\) is free, hence flat.
  Tensoring \(f\) and \(g\) with \(K\) and with \(F_0\), and tensoring
  \(\iota\) and \(\pi\) with \(A\), \(B\), \(C\), assembles a commutative
  \(3 \times 3\) diagram with exact rows
  \(A \otimes K \to A \otimes F_0 \to A \otimes N \to 0\) (and likewise for
  \(B\), \(C\)), by right-exactness of \(- \otimes_R -\); the two right-hand
  columns \(A \otimes K \to B \otimes K \to C \otimes K\) and
  \(A \otimes F_0 \to B \otimes F_0 \to C \otimes F_0\) are exact for the
  same reason, and moreover their first arrow is injective because \(F_0\)
  and \(C\) are flat.

  Given \(x \in A \otimes N\) with \((f \otimes \mathrm{id}_N)(x) = 0\),
  surjectivity of \(\pi\) lifts \(x\) to some \(y \in A \otimes F_0\) with
  \((\mathrm{id}_A \otimes \pi)(y) = x\). Commutativity sends the image of
  \(y\) in \(B \otimes F_0\) to \((f \otimes \mathrm{id}_N)(x) = 0\) in
  \(B \otimes N\), so by exactness of the middle row at \(B \otimes F_0\)
  this image equals the image of some \(z \in B \otimes K\). Pushing \(z\)
  down to \(C \otimes K\) and then right to \(C \otimes F_0\) gives \(0\)
  (the image of \(y\) in \(C \otimes F_0\) vanishes because
  \(g \circ f = 0\)), so injectivity of \(C \otimes K \to C \otimes F_0\)
  (flatness of \(C\)) forces the image of \(z\) in \(C \otimes K\) to
  already vanish, and exactness of the left column at \(B \otimes K\)
  produces \(w \in A \otimes K\) mapping to \(z\). Replacing \(y\) by
  \(y - (\iota \otimes \mathrm{id}_A)(w)\) leaves its image \(x\) under
  \(\mathrm{id}_A \otimes \pi\) unchanged and now maps to \(0\) in
  \(B \otimes F_0\) (it differed from an element mapping to \(z\) exactly by
  the image of \(w\)), so injectivity of \(A \otimes F_0 \to B \otimes F_0\)
  (flatness of \(F_0\)) forces this corrected \(y\) to be \(0\), whence
  \(x = (\mathrm{id}_A \otimes \pi)(y) = 0\).
\end{proof}

\begin{lemma}[Flat cokernel keeps a tensored kernel inclusion injective, left-tensor form]\label{mk-lem:flat_cokernel_lTensor_injective}

  \uses{mk-lem:flat_cokernel_rTensor_injective}
  \dcref{custom}
  \textit{Source: Stacks~\href{https://stacks.math.columbia.edu/tag/00HL}{00HL}.}
  In the situation of \ref{mk-lem:flat_cokernel_rTensor_injective}, the
  left-tensored map \(\mathrm{id}_N \otimes f : N \otimes_R A \to N \otimes_R
  B\) is injective for every \(R\)-module \(N\).
\end{lemma}
\begin{proof}\uses{mk-lem:flat_cokernel_rTensor_injective}
  The left-tensored map is conjugate to the right-tensored map
  \(f \otimes \mathrm{id}_N\) of \ref{mk-lem:flat_cokernel_rTensor_injective} by
  the two tensor-commutativity isomorphisms \(N \otimes A \cong A \otimes N\)
  and \(N \otimes B \cong B \otimes N\), hence injective as a composite of
  injective maps.
\end{proof}

\begin{lemma}[Base change of the ideal of a relative effective divisor]\label{mk-lem:relative_divisor_base_change}

  \uses{mk-def:coherent_sheaf_flat, mk-def:IsLocallyTrivial}
  \dcref{custom}
  \textit{Source: \cite{kleiman-picard}, \S 3, the functoriality argument following
  Def.~df:div.}
  Let
  \[
    \begin{array}{ccc}
      X' & \xrightarrow{\ g'\ } & X \\
      \downarrow f' & & \downarrow f \\
      S' & \xrightarrow{\ g\ } & S
    \end{array}
  \]
  be a cartesian square of schemes, \(q : E \twoheadrightarrow F\) an
  epimorphism of \(\mathcal O_X\)-modules with \(E\) quasi-coherent and
  \(F\) finitely presented and flat over \(S\)
  (\ref{mk-def:coherent_sheaf_flat}), and suppose the kernel
  \(\mathcal K := \ker q\) is locally trivial of rank one
  (\ref{mk-def:IsLocallyTrivial}).  Then \(\ker (g'^* q)\) is locally trivial
  of rank one; in fact the natural map \(g'^* \mathcal K \to \ker(g'^* q)\)
  is an isomorphism.
\end{lemma}
\begin{proof}
  \uses{mk-lem:IsLocallyTrivial_pullback, mk-lem:locally_trivial_of_iso,
    mk-lem:IsLocallyTrivial_trivialization_of_le,
    mk-lem:flat_cokernel_rTensor_injective}
  All three modules in the short exact sequence
  \(0 \to \mathcal K \to E \to F \to 0\) (exact because \(q\) is an
  epimorphism with kernel \(\mathcal K\)) are quasi-coherent: \(\mathcal K\)
  is locally isomorphic to the structure sheaf, \(F\) is finitely presented,
  and \(E\) is quasi-coherent by hypothesis.

  \emph{Charts.}  The claim is local on \(X'\).  Given \(x' \in X'\), choose
  an affine open \(U \subseteq S\) containing the common image
  \(f(g'(x')) = g(f'(x'))\), a trivialising chart of \(\mathcal K\) around
  \(g'(x')\) and inside it an affine open \(V \subseteq f^{-1}(U)\) around
  \(g'(x')\) (a basic open of the chart; \(\mathcal K|_V \cong \mathcal
  O_V\) persists under shrinking,
  \ref{mk-lem:IsLocallyTrivial_trivialization_of_le}), and an affine open
  \(U_t \subseteq g^{-1}(U)\) around \(f'(x')\).  As in
  \ref{mk-lem:coherent_flat_base_change}, the piece
  \(W := g'^{-1}(V) \cap f'^{-1}(U_t)\) is an affine open neighbourhood of
  \(x'\), namely the fibre product \(V \times_U U_t\) of the restricted
  cartesian square; writing \(R := \Gamma(S, U)\),
  \(R' := \Gamma(S', U_t)\) and \(A := \Gamma(X, V)\), it satisfies
  \(\Gamma(X', W) = A \otimes_R R'\), and for every quasi-coherent module
  \(M\) on \(X\),
  \(\Gamma(g'^* M, W) = \Gamma(M, V) \otimes_A \Gamma(X', W)
    = \Gamma(M, V) \otimes_R R'\)
  (the last identification by associativity of the tensor product).

  \emph{Exactness after base change.}  Sections over the affine \(V\) are
  exact on quasi-coherent modules (higher cohomology of a quasi-coherent
  module on an affine scheme vanishes), so
  \(0 \to \Gamma(\mathcal K, V) \to \Gamma(E, V) \to \Gamma(F, V) \to 0\)
  is an exact sequence of \(A\)-modules.  Tensoring over \(R\) with \(R'\)
  and splicing the long exact Tor sequence gives the exact sequence
  \[
    \operatorname{Tor}_1^R\bigl(\Gamma(F, V), R'\bigr)
      \to \Gamma(\mathcal K, V) \otimes_R R'
      \to \Gamma(E, V) \otimes_R R'
      \to \Gamma(F, V) \otimes_R R'
      \to 0,
  \]
  and the Tor term vanishes, equivalently the left-hand map stays injective
  after tensoring (\ref{mk-lem:flat_cokernel_rTensor_injective}), because
  \(\Gamma(F, V)\) is \(R\)-flat (\ref{mk-def:coherent_sheaf_flat} at the
  affine pair \((U, V)\)).  By the chart identification, this says that
  \(0 \to \Gamma(g'^* \mathcal K, W) \to \Gamma(g'^* E, W)
    \to \Gamma(g'^* F, W) \to 0\)
  is exact.

  \emph{Kernel identification.}  The pullback functor \(g'^*\) is right
  exact, so \(g'^* \mathcal K \to g'^* E \to g'^* F \to 0\) is exact, and
  the induced map \(g'^* \mathcal K \to \ker(g'^* q)\) is an epimorphism.
  On each chart \(W\) both sides are quasi-coherent modules on an affine
  scheme, and on sections over \(W\) the map is the isomorphism
  \(\Gamma(g'^* \mathcal K, W) \cong
    \ker\bigl(\Gamma(g'^* E, W) \to \Gamma(g'^* F, W)\bigr)
    = \Gamma(\ker (g'^* q), W)\)
  from the previous paragraph (the second identification is left exactness
  of sections).  A morphism of quasi-coherent modules that is an
  isomorphism on sections over each member of an affine open cover is an
  isomorphism, so \(g'^* \mathcal K \cong \ker(g'^* q)\).

  Finally \(g'^* \mathcal K\) is locally trivial of rank one by
  \ref{mk-lem:IsLocallyTrivial_pullback}, hence so is \(\ker(g'^* q)\) by
  \ref{mk-lem:locally_trivial_of_iso}.
\end{proof}

\begin{remark}
  \label{mk-rem:relative_divisor_base_change_qcoh_removable}
  \dcref{custom}
  \uses{mk-lem:relative_divisor_base_change}
  The hypothesis that \(E\) is quasi-coherent in
  \ref{mk-lem:relative_divisor_base_change} is removable: \(E\) is an
  extension of the quasi-coherent \(F\) by the quasi-coherent
  \(\mathcal K = \ker q\) (\(\mathcal K\) is locally isomorphic to the
  structure sheaf, hence quasi-coherent), and quasi-coherent modules form a
  weak Serre subcategory of all \(\mathcal O_X\)-modules, in particular one
  closed under extensions (Stacks 01LA). Once extension-closure of
  quasi-coherence is available, the hypothesis on \(E\) can be dropped in
  favour of this derivation.
\end{remark}

\begin{definition}
[Kernel--pullback comparison map]
  \label{mk-def:pullback_kernel_comparison}
  \dcref{custom}

  \uses{mk-lem:modules_pullback_mathlib}
  Let \(g' : X' \to X\) be a morphism of schemes and \(q : E \to F\) a
  morphism of \(\mathcal O_X\)-modules. Write \(\iota_q : \ker q \to E\) for
  the kernel inclusion. The \emph{kernel--pullback comparison map}
  \[
    \kappa_{g',q} : g'^{*}(\ker q) \longrightarrow \ker(g'^{*} q)
  \]
  is the unique morphism with
  \(\kappa_{g',q} \circ \iota_{g'^{*}q} = g'^{*}(\iota_q)\); it exists
  because \(g'^{*}\) preserves the zero morphism, being a left adjoint
  (\ref{mk-lem:modules_pullback_mathlib}).
\end{definition}

\begin{lemma}[Local triviality transports along an isomorphic kernel comparison]\label{mk-lem:pullback_kernel_isLocallyTrivial_of_isIso_kernelComparison}
  \dcref{custom}

  \uses{mk-def:pullback_kernel_comparison, mk-lem:IsLocallyTrivial_pullback,
    mk-lem:locally_trivial_of_iso}
  Let \(g' : X' \to X\) and \(q : E \to F\) be as in
  \ref{mk-def:pullback_kernel_comparison}, and suppose the comparison map
  \(\kappa_{g',q}\) is an isomorphism. If \(\ker q\) is locally trivial of
  rank one, then so is \(\ker(g'^{*} q)\).
\end{lemma}
\begin{proof}\uses{mk-lem:IsLocallyTrivial_pullback, mk-lem:locally_trivial_of_iso}
  By \ref{mk-lem:IsLocallyTrivial_pullback}, \(g'^{*}(\ker q)\) is locally
  trivial of rank one; transporting along the isomorphism \(\kappa_{g',q}\)
  via \ref{mk-lem:locally_trivial_of_iso} gives the same for
  \(\ker(g'^{*} q)\).
\end{proof}

\begin{lemma}[Basis-local criterion for a monomorphism of sheaves of modules]\label{mk-lem:modules_mono_of_injective_app_of_isBasis}
  \dcref{custom}

  Let \(B\) be a basis of opens of a scheme \(X\) and let
  \(\varphi : M \to N\) be a morphism of sheaves of \(\mathcal O_X\)-modules
  such that \(\varphi\) is injective on sections over every member of \(B\).
  Then \(\varphi\) is a monomorphism.
\end{lemma}
\begin{proof}Injectivity on a basis of opens gives injectivity of \(\varphi\) on
  stalks, hence the underlying morphism of sheaves of abelian groups is a
  monomorphism; the faithful forgetful functor from sheaves of modules to
  sheaves of abelian groups reflects monomorphisms.
\end{proof}

\begin{lemma}[The kernel--pullback comparison is an epimorphism when \(q\) is]\label{mk-lem:pullback_kernel_comparison_epi}
  \dcref{custom}

  \uses{mk-def:pullback_kernel_comparison}
  Let \(g' : X' \to X\) and \(q : E \twoheadrightarrow F\) be as in
  \ref{mk-def:pullback_kernel_comparison}, with \(q\) an epimorphism. Then the
  comparison map \(\kappa_{g',q} : g'^{*}(\ker q) \to \ker(g'^{*} q)\) is an
  epimorphism.
\end{lemma}
\begin{proof}\(q\) epi identifies \(F\) with the cokernel of the kernel inclusion
  \(\ker q \hookrightarrow E\); the pullback functor \(g'^{*}\) is a left
  adjoint (\ref{mk-lem:modules_pullback_mathlib}), hence preserves cokernels,
  so \(g'^{*}(\ker q) \to g'^{*}E \to g'^{*}F \to 0\) stays exact. Exactness
  at \(g'^{*}E\) says exactly that the induced kernel-comparison lift
  \(\kappa_{g',q}\) is an epimorphism.
\end{proof}

\begin{lemma}[The kernel--pullback comparison is an isomorphism once it is monic]\label{mk-lem:pullback_kernel_comparison_iso_of_mono}
  \dcref{custom}

  \uses{mk-def:pullback_kernel_comparison, mk-lem:pullback_kernel_comparison_epi}
  Let \(g' : X' \to X\) and \(q : E \twoheadrightarrow F\) be as in
  \ref{mk-def:pullback_kernel_comparison}, with \(q\) an epimorphism, and
  suppose the pullback \(g'^{*}(\iota_q) : g'^{*}(\ker q) \to g'^{*}E\) of
  the kernel inclusion is a monomorphism. Then the comparison map
  \(\kappa_{g',q}\) is an isomorphism.
\end{lemma}
\begin{proof}\uses{mk-lem:pullback_kernel_comparison_epi}
  The comparison satisfies \(\kappa_{g',q} \circ \iota_{g'^{*}q} =
  g'^{*}(\iota_q)\) (\ref{mk-def:pullback_kernel_comparison}); since the
  right-hand side is monic by hypothesis and \(\iota_{g'^{*}q}\) is (as a
  kernel inclusion) automatically monic, \(\kappa_{g',q}\) is monic.
  Combined with the epimorphism half of
  \ref{mk-lem:pullback_kernel_comparison_epi}, monic \(+\) epic \(=\) iso in
  the abelian category of \(\mathcal O_{X'}\)-modules.
\end{proof}

\begin{remark}
  \label{mk-rem:relative_divisor_base_change_kernel_comparison}
  \dcref{custom}
  \uses{mk-def:pullback_kernel_comparison,
    mk-lem:pullback_kernel_isLocallyTrivial_of_isIso_kernelComparison,
    mk-lem:flat_cokernel_rTensor_injective, mk-lem:flat_cokernel_lTensor_injective,
    mk-lem:modules_mono_of_injective_app_of_isBasis,
    mk-lem:pullback_kernel_comparison_epi,
    mk-lem:pullback_kernel_comparison_iso_of_mono}
  The conclusion of \ref{mk-lem:relative_divisor_base_change} factors through
  the comparison map of \ref{mk-def:pullback_kernel_comparison}: local
  triviality of \(\ker(g'^{*} q)\) follows from
  \ref{mk-lem:pullback_kernel_isLocallyTrivial_of_isIso_kernelComparison} as
  soon as \(\kappa_{g',q}\) is known to be an isomorphism, and this
  factoring is now itself proved rather than merely anticipated:
  \(\kappa_{g',q}\) is automatically an epimorphism whenever \(q\) is
  (\ref{mk-lem:pullback_kernel_comparison_epi}, right exactness of
  \(g'^{*}\)), so by \ref{mk-lem:pullback_kernel_comparison_iso_of_mono} the
  entire remaining obligation is the single monomorphism fact that
  \(g'^{*}\) keeps the kernel inclusion \(\ker q \hookrightarrow E\)
  injective --- exactly the Tor-vanishing content isolated in the
  ``Exactness after base change'' paragraph of the proof of
  \ref{mk-lem:relative_divisor_base_change}
  (\ref{mk-lem:flat_cokernel_rTensor_injective} and
  \ref{mk-lem:flat_cokernel_lTensor_injective}).
  That monomorphism fact is established scheme-theoretically via the
  basis-local criterion \ref{mk-lem:modules_mono_of_injective_app_of_isBasis}:
  the affine pieces \(W = g'^{-1}(V) \cap f'^{-1}(U_t)\) of the proof of
  \ref{mk-lem:relative_divisor_base_change} cover \(X'\) but need not form a
  topological basis (the scheme-theoretic fibre product carries a
  topology finer than \(|X| \times_{|S|} |S'|\)); descending the
  per-piece section injectivity to the basic opens of each piece --- which
  do form a basis of \(X'\), by flatness of the localisation at a section
  --- completes the reduction, so \(\kappa_{g',q}\) is an isomorphism in
  general (\ref{mk-lem:relative_divisor_base_change_comparison_iso}).
\end{remark}

\begin{lemma}[The kernel--pullback comparison is an isomorphism under base-flatness]\label{mk-lem:relative_divisor_base_change_comparison_iso}
  \dcref{custom}

  \uses{mk-def:pullback_kernel_comparison, mk-lem:relative_divisor_base_change}
  In the situation of \ref{mk-lem:relative_divisor_base_change}, the
  kernel--pullback comparison map
  \(\kappa_{g',q} : g'^{*}(\ker q) \to \ker(g'^{*} q)\)
  (\ref{mk-def:pullback_kernel_comparison}) is itself an isomorphism.
\end{lemma}
\begin{proof}\uses{mk-lem:pullback_kernel_comparison_epi,
    mk-lem:pullback_kernel_comparison_iso_of_mono}
  The comparison map \(\kappa_{g',q}\) is an epimorphism whenever \(q\)
  is (\ref{mk-lem:pullback_kernel_comparison_epi}, right exactness of
  \(g'^{*}\)); the proof of \ref{mk-lem:relative_divisor_base_change} shows
  that \(g'^{*}\) keeps the kernel inclusion \(\ker q \hookrightarrow E\)
  monic, so \ref{mk-lem:pullback_kernel_comparison_iso_of_mono} upgrades
  the epimorphism \(\kappa_{g',q}\) to an isomorphism. This is the
  isomorphism form of \ref{mk-lem:relative_divisor_base_change} consumed by
  the Abel map and by the pullback of families of relative effective
  divisors (\ref{mk-def:div_functor}).
\end{proof}

\begin{lemma}[The structure sheaf is quasi-coherent]\label{mk-lem:unit_isQuasicoherent}
  \dcref{custom}

  For every scheme \(X\), the unit module \(\mathcal O_X\) (the unit object
  of \(X.\mathrm{Modules}\)) is quasi-coherent.
\end{lemma}
\begin{proof}The unit module carries the canonical one-generator, no-relation
  presentation given by the identity on \(\mathcal O_X\) over the
  one-member cover \(\{X\}\); this is a quasi-coherence datum.
\end{proof}

\begin{definition}[Relative-divisor functor]\label{mk-def:div_functor}

  \uses{mk-def:div_family, mk-lem:relative_divisor_base_change,
    mk-lem:pullback_finite_presentation, mk-lem:coherent_flat_base_change,
    mk-lem:proper_support_base_change, mk-lem:unit_isQuasicoherent,
    mk-lem:pullback_isQuasicoherent_hom}
  \dcref{custom}
  \textit{Source: \cite{kleiman-picard}, \S 3, Def.~df:div (FGA Explained Ch.~9).}
  Let \(\pi : X \to S\) be a morphism of schemes.  The
  \emph{relative-divisor functor}
  \[
    \Div_{X/S} : (\mathrm{Sch}/S)^{op} \longrightarrow \mathbf{Set}
  \]
  sends an \(S\)-scheme \(T \to S\) to the set of equivalence classes of
  families of relative effective divisors on \(X_T/T\)
  (\ref{mk-def:div_family}) --- equivalently, the set of relative effective
  divisors on \(X_T/T\) --- and an \(S\)-morphism \(T' \to T\) to the
  pullback of families along the induced morphism
  \(g' : X_{T'} \to X_T\).  The action is well defined: pullback of modules
  preserves surjectivity (it is right exact), finite presentation
  (\ref{mk-lem:pullback_finite_presentation}), flatness over the base
  (\ref{mk-lem:coherent_flat_base_change}) and proper support
  (\ref{mk-lem:proper_support_base_change}), and it preserves the
  invertibility of the kernel ideal because the pulled-back kernel is the
  kernel of the pulled-back quotient (\ref{mk-lem:relative_divisor_base_change},
  applicable since the source of \(q\) is always a pullback of the unit
  module \(\mathcal O_X\), quasi-coherent by \ref{mk-lem:unit_isQuasicoherent}
  together with pullback-stability of quasi-coherence
  (\ref{mk-lem:pullback_isQuasicoherent_hom})).  Functoriality holds because
  the pullback of modules is pseudofunctorial in the scheme morphism.
\end{definition}

\section{The Grassmannian scheme}
\label{mk-sec:grassmannian_scheme}

The Quot construction proceeds by embedding \(\Quot^{\Phi,L}_{E/X/S}\) as a
locally closed subfunctor of a Grassmannian over \(S\). Since Mathlib carries
no Grassmannian \emph{scheme} (only a finite-rank-quotient variety in
linear algebra), we package the construction here.

\begin{definition}[Grassmannian scheme]\label{mk-def:grassmannian_scheme}
  \dcref{custom}

  \uses{mk-def:is_locally_free_of_rank}
  \textit{Source: \cite{nitsure-hilbert-quot}, \S 1, Exercise (2) and "Construction of
  Grassmannian" (FGA Explained Ch.~5).}
  Let \(S\) be a locally noetherian scheme and let \(V\) be a locally free
  \(\mathcal{O}_S\)-module of rank \(r\). For an integer \(1 \le d \le r\), the
  \emph{Grassmannian of rank-\(d\) quotients of \(V\)} is the functor
  \[
    \mathrm{Grass}(V, d) : (\mathrm{Sch}/S)^{op} \longrightarrow \mathbf{Set},
    \qquad
    T \longmapsto \{ \langle \mathcal{F}, q\rangle :
    q : V_T \twoheadrightarrow \mathcal{F},\; \mathcal{F}
    \text{ locally free of rank } d \}
  \]
  with equivalence \(\ker(q) = \ker(q')\); the rank-\(d\) local freeness of
  \(\mathcal{F}\) is as in \ref{mk-def:is_locally_free_of_rank}. Over a locally
  noetherian base this functor agrees with the Quot functor
  (\ref{mk-def:quot_functor}) in the special case \(X = S\), \(E = V\), constant
  Hilbert polynomial \(\Phi = d\): a flat finitely presented module whose fibre
  dimension is \(d\) at every point is locally free of rank \(d\).
\end{definition}

\begin{lemma}[Representability predicate]
  \label{mk-lem:functor_is_representable_mathlib}
  \dcref{custom}

  Let \(\mathcal{C}\) be a category and \(F : \mathcal{C}^{op} \to \mathbf{Set}\)
  a presheaf of sets. The predicate \(F.\mathrm{IsRepresentable}\) asserts that
  \(F\) is representable, i.e.\ that there exists an object \(Y\) of
  \(\mathcal{C}\) together with a natural bijection
  \(\mathrm{Hom}_{\mathcal{C}}(-, Y) \cong F\) (the Yoneda bijection).
\end{lemma}

\begin{lemma}[Local representability of Zariski sheaves]
  \label{mk-lem:local_representability_mathlib}
  \dcref{custom}

  Let \(G\) be a set-valued sheaf on the big Zariski site of schemes,
  \(\{X_i\}\) a family of schemes and \(f_i : h_{X_i} \to G\) a family of
  relatively representable open immersions of presheaves that is jointly
  locally surjective.  Then \(G\) is representable, by the scheme glued from
  the \(X_i\) along the pairwise pullbacks \(h_{X_i} \times_G h_{X_j}\).
\end{lemma}

\begin{lemma}[Representability of the Grassmannian, trivialised case]\label{mk-lem:grassmannian_representable_free}
  \dcref{custom}

  \uses{mk-def:grassmannian_scheme, mk-thm:grassmannian_universal_property,
    mk-lem:functor_is_representable_mathlib}
  \textit{Source: \cite{nitsure-hilbert-quot}, \S 1, Exercise (2).}
  Let \(S\) be a locally noetherian scheme, \(V\) an \(\mathcal{O}_S\)-module
  equipped with a global trivialisation \(V \cong \mathcal{O}_S^{\oplus r}\),
  and \(1 \le d \le r\). Then the functor \(\mathrm{Grass}(V, d)\) of
  \ref{mk-def:grassmannian_scheme} is representable by the base change
  \(S \times_{\Spec \mathbb{Z}} \mathrm{Gr}(d,r)\) of the absolute
  Grassmannian, with structure morphism the first projection.
\end{lemma}

\begin{proof}\uses{mk-thm:grassmannian_universal_property, mk-def:grassmannian_scheme}
  Since \(\Spec \mathbb{Z}\) is the terminal scheme, the base change
  \(S \times_{\Spec\mathbb{Z}} \mathrm{Gr}(d,r)\) is the binary product, and
  for an \(S\)-scheme \(T \to S\) the \(S\)-morphisms
  \(T \to S \times \mathrm{Gr}(d,r)\) (over the first projection) correspond,
  by the universal property of the product, to the bare scheme morphisms
  \(T \to \mathrm{Gr}(d,r)\); these are in natural bijection with
  \(\mathrm{Grass}(r,d)(T)\) by \ref{mk-thm:grassmannian_universal_property}.
  It therefore suffices to identify \(\mathrm{Grass}(V,d)(T)\) with
  \(\mathrm{Grass}(r,d)(T)\), naturally in \(T\).  A \(T\)-point of
  \(\mathrm{Grass}(V,d)\) is a rank-\(d\) locally free quotient
  \(q : V_T \twoheadrightarrow \mathcal{F}\); composing with the pullback of
  the trivialisation and with the canonical identification
  \((\mathcal{O}_S^{\oplus r})_T \cong \mathcal{O}_T^{\oplus r}\) of the
  pullback of a free module re-presents it as a rank-\(d\) quotient of
  \(\mathcal{O}_T^{\oplus r}\), i.e.\ a \(T\)-point of
  \(\mathrm{Grass}(r,d)\); the inverse composes with the inverse
  identifications, and both directions preserve the equivalence
  \(\ker q = \ker q'\).  Naturality in \(T\) is the compatibility of the
  free-module identification with composite pullbacks (the pseudofunctor
  coherence of the module pullback on free modules), and the equivalence
  classes match because the re-presentation composes with isomorphisms only.
\end{proof}

\begin{definition}[Zariski sheaf axiom on \(\mathrm{Sch}/S\)]\label{mk-def:zariski_sheaf_over}
  \dcref{custom}

  A presheaf \(F : (\mathrm{Sch}/S)^{op} \to \mathbf{Set}\) \emph{satisfies
  the Zariski sheaf axiom} when for every \(S\)-scheme \(T \to S\) and every
  open cover \(\{W_k\}\) of \(T\), every family of sections
  \(x_k \in F(T|_{W_k})\) that agrees on the pairwise intersections
  (\(x_k|_{W_k \cap W_l} = x_l|_{W_k \cap W_l}\) for all \(k, l\)) glues to a
  unique section \(x_0 \in F(T)\) with \(x_0|_{W_k} = x_k\).  This is the
  sheaf condition for the big Zariski site of \(S\), stated against open
  covers of the total space; every representable presheaf satisfies it.
\end{definition}

\begin{theorem}[Zariski descent of representability]\label{mk-thm:representability_zariski_descent}
  \dcref{custom}

  \uses{mk-def:zariski_sheaf_over, mk-lem:functor_is_representable_mathlib,
    mk-lem:local_representability_mathlib}
  Let \(S\) be a scheme, \(F : (\mathrm{Sch}/S)^{op} \to \mathbf{Set}\) a
  presheaf satisfying the Zariski sheaf axiom
  (\ref{mk-def:zariski_sheaf_over}), and \(S = \bigcup_i U_i\) an open cover
  such that for every \(i\) the restriction of \(F\) to
  \(\mathrm{Sch}/U_i\) (along \(T \mapsto (T \to U_i \hookrightarrow S)\))
  is representable by a \(U_i\)-scheme \(Y_i\).  Then \(F\) is representable
  by an \(S\)-scheme.
\end{theorem}

\begin{proof}\uses{mk-def:zariski_sheaf_over, mk-lem:local_representability_mathlib}
  Form the \emph{total functor} \(G\) of \(F\) on the absolute category of
  schemes: a point of \(G(T)\) is a morphism \(a : T \to S\) together with
  a family of morphisms \(\gamma_i : T|_{a^{-1}(U_i)} \to Y_i\) over
  \(U_i\) whose classified sections
  \(x_i \in F(T|_{a^{-1}(U_i)})\) agree on pairwise overlaps.  By the
  sheaf axiom of \(F\) applied to the cover \(\{a^{-1}(U_i)\}\) of \(T\),
  such a datum is the same as a pair \((a, x)\) with \(x \in F(T, a)\)
  (the compatible classified sections glue to a unique \(x\), and \(x\)
  restricts back to the \(x_i\)); the encoding by classifying morphisms
  realises \(G\) as a set-valued functor, restriction acting by composition
  on \(a\) and by restriction of the \(\gamma_i\).

  \(G\) is a sheaf for the big Zariski topology: a covering family may be
  replaced by the family of image opens (both generate the same sieve), and
  for a cover of \(T\) by opens \(\{W_k\}\) with compatible points
  \((a_k, x_k) \in G(W_k)\), the base morphisms glue to \(a : T \to S\)
  because morphisms of schemes glue along open covers, and the sections
  glue by the sheaf axiom of \(F\) at \((T, a)\); uniqueness follows from
  locality of morphism equality and the uniqueness half of the sheaf axiom.

  Each \(Y_i\) maps to \(G\) by the \emph{chart}
  \(c_i : h_{Y_i} \to G\) sending \(t : T \to Y_i\) to the point with
  base \(T \to Y_i \to U_i \subseteq S\) classified by the section that
  \(t\) represents.  The pullback of \(c_i\) against an arbitrary point
  \((a, x) : h_T \to G\) is the open subscheme \(a^{-1}(U_i) \subseteq T\)
  (a morphism \(V \to T\) lifts along \(c_i\) iff its base lands in
  \(U_i\), and then the lift is the unique classifying morphism of the
  restricted section), so the \(c_i\) are relatively representable open
  immersions; they are jointly locally surjective because every point of
  \(G(T)\) restricts on the cover \(\{a^{-1}(U_i)\}\) to the chart image
  of its classifying morphisms.  By
  \ref{mk-lem:local_representability_mathlib} the functor \(G\) is
  representable by a scheme \(\hat{Y}\).

  Finally, the tautological point of \(G(\hat{Y})\) corresponding to
  \(\mathrm{id}_{\hat{Y}}\) has base \(\hat{a} : \hat{Y} \to S\), and for
  \(T \to S\) the identification \(G(T) \cong \coprod_{a} F(T, a)\)
  restricts to a natural bijection
  \(\mathrm{Hom}_{\mathrm{Sch}/S}(T, (\hat{Y}, \hat{a})) \cong F(T)\):
  a morphism \(u : T \to \hat{Y}\) over \(S\) corresponds to a point of
  \(G(T)\) with base \(T \to S\), i.e.\ to a section of \(F(T)\), and
  naturality is the restriction-compatibility of the glued sections.  Hence
  \(F\) is representable by the \(S\)-scheme \((\hat{Y}, \hat{a})\).
\end{proof}

\begin{lemma}[The Grassmannian functor is a Zariski sheaf]\label{mk-lem:grassmannian_zariski_sheaf}
  \dcref{custom}

  \uses{mk-def:grassmannian_scheme, mk-def:zariski_sheaf_over,
    mk-def:scheme_modules_glue}
  Let \(S\) be locally noetherian, \(V\) an \(\mathcal{O}_S\)-module and
  \(d \ge 0\).  The functor \(\mathrm{Grass}(V, d)\) of
  \ref{mk-def:grassmannian_scheme} satisfies the Zariski sheaf axiom
  (\ref{mk-def:zariski_sheaf_over}).
\end{lemma}

\begin{proof}\uses{mk-def:scheme_modules_glue, mk-def:grassmannian_scheme,
    mk-def:gr_modules_glueRestrictionIso, mk-lem:gr_modules_glue_unique}
  Let \(T \to S\), let \(\{W_k\}\) be an open cover of \(T\), and let
  \(\langle \mathcal{F}_k, q_k \rangle\) be rank-\(d\) locally free
  quotients of \(V_{T|_{W_k}}\) agreeing on overlaps.  Agreement means:
  there are isomorphisms
  \(f_{kl} : \mathcal{F}_k|_{W_k \cap W_l} \cong \mathcal{F}_l|_{W_k \cap W_l}\)
  commuting with the restricted quotient maps.  Such an isomorphism is
  \emph{unique} when it exists (it is determined on the image of the
  epimorphism \(q_k\)); consequently the \(f_{kl}\) automatically satisfy
  the cocycle condition on triple overlaps -- pulled back to the
  scheme-theoretic overlaps \(W_k \times_T W_l\) of the glue datum attached
  to the cover, both composites of transition transports commute with the
  restricted quotient maps, and gluing isomorphisms are unique.  Descent of
  sheaves of modules along the glue datum of the cover
  (\ref{mk-def:scheme_modules_glue}) glues the \(\mathcal{F}_k\) to a sheaf
  \(\mathcal{F}\) on the glued scheme, whose chart restrictions recover the
  \(\mathcal{F}_k\) (\ref{mk-def:gr_modules_glueRestrictionIso}); the
  transposed quotient maps satisfy the descent-equalizer compatibility, so
  they lift to \(q : V \to \mathcal{F}\) out of the pulled-back test module,
  with chart restrictions the \(q_k\).  Epimorphy is detected on the chart
  cover by joint faithfulness of the chart restrictions
  (\ref{mk-lem:gr_modules_glue_unique}), and rank-\(d\) local freeness is
  local, so transporting along the canonical isomorphism between the glued
  scheme and \(T\) produces the glued point of
  \(\mathrm{Grass}(V,d)(T)\); it restricts to the given classes by the
  chart recovery.  Uniqueness: if \(\langle \mathcal{F}, q \rangle\) and
  \(\langle \mathcal{F}', q' \rangle\) both restrict to
  \(\langle \mathcal{F}_k, q_k \rangle\) up to equivalence, the two kernels
  annihilate each other's quotient map locally on the cover, hence globally
  (the sections of a sheaf are detected on a cover), and the two
  epimorphism descents \(\mathcal{F} \to \mathcal{F}'\),
  \(\mathcal{F}' \to \mathcal{F}\) are mutually inverse by
  epi-cancellation, so the two glued points are equivalent.
\end{proof}

\begin{theorem}[Representability of the Grassmannian]\label{mk-thm:grassmannian_representable}
  \dcref{custom}

  \uses{mk-def:grassmannian_scheme, mk-def:is_locally_free_of_rank,
    mk-lem:grassmannian_representable_free,
    mk-thm:representability_zariski_descent, mk-lem:grassmannian_zariski_sheaf,
    mk-lem:functor_is_representable_mathlib, mk-def:gr_glued_scheme, mk-lem:gr_separated,
    mk-lem:gr_proper}
  \textit{Source: \cite{nitsure-hilbert-quot}, \S 1, "Construction of Grassmannian" (FGA
  Explained Ch.~5).}
  Let \(S\) be a noetherian scheme, \(V\) a locally free \(\mathcal{O}_S\)-module
  of rank \(r\), and \(1 \le d \le r\). The functor \(\mathrm{Grass}(V,d)\) of
  \ref{mk-def:grassmannian_scheme} is representable by a smooth projective
  \(S\)-scheme \(\mathrm{Gr}_S(V, d) \to S\) of relative dimension \(d(r-d)\),
  equipped with a tautological quotient
  \(\pi^* V \twoheadrightarrow \mathcal{U}\) of rank \(d\). The determinant
  line bundle \(\det(\mathcal{U})\) is relatively very ample, and the
  associated linear system gives a closed embedding
  \(\mathrm{Gr}_S(V, d) \hookrightarrow
  \mathbb{P}_S\!\left(\bigwedge^d V\right)\) -- the \emph{Pl\"ucker embedding}.

\end{theorem}
\begin{proof}\uses{mk-lem:grassmannian_representable_free,
    mk-thm:representability_zariski_descent, mk-lem:grassmannian_zariski_sheaf}
  Over \(S = \Spec \mathbb{Z}\) and \(V = \mathcal{O}_S^{\oplus r}\) the
  Grassmannian \(\mathrm{Gr}(r, d)\) is constructed by gluing
  \(\binom{r}{d}\) affine charts \(U^I = \Spec\mathbb{Z}[X^I]\) indexed by
  size-\(d\) subsets \(I \subseteq \{1, \dots, r\}\). The chart \(U^I\)
  carries the matrix \(X^I\) of size \(d \times r\) whose \(I\)-th \(d \times d\)
  minor is the identity, with the remaining \(d(r-d)\) entries free
  variables (so \(U^I \cong \mathbb{A}^{d(r-d)}_{\mathbb{Z}}\)). The
  transition maps \(\theta_{I,J} : U^I \cap U^J \to U^J \cap U^I\) defined
  by \(X^J \mapsto (X^I_J)^{-1} X^I\) on the locus where \(\det X^I_J\) is
  invertible satisfy the cocycle condition
  \(\theta_{I,K} = \theta_{I,J} \circ \theta_{J,K}\). The glued
  \(\mathrm{Gr}(r,d) \to \Spec\mathbb{Z}\) is smooth of relative dimension
  \(d(r-d)\). Separatedness is checked from the diagonal
  \(\Delta_{I,J} \subset U^I \times U^J\) cut out by \(X^J_I X^I - X^J = 0\).
  Properness follows from the valuative criterion for DVRs: given
  \(\varphi : \Spec K \to \mathrm{Gr}(r,d)\) over a DVR \(R\) with fraction
  field \(K\), factoring through some chart \(U^I\) via a ring map
  \(f : \mathbb{Z}[X^I] \to K\), choose \(J\) minimising the valuation
  \(\nu(f(P^I_J))\) of the minor; on the chart \(U^J\) all entries have
  non-negative valuation, so \(f\) extends uniquely to \(\Spec R\).

  The universal quotient \(u : \mathcal{O}_{\mathrm{Gr}(r,d)}^{\oplus r}
  \twoheadrightarrow \mathcal{U}\) is given on \(U^I\) by the matrix \(X^I\)
  and glued by \(g_{I,J} = (X^I_J)^{-1} \in GL_d(U^I \cap U^J)\). The
  determinant line bundle \(\det\mathcal{U}\) has transition functions
  \(\det g_{I,J} = 1/P^I_J\); the global sections
  \(\sigma_I \in \Gamma(\mathrm{Gr}(r,d), \det\mathcal{U})\) defined by
  \(\sigma_I|_{U^J} = P^J_I\) form a base-point-free linear system separating
  points over \(\Spec\mathbb{Z}\), giving an embedding into
  \(\mathbb{P}^{\binom{r}{d}-1}_{\mathbb{Z}}\), which is closed by
  properness. Representability of \(\mathrm{Grass}(r,d)\) is then verified
  on the universal pair \((\mathrm{Gr}(r,d), u)\): a quotient
  \(q : \mathcal{O}_T^{\oplus r} \twoheadrightarrow \mathcal{F}\) of
  rank \(d\) on \(T\) is everywhere locally trivialised; on the locus where
  the columns indexed by \(I\) are a basis the matrix of \(q\) has the form
  \(X^I\), defining the classifying morphism \(T \to U^I \subset
  \mathrm{Gr}(r,d)\). Base change from \(\mathbb{Z}\) to a general noetherian
  base \(S\) with a vector bundle \(V\) is direct: trivialise \(V\) on an open
  cover and glue. For non-trivial \(V\) the same construction produces
  \(\mathrm{Gr}_S(V, d)\) with Pl\"ucker embedding into
  \(\mathbb{P}_S(\bigwedge^d V)\).
\end{proof}
\section{Projective space, the Serre twist, and projective morphisms}
\label{mk-sec:projective_vocabulary}

Representability of the Quot functor holds for a \emph{projective} structural
morphism carrying a relatively very ample line bundle
(\cite{nitsure-hilbert-quot} \S 5); properness alone does not suffice.  This
section builds the ambient vocabulary: relative projective space, the Serre
twisting sheaf, and the projective-morphism predicate together with its
stability properties.

\begin{definition}
[Relative projective space]
  \label{mk-def:projective_space}
  \dcref{custom}
  Let \(S\) be a scheme and let \(n\) be an index set.  The \emph{relative
  projective space} \(\mathbb{P}(n; S)\) is the fibre product
  \[
    \mathbb{P}(n; S) \;=\; S \times_{\operatorname{Spec} \mathbb{Z}}
      \operatorname{Proj} \mathbb{Z}[X_i : i \in n],
  \]
  where the polynomial ring carries the total-degree grading.  Its
  \emph{structural morphism} \(\mathbb{P}(n; S) \to S\) is the first
  projection.  For \(n = \{0, \dots, d\}\) this is the usual projective
  \(d\)-space \(\mathbb{P}^d_S\).
\end{definition}

\begin{lemma}
[Degree-zero part of the homogeneous coordinate ring]
  \label{mk-lem:projective_space_gradeZero}
  \dcref{custom}
  The degree-zero part of the total-degree grading on
  \(\mathbb{Z}[X_i : i \in n]\) is the subring of constants: the structural
  map \(\mathbb{Z} \to (\mathbb{Z}[X_i : i \in n])_0\) is a ring isomorphism.
  Consequently \(\operatorname{Spec}\) of the degree-zero part is a terminal
  scheme.
\end{lemma}

\begin{proof}
  A polynomial is homogeneous of degree \(0\) exactly when it is a constant,
  so the degree-zero submodule is the image of \(\mathbb{Z}\) under the
  (injective) constants map.
\end{proof}

\begin{lemma}
[Properness of the integral model]
  \label{mk-lem:proj_toSpecZero_proper}
  \dcref{custom}
  For a \(\mathbb{N}\)-graded ring \(A\) of finite type over its degree-zero
  part \(A_0\), the structural morphism
  \(\operatorname{Proj} A \to \operatorname{Spec} A_0\) is proper (separated,
  universally closed, and of finite type).
\end{lemma}

\begin{lemma}
[Properness of relative projective space]
  \label{mk-lem:projective_space_proper}
  \dcref{custom}
  \uses{mk-def:projective_space, mk-lem:projective_space_gradeZero,
    mk-lem:proj_toSpecZero_proper}
  For a finite index set \(n\), the structural morphism
  \(\mathbb{P}(n; S) \to S\) is proper (hence separated, universally closed,
  quasi-compact, and locally of finite type).
\end{lemma}

\begin{proof}
  \uses{mk-lem:projective_space_gradeZero, mk-lem:proj_toSpecZero_proper}
  For finite \(n\), the polynomial ring \(\mathbb{Z}[X_i : i \in n]\) is of
  finite type over its degree-zero part, which is \(\mathbb{Z}\) by
  \ref{mk-lem:projective_space_gradeZero}; hence the integral model
  \(\operatorname{Proj} \mathbb{Z}[X_i]\) is proper over
  \(\operatorname{Spec} \mathbb{Z}\) by \ref{mk-lem:proj_toSpecZero_proper},
  which is the terminal scheme.  Properness is stable under base change, and
  the structural morphism of \(\mathbb{P}(n; S)\) is the base change of
  \(\operatorname{Proj} \mathbb{Z}[X_i] \to \operatorname{Spec} \mathbb{Z}\)
  along \(S \to \operatorname{Spec} \mathbb{Z}\).
\end{proof}

\begin{lemma}
[Base change of relative projective space]
  \label{mk-lem:projective_space_base_change}
  \dcref{custom}
  \uses{mk-def:projective_space}
  For every morphism \(g : S' \to S\) the canonical square
  \[
    \begin{array}{ccc}
      \mathbb{P}(n; S') & \longrightarrow & \mathbb{P}(n; S) \\
      \downarrow & & \downarrow \\
      S' & \xrightarrow{\;g\;} & S
    \end{array}
  \]
  is cartesian.
\end{lemma}

\begin{proof}
  Both \(\mathbb{P}(n; S')\) and \(\mathbb{P}(n; S) \times_S S'\) are fibre
  products of \(S' \to \operatorname{Spec}\mathbb{Z} \leftarrow
  \operatorname{Proj} \mathbb{Z}[X_i]\), by the pasting law for cartesian
  squares applied to the defining square of \(\mathbb{P}(n; S)\) over the
  terminal scheme.
\end{proof}

\begin{lemma}
[The transition units of the Serre twist]
  \label{mk-lem:serre_twist_transition_unit}
  \dcref{custom}
  In the degree-zero part of the localization
  \(\mathbb{Z}[X_i : i \in n]_{X_i X_j}\), the fraction
  \(X_i/X_j = X_i^2/(X_i X_j)\) is a unit, with inverse
  \(X_j/X_i = X_j^2/(X_i X_j)\).
\end{lemma}

\begin{proof}
  The product of the two fractions is
  \(X_i^2 X_j^2 / (X_i X_j)^2 = 1\).
\end{proof}

\begin{lemma}
[Cocycle property of the Serre-twist transitions]
  \label{mk-lem:serre_twist_cocycle}
  \dcref{custom}
  \uses{mk-lem:serre_twist_transition_unit}
  On \(\operatorname{Proj} \mathbb{Z}[X_i : i \in n]\), for each \(m \ge 0\)
  the multiplication maps by \((X_i/X_j)^m\) on the trivial line bundles over
  the basic opens \(D_+(X_i)\), transported to the pairwise scheme-theoretic
  overlaps of the cover, satisfy the descent cocycle condition on triple
  overlaps: over \(D_+(X_i) \cap D_+(X_j) \cap D_+(X_k)\),
  \[
    \left(\frac{X_j}{X_k}\right)^{m}
    \left(\frac{X_i}{X_j}\right)^{m}
    \;=\; \left(\frac{X_i}{X_k}\right)^{m},
  \]
  compatibly with the identifications of the restricted trivial bundles.
\end{lemma}

\begin{proof}
  \uses{mk-lem:serre_twist_transition_unit}
  In the degree-zero part of the localization
  \(\mathbb{Z}[X_i]_{X_i X_j X_k}\), the three fractions satisfy
  \((X_i/X_j)(X_j/X_k) = X_i/X_k\) since
  \(X_i^2 X_j^2 X_k^2 = (X_i X_j)(X_j X_k)(X_i X_k)\); raising to the
  \(m\)-th power preserves the identity.  The multiplication maps by these
  units on the trivial line bundles commute with restriction to smaller
  opens, because multiplication by a section commutes with the restriction
  maps of the structure sheaf; and the trivialization comparisons between
  the restricted bundles compose transitively.  Hence the transported
  transition isomorphisms compose as claimed.
\end{proof}

\begin{definition}
[The Serre twisting sheaf on the integral model]
  \label{mk-def:serre_twist}
  \uses{mk-lem:serre_twist_transition_unit, mk-lem:serre_twist_cocycle}
  \dcref{custom}
  For \(m \ge 0\), the \emph{Serre twisting sheaf} \(\mathcal{O}(m)\) on
  \(\operatorname{Proj} \mathbb{Z}[X_i : i \in n]\) is the sheaf of modules
  glued from the trivial line bundles \(\mathcal{O}_{D_+(X_i)}\) on the
  basic opens of the standard affine cover along the transition
  isomorphisms ``multiplication by \((X_i/X_j)^m\)'' of
  \ref{mk-lem:serre_twist_transition_unit}, which satisfy the descent cocycle
  condition by \ref{mk-lem:serre_twist_cocycle}.
\end{definition}

\begin{remark}
[Sign convention]
  \label{mk-rem:serre_twist_sign}
  \dcref{custom}
  The gluing datum realises \(\mathcal{O}(m)\), not \(\mathcal{O}(-m)\).  On
  \(D_+(X_i)\) the trivialisation identifies a local section \(s\) of the twist
  with \(s/X_i^m\); the nowhere-vanishing local frame is thus \(X_i^m\), and its
  expression in the \(j\)-frame is \(X_i^m/X_j^m = (X_i/X_j)^m\) — exactly the
  transition ``multiplication by \((X_i/X_j)^m\)'' of the definition.  Hence a
  global section is a compatible family \((f_i)\) with \(f_i = s/X_i^m\) and
  \(f_i X_i^m = f_j X_j^m =: s\) homogeneous of degree \(m\), so
  \(\Gamma(\operatorname{Proj}\mathbb{Z}[X_i], \mathcal{O}(m))\) is the group of
  degree-\(m\) homogeneous polynomials; in particular \(\mathcal{O}(1)\) has the
  \(X_i\) as sections and is very ample.  The special case \(m = 0\), where the
  factor \((X_i/X_j)^0 = 1\) makes the descent datum that of the structure sheaf,
  is recorded in the formal proof as \texttt{twistTransition\_zero}.
\end{remark}

\begin{definition}
[The Serre twist on relative projective space]
  \label{mk-def:twisting_sheaf}
  \dcref{custom}
  \uses{mk-def:projective_space, mk-def:serre_twist}
  The \emph{Serre twist} \(\mathcal{O}(m)\) on \(\mathbb{P}(n; S)\) is the
  pullback of the twisting sheaf of \ref{mk-def:serre_twist} along the
  projection \(\mathbb{P}(n; S) \to \operatorname{Proj} \mathbb{Z}[X_i]\).
\end{definition}

\begin{lemma}
[Base change of the Serre twist]
  \label{mk-lem:twisting_sheaf_base_change}
  \dcref{custom}
  \uses{mk-def:twisting_sheaf, mk-lem:projective_space_base_change}
  For \(g : S' \to S\), the pullback of \(\mathcal{O}(m)\) on
  \(\mathbb{P}(n; S)\) along the base-change morphism
  \(\mathbb{P}(n; S') \to \mathbb{P}(n; S)\) is canonically isomorphic to
  \(\mathcal{O}(m)\) on \(\mathbb{P}(n; S')\).
\end{lemma}

\begin{proof}
  Both sheaves are pullbacks of the twisting sheaf on the integral model,
  along the two legs of the commuting triangle formed by the base-change
  morphism and the two projections to
  \(\operatorname{Proj} \mathbb{Z}[X_i]\).
\end{proof}

\begin{lemma}[The Serre twist is a locally trivial line bundle]\label{mk-lem:serreTwist_isLocallyTrivial}
  \dcref{custom}

  \uses{mk-def:serre_twist, mk-def:twisting_sheaf, mk-def:glueRestrictionIso,
    mk-def:IsLocallyTrivial}
  For every \(m \ge 0\), the Serre twisting sheaf \(\mathcal{O}(m)\) on
  the integral model \(\operatorname{Proj} \mathbb{Z}[X_i : i \in n]\)
  and its base change \(\mathcal{O}(m)\) on \(\mathbb{P}(n; S)\)
  (\ref{mk-def:twisting_sheaf}) are locally trivial line bundles
  (\ref{mk-def:IsLocallyTrivial}): rank-one free on each member of the
  standard affine cover \(\{D_+(X_i)\}\).
\end{lemma}
\begin{proof}On the basic-open chart \(D_+(X_i)\) the gluing datum of
  \ref{mk-def:serre_twist} trivialises the twist by construction: the
  chart restriction of the twist is the trivial line bundle
  \(\mathcal{O}_{D_+(X_i)}\), realised through the descent-gluing
  comparison isomorphism of \ref{mk-def:glueRestrictionIso}. Since
  \(\{D_+(X_i)\}\) covers \(\operatorname{Proj} \mathbb{Z}[X_i]\) and
  each \(D_+(X_i)\) is affine (a basic open of a homogeneous element of
  positive degree), \(\mathcal{O}(m)\) is locally trivial. Local
  triviality is stable under pullback of sheaves of modules, so the
  twisting sheaf on \(\mathbb{P}(n; S)\) --- the pullback of
  \(\mathcal{O}(m)\) along the projection to the integral model
  (\ref{mk-def:twisting_sheaf}) --- is locally trivial as well.
\end{proof}

\begin{corollary}[Finite presentation and quasi-coherence of the Serre twist]\label{mk-cor:serreTwist_isFinitePresentation}
  \dcref{custom}

  \uses{mk-lem:serreTwist_isLocallyTrivial, mk-thm:lbc_isFinitePresentation}
  The Serre twisting sheaf \(\mathcal{O}(m)\) on the integral model and
  its base change on \(\mathbb{P}(n; S)\) are finitely presented, hence
  quasi-coherent.
\end{corollary}
\begin{proof}Immediate from \ref{mk-lem:serreTwist_isLocallyTrivial} and
  \ref{mk-thm:lbc_isFinitePresentation} (a locally trivial line bundle is
  finitely presented); quasi-coherence then follows from finite
  presentation by \ref{mk-cor:lbc_isFiniteType}.
\end{proof}

\begin{theorem}[\(H^0\) base change of the Serre twist over a flat affine base]\label{mk-thm:serreTwist_sections_baseChange}
  \dcref{custom}

  \uses{mk-thm:quot_pullback_baseMap_sectionEquiv_qc,
    mk-cor:serreTwist_isFinitePresentation, mk-lem:projective_space_base_change,
    mk-lem:proj_toSpecZero_proper}
  Let \(S\) be an affine scheme flat over \(\operatorname{Spec}
  \mathbb{Z}\) and let \(n\) be a finite index set. For every \(m \ge
  0\) there is a \(\Gamma(S, \mathcal{O}_S)\)-linear equivalence
  \[
    \Gamma(S, \mathcal{O}_S) \otimes_{\mathbb{Z}}
      \Gamma\bigl(\operatorname{Proj} \mathbb{Z}[X_i : i \in n],
      \mathcal{O}(m)\bigr)
    \;\xrightarrow{\ \sim\ }\;
    \Gamma\bigl(\mathbb{P}(n; S), \mathcal{O}(m)\bigr).
  \]
\end{theorem}
\begin{proof}Apply \ref{mk-thm:quot_pullback_baseMap_sectionEquiv_qc} (\(H^0\) flat
  base change) to the cartesian square defining \(\mathbb{P}(n; S)\)
  (\ref{mk-lem:projective_space_base_change} for the base-change map
  \(S \to \operatorname{Spec} \mathbb{Z}\)), at \(V = U = \top\): the
  structural morphism
  \(\operatorname{Proj} \mathbb{Z}[X_i] \to \operatorname{Spec}
  \mathbb{Z}\) is proper (\ref{mk-lem:proj_toSpecZero_proper}), hence
  quasi-compact and quasi-separated; the base-change morphism
  \(S \to \operatorname{Spec} \mathbb{Z}\) is flat by hypothesis; and
  \(\mathcal{O}(m)\) is quasi-coherent
  (\ref{mk-cor:serreTwist_isFinitePresentation}). This identifies
  \(\Gamma(S, \mathcal{O}_S) \otimes_{\mathbb{Z}} \Gamma\bigl(
  \operatorname{Proj} \mathbb{Z}[X_i], \mathcal{O}(m)\bigr)\) with
  \(\Gamma(\mathbb{P}(n; S), \mathcal{O}(m))\), \(\Gamma(S,
  \mathcal{O}_S)\)-linearly.
\end{proof}

\begin{definition}
[Projective morphism]
  \label{mk-def:projective_morphism}
  \dcref{custom}
  \uses{mk-def:projective_space}
  A morphism of schemes \(\pi : X \to S\) is \emph{projective} if there are a
  finite set \(I\) and a closed immersion
  \(i : X \hookrightarrow \mathbb{P}(I;S)\) over \(S\).
\end{definition}

\begin{definition}
[Projective morphism carrying a very ample line bundle]
  \label{mk-def:projective_with}
  \uses{mk-def:projective_space, mk-def:twisting_sheaf}
  \dcref{custom}
  A morphism of schemes \(\pi : X \to S\) is \emph{projective carrying the
  line bundle \(L\)} if there are \(d \ge 0\) and a closed immersion
  \(i : X \hookrightarrow \mathbb{P}(\{0,\dots,d\}; S)\) over \(S\) with
  \(L \cong i^* \mathcal{O}(1)\).  In this situation \(L\) is relatively
  very ample and \(\pi\) is projective in the sense of \cite{nitsure-hilbert-quot} \S 5.
\end{definition}

\begin{lemma}
[Forgetting the chosen very ample line bundle]
  \label{mk-lem:projective_with_forget}
  \dcref{custom}
  \uses{mk-def:projective_morphism, mk-def:projective_with}
  A morphism that is projective carrying a very ample line bundle is
  projective.
\end{lemma}

\begin{proof}
  Forget the isomorphism between the chosen line bundle and the pullback of
  \(\mathcal{O}(1)\); the same closed immersion witnesses projectivity.
\end{proof}

\begin{lemma}
[Projective morphisms are proper]
  \label{mk-lem:projective_with_proper}
  \dcref{custom}
  \uses{mk-def:projective_with, mk-lem:projective_space_proper}
  A projective morphism carrying a line bundle is proper.
\end{lemma}

\begin{proof}
  \uses{mk-lem:projective_space_proper}
  A closed immersion is proper, the structural morphism of projective space
  is proper (\ref{mk-lem:projective_space_proper}), and properness is stable
  under composition.
\end{proof}

\begin{lemma}
[Structural consequences of projectivity]
  \label{mk-lem:projective_with_structural}
  \dcref{custom}
  \uses{mk-lem:projective_with_proper}
  A projective morphism carrying a line bundle is locally of finite type,
  separated, and universally closed.
\end{lemma}

\begin{proof}
  \uses{mk-lem:projective_with_proper}
  Immediate from \ref{mk-lem:projective_with_proper}: properness of a morphism
  entails that it is locally of finite type, separated, and universally
  closed.
\end{proof}

\begin{lemma}
[Invariance under isomorphism of the bundle]
  \label{mk-lem:projective_with_iso}
  \dcref{custom}
  \uses{mk-def:projective_with}
  If \(\pi : X \to S\) is projective carrying \(L\) and \(L \cong L'\), then
  \(\pi\) is projective carrying \(L'\).
\end{lemma}

\begin{proof}
  The witnessing closed immersion and its comparison
  \(L \cong i^* \mathcal{O}(1)\) are unchanged; precompose the comparison with
  the given isomorphism \(L' \cong L\) to obtain \(L' \cong i^*
  \mathcal{O}(1)\).
\end{proof}

\begin{lemma}
[Composition with closed immersions]
  \label{mk-lem:projective_with_closed_immersion}
  \dcref{custom}
  \uses{mk-def:projective_with}
  If \(\pi : X \to S\) is projective carrying \(L\) and \(j : Y \to X\) is a
  closed immersion, then \(\pi \circ j\) is projective carrying \(j^* L\).
\end{lemma}

\begin{proof}
  Compose the given closed immersion \(i : X \hookrightarrow
  \mathbb{P}(\{0,\dots,d\}; S)\) with \(j\): a composite of closed
  immersions is a closed immersion, the triangle over \(S\) still commutes,
  and \(j^* L \cong j^* i^* \mathcal{O}(1) \cong (i \circ j)^*
  \mathcal{O}(1)\).
\end{proof}

\begin{lemma}
[Base change of projective morphisms]
  \label{mk-lem:projective_with_base_change}
  \dcref{custom}
  \uses{mk-def:projective_with, mk-lem:projective_space_base_change,
    mk-lem:twisting_sheaf_base_change}
  If \(\pi : X \to S\) is projective carrying \(L\) and \(g : S' \to S\) is
  any morphism, then the base change \(X \times_S S' \to S'\) is projective
  carrying the pullback of \(L\) along \(X \times_S S' \to X\).
\end{lemma}

\begin{proof}
  \uses{mk-lem:projective_space_base_change, mk-lem:twisting_sheaf_base_change}
  Base-change the closed immersion \(i\) along
  \(\mathbb{P}(\{0,\dots,d\}; S') \to \mathbb{P}(\{0,\dots,d\}; S)\)
  (\ref{mk-lem:projective_space_base_change}): closed immersions are stable
  under base change, the induced comparison morphism
  \(X \times_S S' \to \mathbb{P}(\{0,\dots,d\}; S')\) lies over \(S'\), and
  the pulled-back bundle is identified with \(\mathcal{O}(1)\) of the
  base-changed projective space by \ref{mk-lem:twisting_sheaf_base_change}.
\end{proof}

\begin{lemma}
[Projective space is projective over its base]
  \label{mk-lem:projective_space_is_projective}
  \dcref{custom}
  \uses{mk-def:projective_with, mk-def:twisting_sheaf}
  The structural morphism \(\mathbb{P}(\{0,\dots,d\}; S) \to S\) is projective
  carrying the Serre twist \(\mathcal{O}(1)\); the witnessing closed immersion
  is the identity.  In particular the projective-with predicate is
  non-vacuous.
\end{lemma}

\begin{proof}
  Take the identity closed immersion
  \(\mathbb{P}(\{0,\dots,d\}; S) \to \mathbb{P}(\{0,\dots,d\}; S)\); the
  triangle over \(S\) is trivial and \(\mathcal{O}(1) \cong
  \mathrm{id}^* \mathcal{O}(1)\).
\end{proof}

\begin{lemma}
[Projectivity descends to scheme-theoretic fibres]
  \label{mk-lem:projective_with_fiber}
  \dcref{custom}
  \uses{mk-def:projective_with, mk-lem:projective_with_base_change}
  Let \(\pi : X \to S\) be projective carrying \(L\) and let \(s \in S\).  Then
  the structural morphism of the scheme-theoretic fibre
  \(X_s = X \times_S \operatorname{Spec}\kappa(s)\), namely
  \(X_s \to \operatorname{Spec}\kappa(s)\), is projective carrying the
  restriction \(L_s = L|_{X_s}\).
\end{lemma}

\begin{proof}
  \uses{mk-lem:projective_with_base_change}
  The fibre \(X_s\) is the base change of \(\pi\) along the canonical morphism
  \(\operatorname{Spec}\kappa(s) \to S\), and \(X_s \to
  \operatorname{Spec}\kappa(s)\) is the second projection.  Apply
  \ref{mk-lem:projective_with_base_change} to that morphism: the base change is
  projective carrying the pullback of \(L\) along the fibre embedding
  \(X_s \to X\), which is exactly \(L_s\).
\end{proof}

\begin{lemma}
[Finite presentation descends to fibres]
  \label{mk-lem:fiberModule_finitePresentation}
  \dcref{custom}
  \uses{mk-def:fiber_module_restriction}
  If \(F\) is a finitely presented sheaf of modules on \(X\), then its
  restriction \(F_s = F|_{X_s}\) to the fibre over \(s \in S\)
  (\ref{mk-def:fiber_module_restriction}) is finitely presented.
\end{lemma}

\begin{proof}
  \(F_s\) is the pullback of \(F\) along the fibre embedding \(X_s \to X\), and
  finite presentation of a sheaf of modules is stable under pullback along an
  arbitrary morphism of schemes.
\end{proof}

\begin{lemma}
[Serre finiteness along a fibre]
  \label{mk-lem:sectionGradedModule_fg_fiber}
  \dcref{custom}
  \uses{mk-lem:sectionGradedModule_fg, mk-lem:projective_with_fiber,
    mk-lem:fiberModule_finitePresentation, mk-def:sectionGradedRing,
    mk-def:sectionGradedModule}
  Let \(\pi : X \to S\) be projective carrying \(L\), let \(F\) be finitely
  presented, and let \(s \in S\).  Then the section graded ring \(R(X_s, L_s)\)
  (\ref{mk-def:sectionGradedRing}) is Noetherian and the section graded module
  \(M(X_s, F_s, L_s)\) (\ref{mk-def:sectionGradedModule}) is a finitely generated
  \(R(X_s, L_s)\)-module.
\end{lemma}

\begin{proof}
  \uses{mk-lem:sectionGradedModule_fg, mk-lem:projective_with_fiber,
    mk-lem:fiberModule_finitePresentation}
  The fibre \(X_s\) is projective over the field \(\kappa(s)\) carrying \(L_s\)
  (\ref{mk-lem:projective_with_fiber}), and \(F_s\) is coherent
  (\ref{mk-lem:fiberModule_finitePresentation}); apply
  \ref{mk-lem:sectionGradedModule_fg} on \(X_s\), retaining its first two
  conclusions.
\end{proof}

\begin{lemma}
[The graded Hilbert function is a genuine dimension on projective fibres]
  \label{mk-lem:hilbertFunction_finiteDimensional}
  \dcref{custom}
  \uses{mk-lem:sectionGradedModule_fg, mk-lem:projective_with_fiber,
    mk-lem:fiberModule_finitePresentation, mk-def:hilbert_polynomial}
  Let \(\pi : X \to S\) be projective carrying \(L\), let \(F\) be finitely
  presented, and let \(s \in S\).  Then for every \(m\) the space of global
  sections \(\Gamma(X_s, F_s \otimes L_s^{\otimes m})\) is a finite-dimensional
  \(\kappa(s)\)-vector space.  In particular the graded Hilbert function
  \(m \mapsto \dim_{\kappa(s)} \Gamma(X_s, F_s \otimes L_s^{\otimes m})\)
  (\ref{mk-def:hilbert_polynomial}) takes finite values, so it is not the
  \(\operatorname{finrank}\) junk value at any \(m\).
\end{lemma}

\begin{proof}
  \uses{mk-lem:sectionGradedModule_fg, mk-lem:projective_with_fiber,
    mk-lem:fiberModule_finitePresentation}
  This is the third conclusion of \ref{mk-lem:sectionGradedModule_fg} applied on
  the fibre \(X_s\), which is projective over \(\kappa(s)\) carrying \(L_s\)
  (\ref{mk-lem:projective_with_fiber}) with \(F_s\) coherent
  (\ref{mk-lem:fiberModule_finitePresentation}): finite dimensionality over the
  residue field of each graded component
  \(\Gamma(X_s, F_s \otimes L_s^{\otimes m})\).
\end{proof}

\subsection{Global sections of a glued sheaf of modules}
\label{mk-subsec:glue_gamma_compatible}

The Serre-twist sheaf of \ref{mk-def:serre_twist} is built by gluing
(\ref{mk-def:scheme_modules_glue}); the following computes the global sections
of any glued sheaf of modules as an explicit submodule of compatible chart
sections, and specialises it to the Serre twist.

\begin{definition}
[Compatible families of chart sections]
  \label{mk-def:glue_gamma_compatible}
  \dcref{custom}
  \uses{mk-def:scheme_modules_glue}
  Let \(D\) be a glue datum of schemes with chart family
  \((\mathcal M_i)_{i \in J}\) of sheaves of \(\mathcal O_{U_i}\)-modules and
  transition isomorphisms \((g_{ij})\), as in \ref{mk-def:scheme_modules_glue}.
  The \(\Gamma(D.\mathrm{glued}, X)\)-submodule of \emph{compatible families}
  \[
    \mathrm{Compat}(D, M, g) \;\subseteq\; \prod_{i \in J} \Gamma(U_i, \mathcal M_i)
  \]
  consists of the families \((s_i)_i\) such that for every pair \((i,j)\),
  the restriction of \(s_i\) to the overlap \(V_{ij}\) along \(f_{ij}\)
  agrees, through the transition isomorphism \(g_{ij}\), with the
  restriction of \(s_j\) to \(V_{ij}\) along \(t_{ij} \circ f_{ji}\).
\end{definition}

\begin{lemma}
[Global sections of a product of sheaves of modules]
  \label{mk-lem:gamma_pi_sections}
  \dcref{custom}
  Let \((\mathcal N_i)_{i \in I}\) be a family of sheaves of
  \(\mathcal O_X\)-modules on a scheme \(X\) and let
  \(\textstyle\prod^{c}_i \mathcal N_i\) be their categorical product in
  \(X.\mathrm{Modules}\). The family of projections
  \(\pi_i : \prod^{c}_i \mathcal N_i \to \mathcal N_i\) induces, on global
  sections, a bijection
  \[
    \Gamma\Bigl(\textstyle\prod^{c}_i \mathcal N_i,\, X\Bigr)
      \;\xrightarrow{\ \sim\ }\; \prod_i \Gamma(\mathcal N_i, X), \qquad
    x \mapsto \bigl(\pi_i(x)\bigr)_i;
  \]
  in particular the projections are jointly injective, and every family
  \((s_i)_i \in \prod_i \Gamma(\mathcal N_i, X)\) is the family of
  projections of some global section of the product.
\end{lemma}
\begin{proof}Evaluation of sheaves of modules at the maximal open computes global
  sections and preserves all limits, so it carries the categorical product
  to the product, in \(\Gamma(\mathcal O_X, X)\)-modules, of the global
  sections, with components the section-level projections; the resulting
  comparison map is therefore invertible. Injectivity of the family of
  projections and surjectivity onto arbitrary families of sections are the
  two halves of this bijection.
\end{proof}

\begin{lemma}
[Global sections of a glued sheaf are the compatible families]
  \label{mk-lem:glue_sections_equalizer}
  \dcref{custom}
  \uses{mk-def:scheme_modules_glue, mk-def:glue_gamma_compatible,
    mk-lem:gamma_pi_sections}
  With the data and module cocycle conditions (C1)(C2) of
  \ref{mk-def:scheme_modules_glue}, restriction to the charts induces a
  \(\Gamma(D.\mathrm{glued}, X)\)-linear isomorphism
  \[
    \Gamma\bigl(\mathrm{glue}(D,M,g),\, X\bigr) \;\xrightarrow{\ \sim\ }\;
      \mathrm{Compat}(D, M, g)
  \]
  (\ref{mk-def:glue_gamma_compatible}), sending a global section to the family
  of its restrictions to the charts \(U_i\), with inverse sending a
  compatible family to the unique global section restricting to it on every
  chart. In particular global sections of the glued sheaf are jointly
  detected by their restrictions to the charts.
\end{lemma}
\begin{proof}\uses{mk-lem:gamma_pi_sections}
  The glued sheaf is by construction the equalizer of the two parallel maps
  \[
    \textstyle\prod^{c}_i (\iota_i)_{*} \mathcal M_i \;\rightrightarrows\;
      \prod^{c}_{(i,j)} (f_{ij} \circ \iota_i)_{*} (f_{ij}^{*} \mathcal M_i)
  \]
  encoding the two descent legs of \ref{mk-def:scheme_modules_glue}. Taking
  global sections preserves this equalizer (it is a limit-preserving
  functor), yielding, in \(\Gamma(D.\mathrm{glued}, X)\)-modules, the
  equalizer of the corresponding section-level maps; by
  \ref{mk-lem:gamma_pi_sections} the section functor identifies the two outer
  products with the products of sections, and global sections of a
  pushforward at the ambient open agree with the global sections of the
  original sheaf, so the section-level equalizer is exactly the submodule of
  families on which the two descent legs agree componentwise, i.e.\
  \(\mathrm{Compat}(D,M,g)\). Under this identification the chart-restriction
  map corresponds to the equalizer inclusion, giving the stated isomorphism;
  joint detection by chart restrictions is injectivity of the equalizer
  inclusion.
\end{proof}

\begin{lemma}
[Pullback along an isomorphism is pushforward of the inverse]
  \label{mk-lem:pullback_iso_pushforward_inv}
  \dcref{custom}
  Let \(\varphi : X \to Y\) be an isomorphism of schemes. Then the pullback
  functor \(\varphi^{*} : Y.\mathrm{Modules} \to X.\mathrm{Modules}\) is
  naturally isomorphic to the pushforward functor
  \((\varphi^{-1})_{*} : Y.\mathrm{Modules} \to X.\mathrm{Modules}\) along
  the inverse isomorphism.
\end{lemma}
\begin{proof}Pushforward along \(\varphi\) and pushforward along \(\varphi^{-1}\) are
  mutually inverse equivalences between \(X.\mathrm{Modules}\) and
  \(Y.\mathrm{Modules}\), by pseudofunctoriality of pushforward: the
  composite pushforward, in either order, is naturally isomorphic to the
  identity, since \(\varphi \circ \varphi^{-1}\) and
  \(\varphi^{-1} \circ \varphi\) are identities. In particular pushforward
  along \(\varphi^{-1}\) is left adjoint to pushforward along \(\varphi\);
  since \(\varphi^{*}\) is also left adjoint to pushforward along
  \(\varphi\), uniqueness of left adjoints up to natural isomorphism
  identifies the two functors.
\end{proof}

\begin{lemma}
[The first descent leg on unit modules is restriction along the overlap
immersion]
  \label{mk-lem:glue_legA_unit_restriction}
  \dcref{custom}
  \uses{mk-def:scheme_modules_glue}
  Let \(D\) be a glue datum of schemes and instantiate the chart family of
  \ref{mk-def:scheme_modules_glue} at the unit modules
  \(\mathcal M_i = \mathcal O_{U_i}\). Then, through the canonical
  trivialisation of the pulled-back structure sheaf, the \((i,j)\)-component
  of the first descent leg acts on a section \(x \in \Gamma(U_i, \mathcal
  O)\) as its restriction along the overlap immersion
  \(f_{ij} : V_{ij} \to U_i\):
  \[
    a_{ij}(x) = x|_{V_{ij}} \in \Gamma(V_{ij}, \mathcal O).
  \]
\end{lemma}
\begin{proof}The first descent leg is the pushforward, along \(f_{ij}\) followed by
  \(\iota_i\), of the unit of the pullback--pushforward adjunction applied to
  \(\mathcal O_{U_i}\), post-composed with the pushforward-composition
  comparison. Composed with the trivialisation of the pulled-back unit
  sheaf, naturality of the pushforward-composition comparison and the
  identification of the adjunction unit at the unit sheaf with the
  canonical comorphism of \(f_{ij}\) rewrite this composite as the
  pushforward of that comorphism; on global sections the pushforward of a
  comorphism acts by the ring-restriction map, giving the stated formula.
\end{proof}

\begin{lemma}
[Global sections of the glued Serre twist]
  \label{mk-lem:serre_twist_glued_sections}
  \dcref{custom}
  \uses{mk-lem:glue_sections_equalizer, mk-def:serre_twist, mk-lem:serre_twist_cocycle}
  For the descent datum of \ref{mk-def:serre_twist} (the standard affine cover
  of \(\operatorname{Proj} \mathbb Z[X_i : i \in n]\) by the \(D_+(X_i)\),
  unit chart modules, and transition isomorphisms ``multiplication by
  \((X_i/X_j)^m\)''), \ref{mk-lem:glue_sections_equalizer} specialises to a
  \(\Gamma(\operatorname{Proj}\mathbb Z[X_i], \mathcal O)\)-linear
  isomorphism between the global sections of the glued twist and the
  compatible families of chart sections for that datum.
\end{lemma}
\begin{proof}\uses{mk-lem:glue_sections_equalizer}
  Apply \ref{mk-lem:glue_sections_equalizer} with \(M_i = \mathcal
  O_{D_+(X_i)}\) and \(g_{ij}\) multiplication by \((X_i/X_j)^m\); the
  module cocycle conditions (C1)(C2) required by
  \ref{mk-def:scheme_modules_glue} hold: (C1) is the identity of the
  transition at \(i = j\), since \((X_i/X_i)^m = 1\), and (C2) is exactly
  \ref{mk-lem:serre_twist_cocycle}.
\end{proof}

\begin{lemma}
[Global sections of the Serre twist are the compatible chart families]
  \label{mk-lem:serre_twist_sections_compatible}
  \dcref{custom}
  \uses{mk-lem:serre_twist_glued_sections, mk-lem:pullback_iso_pushforward_inv,
    mk-def:serre_twist}
  The Serre twisting sheaf \(\mathcal O(m)\) on
  \(\operatorname{Proj} \mathbb Z[X_i : i \in n]\) (\ref{mk-def:serre_twist})
  is isomorphic to the pushforward, along the canonical comparison from the
  abstractly glued model of \ref{mk-lem:serre_twist_glued_sections} to
  \(\operatorname{Proj}\), of the glued twist
  (\ref{mk-lem:pullback_iso_pushforward_inv}, applied to the inverse of that
  comparison isomorphism); consequently its global sections form an
  additive group isomorphic to the compatible families of
  \ref{mk-lem:serre_twist_glued_sections}:
  \[
    \Gamma\bigl(\operatorname{Proj}\mathbb Z[X_i],\, \mathcal O(m)\bigr)
      \;\cong\; \mathrm{Compat}\bigl(D, \mathcal O, ((X_i/X_j)^m)\bigr).
  \]
  (The sign convention of \ref{mk-rem:serre_twist_sign} governs the
  orientation of the transition isomorphisms entering this identification.)
\end{lemma}
\begin{proof}\uses{mk-lem:serre_twist_glued_sections, mk-lem:pullback_iso_pushforward_inv}
  The gluing construction realises \(\operatorname{Proj}\mathbb Z[X_i]\) as
  the target of an isomorphism from the abstractly glued scheme of the
  standard affine cover; pulling the twisting sheaf back along that
  isomorphism and applying \ref{mk-lem:pullback_iso_pushforward_inv}
  identifies it with a pushforward of the glued twist, and global sections
  of a pushforward at the ambient open agree definitionally with those of
  the original sheaf. Composing this identification of global sections with
  \ref{mk-lem:serre_twist_glued_sections} gives the stated isomorphism.
\end{proof}

\subsection{The compatible-family condition, concretely}
\label{mk-subsec:serre_twist_compatible_condition}

\ref{mk-lem:serre_twist_sections_compatible} identifies global sections of
\(\mathcal O(m)\) with the abstract compatible-family submodule
\(\mathrm{Compat}(D, \mathcal O, ((X_i/X_j)^m))\)
(\ref{mk-def:glue_gamma_compatible}) of the descent datum instantiated at the
transition isomorphisms ``multiplication by \((X_i/X_j)^m\)''. The next
lemmas make this compatibility condition concrete: through the chart
identification with the degree-zero localisation \(\mathrm{Away}(X_i)\), a
compatible family becomes a family of away-fractions related on overlaps by
the transition unit, and a single degree-\(m\) form gives such a family.

\begin{lemma}[The second descent leg of the Serre-twist datum, concretely]\label{mk-lem:serre_twist_b_leg}
  \dcref{custom}

  \uses{mk-lem:glue_legA_unit_restriction, mk-def:scheme_modules_glue, mk-def:serre_twist}
  Instantiate the chart family of \ref{mk-def:scheme_modules_glue} at the unit
  modules and the descent transition of \ref{mk-def:serre_twist}. Through the
  \(i\)-chart trivialisation, the second (``\(b\)'') descent leg sends a
  \(j\)-chart section \(y \in \Gamma(U_j, \mathcal O)\) to its restriction
  along the overlap immersion \(t_{ij} \circ f_{ji} : V_{ij} \to U_j\), scaled
  by the inverse transition unit \((X_j/X_i)^m\):
  \[
    b_{ij}(y) = y|_{V_{ij}} \cdot (X_j/X_i)^m \in \Gamma(V_{ij}, \mathcal O).
  \]
  This is the ``\(b\)-side'' companion of the \(a\)-side restriction formula
  of \ref{mk-lem:glue_legA_unit_restriction}.
\end{lemma}
\begin{proof}\uses{mk-lem:glue_legA_unit_restriction}
  The second descent leg is, before trivialisation, the pushforward of the
  canonical comorphism of \(t_{ij} \circ f_{ji}\) followed by the scalar
  automorphism ``multiplication by \((X_j/X_i)^m\)'' and a reindexing along
  the glue-datum's cocycle identity; composing with the \(i\)-chart
  trivialisation and cancelling it via the shared abstract unit-leg
  trivialisation identity (used identically for the \(a\)-leg in
  \ref{mk-lem:glue_legA_unit_restriction}) leaves exactly the comorphism of
  \(t_{ij} \circ f_{ji}\) composed with the scalar action. On global
  sections the pushforward of a comorphism acts by ring restriction and the
  pushforward of a scalar automorphism acts by multiplication (both
  through the definitional identification \(\Gamma(\rho_* \mathcal O, \top) =
  \Gamma(\mathcal O, \top)\)), and the reindexing comparison acts as the
  identity on sections; composing these gives the stated formula.
\end{proof}

\begin{lemma}[The concrete Serre-twist compatible-family condition]\label{mk-lem:serre_twist_compatible_condition}
  \dcref{custom}

  \uses{mk-lem:glue_legA_unit_restriction, mk-lem:serre_twist_b_leg,
    mk-def:glue_gamma_compatible}
  A chart-section family \((s_i)_i\), \(s_i \in \Gamma(D_+(X_i), \mathcal O)\),
  is a compatible family for the descent datum of \(\mathcal O(m)\)
  (\ref{mk-def:glue_gamma_compatible}) if and only if, on every overlap
  \(D_+(X_i) \cap D_+(X_j)\),
  \[
    s_i|_{V_{ij}} = s_j|_{V_{ij}} \cdot (X_j/X_i)^m .
  \]
  This is the section-level realisation of the sign convention of
  \ref{mk-rem:serre_twist_sign}: the compatible-family condition scales by
  \((X_j/X_i)^m\), not \((X_i/X_j)^m\).
\end{lemma}
\begin{proof}\uses{mk-lem:glue_legA_unit_restriction, mk-lem:serre_twist_b_leg}
  Unfold membership in \(\mathrm{Compat}\) (\ref{mk-def:glue_gamma_compatible})
  and transport the \(a\)-leg and \(b\)-leg conditions through the
  \(i\)-chart trivialisation via \ref{mk-lem:glue_legA_unit_restriction} and
  \ref{mk-lem:serre_twist_b_leg} respectively; the trivialisation is injective
  on sections (being the global-section evaluation of an isomorphism of
  sheaves of modules), so the transported condition is equivalent to the
  original one.
\end{proof}

\begin{definition}[Chart sections as the degree-zero localisation]\label{mk-def:chart_sections_away}
  \dcref{custom}

  \uses{mk-def:serre_twist}
  The structure-sheaf sections \(\Gamma(D_+(X_i), \mathcal O)\) of the
  \(i\)-th chart of \(\operatorname{Proj} \mathbb Z[X_i : i \in n]\) are
  identified with the degree-zero part of the localisation at \(X_i\),
  \(\mathrm{Away}(X_i) := \bigl(\mathbb Z[X_i]_{X_i}\bigr)_0\), via Mathlib's
  chart isomorphism for \(\operatorname{Proj}\) (valid since \(X_i\) is
  homogeneous of positive degree \(1\)), transported to the chart's global
  sections through the open-subscheme identification of the chart with its
  own affine spectrum.
\end{definition}
\begin{proof}Composite of Mathlib's \(\operatorname{Proj}\) basic-open-versus-away
  isomorphism with the canonical global-sections identification of an open
  subscheme with its \(\top\)-sections.
\end{proof}

\subsection{The away-fraction bridge and the degree-\(m\) forms}
\label{mk-subsec:serre_twist_away_bridge}

Through the chart identification \ref{mk-def:chart_sections_away}, the two
descent legs of \ref{mk-lem:serre_twist_compatible_condition} become explicit
maps between localisation rings; this section records that comparison
(the ``away-fraction bridge''), the resulting compatibility condition on a
single homogeneous form, its consequence that the degree-\(m\) forms embed
into \(\Gamma(\mathcal O(m))\), and the coordinate sections of
\(\mathcal O(1)\).

\begin{lemma}[The away-fraction bridge]\label{mk-lem:serre_twist_away_bridge}
  \dcref{custom}

  \uses{mk-lem:serre_twist_compatible_condition, mk-def:chart_sections_away}
  Through the chart identification of \ref{mk-def:chart_sections_away}, the
  two descent legs restricting to the overlap \(D_+(X_i X_j)\) are the
  localisation-algebra maps \(\mathrm{Away}(X_i) \to \mathrm{Away}(X_i X_j)\)
  and \(\mathrm{Away}(X_j) \to \mathrm{Away}(X_i X_j)\) sending a fraction to
  its image under the further localisation at the other variable, followed
  by the canonical identification of \(\mathrm{Away}(X_i X_j)\) with itself
  along the two orderings of the product.
\end{lemma}
\begin{proof}\uses{mk-lem:serre_twist_compatible_condition}
  On global sections, restriction of a chart section along the overlap
  immersion into \(D_+(X_i X_j) \subseteq D_+(X_i)\) (resp.\ \(D_+(X_j)\)) is,
  after the chart identification, exactly the further-localisation
  (away-map) comparison of homogeneous localisations at the common
  degree-zero ring \(\mathrm{Away}(X_i X_j)\).
\end{proof}

\begin{lemma}[A single form gives a compatible family]\label{mk-lem:serre_twist_form_compatible}
  \dcref{custom}

  \uses{mk-lem:serre_twist_away_bridge, mk-lem:serre_twist_compatible_condition}
  For a degree-\(m\) homogeneous form \(F\), the chart family
  \(\bigl(F/X_i^m\bigr)_i\) --- the away-fraction of \(F\) at each chart,
  through \ref{mk-def:chart_sections_away} --- is a compatible family for the
  Serre-twist descent datum (\ref{mk-lem:serre_twist_compatible_condition}).
\end{lemma}
\begin{proof}\uses{mk-lem:serre_twist_away_bridge}
  In \(\mathrm{Away}(X_i X_j)\), the two away-map images of \(F/X_i^m\) and
  \(F/X_j^m\) differ by exactly \((X_j/X_i)^m\) (a direct fraction identity:
  both sides equal \(F/(X_i X_j)^m\) up to the appropriate unit power).
  Transporting through the away-fraction bridge
  (\ref{mk-lem:serre_twist_away_bridge}) turns this into the compatibility
  condition of \ref{mk-lem:serre_twist_compatible_condition}.
\end{proof}

\begin{lemma}[The degree-\(m\) forms embed as global sections of \(\mathcal O(m)\)]\label{mk-lem:serre_twist_forms_embed}
  \dcref{custom}

  \uses{mk-lem:serre_twist_form_compatible, mk-lem:serre_twist_sections_compatible}
  There is an additive homomorphism from the degree-\(m\) homogeneous forms
  of \(\mathbb Z[X_i : i \in n]\) to \(\Gamma\bigl(\operatorname{Proj}
  \mathbb Z[X_i], \mathcal O(m)\bigr)\), sending \(F\) to the global section
  corresponding, under \ref{mk-lem:serre_twist_sections_compatible}, to the
  compatible chart family of \ref{mk-lem:serre_twist_form_compatible}. When the
  index set \(n\) is nonempty, this homomorphism is injective: a global
  section is determined by any one of its chart fractions, and the
  \(i\)-chart fraction of the image of \(F\) recovers \(F\) since \(X_i\) is
  a nonzerodivisor of the domain \(\mathbb Z[X_i]\).
\end{lemma}
\begin{proof}\uses{mk-lem:serre_twist_form_compatible}
  Compose \ref{mk-lem:serre_twist_form_compatible} (additively in \(F\), since
  the chart away-fraction \(F \mapsto F/X_i^m\) and the chart identification
  are additive) with the inverse of the compatible-family isomorphism of
  \ref{mk-lem:serre_twist_sections_compatible}. For injectivity, fix an index
  \(i\) (using nonemptiness of \(n\)): if two forms have equal images, their
  \(i\)-chart fractions agree by naturality of the isomorphisms involved,
  and equality of away-fractions \(F/X_i^m = F'/X_i^m\) in a domain forces
  \(F = F'\) after clearing the nonzerodivisor \(X_i^m\).
\end{proof}

\begin{definition}[Coordinate global sections of \(\mathcal O(1)\)]\label{mk-def:coordinate_section}
  \dcref{custom}

  \uses{mk-lem:serre_twist_forms_embed}
  The \emph{coordinate global section} \(x_j \in \Gamma\bigl(
  \operatorname{Proj} \mathbb Z[X_i], \mathcal O(1)\bigr)\) is the image of
  the degree-one form \(X_j\) under the embedding of
  \ref{mk-lem:serre_twist_forms_embed}. On the \(i\)-th chart it restricts to
  the away-fraction \(X_j/X_i\).
\end{definition}
\begin{proof}\uses{mk-lem:serre_twist_forms_embed}
  Direct instantiation of \ref{mk-lem:serre_twist_forms_embed} at
  \(m = 1\), \(F = X_j\); the chart restriction identity is the \(i\)-chart
  value of \ref{mk-lem:serre_twist_form_compatible} applied to that same form.
\end{proof}

\begin{lemma}[Global sections of \(\mathcal O(m)\) are the degree-\(m\) forms]\label{mk-lem:serre_twist_sections}
  \dcref{custom}

  \uses{mk-lem:serre_twist_forms_embed, mk-lem:serre_twist_away_bridge}
  For \(n\) nonempty, the embedding of \ref{mk-lem:serre_twist_forms_embed} of
  the degree-\(m\) homogeneous forms of \(\mathbb Z[X_i : i \in n]\) into
  \(\Gamma\bigl(\operatorname{Proj}\mathbb Z[X_i], \mathcal O(m)\bigr)\) is
  onto: every global section of \(\mathcal O(m)\) arises from a (unique)
  degree-\(m\) form, i.e.
  \(\Gamma\bigl(\operatorname{Proj}\mathbb Z[X_i], \mathcal O(m)\bigr)
  \cong (\mathbb Z[X_i])_m\).
\end{lemma}
\begin{proof}\uses{mk-lem:serre_twist_sections_compatible, mk-lem:serre_twist_form_compatible}
  Injectivity is \ref{mk-lem:serre_twist_forms_embed}; it remains to prove
  surjectivity. A global section corresponds, via
  \ref{mk-lem:serre_twist_sections_compatible} and the away-fraction bridge
  \ref{mk-lem:serre_twist_away_bridge}, to a compatible family
  \((a_i)_{i \in n}\) of chart fractions, each \(a_i\) an element of the
  degree-zero part of \(\mathrm{Away}(X_i)\), compatible in the sense
  that the images of \(a_i\) and \(a_j\) in \(\mathrm{Away}(X_i X_j)\)
  agree. Fix any index \(i_0\) (possible since \(n\) is nonempty) and
  write \(a_{i_0} = N_0/X_{i_0}^{k_0}\) in lowest terms, with \(N_0\)
  homogeneous of degree \(k_0 + m\). The ring \(\mathbb Z[X_i : i \in n]\)
  is a unique factorisation domain in which every variable \(X_i\) is a
  prime element; comparing \(a_{i_0}\) against every other \(a_i\) on the
  pairwise overlaps forces \(X_{i_0}^{k_0}\) to divide \(X_{i_0}^m N_0\),
  and primality of \(X_{i_0}\) upgrades this to the existence of a
  homogeneous form \(F\) of degree \(m\) with
  \(X_{i_0}^{k_0} F = X_{i_0}^m N_0\), i.e.\ \(F/X_{i_0}^m = a_{i_0}\).
  For every other index \(i\), the same cross relation applied to the
  pair \((i_0, i)\), followed by cancellation of the nonzerodivisor
  \(X_{i_0}^{k_0}\), gives \(F X_i^{k_i} = N_i X_i^m\), i.e.\
  \(F/X_i^m = a_i\). Hence the compatible family \((a_i)\) is exactly the
  chart family of \(F\) (\ref{mk-lem:serre_twist_form_compatible}), so the
  section it represents lies in the image of the embedding of
  \ref{mk-lem:serre_twist_forms_embed}.
\end{proof}

\subsection{The graded Hilbert--Serre bridge}
\label{mk-subsec:graded_hilbert_bridge}

The rationality engine \ref{mk-lem:gradedHilbertSerre_rational} consumes an
abstract finitely generated \(\mathbb{N}\)-graded \(\kappa\)-vector space
carrying \(r\) commuting degree-one endomorphisms.  We record its input as a
package \(\mathrm{GradedHilbert}\), factor Serre finiteness in this engine-ready
form as a single leaf (\ref{mk-lem:gradedHilbert_fiber}, a companion to
\ref{mk-lem:sectionGradedModule_fg}), and derive the rationality of the fibre
graded Hilbert function --- hence the genuineness of the Hilbert polynomial
(\ref{mk-def:hilbert_polynomial}) --- as sorry-free glue.  This is the
\(\mathrm{DirectSum}\)-decomposition bridge of inbox \texttt{I-0109}.

\begin{definition}
[Engine-ready packaging of a Hilbert function]
  \label{mk-def:gradedHilbert}
  \dcref{custom}
  \uses{mk-lem:gradedHilbertSerre_rational, mk-def:ratHilb}
  Fix a field \(\kappa\) and a function \(f : \mathbb{N} \to \mathbb{N}\).  A
  \(\mathrm{GradedHilbert}\) datum for \((\kappa, f)\) is a natural number \(r\),
  an \(\mathbb{N}\)-graded \(\kappa\)-vector space \(M = \bigoplus_n M_n\) (an
  internal grading \(\mathcal{M} : \mathbb{N} \to \mathrm{Submodule}_\kappa M\)
  with \(\mathrm{DirectSum.Decomposition}\)) whose graded pieces \(M_n\) are all
  finite-dimensional, together with \(r\) pairwise-commuting \(\kappa\)-linear
  endomorphisms \(g_0, \dots, g_{r-1}\) that each raise the grading degree by one,
  such that \(\dim_\kappa M_n = f(n)\) for all \(n\) and \(M\) is a finite module
  over the polynomial ring \(\kappa[X_0, \dots, X_{r-1}]\) acting through the
  \(g_i\).  This bundles exactly the data and hypotheses consumed by the graded
  Hilbert--Serre rationality engine \ref{mk-lem:gradedHilbertSerre_rational}.
\end{definition}

\begin{lemma}
[Rationality from a \(\mathrm{GradedHilbert}\) datum]
  \label{mk-lem:gradedHilbert_isRatHilb}
  \dcref{custom}
  \uses{mk-def:gradedHilbert, mk-lem:gradedHilbertSerre_rational, mk-def:ratHilb}
  If \((\kappa, f)\) admits a \(\mathrm{GradedHilbert}\) datum of order \(r\)
  (\ref{mk-def:gradedHilbert}), then \(f\) is a rational Hilbert function of order
  \(r\) (\ref{mk-def:ratHilb}): there are \(p \in \mathbb{Q}[X]\) and \(N\) with
  \(f(n) = [X^n]\bigl(p \cdot (1-X)^{-r}\bigr)\) for all \(n > N\).
\end{lemma}

\begin{proof}
  \uses{mk-lem:gradedHilbertSerre_rational}
  Apply the graded Hilbert--Serre rationality engine
  (\ref{mk-lem:gradedHilbertSerre_rational}) to the graded module \(M\), its
  grading \(\mathcal{M}\), and the degree-one generators \(g_i\) of the datum;
  its hypotheses (finite-dimensional components, finiteness over
  \(\kappa[X_0,\dots,X_{r-1}]\)) are the datum's fields.  This yields that
  \(n \mapsto \dim_\kappa M_n\) is a rational Hilbert function of order \(r\), and
  \(\dim_\kappa M_n = f(n)\) identifies it with \(f\).
\end{proof}

\begin{lemma}
[Serre finiteness in engine-ready form]
  \label{mk-lem:gradedHilbert_fiber}

  \uses{mk-def:gradedHilbert, mk-lem:sectionGradedModule_fg, mk-lem:projective_with_fiber,
    mk-def:projective_with, mk-def:hilbert_polynomial}
  \dcref{custom}
  \textit{Source: \cite{nitsure-hilbert-quot}, \S 1 (FGA Explained Ch.~5); Serre's theorem, as in
  \ref{mk-lem:sectionGradedModule_fg}.}
  Let \(\pi : X \to S\) be projective carrying the very ample line bundle \(L\),
  let \(F\) be finitely presented, and let \(s \in S\).  Then the graded Hilbert
  function \(m \mapsto \dim_{\kappa(s)} \Gamma(X_s, F_s \otimes L_s^{\otimes m})\)
  (\ref{mk-def:hilbert_polynomial}) admits a \(\mathrm{GradedHilbert}\) datum over
  the residue field \(\kappa(s)\) (\ref{mk-def:gradedHilbert}).

\end{lemma}

\begin{proof}
  \uses{mk-lem:sectionGradedModule_fg, mk-lem:projective_with_fiber, mk-def:projective_with}
  The fibre \(X_s\) is projective over \(\kappa(s)\) carrying \(L_s\)
  (\ref{mk-lem:projective_with_fiber}), so it admits a closed immersion
  \(i : X_s \hookrightarrow \mathbb{P}^d_{\kappa(s)}\) with
  \(L_s \cong i^*\mathcal{O}(1)\), and \(F_s\) is coherent.  By Serre
  (\ref{mk-lem:sectionGradedModule_fg}) the section module \(M(X_s, F_s, L_s)\) is
  a finitely generated graded module over \(R(X_s, L_s)\) with
  finite-dimensional components.  The homogeneous coordinates \(x_0, \dots, x_d\)
  restrict along \(i\) to \(d+1\) sections of \(L_s\) generating \(R(X_s, L_s)\)
  in large degrees, so \(M(X_s, F_s, L_s)\) is finite over the polynomial ring
  \(\kappa(s)[x_0, \dots, x_d]\), with \(x_j\) acting by multiplication (raising
  degree by one).  Its \(m\)-th graded piece \(\Gamma(X_s, F_s \otimes
  L_s^{\otimes m})\) has \(\kappa(s)\)-dimension the graded Hilbert function, so
  this data is a \(\mathrm{GradedHilbert}\) datum.
\end{proof}

\begin{lemma}
[Rationality of the fibre graded Hilbert function]
  \label{mk-lem:hilbertFunction_isRatHilb}
  \dcref{custom}

  \uses{mk-lem:gradedHilbert_fiber, mk-lem:gradedHilbert_isRatHilb, mk-def:ratHilb}
  Let \(\pi : X \to S\) be projective carrying \(L\), let \(F\) be finitely
  presented, and let \(s \in S\).  Then the graded Hilbert function
  \(m \mapsto \dim_{\kappa(s)} \Gamma(X_s, F_s \otimes L_s^{\otimes m})\) is a
  rational Hilbert function of some order \(d\) (\ref{mk-def:ratHilb}).
\end{lemma}

\begin{proof}
  \uses{mk-lem:gradedHilbert_fiber, mk-lem:gradedHilbert_isRatHilb}
  By \ref{mk-lem:gradedHilbert_fiber} the graded Hilbert function admits a
  \(\mathrm{GradedHilbert}\) datum over \(\kappa(s)\); apply
  \ref{mk-lem:gradedHilbert_isRatHilb}.
\end{proof}

\begin{lemma}
[The Hilbert polynomial computes the Hilbert function on projective fibres]
  \label{mk-lem:hilbertPolynomial_eval_eventually_of_projective}
  \dcref{custom}

  \uses{mk-lem:hilbertFunction_isRatHilb, mk-cor:hilbertPolynomial_eval_of_isRatHilb,
    mk-def:hilbert_polynomial}
  Let \(S\) be locally Noetherian, let \(\pi : X \to S\) be projective carrying
  \(L\) and locally of finite type, let \(F\) be finitely presented, and let
  \(s \in S\).  Then for all \(m \gg 0\) the Hilbert polynomial of
  \ref{mk-def:hilbert_polynomial} satisfies
  \[
    \Phi_{F,s}(m) \;=\; \dim_{\kappa(s)}
      \Gamma\bigl(X_s,\, F_s \otimes L_s^{\otimes m}\bigr).
  \]
  This is the projective-fibre realisation of
  \ref{mk-thm:hilbertPoly_of_sectionModule}: it is the point at which
  \(\Phi_{F,s}\) ceases to be a formal choice and becomes the genuine
  Snapper/Hilbert polynomial of the fibre.
\end{lemma}

\begin{proof}
  \uses{mk-lem:hilbertFunction_isRatHilb, mk-cor:hilbertPolynomial_eval_of_isRatHilb}
  By \ref{mk-lem:hilbertFunction_isRatHilb} the graded Hilbert function is a
  rational Hilbert function of some order \(d\); conclude by the
  polynomial-extraction corollary
  \ref{mk-cor:hilbertPolynomial_eval_of_isRatHilb}.
\end{proof}

\section{Representability of the Quot scheme}
\label{mk-sec:quot_representable}

\begin{theorem}[Representability of the Quot functor]\label{mk-thm:quot_representable}
  \dcref{custom}

  \uses{mk-def:quot_functor, mk-lem:quot_boundedness, mk-lem:quot_alpha_injective,
    mk-lem:quot_reduction_to_pi_star_W, mk-lem:quot_valuative_criterion,
    mk-thm:grassmannian_representable, mk-thm:flattening_stratification_exists,
    mk-thm:flat_base_change_cohomology}
  \textit{Source: \cite{nitsure-hilbert-quot}, \S 5, Theorem (Grothendieck), Theorem
  (Altman--Kleiman) (FGA Explained Ch.~5); cf.\ Grothendieck, FGA TDTE-IV.}
  Let \(S\) be a noetherian scheme, \(\pi : X \to S\) a projective morphism,
  \(L\) a relatively very ample line bundle on \(X\), \(E\) a coherent
  \(\mathcal{O}_X\)-module, and \(\Phi \in \mathbb{Q}[\lambda]\). Then the
  functor \(\Quot^{\Phi,L}_{E/X/S}\) of \ref{mk-def:quot_functor} is
  representable by a projective \(S\)-scheme \(\Quot^{\Phi,L}_{E/X/S}\),
  equipped with a universal quotient
  \[
    q^{\mathrm{univ}} :
    \pi^*_{\Quot} E \twoheadrightarrow \mathcal{F}^{\mathrm{univ}}
    \quad \text{on } X \times_S \Quot^{\Phi,L}_{E/X/S},
  \]
  flat over \(\Quot^{\Phi,L}_{E/X/S}\) with constant Hilbert polynomial
  \(\Phi\) on every fiber. The Hilbert scheme is the special case
  \(E = \mathcal{O}_X\): \(\Hilb^{\Phi,L}_{X/S} = \Quot^{\Phi,L}_{\mathcal{O}_X/X/S}\)
  represents the functor of \(T\)-flat closed subschemes
  \(Y \subset X_T\) with fiberwise Hilbert polynomial \(\Phi\).
\end{theorem}
\begin{proof}

  The proof has four steps; each step is a sub-lemma in the next
  sub-sections.
  \begin{enumerate}
    \item \textbf{(Boundedness, \ref{mk-lem:quot_boundedness}.)} Pick the
      Castelnuovo--Mumford regularity bound \(m = m(\mathrm{rank}\,V,
      \mathrm{rank}\,W, \Phi)\) uniformly across all fibers (by Mumford's
      \(m\)-regularity theorem applied to \(E_s\) and to all coherent
      quotients of \(E_s\) with Hilbert polynomial \(\Phi\)). For \(r \ge m\)
      and any \(T \to S\) and any \(T\)-flat quotient
      \(q : E_T \twoheadrightarrow \mathcal{F}\) with Hilbert polynomial
      \(\Phi\), the higher direct images \(R^i (\pi_T)_* \mathcal{F}(r)\)
      and \(R^i (\pi_T)_* E_T(r)\) and \(R^i (\pi_T)_* \mathcal{G}(r)\)
      vanish for \(i \ge 1\) (where \(\mathcal{G} = \ker q\)), and the
      direct images \((\pi_T)_* \mathcal{F}(r)\), \((\pi_T)_* E_T(r)\),
      \((\pi_T)_* \mathcal{G}(r)\) are locally free of ranks fixed by the
      data; the canonical maps \((\pi_T)^* (\pi_T)_*(-)(r) \to (-)(r)\)
      are surjective. Cohomology-and-base-change (\ref{mk-thm:flat_base_change_cohomology},
      Stacks tag 02KH) propagates these properties across \(T\)-base change.
    \item \textbf{(Grassmannian embedding, \ref{mk-lem:quot_alpha_injective}.)}
      Reduce by \ref{mk-lem:quot_reduction_to_pi_star_W} to the case
      \(X = \mathbb{P}_S(V)\), \(E = \pi^* W\) with \(V, W\) locally free
      on \(S\), \(L = \mathcal{O}_{\mathbb{P}(V)}(1)\). Fix \(r \ge m\).
      The Grassmannian embedding
      \[
        \alpha : \Quot^{\Phi,L}_{E/X/S} \longrightarrow
        \mathrm{Grass}\!\left(
          W \otimes_{\mathcal{O}_S} \mathrm{Sym}^r V,\; \Phi(r)
        \right)
      \]
      sends \(\langle \mathcal{F}, q\rangle\) on \(T\) to the locally free
      quotient \((\pi_T)_*(q(r)) : (\pi_T)_* E_T(r) \to (\pi_T)_* \mathcal{F}(r)\).
      This is injective on \(T\)-valued points: from the map \(T \to
      \mathrm{Gr}_S(W \otimes \mathrm{Sym}^r V, \Phi(r))\) one recovers
      the kernel-pull-back \((\pi_T)^* (\pi_T)_* \mathcal{G}(r) \to E_T(r)\),
      whose cokernel is \(q(r) : E_T(r) \to \mathcal{F}(r)\);
      twisting back by \(-r\) recovers \(q\).
    \item \textbf{(Locally closed in the Grassmannian via flattening
      stratification.)} Let \(T_0 = \mathrm{Gr}_S(W \otimes \mathrm{Sym}^r V,
      \Phi(r))\) with universal quotient
      \(f : (W \otimes \mathrm{Sym}^r V)_{T_0} \twoheadrightarrow \mathcal{J}\)
      and kernel \(\mathcal{K}\). Form the composite
      \(h : \pi_{T_0}^* \mathcal{K} \hookrightarrow
      \pi_{T_0}^* (W \otimes \mathrm{Sym}^r V) = \pi_{T_0}^* (\pi_{T_0})_* E_{T_0}
      \twoheadrightarrow E_{T_0}\)
      and let \(q := \mathrm{coker}(h) : E_{T_0} \to \mathcal{F}\). By the
      flattening stratification (\ref{mk-thm:flattening_stratification_exists}
      of \ref{mk-chap:Picard_FlatteningStratification}) applied to \(\mathcal{F}\)
      on \(X_{T_0}\), there exists a locally closed subscheme
      \(T_0^\Phi \subseteq T_0\) such that for any \(\phi : Y \to T_0\), the
      pulled-back cokernel \(\mathcal{F}_Y\) is \(Y\)-flat with fiberwise
      Hilbert polynomial \(\Phi\) iff \(\phi\) factors through \(T_0^\Phi\).
      The pair \((T_0^\Phi, \text{universal cokernel})\) represents
      \(\Quot^{\Phi,L}_{E/X/S}\).
    \item \textbf{(Closed, hence projective,
      \ref{mk-lem:quot_valuative_criterion}.)} The functor
      \(\Quot^{\Phi,L}_{E/X/S}\) satisfies the valuative criterion of
      properness with respect to DVRs (Exercise (7) of \S 1, proved
      using that flatness over a DVR is torsion-freeness). Combined
      with the locally closed embedding of step 3 into the projective
      \(\mathrm{Gr}_S \hookrightarrow \mathbb{P}_S(\bigwedge^{\Phi(r)}
      (W \otimes \mathrm{Sym}^r V))\), this upgrades to a \emph{closed}
      embedding. Hence \(\Quot^{\Phi,L}_{E/X/S}\) is a projective \(S\)-scheme.
  \end{enumerate}
  Finally, by \ref{mk-lem:quot_reduction_to_pi_star_W}, the general case
  (arbitrary \(X\) projective over \(S\) and arbitrary coherent \(E\)) follows
  from the case \(X = \mathbb{P}(V)\), \(E = \pi^* W\): pick a closed
  embedding \(X \subset \mathbb{P}(V)\) and write \(E\) as a coherent quotient
  of \(\pi^*(W)(\nu)\) for some vector bundle \(W\) on \(S\) and integer \(\nu\);
  the Quot scheme of the original data is then a closed subscheme of the
  Quot scheme of the universal data, by the closed-embedding clause of
  \ref{mk-lem:quot_reduction_to_pi_star_W}.
\end{proof}

\section{Sub-lemmas of the construction}
\label{mk-sec:quot_sub_lemmas}

\begin{lemma}[Reduction to the universal case]
  \label{mk-lem:quot_reduction_to_pi_star_W}
  \dcref{custom}

  \uses{mk-def:quot_functor}
  \textit{Source: \cite{nitsure-hilbert-quot}, \S 5 (FGA Explained Ch.~5).}
  Let \(S\), \(X\), \(L\), \(E\), \(\Phi\) be as in \ref{mk-def:quot_functor}.
  \begin{enumerate}
    \item (Twist-shift.) For \(\nu \in \mathbb{Z}\), tensoring by
      \(L^{\otimes \nu}\) gives an isomorphism of functors
      \(\Quot^{\Phi,L}_{E/X/S} \xrightarrow{\sim}
      \Quot^{\Psi,L}_{E(\nu)/X/S}\) with \(\Psi(\lambda) = \Phi(\lambda + \nu)\).
    \item (Surjection induces closed embedding.) Given a surjection
      \(\phi : E \twoheadrightarrow G\) of coherent sheaves on \(X\), the
      induced natural transformation
      \(\Quot^{\Phi,L}_{G/X/S} \to \Quot^{\Phi,L}_{E/X/S}\) is a closed
      embedding of functors.
  \end{enumerate}
  Consequently, if \(\Quot^{\Phi,L}_{\pi^* W / \mathbb{P}(V)/S}\) is
  representable for all \(V\), \(W\) vector bundles on \(S\), then so is
  \(\Quot^{\Phi,L}_{E/X/S}\) for every \(X \subset \mathbb{P}(V)\) closed and
  every coherent \(E\) on \(X\) that is a quotient of \(\pi^*(W)(\nu)|_X\) for
  some \(W\) and \(\nu\).
\end{lemma}
\begin{proof}

  (i) is direct from the definitions: tensoring with \(L^{\otimes \nu}\) is
  an autoequivalence of coherent sheaves on \(X\) (and on every \(X_T\)),
  preserving surjectivity, \(T\)-flatness, and properness of support, and
  shifting the Hilbert polynomial by \(\nu\) in the variable.

  (ii) Given a family \(\langle \mathcal{F}, q\rangle\) on \(T\) in
  \(\Quot^{\Phi,L}_{E/X/S}\), the locus where \(q\) factors through
  \(\phi : E \twoheadrightarrow G\) is cut out by the vanishing of the
  composite \(\ker(\phi)|_{X_T} \hookrightarrow E_T \xrightarrow{q}
  \mathcal{F}\). Because \(\mathcal{F}\) is \(T\)-flat, the schematic image of
  this composite descends to a closed subscheme \(T' \subset T\) such that
  \(f : U \to T\) factors through \(T'\) iff \(q_U\) factors through \(\phi_U\).
  Hence \(\Quot^{\Phi,L}_{G/X/S} \hookrightarrow \Quot^{\Phi,L}_{E/X/S}\)
  is a closed embedding of subfunctors.
\end{proof}

\begin{lemma}[Uniform Castelnuovo--Mumford regularity bound]
  \label{mk-lem:quot_boundedness}
  \dcref{custom}

  \uses{mk-thm:flat_base_change_cohomology}
  \textit{Source: \cite{nitsure-hilbert-quot}, \S 5, "Use of \(m\)-Regularity" (FGA Explained
  Ch.~5).}
  Let \(X = \mathbb{P}_S(V)\), \(E = \pi^* W\) with \(V\), \(W\) vector bundles on
  \(S\) of ranks \(n+1\), \(p\), \(L = \mathcal{O}_{\mathbb{P}(V)}(1)\), and
  \(\Phi \in \mathbb{Q}[\lambda]\). There exists an integer
  \(m = m(n, p, \Phi)\), depending only on \(n\), \(p\), and \(\Phi\), such that
  for every geometric point \(s\) of \(S\) and every coherent quotient
  \(q : E_s \twoheadrightarrow \mathcal{F}\) with Hilbert polynomial
  \(\Phi\), the sheaves \(E_s\), \(\mathcal{F}\), and \(\ker q\) on \(X_s\) are all
  \(m\)-regular in the sense of Castelnuovo--Mumford. Consequently, for
  every \(T \to S\) and every \(T\)-flat quotient
  \(q : E_T \twoheadrightarrow \mathcal{F}\) on \(X_T\) with Hilbert
  polynomial \(\Phi\) and every \(r \ge m\), the conclusions \textbf{(*)} and
  \textbf{(**)} of \cite[\S 5, ``Use of \(m\)-Regularity'']{nitsure-hilbert-quot}
  hold: the direct images \((\pi_T)_* \mathcal{F}(r)\), \((\pi_T)_* E_T(r)\),
  \((\pi_T)_* \mathcal{G}(r)\) are locally free of fixed ranks; the
  canonical surjections \((\pi_T)^* (\pi_T)_*(-)(r) \twoheadrightarrow (-)(r)\)
  hold; and the higher direct images \(R^i (\pi_T)_* (-)(r)\) vanish for
  \(i \ge 1\).
\end{lemma}
\begin{proof}

  Each geometric fiber \(E_s\) is a trivial bundle on \(\mathbb{P}^n_{\kappa(s)}\)
  of rank \(p\), so \(E_s\) is \(0\)-regular. The Castelnuovo--Mumford
  \(m\)-regularity theorem of \cite{nitsure-hilbert-quot} \S 2 produces, for any fixed Hilbert
  polynomial \(\Phi\), a uniform integer \(m = m(n, p, \Phi)\) above which
  every coherent quotient of \(E_s\) with Hilbert polynomial \(\Phi\) -- and
  its kernel -- is \(m\)-regular. The Castelnuovo Lemma of \cite{nitsure-hilbert-quot} \S 2
  converts \(m\)-regularity into vanishing of \(H^i((-) (r))\) for \(i \ge 1\)
  and \(r \ge m\), and into global generation of \(H^0((-)(r))\). To
  propagate the fiber-wise statement to families over \(T\), apply the
  flat-base-change theorem (\ref{mk-thm:flat_base_change_cohomology}, Stacks
  tag 02KH) along \(T \to S\): cohomology of flat families commutes with
  base change, so the direct images \((\pi_T)_*(-)(r)\) are locally free of
  the rank dictated by the fiber Hilbert polynomials, and the canonical
  maps \((\pi_T)^* (\pi_T)_*(-)(r) \to (-)(r)\) inherit fiberwise
  surjectivity.
\end{proof}

\begin{lemma}[Grassmannian embedding \(\alpha\) is injective]
  \label{mk-lem:quot_alpha_injective}
  \dcref{custom}

  \uses{mk-lem:quot_boundedness}
  \textit{Source: \cite{nitsure-hilbert-quot}, \S 5, "Embedding Quot into Grassmannian"
  (FGA Explained Ch.~5).}
  In the universal setup \(X = \mathbb{P}_S(V)\), \(E = \pi^* W\), fix
  \(r \ge m\) from \ref{mk-lem:quot_boundedness}. The natural transformation
  \[
    \alpha :
    \Quot^{\Phi,L}_{E/X/S} \longrightarrow
    \mathrm{Grass}\!\left(
      W \otimes_{\mathcal{O}_S} \mathrm{Sym}^r V,\; \Phi(r)
    \right),
    \qquad
    \langle \mathcal{F}, q\rangle \longmapsto
    \langle (\pi_T)_*\mathcal{F}(r),\, (\pi_T)_*(q(r))\rangle
  \]
  is well-defined (the target is locally free of rank \(\Phi(r)\), and the
  canonical isomorphism \(\pi_* E(r) \cong W \otimes_{\mathcal{O}_S}
  \mathrm{Sym}^r V\) gives the source) and \emph{injective} on \(T\)-points:
  the data of the locally free quotient at level \(r\) determines
  \(q : E_T \to \mathcal{F}\).
\end{lemma}

\begin{proof}

  Suppose \(\alpha(T)\) sends both \(\langle \mathcal{F}_1, q_1\rangle\) and
  \(\langle \mathcal{F}_2, q_2\rangle\) to the same \(T\)-point of the
  Grassmannian. Then the pulled-back kernel
  \((\pi_T)^* (\pi_T)_* \mathcal{G}_i(r) \to (\pi_T)^* (\pi_T)_* E_T(r)
  = \pi_{T_0}^* (W \otimes \mathrm{Sym}^r V)_T\)
  is the same. The diagram \textbf{(**)} of
  \ref{mk-lem:quot_boundedness} gives a right-exact sequence
  \[
    (\pi_T)^* (\pi_T)_* \mathcal{G}_i(r) \xrightarrow{h}
    E_T(r) \xrightarrow{q_i(r)} \mathcal{F}_i(r) \to 0,
  \]
  so \(q_i(r)\) is recovered as the cokernel of \(h\); twisting by \(-r\)
  recovers \(q_i\), giving \(q_1 = q_2\) and hence
  \(\langle \mathcal{F}_1, q_1\rangle = \langle \mathcal{F}_2, q_2\rangle\).
\end{proof}

\begin{lemma}[Valuative criterion of properness for \(\Quot\)]
  \label{mk-lem:quot_valuative_criterion}
  \dcref{custom}

  \textit{Source: \cite{nitsure-hilbert-quot}, \S 5, "Valuative Criterion for Properness"
  (FGA Explained Ch.~5); EGA IV(2) 2.8.1.}
  The functor \(\Quot^{\Phi,L}_{E/X/S}\) satisfies the valuative criterion
  for properness over \(S\): for every discrete valuation ring \(R\) with
  fraction field \(K\) and every morphism \(\Spec R \to S\), the restriction
  map
  \[
    \Quot^{\Phi,L}_{E/X/S}(\Spec R)
    \longrightarrow
    \Quot^{\Phi,L}_{E/X/S}(\Spec K)
  \]
  is bijective.
\end{lemma}

\begin{proof}

  Given a \(\Spec K\)-point \(\langle \mathcal{F}_K, q_K\rangle\) on \(X_K\),
  let \(j : X_K \hookrightarrow X_R\) be the open inclusion. Define
  \(\overline{\mathcal{F}}\) as the image of the composite
  \(E_R \to j_* E_K \to j_* \mathcal{F}_K\), with induced surjection
  \(\overline{q} : E_R \to \overline{\mathcal{F}}\). The sheaf
  \(\overline{\mathcal{F}}\) is torsion-free over \(R\) (being a subsheaf of
  \(j_* \mathcal{F}_K\)), hence \(R\)-flat (flatness over a DVR
  \(=\) torsion-freeness), and restricts on the generic fiber to
  \(\mathcal{F}_K\). The pair \((\overline{\mathcal{F}}, \overline{q})\) is
  the unique extension of \(\langle \mathcal{F}_K, q_K\rangle\) to a
  \(\Spec R\)-point (uniqueness: any other extension agrees on the dense
  open \(X_K \subset X_R\) and is torsion-free, hence equals the image
  construction). The Hilbert polynomial is constant on connected
  \(\Spec R\) by flatness.
\end{proof}

\section{Cohomology and base change}
\label{mk-sec:quot_iter195_bc_substrate}
\label{mk-sec:cohomology_base_change}


\begin{theorem}[Degree-zero flat base change]\label{mk-thm:flat_base_change_cohomology}
  \dcref{02KH,custom}

  \uses{mk-def:quot_canonical_basechange_map}
  \textit{Source: \cite{stacks-project}, tag 02KH (lemma-flat-base-change-cohomology).}
  Let
  \[
    \begin{array}{ccc}
      X' & \xrightarrow{g'} & X \\
      \downarrow f' &  & \downarrow f \\
      S' & \xrightarrow{g} & S
    \end{array}
  \]
  be a cartesian diagram of schemes with \(g\) flat and \(f\) quasi-compact and
  quasi-separated. Let \(\mathcal{F}\) be a quasi-coherent
  \(\mathcal{O}_X\)-module with pull-back
  \(\mathcal{F}' = (g')^* \mathcal{F}\). Then the canonical base-change map
  \[
    g^* f_* \mathcal{F} \longrightarrow f'_* \mathcal{F}'
  \]
  is an isomorphism.
\end{theorem}
\begin{proof}\uses{mk-thm:quot_canonical_basechange_isIso}
  The displayed morphism is the canonical base-change map of
  \ref{mk-def:quot_canonical_basechange_map}. Its invertibility under the stated
  hypotheses is \ref{mk-thm:quot_canonical_basechange_isIso}; taking the associated
  isomorphism proves the claim.
\end{proof}

\subsection{Project-side base-change substrate}
\label{mk-subsec:quot_project_basechange}


\begin{lemma}
[Injectivity along a finite cover, by flat base change]
  \label{mk-lem:section_base_change_injective_cover}
  \dcref{custom}
  Let \(A \to B\) be a flat homomorphism of commutative rings, let \(M\) be
  an \(A\)-module and \(P\) a \(B\)-module, and let \((M_i)_{i \in \iota}\),
  \((P_i)_{i \in \iota}\) be finite families of modules with \(A\)-linear
  maps \(\mathrm{res}_i : M \to M_i\) jointly injective and \(B\)-linear
  maps \(\mathrm{res}'_i : P \to P_i\). Suppose \(s : B \otimes_A M \to P\)
  is \(B\)-linear and, for each \(i\), there is an injective \(B\)-linear
  comparison \(\varepsilon_i : B \otimes_A M_i \to P_i\) with
  \(\mathrm{res}'_i(s(1 \otimes m)) = \varepsilon_i(1 \otimes
  \mathrm{res}_i(m))\) for all \(m \in M\). Then \(s\) is injective.
\end{lemma}
\begin{proof}An element \(z\) of \(\ker s\) has, for every \(i\), image \(0\) under
  \(\varepsilon_i \circ (\mathrm{id}_B \otimes \mathrm{res}_i)\) (by
  \(B\)-linearity, the intertwining hypothesis extends from the generators
  \(1 \otimes m\) to all of \(B \otimes_A M\)), hence lies in the kernel of
  \(\mathrm{id}_B \otimes \mathrm{res}_i\) by injectivity of \(\varepsilon_i\).
  Joint injectivity of the \(\mathrm{res}_i\) on \(M\) transports, via
  flatness of \(B\) over \(A\), to injectivity of the induced map
  \(B \otimes_A M \to \prod_i (B \otimes_A M_i)\) (flatness of \(B\) means
  tensoring with \(B\) preserves injective \(A\)-linear maps, applied to
  \(\prod_i \mathrm{res}_i\)); since \(z\) lies in the kernel of every
  component of this map, \(z = 0\).
\end{proof}

\begin{lemma}
[The Stacks 02KE equalizer ladder, abstract form]
  \label{mk-lem:section_base_change_gluing_ladder}
  \uses{mk-lem:section_base_change_injective_cover}
  \dcref{custom}
  \textit{Source: Stacks~\href{https://stacks.math.columbia.edu/tag/02KE}{02KE},
  module dialect.}
  Let \(A \to B\) be a flat homomorphism of commutative rings. Let \(M\) be
  an \(A\)-module equipped with a finite sheaf-style gluing datum: finite
  index sets \(\iota\) (pieces) and \(\kappa\) (overlaps), \(A\)-linear maps
  \(\mathrm{res}_i : M \to M_i\) jointly injective, and \(A\)-linear maps
  \(l_k, r_k : \prod_i M_i \to M_k\) such that \(l_k(\mathrm{res}(m)) =
  r_k(\mathrm{res}(m))\) for all \(m \in M\), \(k \in \kappa\)
  (compatibility), and every family \((s_i) \in \prod_i M_i\) with
  \(l_k(s) = r_k(s)\) for all \(k\) equals \(\mathrm{res}(m)\) for some
  \(m \in M\) (gluing). Let \(P\) be a \(B\)-module with an analogous
  gluing datum \((P_i, P_k, \mathrm{res}', l', r')\). Suppose given, for
  each \(i\), a bijective \(B\)-linear comparison
  \(\varepsilon_i : B \otimes_A M_i \to P_i\), and for each \(k\), an
  injective \(B\)-linear comparison \(\mu_k : B \otimes_A M_k \to P_k\),
  both intertwining the restriction and overlap maps on tensor generators.
  Then there is a \(B\)-linear equivalence
  \[
    \Phi : B \otimes_A M \;\xrightarrow{\ \sim\ }\; P
  \]
  compatible with all the restrictions:
  \(\mathrm{res}'_i(\Phi(z)) = \varepsilon_i((\mathrm{id}_B \otimes
  \mathrm{res}_i)(z))\) for every \(z\) and \(i\).
\end{lemma}
\begin{proof}\uses{mk-lem:section_base_change_injective_cover}
  Assemble the maps \(\mathrm{res}_i\) into \(T : M \to \prod_i M_i\) and
  the overlap maps into \(L, R : \prod_i M_i \to \prod_k M_k\); the gluing
  hypothesis on \(M\) says \(T\) is an isomorphism onto the equalizer of
  \(L\) and \(R\), and likewise \(T' : P \to \prod_i P_i\) is an
  isomorphism onto the equalizer of \(L'\) and \(R'\). Tensoring the
  equalizer presentation of \(M\) with the flat \(A\)-algebra \(B\)
  preserves it (a flat module is one for which tensoring preserves finite
  limits of modules, in particular equalizers), giving
  \(B \otimes_A M \cong \mathrm{eq}(B \otimes_A L,\, B \otimes_A R)\).
  Composing the bijective piece comparisons \(\varepsilon_i\) into a
  bijection \(\prod_i(B \otimes_A M_i) \cong \prod_i P_i\), and using
  injectivity of the overlap comparisons \(\mu_k\) to check that this
  bijection carries the equalizer condition for \((L,R)\) exactly onto the
  equalizer condition for \((L',R')\) (\ref{mk-lem:section_base_change_injective_cover}
  supplies the injectivity input for this check), transports the equalizer
  presentation of \(B \otimes_A M\) isomorphically onto the equalizer
  presentation of \(P\), i.e.\ onto \(P\) itself via \(T'\). The resulting
  equivalence \(\Phi\) is compatible with the restrictions by construction;
  this compatibility, together with the fact that the generators
  \(1 \otimes m\) generate \(B \otimes_A M\) as a \(B\)-module, determines
  \(\Phi\) uniquely.
\end{proof}

\begin{definition}[Canonical base-change map on sheaves of modules]\label{mk-def:quot_canonical_basechange_map}
  \dcref{custom}

  Given a cartesian square of schemes
  \(\square = (g' : X' \to X,\ f' : X' \to S',\ f : X \to S,\ g : S' \to S)\)
  with \(g \circ f' = f \circ g'\), the
  \emph{canonical base-change map} on sheaves of modules is the natural
  transformation
  \[
    g^* f_* \mathcal F \longrightarrow f'_* g'^* \mathcal F.
  \]
  It is the mate, under the pullback--pushforward adjunctions, of the canonical
  identification \(f'^*g^*\cong g'^*f^*\) supplied by the cartesian square. Its
  construction requires no flatness hypothesis.
\end{definition}

\begin{theorem}[Section-wise base-change isomorphism]\label{mk-thm:quot_canonical_basechange_app_app_isIso}
  \dcref{custom}

  \uses{mk-def:quot_canonical_basechange_map}
  Let \(\square\) be a cartesian square as in
  \ref{mk-def:quot_canonical_basechange_map} with \(f\) quasi-compact and
  quasi-separated, \(g\) flat, and \(\mathcal F\) quasi-coherent. For every
  open \(U \subseteq S'\), the section over \(U\) of the canonical
  base-change map
  \(g^* f_* \mathcal F \to f'_* g'^* \mathcal F\)
  is an isomorphism (Stacks 02KH(1) at \(i = 0\)).
\end{theorem}
\begin{proof}\uses{mk-thm:quot_canonical_basechange_app_app_isIso_of_affineCover,
    mk-thm:quot_canonical_basechange_app_app_isIso_of_isAffineOpen}
  Combine \ref{mk-thm:quot_canonical_basechange_app_app_isIso_of_isAffineOpen}
  (the section map is an isomorphism over every affine open of \(S'\)) with
  the basis-locality reduction
  \ref{mk-thm:quot_canonical_basechange_app_app_isIso_of_affineCover}.
\end{proof}

\begin{lemma}
[Base map at a reflexive inclusion]
  \label{mk-lem:quot_baseMap_le_refl}
  \dcref{custom}

  \uses{mk-def:quot_pullback_app_isoTensor_baseMap,
    mk-def:quot_pullback_app_isoTensor_unitAtV}
  For a morphism \(g : Y \to X\), a sheaf of \(\mathcal O_X\)-modules
  \(\mathcal N\), an open \(V \subseteq X\), and the reflexive inclusion
  \(g^{-1}(V) \subseteq g^{-1}(V)\), the section-level base map
  (\ref{mk-def:quot_pullback_app_isoTensor_baseMap}) coincides with the
  \(V\)-component of the adjunction unit
  (\ref{mk-def:quot_pullback_app_isoTensor_unitAtV}):
  \(\mathrm{baseMap}_g(\mathcal N, g^{-1}V \le g^{-1}V)
    = \Gamma\bigl(V, \eta_{\mathcal N}\bigr)\).
\end{lemma}
\begin{proof}The base map is by construction the unit component followed by the
  presheaf restriction along the inclusion \(g^{-1}(V) \subseteq g^{-1}(V)\);
  that inclusion is the identity of \(g^{-1}(V)\), and a presheaf sends the
  identity to the identity.
\end{proof}

\begin{theorem}
[Elementwise Beck--Chevalley compatibility of the canonical mate]
  \label{mk-thm:quot_canonical_basechange_app_baseMap_compat}
  \dcref{custom}

  \uses{mk-def:quot_canonical_basechange_map,
    mk-def:quot_pullback_app_isoTensor_baseMap}
  Let \(\square = (g', f', f, g)\) be a cartesian square as in
  \ref{mk-def:quot_canonical_basechange_map} (no flatness, compactness or
  affineness hypotheses), \(\mathcal F\) a sheaf of
  \(\mathcal O_X\)-modules, and \(V \subseteq S\), \(U \subseteq S'\) opens
  with \(U \subseteq g^{-1}(V)\), so that
  \(f'^{-1}(U) \subseteq g'^{-1}(f^{-1}(V))\). For every section
  \(t \in \Gamma(f_* \mathcal F, V) = \Gamma(\mathcal F, f^{-1}V)\), the
  \(U\)-section of the canonical base-change map sends the base-map image
  of \(t\) along \(g\) to the base-map image of \(t\) along \(g'\)
  (\ref{mk-def:quot_pullback_app_isoTensor_baseMap}):
  \[
    \Gamma\bigl(U, (\mathrm{canonicalBaseChangeMap}\ \square)_{\mathcal F}\bigr)
      \bigl(\mathrm{baseMap}_{g}(f_*\mathcal F)(t)\bigr)
    \;=\;
    \mathrm{baseMap}_{g'}(\mathcal F)(t)
    \quad\text{in } \Gamma\bigl(g'^*\mathcal F,\ f'^{-1}(U)\bigr)
      = \Gamma\bigl(f'_* g'^*\mathcal F,\ U\bigr).
  \]
\end{theorem}
\begin{proof}\uses{mk-lem:quot_baseMap_le_refl, mk-lem:baseMap_pullbackComp_apply,
    mk-lem:baseMap_pullback_comp_apply, mk-lem:baseMap_pullbackCongr_apply}
  The component at \(\mathcal F\) of the canonical base-change map is, by
  definition of the mate of the pseudofunctor \(2\)-isomorphism
  \(f'^* g^* \cong (f' g)^* = (g' f)^* \cong g'^* f^*\) under the two
  adjunctions \(f^* \dashv f_*\) and \(f'^* \dashv f'_*\), the composite
  \[
    g^* f_* \mathcal F
    \xrightarrow{\ \eta_{f'}\ }
    f'_* f'^* g^* f_* \mathcal F
    \xrightarrow{\ f'_*(\alpha_{f_*\mathcal F})\ }
    f'_* g'^* f^* f_* \mathcal F
    \xrightarrow{\ f'_* g'^*(\varepsilon_f)\ }
    f'_* g'^* \mathcal F .
  \]
  Evaluate on the section \(\mathrm{baseMap}_g(f_*\mathcal F)(t)\) over
  \(U\) and follow the stages. (1) The \(\eta_{f'}\)-stage produces the
  base-map image along \(f'\) at the reflexive inclusion
  (\ref{mk-lem:quot_baseMap_le_refl}); (2) the
  \(f'^* g^* \cong (f' g)^*\) stage composes the two nested base maps into
  \(\mathrm{baseMap}_{f' g}(f_*\mathcal F)(t)\)
  (\ref{mk-lem:baseMap_pullback_comp_apply}); (3) the transport along
  \(f' g = g' f\) turns it into
  \(\mathrm{baseMap}_{g' f}(f_*\mathcal F)(t)\)
  (\ref{mk-lem:baseMap_pullbackCongr_apply}); (4) the inverse of the
  \((g' f)^* \cong g'^* f^*\) stage splits it back as
  \(\mathrm{baseMap}_{g'}\bigl(f^* f_* \mathcal F\bigr)
    \bigl(\mathrm{baseMap}_{f}(f_*\mathcal F)(t)\bigr)\)
  (\ref{mk-lem:baseMap_pullback_comp_apply} again, cancelled against its
  inverse); (5) naturality of the base map in the sheaf argument
  (\ref{mk-lem:baseMap_pullbackComp_apply}) moves the counit stage inside:
  the result is
  \(\mathrm{baseMap}_{g'}(\mathcal F)\bigl(
     \Gamma(f^{-1}V, \varepsilon_f)(\mathrm{baseMap}_{f}(f_*\mathcal F)(t))
   \bigr)\);
  and (6) by \ref{mk-lem:quot_baseMap_le_refl} and the right-triangle
  identity \(f_*(\varepsilon_f) \circ \eta_{f}|_{f_*} = \mathrm{id}\) of the
  adjunction \(f^* \dashv f_*\), the inner section is \(t\) itself.
\end{proof}

\begin{lemma}
[Affineness of the fibre-product cover pieces]
  \label{mk-lem:quot_isAffineOpen_preimage_inf}
  \dcref{custom}

  \uses{mk-def:quot_canonical_basechange_map}
  Let \(\square = (g', f', f, g)\) be a cartesian square as in
  \ref{mk-def:quot_canonical_basechange_map} and let \(V \subseteq S\),
  \(U \subseteq S'\), \(V' \subseteq X\) be affine opens with
  \(U \subseteq g^{-1}(V)\) and \(V' \subseteq f^{-1}(V)\). Then the open
  \(g'^{-1}(V') \cap f'^{-1}(U) \subseteq X'\) is affine.
\end{lemma}
\begin{proof}Restricting the cartesian square to the opens \(V, U, V'\) and
  \(g'^{-1}(V') \cap f'^{-1}(U)\) yields a cartesian square of schemes
  exhibiting \(g'^{-1}(V') \cap f'^{-1}(U)\) as the fibre product
  \(V' \times_V U\). A fibre product of affine schemes over an affine
  scheme is affine, and affineness transports along the comparison
  isomorphism with the categorical pullback.
\end{proof}

\begin{lemma}
[Section algebras of a flat morphism are flat]
  \label{mk-lem:quot_flat_gamma_appLE}
  \dcref{custom}

  Let \(g : S' \to S\) be a flat morphism of schemes and \(V \subseteq S\),
  \(U \subseteq S'\) affine opens with \(U \subseteq g^{-1}(V)\). Then
  \(\Gamma(S', U)\) is a flat \(\Gamma(S, V)\)-algebra via the induced ring
  map.
\end{lemma}
\begin{proof}Flatness of a morphism of schemes is affine-local: it holds if and only
  if for every pair of affine opens \(U \subseteq g^{-1}(V)\) the induced
  ring map \(\Gamma(S, V) \to \Gamma(S', U)\) is flat.
\end{proof}

\begin{theorem}
[\(H^0\) flat base change over a compatible affine pair]
  \label{mk-thm:quot_pullback_baseMap_sectionEquiv_qc}
  \dcref{custom}

  \uses{mk-def:quot_canonical_basechange_map,
    mk-def:quot_pullback_app_isoTensor_baseMap, mk-lem:quot_flat_gamma_appLE}
  \textit{Source: \cite{stacks-project}, tags 02KE and 02KH.}
  Let \(\square = (g', f', f, g)\) be a cartesian square as in
  \ref{mk-def:quot_canonical_basechange_map} with \(f\) quasi-compact and
  quasi-separated, \(g\) flat, and \(\mathcal F\) quasi-coherent on \(X\).
  Let \(V \subseteq S\) and \(U \subseteq S'\) be affine opens with
  \(U \subseteq g^{-1}(V)\), and write \(A := \Gamma(S, V)\),
  \(B := \Gamma(S', U)\), a flat \(A\)-algebra. Then there is a
  \(B\)-linear equivalence
  \[
    B \otimes_A \Gamma\bigl(\mathcal F,\ f^{-1}(V)\bigr)
    \;\xrightarrow{\ \sim\ }\;
    \Gamma\bigl(g'^* \mathcal F,\ f'^{-1}(U)\bigr)
  \]
  sending \(1 \otimes x\) to the canonical base-map image
  \(\mathrm{baseMap}_{g'}(\mathcal F)(x)\)
  (\ref{mk-def:quot_pullback_app_isoTensor_baseMap}); this condition
  determines the equivalence on the \(B\)-module generators, hence
  uniquely.
\end{theorem}
\begin{proof}\uses{mk-def:quot_pullback_app_isoTensor, mk-lem:quot_baseMap_le_refl,
    mk-lem:baseMap_pullbackComp_apply, mk-lem:baseMap_pullback_comp_apply,
    mk-lem:quot_isAffineOpen_preimage_inf, mk-lem:quot_flat_gamma_appLE}
  Since \(V\) is affine and \(f\) is quasi-compact, \(Y := f^{-1}(V)\) is a
  quasi-compact open of \(X\); choose finitely many affine opens
  \(V_1, \dots, V_n \subseteq X\) with \(Y = \bigcup_i V_i\). Because \(f\)
  is quasi-separated, each intersection \(V_i \cap V_j\) is quasi-compact,
  hence a finite union of affine opens \(V_{ijk}\). The sheaf axiom for
  \(\mathcal F\) presents \(\Gamma(\mathcal F, Y)\) as the equalizer of the
  two restriction maps
  \(\prod_i \Gamma(\mathcal F, V_i) \rightrightarrows
    \prod_{i,j,k} \Gamma(\mathcal F, V_{ijk})\),
  an equalizer of \(A\)-modules with finitely many factors. Since \(B\) is
  flat over \(A\) and the products are finite, applying
  \(B \otimes_A (-)\) preserves this equalizer:
  \(B \otimes_A \Gamma(\mathcal F, Y)\) is the equalizer of
  \(\prod_i B \otimes_A \Gamma(\mathcal F, V_i) \rightrightarrows
    \prod_{i,j,k} B \otimes_A \Gamma(\mathcal F, V_{ijk})\).

  On the primed side set \(W_i := g'^{-1}(V_i) \cap f'^{-1}(U)\). Then
  \(W_i\), as a scheme, is the fibre product \(V_i \times_V U\) of affine
  schemes over the affine \(V\), hence is an affine open of \(X'\)
  (\ref{mk-lem:quot_isAffineOpen_preimage_inf}); the
  \(W_i\) cover \(f'^{-1}(U)\) because the \(V_i\) cover \(f^{-1}(V)\) and
  \(f'^{-1}(U) \subseteq g'^{-1}(f^{-1}(V))\). The sheaf axiom for
  \(g'^* \mathcal F\) presents \(\Gamma(g'^*\mathcal F, f'^{-1}(U))\) as the
  equalizer of
  \(\prod_i \Gamma(g'^*\mathcal F, W_i) \rightrightarrows
    \prod_{i,j,k} \Gamma(g'^*\mathcal F, W_{ijk})\)
  with \(W_{ijk} := g'^{-1}(V_{ijk}) \cap f'^{-1}(U)\) (affine opens
  covering \(W_i \cap W_j\)).

  Per index \(i\): since \(W_i = V_i \times_V U\) is an affine fibre
  product, \(\Gamma(X', W_i) = \Gamma(X, V_i) \otimes_A B\) (global sections
  of an affine fibre product form the pushout of the section rings), and the
  affine section formula for the module pullback
  (\ref{mk-def:quot_pullback_app_isoTensor}, applied to \(g'\) and the affine
  pair \(W_i \subseteq g'^{-1}(V_i)\)) gives
  \[
    \Gamma(g'^*\mathcal F, W_i)
    \cong \Gamma(X', W_i) \otimes_{\Gamma(X, V_i)} \Gamma(\mathcal F, V_i)
    \cong \bigl(\Gamma(X, V_i) \otimes_A B\bigr)
      \otimes_{\Gamma(X, V_i)} \Gamma(\mathcal F, V_i)
    \cong B \otimes_A \Gamma(\mathcal F, V_i),
  \]
  the last step by tensor cancellation; the composite sends
  \(b \otimes x \mapsto b \cdot \mathrm{baseMap}_{g'}(x|_{V_i})\), so it
  commutes with the restriction maps on both sides (the base map is
  natural in the open pair; \ref{mk-lem:baseMap_pullbackComp_apply},
  \ref{mk-lem:quot_baseMap_le_refl}). The same computation applies verbatim
  to each affine \(W_{ijk}\) over \(V_{ijk}\), giving on the overlap
  columns \emph{injective} comparisons
  \(B \otimes_A \Gamma(\mathcal F, V_i \cap V_j) \to
    \Gamma(g'^*\mathcal F, W_i \cap W_j)\):
  indeed \(\Gamma(\mathcal F, V_i \cap V_j)\) embeds by sheaf separatedness
  into \(\prod_k \Gamma(\mathcal F, V_{ijk})\), flatness of \(B\) over \(A\)
  preserves the injection, and on each \(V_{ijk}\)-factor the comparison is
  the isomorphism above.

  The two equalizer presentations are therefore connected by a commuting
  ladder whose product-column rungs are isomorphisms and whose
  discriminating-column rungs are injections; hence the induced map of
  equalizers
  \(B \otimes_A \Gamma(\mathcal F, Y) \to
    \Gamma(g'^*\mathcal F, f'^{-1}(U))\)
  is an isomorphism (an equalizer of two diagrams compared by isomorphisms
  on the product column and injections on the target column is preserved:
  surjectivity of the equalizer comparison uses injectivity of the overlap
  rung to check the equalizing condition upstairs). Finally, chasing the
  generator \(1 \otimes x\) through the ladder: its image in the \(i\)-th
  factor is \(\mathrm{baseMap}_{g'}(x|_{V_i})\), which is the
  \(W_i\)-restriction of \(\mathrm{baseMap}_{g'}(x)\) by compatibility of
  the base map with restrictions; since the equalizer includes into the
  product, the equivalence sends \(1 \otimes x\) to
  \(\mathrm{baseMap}_{g'}(\mathcal F)(x)\), as claimed.
\end{proof}

\begin{theorem}
[Affine-pair workhorse for the base-change isomorphism]
  \label{mk-thm:quot_canonical_basechange_app_app_isIso_of_le_preimage}
  \dcref{custom}

  \uses{mk-def:quot_canonical_basechange_map}
  Let \(\square\) be a cartesian square as in
  \ref{mk-def:quot_canonical_basechange_map} with \(f\) quasi-compact and
  quasi-separated, \(g\) flat, and \(\mathcal F\) quasi-coherent. For every
  pair of affine opens \(V \subseteq S\), \(U \subseteq S'\) with
  \(U \subseteq g^{-1}(V)\), the section over \(U\) of the canonical
  base-change map \(g^* f_* \mathcal F \to f'_* g'^* \mathcal F\) is an
  isomorphism.
\end{theorem}
\begin{proof}\uses{mk-def:pullback_app_isoTensor_sigma, mk-lem:pushforward_isQuasicoherent,
    mk-thm:quot_pullback_baseMap_sectionEquiv_qc,
    mk-thm:quot_canonical_basechange_app_baseMap_compat}
  Write \(A := \Gamma(S, V)\), \(B := \Gamma(S', U)\). The pushforward
  \(f_* \mathcal F\) is quasi-coherent
  (\ref{mk-lem:pushforward_isQuasicoherent}), so the affine section formula
  (\ref{mk-def:pullback_app_isoTensor_sigma}, applied to \(g\) and
  \(f_*\mathcal F\)) provides a \(B\)-linear equivalence
  \(e_L : B \otimes_A \Gamma(\mathcal F, f^{-1}V) \cong
    \Gamma(g^* f_* \mathcal F, U)\) with
  \(e_L(1 \otimes t) = \mathrm{baseMap}_{g}(f_*\mathcal F)(t)\), and the
  \(H^0\) flat-base-change equivalence
  (\ref{mk-thm:quot_pullback_baseMap_sectionEquiv_qc}) provides
  \(e_R : B \otimes_A \Gamma(\mathcal F, f^{-1}V) \cong
    \Gamma(f'_* g'^* \mathcal F, U)\) with
  \(e_R(1 \otimes t) = \mathrm{baseMap}_{g'}(\mathcal F)(t)\).
  The section over \(U\) of the canonical base-change map is
  \(B\)-linear, and by
  \ref{mk-thm:quot_canonical_basechange_app_baseMap_compat} it carries
  \(e_L(1 \otimes t)\) to \(e_R(1 \otimes t)\) for every \(t\); since the
  elements \(1 \otimes t\) generate
  \(B \otimes_A \Gamma(\mathcal F, f^{-1}V)\) as a \(B\)-module and all
  three maps are \(B\)-linear, it equals the composite
  \(e_R \circ e_L^{-1}\), hence is bijective.
\end{proof}

\begin{theorem}[Affine-base case of the section-wise base-change isomorphism]\label{mk-thm:quot_canonical_basechange_app_app_isIso_of_isAffineBase}
  \dcref{custom}

  \uses{mk-def:quot_canonical_basechange_map}
  Assume in addition that the base \(S\) is affine, that \(f\) is
  quasi-compact and quasi-separated, that \(g\) is flat, and that
  \(\mathcal{F}\) is quasi-coherent. Then for every affine open
  \(U \subseteq S'\) the section of the canonical base-change map
  (\ref{mk-def:quot_canonical_basechange_map}) over \(U\) is an isomorphism.
\end{theorem}
\begin{proof}\uses{mk-thm:quot_canonical_basechange_app_app_isIso_of_le_preimage}
  Take \(V := S\), affine by hypothesis; every affine open
  \(U \subseteq S'\) satisfies \(U \subseteq g^{-1}(S)\), so the affine-pair
  workhorse
  \ref{mk-thm:quot_canonical_basechange_app_app_isIso_of_le_preimage}
  applies to the pair \((V, U)\).
\end{proof}

\begin{theorem}[Affine-open case of the section-wise base-change isomorphism]\label{mk-thm:quot_canonical_basechange_app_app_isIso_of_isAffineOpen}
  \dcref{custom}

  \uses{mk-def:quot_canonical_basechange_map}
  With \(f\) quasi-compact and quasi-separated, \(g\) flat, and
  \(\mathcal{F}\) quasi-coherent (but \(S\) now arbitrary), the section of
  the canonical base-change map (\ref{mk-def:quot_canonical_basechange_map})
  over an affine open \(U \subseteq S'\) is an isomorphism.
\end{theorem}
\begin{proof}\uses{mk-thm:quot_canonical_basechange_app_app_isIso_of_le_preimage,
    mk-lem:modules_isIso_of_isBasis}
  The affine opens \(W \subseteq S'\) that are \emph{compatible} with some
  affine open \(V \subseteq S\) (that is, \(W \subseteq g^{-1}(V)\)) form a
  basis of the topology of \(S'\): given \(x \in O\) with \(O\) open, pick
  an affine open \(V \ni g(x)\) of \(S\), then an affine open
  \(W \subseteq O \cap g^{-1}(V)\) containing \(x\). On every such pair the
  section of the canonical base-change map is an isomorphism by the
  affine-pair workhorse
  (\ref{mk-thm:quot_canonical_basechange_app_app_isIso_of_le_preimage}); by
  basis locality for morphisms of sheaves of modules
  (\ref{mk-lem:modules_isIso_of_isBasis}) the map
  \(g^* f_* \mathcal F \to f'_* g'^* \mathcal F\) is then an isomorphism of
  sheaves, and in particular its section over the given affine \(U\) (in
  fact over every open) is an isomorphism.
\end{proof}

\begin{theorem}[Basis-locality reduction for the section-wise base-change isomorphism]\label{mk-thm:quot_canonical_basechange_app_app_isIso_of_affineCover}
  \dcref{custom}

  \uses{mk-def:quot_canonical_basechange_map}
  Suppose the section of the canonical base-change map
  (\ref{mk-def:quot_canonical_basechange_map}) is an isomorphism over
  \emph{every} affine open of \(S'\). Then it is an isomorphism over every
  open \(U \subseteq S'\). It is the reduction half of the section-wise
  statement \ref{mk-thm:quot_canonical_basechange_app_app_isIso}, whose proof
  feeds the affine-open case
  \ref{mk-thm:quot_canonical_basechange_app_app_isIso_of_isAffineOpen} as the
  affine-open hypothesis.
\end{theorem}
\begin{proof}\uses{mk-lem:modules_isIso_of_isBasis}
  The affine opens form a basis of the topology of \(S'\). A morphism of
  sheaves of \(\mathcal O_{S'}\)-modules whose section map is an
  isomorphism over every member of a basis is an isomorphism of sheaves
  (\ref{mk-lem:modules_isIso_of_isBasis}, via stalks: injectivity and
  surjectivity of each stalk map are checked on germs represented on basic
  opens); its section over every open \(U\) is then an isomorphism.
\end{proof}

\begin{lemma}[Definitional section identity for pushforward of a pullback]\label{mk-lem:quot_pushforward_pullback_section_eq}
  \dcref{custom}

  \uses{mk-def:quot_canonical_basechange_map}
  Let \(f' : X' \to S'\) and \(g' : X' \to X\) be morphisms of schemes,
  \(\mathcal{F}\) a sheaf of \(\mathcal{O}_X\)-modules, and \(U \subseteq S'\)
  an open. The sections over \(U\) of the pushforward along \(f'\) of the
  pullback along \(g'\) of \(\mathcal{F}\) coincide \emph{definitionally}
  with the sections of the pullback \(g'^* \mathcal{F}\) over the preimage
  \(f'^{-1}(U)\):
  \[
    \Gamma\bigl(U,\ f'_* (g'^* \mathcal{F})\bigr)
    \;=\;
    \Gamma\bigl(f'^{-1}(U),\ g'^* \mathcal{F}\bigr).
  \]
  This records step~(3) of the affine-open construction
  (\ref{mk-thm:quot_canonical_basechange_app_app_isIso_of_isAffineOpen}): the
  target of the canonical base-change map
  (\ref{mk-def:quot_canonical_basechange_map}) is read through this identity.
\end{lemma}
\begin{proof}
  By definition of pushforward, the sections of \(f'_*\mathcal G\) over
  \(U\) are the sections of \(\mathcal G\) over \(f'^{-1}(U)\). Apply this
  identity to \(\mathcal G=g'^*\mathcal F\).
\end{proof}

\begin{theorem}[Canonical base-change map is a sheaf isomorphism]\label{mk-thm:quot_canonical_basechange_isIso}
  \dcref{custom}

  \uses{mk-def:quot_canonical_basechange_map}
  With \(f\) quasi-compact and quasi-separated, \(g\) flat, and
  \(\mathcal{F}\) quasi-coherent, the canonical base-change map
  \(g^* f_* \mathcal{F} \to f'_* g'^* \mathcal{F}\) of
  \ref{mk-def:quot_canonical_basechange_map} is an isomorphism of sheaves of
  modules. This is the canonical-morphism form of the degree-zero isomorphism
  in \ref{mk-thm:flat_base_change_cohomology}.
\end{theorem}
\begin{proof}\uses{mk-thm:quot_canonical_basechange_app_app_isIso}
  A morphism of sheaves of modules is an isomorphism if and only if its
  section map over every open is an isomorphism; apply
  \ref{mk-thm:quot_canonical_basechange_app_app_isIso}.
\end{proof}

\begin{definition}[Tensor-presentation of the pullback-of-sections functor]\label{mk-def:quot_pullback_app_isoTensor}
  \dcref{custom}

  \uses{mk-def:quot_canonical_basechange_map, mk-thm:quot_pullback_app_isoTensor_isBaseChange}
  Given a morphism \(g : X \to Y\) of schemes, affine opens
  \(U \subseteq Y\), \(V \subseteq X\) with \(g(V) \subseteq U\), and a
  quasi-coherent sheaf of modules \(\mathcal{N}\) on \(Y\), the
  \emph{tensor-presentation iso} packages the canonical iso
  \(\Gamma(V,\, g^* \mathcal{N}) \;\cong\;
   \Gamma(V,\, \mathcal{O}_X) \otimes_{\Gamma(U,\, \mathcal{O}_Y)}
     \Gamma(U,\, \mathcal{N})\)
  as a linear equivalence. It is induced by the adjunction-unit base map:
  \ref{mk-thm:quot_pullback_app_isoTensor_baseMap_isBaseChange} identifies
  the scalar extension of that map with an isomorphism, and hence with the
  displayed tensor-presentation equivalence.
\end{definition}

\begin{definition}[Adjunction-unit base linear map at a section]\label{mk-def:quot_pullback_app_isoTensor_unitAtV}
  \dcref{custom}

  Given a morphism \(g : Y \to X\) of schemes, a sheaf of
  \(\mathcal{O}_X\)-modules \(\mathcal{N}\), and an open \(V \subseteq X\),
  the unit of the pullback--pushforward adjunction yields a morphism of
  \(\mathcal{O}_X\)-modules \(\mathcal{N} \to g_*(g^* \mathcal{N})\).
  Evaluating its \(V\)-component gives the \(\Gamma(X, V)\)-linear
  \emph{unit-at-section} map
  \[
    \Gamma(\mathcal{N}, V)
    \;\longrightarrow\;
    \Gamma\bigl(g_*(g^* \mathcal{N}),\ V\bigr).
  \]
  This is the adjunction-unit map underlying the tensor-presentation
  isomorphism of \ref{mk-def:quot_pullback_app_isoTensor}.
\end{definition}
\begin{proof}
  Evaluate the unit natural transformation
  \(\mathcal N\to g_*g^*\mathcal N\) on the open \(V\). Naturality with
  respect to multiplication by sections of \(\mathcal O_X\) gives
  \(\Gamma(X,V)\)-linearity.
\end{proof}

\begin{definition}[Section-level base map for the pullback]\label{mk-def:quot_pullback_app_isoTensor_baseMap}
  \dcref{custom}

  \uses{mk-def:quot_pullback_app_isoTensor_unitAtV}
  For a morphism \(g : Y \to X\), a sheaf \(\mathcal{N}\) on \(X\), and
  compatible affine opens \(U \subseteq Y\), \(V \subseteq X\) with
  \(g(U) \subseteq V\), composing the unit-at-section map
  (\ref{mk-def:quot_pullback_app_isoTensor_unitAtV}) with the presheaf
  restriction from \(g^{-1}(V)\) down to \(U\) produces the
  \(\Gamma(X, V)\)-linear \emph{base map}
  \[
    \Gamma(\mathcal{N}, V)
    \;\longrightarrow\;
    \Gamma\bigl(g^* \mathcal{N},\ U\bigr),
  \]
  where the \(\Gamma(X, V)\)-action on the target is restriction of scalars
  along the ring map \(\Gamma(X, V) \to \Gamma(Y, U)\) induced by \(g\). This
  is the candidate linear map whose base-change property
  (\ref{mk-thm:quot_pullback_app_isoTensor_baseMap_isBaseChange}) yields the
  tensor presentation.
\end{definition}
\begin{proof}
  Restrict the section supplied by the adjunction unit from
  \(g^{-1}(V)\) to \(U\). The compatibility of the unit with the structure
  sheaf map makes the composite semilinear for
  \(\Gamma(X,V)\to\Gamma(Y,U)\), hence linear after restriction of scalars.
\end{proof}

\begin{theorem}[The base map is a base change]\label{mk-thm:quot_pullback_app_isoTensor_baseMap_isBaseChange}
  \dcref{custom}

  \uses{mk-def:quot_pullback_app_isoTensor_baseMap, mk-def:pullback_app_isoTensor_sigma}
  For compatible affine opens \(U \subseteq Y\), \(V \subseteq X\) of a
  morphism \(g : Y \to X\) and a quasi-coherent sheaf \(\mathcal{N}\) on
  \(X\), the base map of \ref{mk-def:quot_pullback_app_isoTensor_baseMap}
  exhibits \(\Gamma(g^* \mathcal{N}, U)\) as the base change of
  \(\Gamma(\mathcal{N}, V)\) along the ring map
  \(\Gamma(X, V) \to \Gamma(Y, U)\); i.e.\ the induced
  \(\Gamma(Y, U) \otimes_{\Gamma(X, V)} \Gamma(\mathcal{N}, V)
   \to \Gamma(g^* \mathcal{N}, U)\) is an isomorphism. The proof feeds the
  packaged section-level linear equivalence and its intertwining property
  (\ref{mk-def:pullback_app_isoTensor_sigma}) into the universal
  characterization of a base change.
\end{theorem}
\begin{proof}
  Choose the linear equivalence and its intertwining identity from
  \ref{mk-def:pullback_app_isoTensor_sigma}. The universal property of the
  tensor product identifies the scalar-extension map induced by the base
  map with this equivalence, since they agree on every pure tensor
  \(1\otimes x\). It is therefore an isomorphism, which is precisely the
  asserted base-change property.
\end{proof}

\begin{theorem}[Existence of the tensor-presentation equivalence]\label{mk-thm:quot_pullback_app_isoTensor_isBaseChange}
  \dcref{custom}

  \uses{mk-thm:quot_pullback_app_isoTensor_baseMap_isBaseChange}
  In the same affine setup, there exists a \(\Gamma(Y, U)\)-linear
  isomorphism
  \[
    \Gamma\bigl(g^* \mathcal{N},\ U\bigr)
    \;\cong\;
    \Gamma(Y, U) \otimes_{\Gamma(X, V)} \Gamma(\mathcal{N}, V),
  \]
  obtained from the base-change property
  (\ref{mk-thm:quot_pullback_app_isoTensor_baseMap_isBaseChange}) by passing to
  the induced equivalence. This is the existence statement consumed by the
  tensor-presentation iso \ref{mk-def:quot_pullback_app_isoTensor}.
\end{theorem}
\begin{proof}
  By \ref{mk-thm:quot_pullback_app_isoTensor_baseMap_isBaseChange}, the
  scalar-extension map induced by the base map is an isomorphism. Its
  inverse, viewed as a linear equivalence, gives the displayed
  tensor-presentation equivalence.
\end{proof}

\begin{definition}[\(\Sigma\)-pair packaging: section-level LinearEquiv + intertwining]\label{mk-def:pullback_app_isoTensor_sigma}
  \dcref{custom}

  \uses{mk-def:quot_pullback_app_isoTensor_baseMap, mk-def:quot_canonical_basechange_map,
        mk-lem:pullback_tildeIso, mk-lem:tildeIso_of_isQuasicoherent_isAffineOpen,
        mk-lem:pullback_of_openImmersion_iso_restrict, mk-lem:pushforward_isQuasicoherent}
  Write \(A=\Gamma(X,V)\) and \(B=\Gamma(Y,U)\). The packaged data consist
  of a \(B\)-linear equivalence
  \[
    f : B\otimes_A\Gamma(\mathcal N,V)
      \;\xrightarrow{\sim}\; \Gamma(g^*\mathcal N,U)
  \]
  together with the identity
  \(f(1\otimes x)=\operatorname{baseMap}(x)\) for every
  \(x\in\Gamma(\mathcal N,V)\). To construct \(f\), identify the
  restriction of \(\mathcal N\) to \(\operatorname{Spec} A\) with its
  associated tilde sheaf, pull this sheaf back along
  \(\operatorname{Spec} B\to\operatorname{Spec} A\), and apply the tilde
  base-change isomorphism. Transporting through the affine-open
  identifications gives the displayed equivalence. The naturality,
  composition, and open-immersion identities of
  \ref{mk-sec:quot_iter195_bc_substrate} show that this composite sends
  \(1\otimes x\) to the canonical base-map section.
\end{definition}

\subsection{Adjunction substrate for pullback of tilde sheaves}
\label{mk-sec:quot_analogist_gaps}

The affine pullback formula rests on the tilde--global-sections adjunction,
the extension--restriction-of-scalars adjunction, and the
pullback--pushforward adjunction. The following comparison identifies global
sections of a pushforward with restricted scalars; uniqueness of left
adjoints then yields the pullback formula for tilde sheaves used above.

\begin{lemma}[Tilde--global-sections adjunction on an affine scheme]
  \label{mk-lem:tilde_gamma_adjunction_mathlib}
  \dcref{custom}

  For a commutative ring \(R\), the tilde functor
  \(\widetilde{(\cdot)} : \mathrm{Mod}_{R} \to (\Spec R).\mathrm{Modules}\)
  is left adjoint to the global-sections functor
  \(\Gamma(\Spec R, -) : (\Spec R).\mathrm{Modules} \to \mathrm{Mod}_{R}\)
  (\texttt{moduleSpecΓFunctor}); the unit is the isomorphism
  \(M \to \Gamma(\widetilde M, \top)\) and the counit is the
  tilde--\(\Gamma\) comparison
  \(\widetilde{\Gamma(\mathcal N, \top)} \to \mathcal N\)
  (\texttt{fromTildeΓ}). Mathlib also supplies the two companion
  adjunctions used with it below: extension--restriction of scalars
  \(B \otimes_A - \dashv \mathrm{res}_\varphi\)
  (\texttt{ModuleCat.extendRestrictScalarsAdj}) and
  pullback--pushforward \(\pi^* \dashv \pi_*\) for sheaves of modules
  over schemes (\texttt{Scheme.Modules.pullbackPushforwardAdjunction}).
\end{lemma}

\begin{lemma}[Global sections of the pushforward as restricted scalars]
  \label{mk-lem:pullbackTilde_gammaBridge}
  \dcref{custom}

  \uses{mk-lem:tilde_gamma_adjunction_mathlib, mk-lem:modules_pullback_mathlib}
  For a ring homomorphism \(\varphi : A \to B\) with
  \(\pi = \Spec\varphi : \Spec B \to \Spec A\), the two functors
  \((\Spec B).\mathrm{Modules} \to \mathrm{Mod}_{A}\)
  \[
    \mathcal N \longmapsto \Gamma(\Spec A,\, \pi_{*}\mathcal N)
    \qquad\text{and}\qquad
    \mathcal N \longmapsto \mathrm{res}_{\varphi}\, \Gamma(\Spec B,\, \mathcal N)
  \]
  (global sections of the pushforward, with \(A\) acting through
  \(\Gamma(\pi^{\sharp})\); respectively global sections, with \(A\)
  acting through \(\varphi\)) are naturally isomorphic.
\end{lemma}
\begin{proof}
  \uses{mk-lem:tilde_gamma_adjunction_mathlib, mk-lem:modules_pullback_mathlib}
  On carriers the component at \(\mathcal N\) is restriction of sections
  along the equality of opens \(\pi^{-1}(\Spec A) = \Spec B\), a
  bijection. It is \(A\)-linear: for \(a \in A\) the action on the
  source is by \(\Gamma(\pi^{\sharp})(\iota_A(a))\) and on the target by
  \(\iota_B(\varphi(a))\), where
  \(\iota_R : R \cong \Gamma(\Spec R, \mathcal O)\) is the canonical
  identification; these agree by naturality of \(\iota\) in \(R\)
  applied to \(\varphi\) (that is,
  \(\Gamma(\pi^{\sharp}) \circ \iota_A = \iota_B \circ \varphi\)),
  combined with semilinearity of section restriction for the module
  sheaf \(\mathcal N\). Naturality in \(\mathcal N\) is naturality of
  section restriction against module-sheaf morphisms.
\end{proof}

\begin{lemma}[Pullback of tilde is tilde of base change (Stacks 01HQ)]
  \label{mk-lem:pullback_tildeIso}
  \dcref{custom}

  \uses{mk-def:quot_pullback_app_isoTensor_baseMap}
  \textit{Source: Stacks~\href{https://stacks.math.columbia.edu/tag/01HQ}{01HQ}
  / \href{https://stacks.math.columbia.edu/tag/0BJ8}{0BJ8} (pullback of
  quasi-coherent modules along a Spec morphism).}
  For a ring homomorphism \(\varphi : A \to B\) and an \(A\)-module
  \(M\), there is a canonical iso between the pullback of the
  tilde-sheaf and the tilde of the base-changed module:
  \[
    (\mathtt{Scheme.Modules.pullback}\, (\mathrm{Spec.map}\, \varphi)).\mathrm{obj}\,
      (\widetilde{M})
    \;\cong\;
    \widetilde{B \otimes_A M}.
  \]
\end{lemma}
\begin{proof}
  \uses{mk-def:quot_pullback_app_isoTensor_baseMap,
    mk-lem:tilde_gamma_adjunction_mathlib, mk-lem:pullbackTilde_gammaBridge}
  Write \(\pi = \Spec\varphi : \Spec B \to \Spec A\). Both composite
  functors
  \[
    \mathrm{Mod}_{A} \xrightarrow{\ \widetilde{(\cdot)}\ }
      (\Spec A).\mathrm{Modules} \xrightarrow{\ \pi^{*}\ }
      (\Spec B).\mathrm{Modules}
    \qquad\text{and}\qquad
    \mathrm{Mod}_{A} \xrightarrow{\ B \otimes_A -\ } \mathrm{Mod}_{B}
      \xrightarrow{\ \widetilde{(\cdot)}\ } (\Spec B).\mathrm{Modules}
  \]
  are left adjoints: the first of
  \(\mathcal N \mapsto \Gamma(\Spec A, \pi_{*}\mathcal N)\) (compose the
  tilde--\(\Gamma\) adjunction on \(\Spec A\),
  \ref{mk-lem:tilde_gamma_adjunction_mathlib}, with
  \(\pi^{*} \dashv \pi_{*}\)), the second of
  \(\mathcal N \mapsto \mathrm{res}_{\varphi}\,\Gamma(\Spec B, \mathcal N)\)
  (compose extension--restriction of scalars along \(\varphi\) with the
  tilde--\(\Gamma\) adjunction on \(\Spec B\)). By
  \ref{mk-lem:pullbackTilde_gammaBridge} the two right adjoints are
  naturally isomorphic, so uniqueness of left adjoints yields a natural
  isomorphism of the two composites; its component at \(M\) is the
  required isomorphism
  \(\pi^{*}\widetilde M \cong \widetilde{B \otimes_A M}\).

  For the section-level identity, apply the unit compatibility of the
  left-adjoint-uniqueness isomorphism at \(m \in M\): the unit of the
  first composed adjunction sends \(m\) to the canonical pullback base map
  on global sections of \(\widetilde M\), while the unit
  of the second sends \(m\) to the global section of
  \(\widetilde{B \otimes_A M}\) attached to \(1 \otimes m\); the
  compatibility square states precisely that the isomorphism carries
  the former to the latter.
\end{proof}

\subsection{qcqs section localization and Stacks 01XJ}
\label{mk-sec:quot_qcqs_section_localization}

The proof of \ref{mk-lem:pushforward_isQuasicoherent} generalizes the affine
section-localization keystone to quasi-compact quasi-separated opens by a
Mayer--Vietoris induction, then runs the tilde--$\Gamma$ comparison backwards
at each affine open of the base.

\begin{lemma}[Affine basic-open section localization for quasi-coherent modules]
  \label{mk-lem:qcoh_section_localization_basicOpen}
  \dcref{II.5.3,custom}

  \uses{mk-lem:qcoh_affine_isIso_fromTildeΓ,
    mk-lem:isLocalizedModule_restrict_of_isIso_fromTildeΓ}
  \textit{Source: Stacks~\href{https://stacks.math.columbia.edu/tag/01IB}{01IB}
  / Hartshorne II.5.3 (\texttt{lemma-invert-f-sections}, affine case).}
  Let $X$ be a scheme, $\mathcal{M}$ a quasi-coherent
  $\mathcal{O}_X$-module, $U \subseteq X$ an affine open, and
  $f \in \Gamma(X, U)$. Then the section restriction
  $\Gamma(\mathcal{M}, U) \to \Gamma(\mathcal{M}, D(f))$ exhibits the target as
  the localization $\Gamma(\mathcal{M}, U)[1/f]$ over $\Gamma(X, U)$.
\end{lemma}
\begin{proof}
  \uses{mk-lem:qcoh_affine_isIso_fromTildeΓ,
    mk-lem:isLocalizedModule_restrict_of_isIso_fromTildeΓ}
  Pull $\mathcal{M}$ back along the affine immersion
  $\operatorname{Spec} \Gamma(X, U) \to X$; the pullback is quasi-coherent
  (pullback along an open immersion preserves quasi-coherence), so its
  tilde--$\Gamma$ counit is an isomorphism
  (\ref{mk-lem:qcoh_affine_isIso_fromTildeΓ}), identifying it with
  $\widetilde{\Gamma(\mathcal{M}, U)}$; the sections of a tilde over a basic
  open are the localization
  (\ref{mk-lem:isLocalizedModule_restrict_of_isIso_fromTildeΓ}), and the
  section comparison transports the statement back to $X$.
\end{proof}

\begin{lemma}[qcqs section localization, torsion half]
  \label{mk-lem:qcqs_section_localization_torsion}
  \dcref{custom}

  \uses{mk-lem:qcoh_section_localization_basicOpen}
  \textit{Source: Stacks~\href{https://stacks.math.columbia.edu/tag/01P0}{01P0}
  (first half).}
  Let $X$ be a scheme, $\mathcal{M}$ a quasi-coherent $\mathcal{O}_X$-module,
  $W \subseteq X$ open, $g \in \Gamma(X, W)$, and $U \le W$ a quasi-compact
  open. If $x \in \Gamma(\mathcal{M}, U)$ restricts to zero on
  $U \cap D(g)$, then $(g|_U)^n \cdot x = 0$ for some $n \in \mathbb{N}$.
\end{lemma}
\begin{proof}
  \uses{mk-lem:qcoh_section_localization_basicOpen}
  Induction on the quasi-compact open $U$, adding one affine open at a time.
  For $U = \emptyset$ the section vanishes by sheaf separation over the empty
  cover. For $U = S \cup V$ with $V$ affine and the claim known on the
  quasi-compact $S$: the restriction of $x$ to $S$ is killed by a power of
  $g|_S$ by induction, and the restriction to $V$ is killed by a power of
  $g|_V$ by the injectivity half of
  \ref{mk-lem:qcoh_section_localization_basicOpen} at the affine $V$ (note
  $D(g|_V) = V \cap D(g)$). Taking the maximum $n$ of the two exponents,
  $(g|_U)^n \cdot x$ restricts to zero on both members of the cover
  $\{S, V\}$, hence vanishes by sheaf separation.
\end{proof}

\begin{lemma}[qcqs section localization, surjectivity half]
  \label{mk-lem:qcqs_section_localization_surj}
  \dcref{custom}

  \uses{mk-lem:qcoh_section_localization_basicOpen, mk-lem:qcqs_section_localization_torsion}
  \textit{Source: Stacks~\href{https://stacks.math.columbia.edu/tag/01P0}{01P0}
  (second half).}
  Let $X$, $\mathcal{M}$, $W$, $g$ be as above, with $W$ quasi-separated, and
  let $U \le W$ be quasi-compact. For every
  $y \in \Gamma(\mathcal{M}, U \cap D(g))$ there are $n \in \mathbb{N}$ and
  $x \in \Gamma(\mathcal{M}, U)$ with
  $x|_{U \cap D(g)} = (g|_{U \cap D(g)})^n \cdot y$.
\end{lemma}
\begin{proof}
  \uses{mk-lem:qcoh_section_localization_basicOpen, mk-lem:qcqs_section_localization_torsion}
  Induction on the quasi-compact $U$ as in
  \ref{mk-lem:qcqs_section_localization_torsion}. In the step $U = S \cup V$
  ($V$ affine), the induction hypothesis and the affine surjectivity half of
  \ref{mk-lem:qcoh_section_localization_basicOpen} produce candidate sections
  $x_S$ on $S$ and $x_V$ on $V$ with a common normalization exponent $n$.
  Their difference on the overlap $S \cap V$ --- quasi-compact because $W$ is
  quasi-separated --- restricts to zero on $(S \cap V) \cap D(g)$, so by the
  torsion half a further power $g^m$ kills it; after multiplying both
  candidates by $g^m$ they agree on the overlap and glue to a section $x$
  over $U$ with $x|_{U \cap D(g)} = g^{n+m} \cdot y$, checked by sheaf
  separation over the cover $\{S \cap D(g), D(g|_V)\}$ of $U \cap D(g)$.
\end{proof}

\begin{lemma}[qcqs section localization]
  \label{mk-lem:qcqs_section_localization}
  \dcref{custom}

  \uses{mk-lem:qcqs_section_localization_torsion, mk-lem:qcqs_section_localization_surj}
  \textit{Source: Stacks~\href{https://stacks.math.columbia.edu/tag/01P0}{01P0}.}
  Let $X$ be a scheme, $\mathcal{M}$ a quasi-coherent $\mathcal{O}_X$-module,
  and $W \subseteq X$ a quasi-compact, quasi-separated open. For every
  $g \in \Gamma(X, W)$ the section restriction
  $\Gamma(\mathcal{M}, W) \to \Gamma(\mathcal{M}, D(g))$ exhibits the target
  as the localization $\Gamma(\mathcal{M}, W)[1/g]$ over $\Gamma(X, W)$.
\end{lemma}
\begin{proof}
  \uses{mk-lem:qcqs_section_localization_torsion, mk-lem:qcqs_section_localization_surj}
  The element $g$ restricts to a unit of $\Gamma(X, D(g))$, so every power of
  $g$ acts invertibly on $\Gamma(\mathcal{M}, D(g))$; the surjectivity and
  torsion (injectivity) conditions of the localization are
  \ref{mk-lem:qcqs_section_localization_surj} and
  \ref{mk-lem:qcqs_section_localization_torsion} at $U := W$.
\end{proof}

\begin{lemma}[Converse tilde--$\Gamma$ transport at an affine open]
  \label{mk-lem:fromTildeGamma_pullback_fromSpec_converse}
  \dcref{custom}

  \uses{mk-lem:qcoh_section_localization_basicOpen}
  Let $Y$ be a scheme, $\mathcal{N}$ an $\mathcal{O}_Y$-module, and
  $U \subseteq Y$ an affine open with canonical immersion
  $j : \operatorname{Spec} \Gamma(Y, U) \to Y$. If for every
  $f' \in \Gamma(Y, U)$ the restriction
  $\Gamma(\mathcal{N}, U) \to \Gamma(\mathcal{N}, D(f'))$ is a localization at
  the powers of $f'$ over $\Gamma(Y, U)$, then the tilde--$\Gamma$ counit of
  the pullback $j^* \mathcal{N}$ is an isomorphism.
\end{lemma}
\begin{proof}
  The sections of $j^* \mathcal{N}$ over an open of the spectrum are the
  sections of $\mathcal{N}$ over its image ($j$ is an open immersion), with
  $j''(\top) = U$ and $j''(D(f')) = D(f')$; under these identifications --
  which are semilinear over the canonical section-ring comparisons, by the
  naturality of the immersion's application isomorphisms -- the hypothesis
  says exactly that all basic-open section restrictions of $j^* \mathcal{N}$
  on the spectrum are localizations, which characterizes the invertibility of
  the tilde--$\Gamma$ counit.
\end{proof}

\begin{lemma}[Pushforward basic-open section localization]
  \label{mk-lem:pushforward_section_localization}
  \dcref{custom}

  \uses{mk-lem:qcqs_section_localization}
  Let $\pi : X \to S$ be a quasi-compact, quasi-separated morphism of
  schemes, $\mathcal{F}$ a quasi-coherent $\mathcal{O}_X$-module,
  $U \subseteq S$ an affine open, and $f' \in \Gamma(S, U)$. Then the section
  restriction
  $\Gamma(\pi_* \mathcal{F}, U) \to \Gamma(\pi_* \mathcal{F}, D(f'))$ is a
  localization at the powers of $f'$ over $\Gamma(S, U)$.
\end{lemma}
\begin{proof}
  \uses{mk-lem:qcqs_section_localization}
  The preimage $W := \pi^{-1}(U)$ is quasi-compact and quasi-separated
  ($\pi$ is qcqs and $U$ is affine), and
  $\pi^{-1}(D_S(f')) = D_X(g)$ for $g := \pi^\sharp f' \in \Gamma(X, W)$. The
  sections of the pushforward are by definition the sections of
  $\mathcal{F}$ over these preimages, with the $\Gamma(S, U)$-action given
  through $\pi^\sharp$, so the statement is
  \ref{mk-lem:qcqs_section_localization} for $\mathcal{F}$ at $g$ over $W$,
  read through $\pi^\sharp$ (a localization over the base ring restricts to a
  localization along any ring map matching the multiplicative sets).
\end{proof}

\begin{lemma}[Pushforward of quasi-coherent modules is quasi-coherent (Stacks 01XJ)]
  \label{mk-lem:pushforward_isQuasicoherent}
  \dcref{custom}

  \textit{Source: Stacks~\href{https://stacks.math.columbia.edu/tag/01XJ}{01XJ}
  (preservation of quasi-coherence under pushforward).}
  Let \(f : X \to Y\) be a quasi-compact, quasi-separated morphism of
  schemes, and let \(\mathcal{F}\) be a quasi-coherent
  \(\mathcal{O}_X\)-module. Then \(f_* \mathcal{F}\) is a quasi-coherent
  \(\mathcal{O}_Y\)-module. Mathlib \texttt{b80f227} ships no
  \texttt{IsQuasicoherent.pushforward} instance at the
  \texttt{Scheme.Modules} layer.
\end{lemma}
\begin{proof}
  \uses{mk-lem:pushforward_section_localization, mk-lem:fromTildeGamma_pullback_fromSpec_converse}
  Quasi-coherence is local on $Y$, so it suffices to produce a presentation
  of the slice of $f_* \mathcal{F}$ over each affine open $U \subseteq Y$.
  By \ref{mk-lem:pushforward_section_localization} all basic-open section
  restrictions of $f_* \mathcal{F}$ at $U$ are localizations, so by
  \ref{mk-lem:fromTildeGamma_pullback_fromSpec_converse} the pullback of
  $f_* \mathcal{F}$ along $\operatorname{Spec} \Gamma(Y, U) \to Y$ is
  isomorphic to the tilde of its global sections. A tilde admits a global
  presentation; transporting it along
  $U \hookrightarrow Y = (\operatorname{Spec}\Gamma(Y,U) \cong U)
  \hookrightarrow Y$ yields a presentation of the geometric restriction of
  $f_* \mathcal{F}$ to $U$, hence of the slice. (The previously sketched
  ``right adjoints preserve quasi-coherence'' argument is not a proof:
  adjointness yields colimit preservation only.)
\end{proof}


The tensor-presentation equivalence uses three comparison isomorphisms:
a quasi-coherent sheaf on an affine open is the tilde sheaf associated to
its sections, tilde sheaves commute with affine base change, and pullback
along the spectrum presentation of an affine open agrees with restriction.
The following results record these ingredients separately.

\begin{lemma}[Quasi-coherent module on an affine open is \texttt{tilde} of its sections (Stacks 01I8)]\label{mk-lem:tildeIso_of_isQuasicoherent_isAffineOpen}
  \dcref{custom}

  \uses{mk-def:quot_canonical_basechange_map, mk-lem:qcoh_affine_isIso_fromTildeΓ}
  \textit{Source: Stacks~\href{https://stacks.math.columbia.edu/tag/01I8}{01I8}
  (a quasi-coherent module on an affine scheme is \(\widetilde M\) for
  some module).}
  Let \(X\) be a scheme, \(\mathcal N\) a quasi-coherent
  \(\mathcal O_X\)-module, and \(V \subset X\) an affine open. Then the
  pullback of \(\mathcal N\) along \texttt{hV.fromSpec} is canonically
  isomorphic to \texttt{tilde} of its sections:
  \[
    (\mathtt{Scheme.Modules.pullback}\,\mathtt{hV.fromSpec}).\mathrm{obj}\,
      \mathcal N
    \;\cong\;
    \widetilde{\,\Gamma(\mathcal N, V)\,}
    \quad\text{(on \(\mathrm{Spec}\,\Gamma(X, V)\)).}
  \]
  Mathlib \texttt{b80f227} ships the analog for \emph{free}
  quasi-coherent modules (\texttt{isIso\_fromTildeΓ\_of\_presentation}
  partial), but the general statement
  ``\([N.IsQuasicoherent] \Rightarrow N|_V \cong \widetilde{\Gamma(N,V)}\)
  on an affine open \(V\)'' (Stacks 01I8 \emph{lemma-affine-qc-module-iff})
  is unowned at the \texttt{Scheme.Modules} layer.
\end{lemma}
\begin{proof}
  \uses{mk-def:quot_canonical_basechange_map, mk-lem:modules_pullback_mathlib,
    mk-lem:modules_restrictFunctorIsoPullback_mathlib,
    mk-lem:qcoh_affine_isIso_fromTildeΓ}
  Write \(j = \mathtt{hV.fromSpec} : \Spec\Gamma(X,V) \to X\) (an open
  immersion with image \(V\)) and \(G = j^{*}\mathcal N\). Pullback along an
  open immersion preserves quasi-coherence, so \(G\) is a
  quasi-coherent module on the affine scheme \(\Spec\Gamma(X,V)\); by the
  affine structure theorem (\ref{mk-lem:qcoh_affine_isIso_fromTildeΓ},
  Stacks 01I8) the tilde--\(\Gamma\)
  counit \(\widetilde{\Gamma(G,\top)} \to G\) is an isomorphism.

  The canonical base map \(b : \Gamma(\mathcal N, V) \to \Gamma(G, \top)\)
  (the \(V\)-sections of the pullback--pushforward adjunction unit followed by
  restriction along \(\top \le j^{-1}V\)) is bijective. Indeed the two
  adjunctions \(\mathrm{restrict} \dashv j_{*}\) and \(j^{*} \dashv j_{*}\)
  share the right adjoint, so their units are identified through the canonical
  natural isomorphism \(\mathrm{restrict} \cong j^{*}\)
  (\ref{mk-lem:modules_restrictFunctorIsoPullback_mathlib}). Under this
  identification, the \(V\)-component of the unit of the restriction
  adjunction is the presheaf restriction of \(\mathcal N\) along the
  \emph{equality} of opens \(j(j^{-1}V) = V\), hence bijective; the component
  of the natural isomorphism is bijective; and the final restriction is along
  the equality \(\top = j^{-1}V\). Thus \(b\) is a composite of three
  bijections.

  Since \(b\) is a bijective \(\Gamma(X,V)\)-linear map (the
  \(\Gamma(X,V)\)-action on \(\Gamma(G,\top)\) being through the canonical
  ring isomorphism \(\Gamma(X,V) \cong \Gamma(\Spec\Gamma(X,V),\top)\)),
  it induces an isomorphism
  \(\widetilde{b} : \widetilde{\Gamma(\mathcal N,V)} \cong
  \widetilde{\Gamma(G,\top)}\). The required isomorphism is the inverse of
  \(\widetilde{b}\) followed by the counit. The characterizing section
  identity holds because \(\widetilde{(-)}\)-functoriality intertwines the
  canonical maps \(M \to \Gamma(\widetilde M, U)\) (naturality of
  \texttt{toOpen}), while the counit composed with \(\mathtt{toOpen}\) at
  \(\top\) is the identity restriction.
\end{proof}

\begin{lemma}[Section-level transport for pullback along an affine open's \texttt{fromSpec}]\label{mk-lem:pullback_of_openImmersion_iso_restrict}
  \dcref{custom}

  \uses{mk-def:quot_canonical_basechange_map}
  \textit{Source: Project-bespoke transport, formed from
  Stacks~\href{https://stacks.math.columbia.edu/tag/01HH}{01HH} +
  \texttt{tilde.isoTop}.}
  Let \(Y\) be a scheme, \(\mathcal N\) an \(\mathcal O_Y\)-module, and
  \(U \subset Y\) an affine open. Equip \(\Gamma((\mathrm{Spec}\,\Gamma(Y, U)),
  \top)\) with the \(\Gamma(Y, U)\)-algebra structure inherited from
  \texttt{Scheme.ΓSpecIso} (the canonical
  \(\Gamma(\mathrm{Spec}\,R, \top) \cong R\) for any commutative ring \(R\)),
  and equip
  \(\Gamma((\text{pullback}\,\mathtt{hU.fromSpec}).\mathrm{obj}\,\mathcal N,
  \top)\) with the \(\Gamma(Y, U)\)-module structure inherited via
  \texttt{Module.compHom}. Then there is a canonical \(\Gamma(Y, U)\)-linear
  isomorphism between this pullback-section module and the original section
  module \(\Gamma(\mathcal N, U)\):
  \[
    \Gamma((\text{pullback}\,\mathtt{hU.fromSpec}).\mathrm{obj}\,\mathcal N,
      \top)
    \;\simeq_{\Gamma(Y, U)}\;
    \Gamma(\mathcal N, U).
  \]
  The map identifies a section of the pulled-back sheaf at \(\top\) of
  \(\mathrm{Spec}\,\Gamma(Y, U)\) with the corresponding section of
  \(\mathcal N\) on the affine open \(U \subset Y\).
\end{lemma}
\begin{proof}
  \uses{mk-def:quot_canonical_basechange_map,
    mk-lem:modules_restrictFunctorIsoPullback_mathlib}
  The underlying additive equivalence is the composite of the
  \(\top\)-section of the natural isomorphism
  \(j^{*} \cong \mathrm{restrict}\,j\) for the open immersion
  \(j = \mathtt{hU.fromSpec}\)
  (\ref{mk-lem:modules_restrictFunctorIsoPullback_mathlib}) with the
  restriction of \(\mathcal N\) along the equality of opens
  \(j(\top) = U\); its \(\Gamma(Y, U)\)-linearity reduces (through
  \texttt{Hom.app\_smul} and semilinearity of section restriction) to
  the ring identity
  \(\iota^{-1} \circ (j^{\sharp}\top)^{-1} \circ
    \mathrm{res}_{U = j(\top)} = \mathrm{id}\), which follows from
  \texttt{fromSpec\_app\_self}.

  The additional \(\Sigma\)-pair characterization (T12, 2026-07-03) —
  the inverse of the equivalence is the canonical base map — is the
  unit-compatibility of \texttt{Adjunction.leftAdjointUniq} for the
  restriction adjunction versus the pullback--pushforward adjunction,
  as spelled out in \ref{mk-lem:baseMap_inv_step3_open_immersion}.
\end{proof}


\subsection{Finite generation of affine sections (Stacks 01PC, finite-type half)}

The chart-level input to generic flatness (Nitsure \S4): over an affine open
contained in a presenting-cover member, the sections of a finitely presented
module sheaf form a finite module. The route runs through the tilde--$\Gamma$
adjunction on the affine model.

\begin{lemma}[Finite generation from a finite generating family of the tilde]
  \label{mk-lem:module_finite_of_tilde_genSections}
  \dcref{custom}

  \textit{Source: Stacks~\href{https://stacks.math.columbia.edu/tag/01PC}{01PC}
  (finite-type modules on affines).}
  Let \(R\) be a commutative ring and \(N\) an \(R\)-module. If the associated
  sheaf \(\widetilde N\) on \(\operatorname{Spec} R\) admits a generating
  family of sections indexed by a finite type \(I\) — that is, an epimorphism
  of \(\mathcal O_{\operatorname{Spec} R}\)-modules
  \(\pi : \mathcal O^{\oplus I} \twoheadrightarrow \widetilde N\) — then \(N\)
  is a finite \(R\)-module.
\end{lemma}
\begin{proof}
  Precomposing \(\pi\) with the canonical isomorphism
  \(\widetilde{R^{\oplus I}} \cong \mathcal O^{\oplus I}\) gives an epimorphism
  \(\pi' : \widetilde{R^{\oplus I}} \twoheadrightarrow \widetilde N\).
  Transpose \(\pi'\) through the adjunction
  \(\widetilde{(\cdot)} \dashv \Gamma\): writing
  \(t : R^{\oplus I} \to \Gamma(\widetilde N)\) for the transpose,
  the counit factorization reads
  \(\widetilde{t} \mathbin{;} \varepsilon_{\widetilde N} = \pi'\).
  The counit \(\varepsilon\) of the tilde--\(\Gamma\) adjunction is invertible
  at every tilde: the triangle identity
  \(\widetilde{\eta_N} \mathbin{;} \varepsilon_{\widetilde N} = \mathrm{id}\)
  exhibits it as inverse to the tilde of the (invertible) unit. Hence
  \(\widetilde t = \pi' \mathbin{;} \varepsilon_{\widetilde N}^{-1}\) is an
  epimorphism; since \(\widetilde{(\cdot)}\) is faithful, \(t\) is an
  epimorphism of \(R\)-modules, i.e.\ surjective. Composing with the inverse
  of the unit isomorphism \(N \cong \Gamma(\widetilde N)\) exhibits \(N\) as a
  quotient of the finite free module \(R^{\oplus I}\).
\end{proof}

\begin{lemma}[Affine sections of a finitely presented module sheaf are
  finitely generated, chart form]
  \label{mk-lem:module_finite_sections_qcdata}
  \dcref{custom}

  \uses{mk-lem:module_finite_of_tilde_genSections,
    mk-lem:tildeIso_of_isQuasicoherent_isAffineOpen}
  \textit{Source: Stacks~\href{https://stacks.math.columbia.edu/tag/01PC}{01PC}.}
  Let \(X\) be a scheme, \(\mathcal F\) a quasi-coherent
  \(\mathcal O_X\)-module, and \(q\) a quasi-coherence datum for
  \(\mathcal F\) all of whose presentations are finite. For every affine open
  \(V \subseteq X\) contained in a cover member \(q.X_i\), the section module
  \(\Gamma(\mathcal F, V)\) is a finite \(\Gamma(X, V)\)-module.
\end{lemma}
\begin{proof}
  The finite presentation of the slice restriction
  \(\mathcal F|_{q.X_i}\) provides a generating family of sections indexed by
  a finite type. Generating families transport along colimit-preserving
  functors carrying the structure sheaf to the structure sheaf, and along
  epimorphisms; apply this three times: across the slice-to-geometric
  equivalence onto the open subscheme \(q.X_i\); along the pullback by the
  open immersion \(k : \operatorname{Spec} \Gamma(X, V) \to q.X_i\) through
  which \(\mathrm{fromSpec}_V\) factors (it factors because
  \(V \subseteq q.X_i\)); and along the composite isomorphism identifying the
  result with the pullback of \(\mathcal F\) along \(\mathrm{fromSpec}_V\),
  which by \ref{mk-lem:tildeIso_of_isQuasicoherent_isAffineOpen} is
  \(\widetilde{\Gamma(\mathcal F, V)}\). The tilde of
  \(\Gamma(\mathcal F, V)\) thus has a finite generating family, and
  \ref{mk-lem:module_finite_of_tilde_genSections} concludes.
\end{proof}

\begin{lemma}[Charts with finitely generated sections]
  \label{mk-lem:affine_finite_sections_nhds}
  \dcref{custom}

  \uses{mk-lem:module_finite_sections_qcdata}
  \textit{Source: Stacks~\href{https://stacks.math.columbia.edu/tag/01PC}{01PC}.}
  Let \(\mathcal F\) be a module sheaf of finite presentation on a scheme
  \(X\). Every point \(x\) of every open \(O \subseteq X\) has an affine open
  neighbourhood \(V\) with \(x \in V \subseteq O\) such that
  \(\Gamma(\mathcal F, V)\) is a finite \(\Gamma(X, V)\)-module.
\end{lemma}
\begin{proof}
  Choose a quasi-coherence datum \(q\) with finite presentations. The cover
  members \(q.X_i\) cover \(X\), so \(x\) lies in some \(q.X_i\); affine opens
  form a basis of the topology, so there is an affine
  \(V \subseteq O \cap q.X_i\) containing \(x\). Conclude by
  \ref{mk-lem:module_finite_sections_qcdata}.
\end{proof}


The base map is obtained from the unit of the pullback--pushforward
adjunction. Its compatibility with the affine tilde comparisons therefore
follows from four identities: naturality in the module, compatibility with
composition of pullbacks, invariance under equality of the underlying
morphisms, and inversion of the comparison for an affine open immersion.
Together these identities identify the composite used in
\ref{mk-def:pullback_app_isoTensor_sigma} with the canonical base map on every
section.

\begin{lemma}
[(N1) baseMap naturality in input sheaf]
  \label{mk-lem:baseMap_pullbackComp_apply}
  \dcref{custom}

  \uses{mk-def:quot_canonical_basechange_map, mk-lem:pullback_tildeIso}
  Let \(g : X \to Y\) be a morphism of schemes, let \(\mathcal N\),
  \(\mathcal N'\) be \(\mathcal O_Y\)-modules, and let
  \(\eta : \mathcal N \to \mathcal N'\) be a morphism of
  \(\mathcal O_Y\)-modules. For affine opens \(U \subseteq Y\) and
  \(V \subseteq X\) with \(g(V) \subseteq U\) and any
  \(x \in \Gamma(\mathcal N, U)\), the section-level component of
  \ref{mk-def:quot_canonical_basechange_map} satisfies
  \[
    \bigl((\text{pullback}\,g).\mathrm{map}\,\eta\bigr).\mathrm{app}\,V\;
    \bigl(\text{baseMap}\,g\,\mathcal N\,e\,x\bigr)
    \;=\;
    \text{baseMap}\,g\,\mathcal N'\,e\,(\eta.\mathrm{app}\,U\,x),
  \]
  i.e.\ \texttt{baseMap} is natural in its module argument. Used in
  Stage~2 of the \ref{mk-def:pullback_app_isoTensor_sigma}
  decomposition to transport \(\text{step1}.\mathrm{inv}\) (a
  morphism of \(\mathcal O_Y\)-modules) across \texttt{baseMap} of
  \(\text{Spec.map}\,\varphi\).
\end{lemma}
\begin{proof}
  \uses{mk-def:quot_canonical_basechange_map}
  Direct from the naturality square of
  \texttt{pullbackPushforwardAdjunction.unit} at the morphism
  \(\eta : \mathcal N \to \mathcal N'\) in the input category
  (giving the unit half), combined with the naturality of the
  section restriction against \((\text{pullback}\,g).\mathrm{map}\,\eta\)
  (giving the restriction half); the two squares chain elementwise.
\end{proof}

\begin{lemma}
[(N2) baseMap composition via pullbackComp]
  \label{mk-lem:baseMap_pullback_comp_apply}
  \dcref{custom}

  \uses{mk-def:quot_canonical_basechange_map, mk-lem:pullback_tildeIso}
  Let \(f : X \to Y\) and \(g : Y \to Z\) be morphisms of schemes and
  let \(\mathcal N\) be an \(\mathcal O_Z\)-module. The
  Mathlib-canonical iso
  \(\text{pullbackComp}\,f\,g :
    (\text{pullback}\,f) \circ (\text{pullback}\,g)
    \cong \text{pullback}\,(f \circ g)\)
  intertwines \texttt{baseMap}: for affine opens \(W \subseteq Z\),
  \(V \subseteq X\) with \(f(g(V)) \subseteq W\) and any
  \(x \in \Gamma(\mathcal N, W)\),
  \[
    ((\text{pullbackComp}\,f\,g)\,\mathcal N).\mathrm{hom}.\mathrm{app}\,V\;
    \bigl(\text{baseMap}\,f\,((\text{pullback}\,g).\mathrm{obj}\,\mathcal N)\,
       e_f\,(\text{baseMap}\,g\,\mathcal N\,e_g\,x)\bigr)
    \;=\;
    \text{baseMap}\,(f \circ g)\,\mathcal N\,e\,x.
  \]
  In words: applying \texttt{baseMap} along \(g\) then along \(f\),
  routed through the canonical composition iso, yields the
  \texttt{baseMap} along the composite \(f \circ g\). Used in Stages~3
  and~5 of the \ref{mk-def:pullback_app_isoTensor_sigma}
  decomposition.
\end{lemma}
\begin{proof}
  \uses{mk-def:quot_canonical_basechange_map}
  The unit of a composition of left adjoints
  \(F \circ G \dashv R \circ R'\) satisfies the canonical composition
  identity \(\eta_{FG} = R'(\eta_F) \circ \eta_G\) (Mathlib
  \texttt{Adjunction.comp}). The composed unit is identified with the
  \(\text{pullbackComp}\)-twist of the unit of the composite by the
  conjugation identity \texttt{unit\_conjugateEquiv} applied to
  Mathlib's \texttt{conjugateEquiv\_pullbackComp\_inv} (the components
  of \texttt{pushforwardComp} are identities). Together with the
  naturality of the \(f\)-unit component against the restriction
  \(V \le g^{-1}U\), the collapse of consecutive section
  restrictions, the naturality of
  \(\text{pullbackComp}.\mathrm{inv}\) against the restriction
  \(T \le (f \circ g)^{-1}U\), and the hom--inv cancellation of
  \(\text{pullbackComp}\) at \(T\)-sections, this yields the stated
  formula elementwise. Used twice (Stages~3 and~5).
\end{proof}

\begin{lemma}
[(N3) baseMap transport via pullbackCongr]
  \label{mk-lem:baseMap_pullbackCongr_apply}
  \dcref{custom}

  \uses{mk-def:quot_canonical_basechange_map, mk-lem:pullback_tildeIso}
  Let \(f_1, f_2 : X \to Y\) be morphisms of schemes with
  \(f_1 = f_2\) as a propositional equality \(h_{eq}\), and let
  \(\mathcal N\) be an \(\mathcal O_Y\)-module. The Mathlib-canonical
  equality-transport iso
  \(\text{pullbackCongr}\,h_{eq} : \text{pullback}\,f_1
    \cong \text{pullback}\,f_2\)
  intertwines \texttt{baseMap}: for affine opens \(U \subseteq Y\),
  \(V \subseteq X\) with \(f_1(V) \subseteq U\) (equivalently
  \(f_2(V) \subseteq U\) via \(h_{eq}\)) and any
  \(x \in \Gamma(\mathcal N, U)\),
  \[
    ((\text{pullbackCongr}\,h_{eq})\,\mathcal N).\mathrm{inv}.\mathrm{app}\,V\;
    \bigl(\text{baseMap}\,f_2\,\mathcal N\,e_2\,x\bigr)
    \;=\;
    \text{baseMap}\,f_1\,\mathcal N\,e_1\,x.
  \]
  I.e.\ \texttt{baseMap} is invariant under the equality-transport iso
  of the morphism argument. Used in Stage~4 of the
  \ref{mk-def:pullback_app_isoTensor_sigma} decomposition to re-associate
  the composite
  \(\text{Spec.map}\,\varphi \circ \mathtt{\_hV.fromSpec}
    = \mathtt{\_hU.fromSpec} \circ g\) under the Σ-pair packaging.
\end{lemma}
\begin{proof}
  \uses{mk-def:quot_canonical_basechange_map}
  By propositional-equality elimination (\texttt{subst} on
  \(h_{eq}\)): after substitution \(\text{pullbackCongr}\,\mathrm{rfl}\)
  is the identity isomorphism by iota-reduction of \texttt{eqToHom},
  and the two \(\le\)-witnesses are proof-irrelevant, so both sides
  are definitionally equal.
\end{proof}

\begin{lemma}
[(N4) step3 inversion identity for an open immersion]
  \label{mk-lem:baseMap_inv_step3_open_immersion}
  \dcref{custom}

  \uses{mk-def:quot_canonical_basechange_map,
        mk-lem:pullback_of_openImmersion_iso_restrict}
  Let \(Y\) be a scheme, \(U \subseteq Y\) an affine open with
  canonical open immersion
  \(\mathtt{\_hU.fromSpec} : \text{Spec}\,\Gamma(Y, U) \to Y\), and
  \(\mathcal N\) an \(\mathcal O_Y\)-module. Then the canonical
  section-transport isomorphism of
  \ref{mk-lem:pullback_of_openImmersion_iso_restrict} is the inverse
  of \texttt{baseMap} along \(\mathtt{\_hU.fromSpec}\) at the
  \(\top\)-section:
  \[
    \text{step3}\bigl(
      \text{baseMap}\,\mathtt{\_hU.fromSpec}\,
        \mathcal N\,
        (\text{le\_of\_eq}\,\mathtt{\_hU.fromSpec\_preimage\_self.symm})\,y
    \bigr) \;=\; y
    \qquad\text{for every } y \in \Gamma(\mathcal N, U).
  \]
  This identifies the inverse of \texttt{step3} with the natural
  \texttt{baseMap} action of the open immersion
  \(\mathtt{\_hU.fromSpec}\). Used in Stage~6 of the
  \ref{mk-def:pullback_app_isoTensor_sigma} decomposition, closing the
  intertwining identity.
\end{lemma}
\begin{proof}
  \uses{mk-def:quot_canonical_basechange_map,
        mk-lem:pullback_of_openImmersion_iso_restrict}
  Realized as the \(\Sigma\)-pair characterization carried by
  \ref{mk-lem:pullback_of_openImmersion_iso_restrict}: the inverse of
  its equivalence IS the canonical base map. The proof is the
  unit-compatibility of \texttt{Adjunction.leftAdjointUniq} for
  \texttt{restrictAdjunction} versus
  \texttt{pullbackPushforwardAdjunction} (whose
  \texttt{leftAdjointUniq} is Mathlib's
  \texttt{restrictFunctorIsoPullback}), combined with the fact that
  the unit of the restriction adjunction is a plain presheaf
  restriction (definitional in Mathlib), the naturality of the
  comparison against the restriction
  \(\top \le j^{-1}U\), and the collapse of consecutive
  restrictions along the equality \(j(\top) = U\). The stated
  inversion identity then follows by applying the equivalence to both
  sides of the characterization.
\end{proof}

\section{Construction overview}
The construction proceeds in three phases.

\paragraph{Phase 1: Hilbert polynomial (\ref{mk-def:hilbert_polynomial}).}
On a locally-of-finite-type morphism \(\pi : X \to S\) with \(S\) locally
noetherian, the fibers determine a function
\(s \mapsto \Phi_{\mathcal{F},s} \in \mathbb{Q}[\lambda]\)
(\ref{mk-def:fiber_module_restriction}, \ref{mk-def:fiber_sections_module},
\ref{mk-def:fiber_hilbert_function}, \ref{mk-lem:hilbertPolynomial_unique},
\ref{mk-lem:hilbertPoly_of_isRatHilb}). The \(H^0\)-only route through the
  graded Hilbert function of the section module reduces polynomiality to
  finite generation (\ref{mk-lem:sectionGradedModule_fg}) and the
  Hilbert--Serre theorem (\ref{mk-lem:gradedHilbertSerre_rational}), yielding
  \ref{mk-thm:hilbertPoly_of_sectionModule}.

\paragraph{Phase 2: Grassmannian scheme (\ref{mk-def:grassmannian_scheme},
\ref{mk-thm:grassmannian_representable}).} The Grassmannian
\(\mathrm{Gr}_S(V, d)\) is obtained by gluing \(\binom{r}{d}\)
affine charts \(U^I\) (each isomorphic to a relative affine space
\(\mathbb{A}^{d(r-d)}_S\)) via the Pl\"ucker cocycle \(\theta_{I,J}\). The
universal quotient sheaf \(\mathcal{U}\) is glued by the transition cocycle
  \(g_{I,J} = (X^I_J)^{-1}\), and the determinant line bundle
  \(\det \mathcal{U}\) gives the Pl\"ucker closed embedding into
  \(\mathbb{P}_S(\bigwedge^d V)\). The universal pair
  \((\mathrm{Gr}_S(V, d), \mathcal{U})\) represents rank-\(d\) locally
  free quotients.

\paragraph{Phase 3: Quot scheme (\ref{mk-thm:quot_representable}).}
The Quot functor is constructed as a locally closed subfunctor of the
Grassmannian
\(\mathrm{Grass}(W \otimes_{\mathcal{O}_S} \mathrm{Sym}^r V, \Phi(r))\), via:
\begin{enumerate}
  \item the uniform \(m\)-regularity bound from
    \ref{mk-lem:quot_boundedness}, using Castelnuovo--Mumford regularity;
  \item the Grassmannian embedding \(\alpha\) from
    \ref{mk-lem:quot_alpha_injective};
  \item the flattening stratification (\ref{mk-thm:flattening_stratification_exists}
    of \ref{mk-chap:Picard_FlatteningStratification}, A.2.a) applied to the
    universal quotient on the Grassmannian, producing the locally closed
    stratum \(T_0^\Phi\);
  \item the valuative-criterion upgrade
    (\ref{mk-lem:quot_valuative_criterion}) to closed embedding, hence
    projective \(S\)-scheme.
\end{enumerate}
The reduction \ref{mk-lem:quot_reduction_to_pi_star_W} from the universal
case to the general case is a closed embedding of representable
subfunctors via the kernel-vanishing scheme.

\paragraph{Cohomology and base change.}
\ref{mk-thm:flat_base_change_cohomology} proves the degree-zero flat-base-change
isomorphism \(g^*f_*\mathcal F\cong f'_*g'^*\mathcal F\) (Stacks tag 02KH).
It is assembled from the canonical base-change map using the affine
tensor/base-map comparison developed above; higher direct images are not
claimed here.

\paragraph{Dependencies summary.}
Relative-spectrum gluing from \ref{mk-chap:Picard_RelativeSpec} constructs
the affine charts, while
\ref{mk-thm:flattening_stratification_exists} from
\ref{mk-chap:Picard_FlatteningStratification} provides the flattening stratum.
Downstream, the special case \(E = \mathcal{O}_C\) of
\ref{mk-thm:quot_representable} on a smooth proper curve \(C/k\) gives the Hilbert
scheme \(\Hilb_{C/k} = \Quot^{\Phi,\mathcal{O}_C(1)}_{\mathcal{O}_C/C/k}\)
and its open subscheme \(\mathrm{Div}_{C/k}\) of relative effective divisors:
the engine of the quotient route to the Picard scheme, which
\ref{mk-sec:fga_pic_setup} sets beside the Milne--Koll\'ar route the formal proof
follows.  The substrate developed in this chapter --- Hilbert polynomial,
Grassmannian scheme, graded algebra --- is consumed by either route.
